%%%%%%%%%%%%%%%%%%%%%%%%%%%%%%%%%%%%%%%%%%%%%%%%%%%%%%%%%%%%%%%%%%%%
%-----------------------------------------------------------------------------
% Beginning of main.tex
%-----------------------------------------------------------------------------
%
% This is a the text file for ``Computing with Modular Forms'.
% Use AMS-LaTeX to typeset it.
%
% This is the driver file.  The body is included below.
% of this file.
%%%%%%%%%%%%%%%%%%%%%%%%%%%%%%%%%%%%%%%%%%%%%%%%%%%%%%%%%%%%%%%%%%%%

\documentclass{gsm-l}
%\documentclass[draft]{gsm-l}
%\usepackage[american]{babel}
% these are for convenience while working with
% the electronic version of the document; then
% can be safely commented out.
%\usepackage[hypertex]{hyperref}
\usepackage[hyperindex,pdfmark]{hyperref}


\bibliographystyle{amsalpha}

\makeindex    

%\includeonly{macros, appendix}
\includeonly{macros, preface, body, solutions, appendix, index}
%\includeonly{macros, solutions}

\numberwithin{section}{chapter}
\numberwithin{equation}{chapter}

\newcommand{\defn}[1]{{\em #1}\index{{\bf Definition}!\protect{#1}}}
%\newcommand{\defn}[1]{{\em #1}\index{\protect{#1}}}
\newcommand{\sym}[1]{$#1$\index{\protect{$#1$}}}

%\newcommand{\sym}[1]{$#1$\index{hi}}
\newcommand{\indexp}[1]{\index{\protect{#1}}}
 

%    Absolute value notation
\newcommand{\abs}[1]{\lvert#1\rvert}

%    Blank box placeholder for figures (to avoid requiring any
%    particular graphics capabilities for printing this document).
\newcommand{\blankbox}[2]{%
  \parbox{\columnwidth}{\centering
%    Set fboxsep to 0 so that the actual size of the box will match the
%    given measurements more closely.
    \setlength{\fboxsep}{0pt}%
    \fbox{\raisebox{0pt}[#2]{\hspace{#1}}}%
  }%
}

\usepackage{amssymb}
\usepackage{fancybox}
\usepackage{mathrsfs}
\usepackage{subfigure}
\usepackage{graphicx}
\usepackage{psfrag}  
\usepackage[all]{xy}
\usepackage{sage}


\usepackage{xspace}  % to allow for text macros that don't eat space 
%\newcommand{\SAGE}{{\sf SAGE}\xspace\index{\protect{\sf SAGE}}}
%\newcommand{\SAGE}{{\sf SAGE}\xspace}
\newcommand{\SAGE}{{\sf SAGE}\index{\protect{\sf SAGE}}}

\newcommand{\sage}{\SAGE}
\newcommand{\sageindex}[1]{\index{\protect{\sf SAGE}!#1}}


\newcommand{\an}{{\rm an}}
\newcommand{\examplesession}[1]{}
\newcommand{\gzero}{\Gamma_0(N)}
\newcommand{\esM}{\overline{\sM}}
\newcommand{\sM}{\mathbb{M}}
\newcommand{\sS}{\mathbb{S}}
\newcommand{\sB}{\mathbb{B}}
\newcommand{\eE}{\mathbb{E}}
\newcommand{\Adual}{A^{\vee}}
\newcommand{\Bdual}{B^{\vee}}


\renewcommand{\div}{\mbox{\rm div}}
\newcommand{\Eis}{E}
\newcommand{\unity}{\mathbf{1}}
\newcommand{\dz}{dz}
\renewcommand{\Re}{\mbox{\rm Re}}
\newcommand{\pfSel}{\widehat{\Sel}}
\newcommand{\qe}{\stackrel{\mbox{\tiny ?}}{=}}
\newcommand{\isog}{\simeq}
\newcommand{\e}{\mathbf{e}}
\newcommand{\bN}{\mathbf{N}}
\newcommand{\set}{=}
\newcommand{\Gam}{\Gamma}
\renewcommand{\Im}{\text{Im}}
\newcommand{\X}{\mathcal{X}}
\newcommand{\cH}{\mathcal{H}}
\newcommand{\cA}{\mathcal{A}}
\newcommand{\cF}{\mathcal{F}}
\newcommand{\cG}{\mathcal{G}}
\newcommand{\cJ}{\mathcal{J}}
\newcommand{\cL}{\mathcal{L}}
\newcommand{\cV}{\mathcal{V}}
\newcommand{\cO}{\mathcal{O}}
\newcommand{\cQ}{\mathcal{Q}}
\newcommand{\cX}{\mathcal{X}}
\newcommand{\ds}{\displaystyle}
\newcommand{\M}{\sM}
\newcommand{\B}{\sB}
\newcommand{\E}{\sE}
\renewcommand{\L}{\mathcal{L}}
\newcommand{\J}{\mathcal{J}}
\newcommand{\dw}{\delta(w)} 
\newcommand{\dwh}{\widehat{\delta(w)}}      
\newcommand{\dlwh}{\widehat{\delta_\l(w)}} 
\newcommand{\dash}{-\!\!\!\!-\!\!\!\!-\!\!\!\!-} 
\newcommand{\Frobl}{\Frob_{\ell}}
\newcommand{\tE}{\tilde{E}}
\renewcommand{\l}{\ell}
\renewcommand{\t}{\tau}
\newcommand{\N}{\mathcal{N}}
\newcommand{\U}{\mathcal{U}}
\newcommand{\Kbar}{\overline{K}}
\newcommand{\Lbar}{\overline{L}}
\newcommand{\gammabar}{\overline{\gamma}}
\newcommand{\q}{\mathbf{q}}
\renewcommand{\star}{\times}
\newcommand{\gM}{\mathfrak{M}}
\newcommand{\gA}{\mathfrak{A}}
\newcommand{\gP}{\mathfrak{P}}
\newcommand{\bmu}{\boldsymbol{\mu}}
\newcommand{\union}{\cup}
\newcommand{\Tl}{T_{\ell}}
\newcommand{\into}{\rightarrow}
\newcommand{\onto}{\rightarrow\!\!\!\!\rightarrow}
\newcommand{\intersect}{\cap}
\newcommand{\meet}{\cap}
\newcommand{\cross}{\times}
\newcommand{\rb}{\overline{\rho}}
\newcommand{\ra}{\rightarrow}
\newcommand{\xra}[1]{\xrightarrow{#1}}
\newcommand{\hra}{\hookrightarrow}
\newcommand{\kr}[2]{\left(\frac{#1}{#2}\right)}
\newcommand{\la}{\leftarrow}
\newcommand{\lra}{\longrightarrow}
\newcommand{\riso}{\xrightarrow{\sim}}
\newcommand{\da}{\downarrow}
\newcommand{\ua}{\uparrow}
\newcommand{\con}{\equiv}
\newcommand{\Gm}{\mathbb{G}_m}
\newcommand{\pni}{\par\noindent}
\newcommand{\iv}{^{-1}}
\newcommand{\alp}{\alpha}
\newcommand{\bq}{\mathbf{q}}
\newcommand{\cpp}{{\tt C++}}
\newcommand{\tensor}{\otimes}
\newcommand{\bg}{{\tt BruceGenus}}
\newcommand{\abcd}[4]{\left(
        \begin{smallmatrix}#1&#2\\#3&#4\end{smallmatrix}\right)}
\newcommand{\mthree}[9]{\left(
        \begin{matrix}#1&#2&#3\\#4&#5&#6\\#7&#8&#9
        \end{matrix}\right)}
\newcommand{\mtwo}[4]{\left(
        \begin{matrix}#1&#2\\#3&#4
        \end{matrix}\right)}
\newcommand{\vtwo}[2]{\left(
        \begin{matrix}#1\\#2
        \end{matrix}\right)}
\newcommand{\smallmtwo}[4]{\left(
        \begin{smallmatrix}#1&#2\\#3&#4
        \end{smallmatrix}\right)}
\newcommand{\smallvtwo}[2]{\left(
        \begin{smallmatrix}#1\\#2
        \end{smallmatrix}\right)}
\newcommand{\twopii}{2\pi{}i{}}  
\newcommand{\eps}{\varepsilon}
\newcommand{\vphi}{\varphi}
\newcommand{\gp}{\mathfrak{p}}
\newcommand{\W}{\mathcal{W}}
\newcommand{\oz}{\overline{z}}
\newcommand{\Zpstar}{\Zp^{\star}}
\newcommand{\Zhat}{\widehat{\Z}}
\newcommand{\Zbar}{\overline{\Z}}
\newcommand{\Zl}{\Z_{\ell}}
\newcommand{\Q}{\mathbb{Q}}
\newcommand{\GQ}{G_{\Q}}
\newcommand{\R}{\mathbb{R}}
\newcommand{\D}{{\mathbf D}}
\newcommand{\cC}{\mathcal{C}}
\newcommand{\cD}{\mathcal{D}}
\newcommand{\cP}{\mathcal{P}}
\newcommand{\cS}{\mathcal{S}}
\newcommand{\Sbar}{\overline{S}}
\newcommand{\K}{{\mathbf K}}
\newcommand{\C}{\mathbb{C}}
\newcommand{\Cp}{{\mathbf C}_p}
\newcommand{\Sets}{\mbox{\rm\bf Sets}}
\newcommand{\bcC}{\boldsymbol{\mathcal{C}}}
\renewcommand{\P}{\mathbb{P}}
\newcommand{\Qbar}{\overline{\Q}}
\newcommand{\kbar}{\overline{k}}
\newcommand{\dual}{\bot}
\newcommand{\T}{\mathbb{T}}
\newcommand{\tT}{\tilde{\T}}
\newcommand{\calT}{\mathcal{T}}
\newcommand{\cT}{\mathcal{T}}
\newcommand{\cbT}{\mathbf{\mathcal{T}}}
\newcommand{\cU}{\mathcal{U}}
\newcommand{\Z}{\mathbb{Z}}
\newcommand{\F}{\mathbb{F}}
\newcommand{\Fl}{\F_{\ell}}
\newcommand{\Fell}{\Fl}
\newcommand{\Flbar}{\overline{\F}_{\ell}}
\newcommand{\Flnu}{\F_{\ell^{\nu}}}
\newcommand{\Fbar}{\overline{\F}}
\newcommand{\Fpbar}{\overline{\F}_p}
\newcommand{\fbar}{\overline{f}}
\newcommand{\Qp}{\Q_p}
\newcommand{\Ql}{\Q_{\ell}}
\newcommand{\Qell}{\Q_{\ell}}
\newcommand{\Qlbar}{\overline{\Q}_{\ell}}
\newcommand{\Qlnr}{\Q_{\ell}^{\text{nr}}}
\newcommand{\Qlur}{\Q_{\ell}^{\text{ur}}}
\newcommand{\Qltm}{\Q_{\ell}^{\text{tame}}}
\newcommand{\Qv}{\Q_v}
\newcommand{\Qpbar}{\Qbar_p}
\newcommand{\Zp}{\Z_p}
\newcommand{\Fp}{\F_p}
\newcommand{\Fq}{\F_q}
\newcommand{\Fqbar}{\overline{\F}_q}
\newcommand{\Ad}{Ad}
\newcommand{\adz}{\Ad^0}
\renewcommand{\O}{\mathcal{O}}
\newcommand{\A}{\mathcal{A}}
\newcommand{\Og}{O_{\gamma}}
\newcommand{\isom}{\cong}
\newcommand{\ncisom}{\approx}   % noncanonical isomorphism
\newcommand{\galq}{\Gal(\Qbar/\Q)}
\newcommand{\rhobar}{\overline{\rho}}
\newcommand{\cM}{\mathcal{M}}
\newcommand{\cB}{\mathcal{B}}
\newcommand{\cE}{\mathcal{E}}
\newcommand{\cR}{\mathcal{R}}
\newcommand{\et}{\text{\rm\'et}}
\newcommand{\hd}[1]{\vspace{1ex}\noindent{\bf #1} }
\newcommand{\nf}[1]{\underline{#1}} 
\newcommand{\cbar}{\overline{c}}
\newcommand{\sltwoz}{\SL_2(\Z)}
\newcommand{\sltwo}{\SL_2}
\newcommand{\gltwoz}{\GL_2(\Z)}
\newcommand{\gltwoq}{\GL_2(\Q)}
\newcommand{\gltwo}{\GL_2}
\newcommand{\gln}{\GL_n}
\newcommand{\psltwoz}{\PSL_2(\Z)}
\newcommand{\psltwo}{\PSL_2}
\newcommand{\h}{\mathfrak{h}}
\renewcommand{\a}{\mathfrak{a}}
\newcommand{\p}{\mathfrak{p}}
\newcommand{\m}{\mathfrak{m}}
\newcommand{\trho}{\tilde{\rho}}
\newcommand{\rhol}{\rho_{\ell}}
\newcommand{\rhoss}{\rho^{\text{ss}}}
\newcommand{\dbd}[1]{\langle #1 \rangle}
\renewcommand{\H}{\HH}
\newcommand{\smallzero}{\left(\begin{smallmatrix}0&0\\0&0
                        \end{smallmatrix}\right)}
\newcommand{\smallone}{\left(\begin{smallmatrix}1&0\\0&1
                        \end{smallmatrix}\right)}
\newcommand{\pari}{{\sc Pari}\index{Pari}}
\newcommand{\python}{{\sc Python}\index{Python}}
\newcommand{\magma}{{\sc Magma}\index{Magma}}
\newcommand{\manin}{{\sc Manin}\index{Manin}}
\newcommand{\zbar}{\overline{z}}
\newcommand{\dzbar}{d\overline{z}}

%%%%%%%%%%%%%%%%%%%%%%%%%%%%%%%%%%%%%%%%%%%%%%%%%%%%%%%%%%%%%%%%%%

%%%% Theoremstyles
\theoremstyle{plain}
\newtheorem{theorem}{Theorem}[chapter]
\newtheorem{proposition}[theorem]{Proposition}
\newtheorem{corollary}[theorem]{Corollary}
\newtheorem{claim}[theorem]{Claim}
\newtheorem{lemma}[theorem]{Lemma}
\newtheorem{conjecture}[theorem]{Conjecture}

\theoremstyle{definition}
\newtheorem{definition}[theorem]{Definition}
\newtheorem{alg}[theorem]{Algorithm}
\newtheorem{question}[theorem]{Question}
\newtheorem{problem}[theorem]{Problem}

\newtheorem{goal}[theorem]{Goal}
\newtheorem{remark}[theorem]{Remark}
\newtheorem{remarks}[theorem]{Remarks}
\newtheorem{example}[theorem]{Example}
\newtheorem{warning}[theorem]{Warning}

\numberwithin{section}{chapter}
\numberwithin{equation}{section}
\numberwithin{figure}{section}
\numberwithin{table}{section}


%%%%%%%%%%%%%%%%%%%%%%%%%%%%%%%%%%%%%%%%%%%%%%%%%%
%% Environments
%%%%%%%%%%%%%%%%%%%%%%%%%%%%%%%%%%%%%%%%%%%%%%%%%%
\newenvironment{algorithm}[1]{%
\begin{alg}[#1]\index{{\bf Algorithm}!#1}\em
}%
{\end{alg}}
\newenvironment{steps}%
{\rm \begin{enumerate}\setlength{\itemsep}{0.1ex}}{\end{enumerate}}

\newcommand{\exref}[2]{Exercise~\ref{#1}.\ref{#2}}

\newenvironment{exercises}
{\section{Exercises}
 \renewcommand{\labelenumi}{\thechapter.\theenumi}
  \begin{enumerate}
}{\end{enumerate}}


%%%%%%%%%%%%%%%%%%%%%%%%%%%%%%%%%%%%%%%%%%%%%%%%%%%%%%%%

\DeclareMathOperator{\rad}{rad}
\DeclareMathOperator{\Ext}{Ext}
\DeclareMathOperator{\Tor}{Tor}
\DeclareMathOperator{\ab}{ab}
\DeclareMathOperator{\Aut}{Aut}
\DeclareMathOperator{\Frob}{Frob}
\DeclareMathOperator{\Fr}{Fr}
\DeclareMathOperator{\Ver}{Ver}
\DeclareMathOperator{\Norm}{Norm}
\DeclareMathOperator{\norm}{norm}
\DeclareMathOperator{\disc}{disc}
\DeclareMathOperator{\ord}{ord}
\DeclareMathOperator{\GL}{GL}
\DeclareMathOperator{\PSL}{PSL}
\DeclareMathOperator{\PGL}{PGL}
\DeclareMathOperator{\Gal}{Gal}
\DeclareMathOperator{\SL}{SL}
\DeclareMathOperator{\SO}{SO}
\DeclareMathOperator{\Stab}{Stab}
\DeclareMathOperator{\tr}{tr}
\DeclareMathOperator{\order}{order}
\DeclareMathOperator{\ur}{ur}
\DeclareMathOperator{\Tr}{Tr}
\DeclareMathOperator{\Hom}{Hom}
\DeclareMathOperator{\Mor}{Mor}
\DeclareMathOperator{\HH}{H}
\DeclareMathOperator{\rref}{rref}
\DeclareMathOperator{\diag}{diag}
\DeclareMathOperator{\her}{her}
\DeclareMathOperator*{\Max}{Max}
\DeclareMathOperator*{\Min}{Min}
\DeclareMathOperator{\support}{supp}
\DeclareMathOperator{\Comm}{comm}
\DeclareMathOperator{\chr}{char}
\DeclareMathOperator{\denom}{denom}
\DeclareMathOperator{\ind}{ind}
\DeclareMathOperator{\im}{im}
\DeclareMathOperator{\oo}{\infty}
\DeclareMathOperator{\lcm}{lcm}
\DeclareMathOperator{\cores}{cores}
\DeclareMathOperator{\coker}{coker}
\DeclareMathOperator{\image}{image}
\DeclareMathOperator{\prt}{part}
\DeclareMathOperator{\proj}{proj}
\DeclareMathOperator{\Br}{Br}
\DeclareMathOperator{\Ann}{Ann}
\DeclareMathOperator{\End}{End}
\DeclareMathOperator{\Tan}{Tan}
\DeclareMathOperator{\Pic}{Pic}
\DeclareMathOperator{\Vol}{Vol}
\DeclareMathOperator{\Vis}{Vis}
\DeclareMathOperator{\Reg}{Reg}
\DeclareMathOperator{\Res}{Res}
\DeclareMathOperator{\rank}{rank}
\DeclareMathOperator{\Sel}{Sel}
\DeclareMathOperator{\Mat}{Mat}
\DeclareMathOperator{\BSD}{BSD}
\DeclareMathOperator{\id}{id}
\DeclareMathOperator{\Selmer}{Selmer}
\DeclareMathOperator{\SU}{SU}
\DeclareMathOperator{\SP}{Sp}
\DeclareMathOperator{\UN}{U}
\DeclareMathOperator{\Id}{Id}
\DeclareMathOperator{\vcd}{vcd}
\DeclareMathOperator{\Trace}{Trace}
\DeclareMathOperator{\new}{new}
\DeclareMathOperator{\Morph}{Morph}
\DeclareMathOperator{\old}{old}
\DeclareMathOperator{\Sym}{Sym}
\DeclareMathOperator{\Symb}{Symb}
\DeclareMathOperator{\tor}{tor}  
\DeclareMathOperator{\cond}{cond}
\DeclareMathOperator{\Spec}{Spec}
\DeclareMathOperator{\Div}{Div}
\DeclareMathOperator{\Jac}{Jac}
\DeclareMathOperator{\res}{res}
\DeclareMathOperator{\Ker}{Ker}
\DeclareMathOperator{\Coker}{Coker}
\DeclareMathOperator{\sep}{sep}
\DeclareMathOperator{\sign}{sign}
\DeclareMathOperator{\unr}{unr}
\DeclareMathOperator{\Char}{char}
\DeclareMathOperator{\md}{mod}
\DeclareMathOperator{\toric}{toric}
\DeclareMathOperator{\tors}{tors}
\DeclareMathOperator{\Frac}{Frac}
\DeclareMathOperator{\corank}{corank}
\DeclareMathOperator{\GSp}{GSp}



\begin{document}

\frontmatter
\title{Modular Forms:\\A Computational Approach}

%    Information for first author
\author{William A. Stein\\(with an appendix by Paul E. Gunnells)}
%    Address of record for the research reported here
% i changed this to the more formal ``Mathematics'' rather than ``Math''
\address{Department of Mathematics, University of Washington}
%    Current address
%\curraddr{Dept. of Mathematics, University of Washington, Seattle, WA 98195-4350}
\email{wstein@math.washington.edu}
%    \thanks will become a 1st page footnote.
 

% suggestion (first and second is a little awkward because the usual
% meaning in mathematics is lexicographic order, not importance of
% contribution, but in this case the
% alphabetical order of our names doesn't correspond to our
% contributions)
% \thanks{Stein was supported in part by NSF Grant DMS 05-55776.
% Gunnells was supported in part by NSF Grants DMS 02-45580 and DMS
% 04-01525.}

%    Information for second author
%\author{}
% our statistics people are a vociferous minority
\address{Department of Mathematics and Statistics, University of Massachusetts}
\email{gunnells@math.umass.edu}
%\thanks{The second author was supported in part by NSF Grant \#.}

\date{June 1, 2006}
\subjclass{Primary 11;\\Secondary 11-04}
\keywords{abelian varieties, cohomology of arithmetic groups,
  computation, elliptic curves,
  Hecke operators,
  modular curves, modular forms,  modular symbols, Manin
  symbols, number theory}


\begin{abstract}
This is a textbook about algorithms for computing with
modular forms.  It is nontraditional in that the primary focus is not
on underlying theory; instead, it answers the question {\em ``how do you 
explicitly compute spaces of modular forms?''}
\end{abstract}

\maketitle

\setcounter{page}{4}

\newpage
To my grandmother, Annette Maurer. 

\tableofcontents

%-----------------------------------------------------------------------------
% Beginning of preface.tex
%-----------------------------------------------------------------------------
%
% AMS-LaTeX 1.2 sample file for a monograph, based on amsbook.cls.
% This is a data file input by chapter.tex.
%%%%%%%%%%%%%%%%%%%%%%%%%%%%%%%%%%%%%%%%%%%%%%%%%%%%%%%%%%%%%%%%%%%%%%%%

\chapter*{Preface}


This is a graduate-level textbook about algorithms for computing with
modular forms.  It is nontraditional in that the primary focus is not
on underlying theory; instead, it answers the question {\em ``how do you use a
computer to explicitly compute spaces of modular forms?''}

This book emerged from notes for a course the author
taught at Harvard University in 2004, a course
at UC San Diego in 2005, and a course at
the University of Washington in 2006. 

The author has spent years trying to find good practical ways to
compute with classical modular forms for congruence subgroups of
$\SL_2(\Z)$ and has implemented most of these algorithms several
times, first in C++ \cite{stein:hecke}, then in MAGMA \cite{magma},
and as part of the free open source computer
algebra system \sage{} (see \cite{sage}).  Much of this work has involved
turning formulas and constructions buried in obscure research papers
into precise computational recipes then testing these 
and eliminating inaccuracies.
%The goal of this book is to explain some of what I
%have learned along the way.
%, and also describe unsolved problems whose
%solution would move the theory forward.

The author is aware of no other textbooks on computing with modular
forms, the closest work being Cremona's book \cite{cremona:algs},
which is about computing with elliptic curves, and Cohen's book
\cite{cohen:course_ant} about algebraic number theory.  
%There are
%several Ph.D. thesis on computing modular forms, including
%\cite{}\edit{me, gabor, lassine, belgium mod p guy, maybe denis
%  charles???}.  
%Though many of the techniques described in this book
%have been around for decades, the field is still not yet mature, and
%there are missing details and potential improvements to many of the
%algorithms, which you the reader might fill in, and which would be
%greatly appreciated by other mathematicians.

In this book we focus on how to compute {\em in practice} the spaces
$M_k(N,\eps)$ of modular forms, where $k\geq 2$ is an integer
and~$\eps$ is a Dirichlet character of modulus~$N$ (the appendix
treats modular forms for higher rank groups).  We spend the most
effort explaining the general algorithms that appear so far to be the best (in
practice!) for such computations.  We will not discuss in any detail
computing with quaternion algebras,
half-integral weight forms, weight 1 forms, forms for
noncongruence subgroups or groups other than $\GL_2$, Hilbert and
Siegel modular forms, trace formulas, $p$-adic modular forms, and
modular abelian varieties, all of which are topics for additional
books.  We also rarely analyze the complexity of the
algorithms, but instead settle for occasional remarks about their
practical efficiency.  

% \edit{Actually, I probably
%will add some less in-depth discussion of each of these topics.}

%Classical modular forms, which are the main focus of this book, are
%more accessible than more general automorphic forms (see the
%appendix).  
For most of this book we assume the reader has some prior exposure to
modular forms (e.g., \cite{diamond-shurman}), though we recall many of
the basic definitions.  We cite standard books for proofs of the
fundamental results about modular forms that we will use.  The reader
should also be familiar with basic algebraic number theory, linear
algebra, complex analysis (at the level of \cite{ahlfors}), and
algorithms (e.g., know what an algorithm is and what big oh notation
means). In some of the examples and applications we assume that the
reader knows about elliptic curves at the level
of~\cite{silverman:aec}.


Chapter~\ref{ch:modform} is foundational for the rest of this book.  It
introduces congruence subgroups of $\SL_2(\Z)$ and modular forms as
functions on the complex upper half plane.  We discuss $q$-expansions,
which provide an important computational handle on modular forms.  We
also study an algorithm for computing with congruence subgroups.  The
chapter ends with a list of applications of modular forms throughout
mathematics.

In Chapter~\ref{ch:level1} we discuss level 1 modular forms in much
more detail.  In particular, we introduce Eisenstein series and the
cusp form $\Delta$ and describe their $q$-expansions and basic
properties.  Then we prove a structure theorem for level 1 modular
forms and use it to deduce dimension formulas and give an algorithm
for explicitly computing a basis.  We next introduce Hecke operators
on level 1 modular forms, prove several results about them, and deduce
multiplicativity of the Ramanujan $\tau$ function as an application.
We also discuss explicit computation of Hecke operators.  
In Section~\ref{sec:fasttau} we make some brief remarks on
recent work on asymptotically fast computation of
values of $\tau$. Finally, we describe computation of constant terms
of Eisenstein series using an analytic algorithm.  We generalize many
of the constructions in this chapter to higher level in subsequent
chapters.

In Chapter~\ref{ch:modform2} we turn to modular forms of higher level
but restrict for simplicity to weight $2$ since much is clearer in
this case.  (We remove the weight restriction later in
Chapter~\ref{ch:modsym}.)  We describe a geometric way of viewing
cuspidal modular forms as differentials on modular curves, which leads
to modular symbols, which are an explicit way to present a
certain homology group.  This chapter closes with methods for
explicitly computing cusp forms of weight 2 using modular symbols,
which we generalize in Chapter~\ref{ch:modformcomp}.

In Chapter~\ref{ch:dirichlet} we introduce Dirichlet characters, which
are important both in explicit construction of Eisenstein series (in 
Chapter~\ref{ch:eisen}) and
in decomposing spaces of modular forms as direct sums of simpler
spaces.  The main focus of this chapter is a detailed study of
how to explicitly represent and compute with Dirichlet characters.

Chapter~\ref{ch:eisen} is about how to explicitly construct the
Eisenstein subspace of modular forms.  First we define generalized
Bernoulli numbers attached to a Dirichlet character and an integer
then explain a new analytic algorithm for computing them (which
generalizes the algorithm in Chapter~\ref{ch:level1}).  Finally we
give without proof an explicit description of a basis of Eisenstein
series, explain how to compute it, and give some examples.

Chapter~\ref{ch:dim} records a wide range of dimension formulas for
spaces of modular forms, along with a few remarks about where they
come from and how to compute them.  

Chapter~\ref{ch:linalg} is about linear algebra over exact fields,
mainly the rational numbers.  This chapter can be read independently
of the others and does not require any background in modular forms.
Nonetheless, this chapter occupies a central position in this book,
because the algorithms in this chapter are of crucial importance to
any actual implementation of algorithms for computing with modular
forms.

Chapter~\ref{ch:modsym} is the most important chapter in this book; it
generalizes Chapter~\ref{ch:modform2} to higher weight and general
level.  The modular symbols formulation described here is central to
general algorithms for computing with modular forms.

Chapter~\ref{ch:modformcomp} applies the algorithms from
Chapter~\ref{ch:modsym} to the problem of computing with modular forms. First
we discuss decomposing spaces of modular forms using Dirichlet
characters, and then explain how to compute a basis of Hecke eigenforms
for each subspace using several approaches.  We also discuss
congruences between modular forms and bounds needed to provably
generate the Hecke algebra. 

Chapter~\ref{ch:periods} is about computing analytic invariants
of modular forms.  It discusses tricks for speeding convergence
of certain infinite series and sketches how to compute every
elliptic curve over $\Q$ with given conductor.

Chapter~\ref{ch:solutions} contains detailed solutions to most of the
exercises in this book.  (Many of these were written by students
in a course taught at the University of Washington.)

Appendix~\ref{ch:appendix} deals with computational techniques
for working with generalizations of modular forms to more general
groups than $\SL_2(\Z)$, such as $\SL_n(\Z)$ for $n\geq 3$.  Some of
this material requires more prerequisites than the rest of the book.
Nonetheless, seeing a natural generalization of the material in the
rest of this book helps to clarify the key ideas.  The
topics in the appendix are directly related to the main themes of
this book: modular symbols, Manin symbols, cohomology
of subgroups of $\SL_2(\Z)$ with various coefficients, explicit
computation of modular forms, etc.

\medskip

%begin_sage
\vspace{2ex}
\noindent{}{\bf Software.} 
We use \sage{}, Software for Algebra and Geometry Experimentation 
(see \cite{sage}), to illustrate how to do
many of the examples.  \sage{} is completely free and packages
together a wide range of open source mathematics software for doing
much more than just computing with modular forms.  \sage{} can be
downloaded and run on your computer or can be used via a web browser over
the Internet.  The reader is encouraged to experiment with many of the
objects in this book using \sage.  We do not describe the basics of
using \sage{} in this book; the reader should read the \sage{}
tutorial (and other documentation) available at the \sage{} website
\cite{sage}.   All examples in this book have been automatically
tested and should work exactly as indicated in \sage{} version
at least 1.5. 

%end_sage

%The text of this book is licensed under the GNU Free Documentation
%License, Version 1.2, November 2002.  This means that you may freely
%copy this book.  For the precise details, see the Appendix.

%\noindent{}{\bf Acknowledgement:} None yet.
% \mbox{}\vspace{2ex} 

% \noindent{}William Stein\\
% {\tt http://modular.fas.harvard.edu}\\
% {\tt was@math.harvard.edu}

\vspace{2ex}
\noindent{}{\bf Acknowledgements.} 
David Joyner and Gabor Wiese carefully read the book and provided
a huge number of helpful comments.

John Cremona and Kevin Buzzard both made many helpful remarks that
were important in the development of the algorithms in this
book.  Much of the mathematics (and some of the writing) in
Chapter~\ref{ch:periods} is joint work with Helena Verrill.

Noam Elkies made remarks about Chapters~\ref{ch:modform} and
\ref{ch:level1}.  S\'andor Kov\'acs provided interesting comments on
Chapter~\ref{ch:modform}.  Allan Steel provided helpful feedback on
Chapter~\ref{ch:linalg}.  Jordi Quer made useful remarks
about Chapter~\ref{ch:dirichlet} and Chapter~\ref{ch:dim}.

The students in the courses that I taught on this material at Harvard,
San Diego, and Washington provided substantial feedback: in
particular, Abhinav Kumar made numerous observations about computing
widths of cusps (see Section~\ref{sec:cuspwidth}) and Thomas James
Barnet-Lamb made helpful remarks about 
how to represent Dirichlet characters.  James
Merryfield made helpful remarks about complex analytic issues and
about convergence in Stirling's formula.  Robert Bradshaw, Andrew
Crites (who wrote \exref{ch:linalg}{ex:hnf}), Michael Goff, Dustin
Moody, and Koopa Koo wrote most of the solutions included in
Chapter~\ref{ch:solutions} and found numerous typos throughout the
book.  Dustin Moody also carefully read through the book and provided
feedback.

H.~Stark suggested using Stirling's formula in Section~\ref{sec:bernnum},
and Mark Watkins and Lynn Walling made comments on
Chapter~\ref{ch:modform2}.

Justin Walker found typos in the first published version of the book.

Parts of Chapter~\ref{ch:modform} follow Serre's beautiful
introduction to modular forms \cite[Ch.~VII]{serre:arithmetic}
closely, though we adjust the notation, definitions, and order of
presentation to be consistent with the rest of this book. 

I would like to acknowledge the partial support of 
NSF Grant DMS 05-55776.
Gunnells was supported in part by NSF Grants DMS 02-45580 and
 DMS 04-01525.


\vspace{2ex}
\noindent{\bf Notation and Conventions.}
We denote canonical isomorphisms by $\isom$
and noncanonical isomorphisms by $\ncisom$. 
If $V$ is a vector space and $s$ denotes
some sort of construction involving $V$,
we let $V_s$ denote the corresponding subspace
and $V^s$ the quotient space.  E.g., if $\iota$
is an involution of $V$, then $V_+$ is $\Ker(\iota-1)$
and $V^+ = V/\Im(\iota-1)$.  If $A$ is a finite
abelian group, then $A_{\tor}$ denotes the torsion
subgroup and $A/{\tor}$ denotes the quotient $A/A_{\tor}$.
We denote right group actions using exponential
notation.  Everywhere in this book,~$N$ is a positive
integer and $k$ is an integer. 

If $N$ is an integer, a \defn{divisor} $t$ of $N$
is a {\em positive} integer such that $N/t$ is an integer.

%-----------------------------------------------------------------------------
% End of preface.tex
%-----------------------------------------------------------------------------


\mainmatter

\chapter{Modular Forms}\label{ch:modform}

This chapter introduces modular forms and congruence subgroups, which
are central objects in this book.  We first introduce the upper half
plane and the group $\SL_2(\Z)$ then recall some definitions from
complex analysis.  Next we define modular forms of level $1$ followed
by modular forms of general level.  In Section~\ref{sec:congsgalg} we
discuss congruence subgroups and explain a
simple way to compute generators for them and determine element
membership.  Section~\ref{sec:apps} lists applications of
modular forms.

We assume familiarity with basic number theory, group theory, and
complex analysis.  For a deeper understanding of modular forms, the
reader is urged to consult the standard books in the field, e.g.,
\cite{lang:modular, serre:arithmetic, diamond-im, miyake,
  shimura:intro, koblitz:cong}.  See also \cite{diamond-shurman},
which is an excellent first introduction to the theoretical
foundations of modular forms.


\section{Basic Definitions}
The group
$$
 \SL_2(\R) = 
\left\{ \mtwo{a}{b}{c}{d} : ad-bc=1\text{ and } a,b,c,d\in\R\right\}
$$
acts on the \defn{complex upper half plane}\index{$\h$}
$$
  \h = \{z \in \C : \Im(z) > 0\}
$$
by \defn{linear fractional transformations}, as follows.
\index{$\SL_2(\Z)$}
If
$\gamma = \abcd{a}{b}{c}{d} \in \SL_2(\R)$,
then for any $z\in\h$ we let
\begin{equation}\label{eqn:lft}
  \gamma(z) = \frac{az + b}{cz + d} \in \h.
\end{equation}
Since the determinant of $\gamma$ is $1$, we have
$$
   \left(\frac{d}{dz}\gamma\right)(z) = \frac{1}{(cz+d)^2}.
$$

\begin{definition}[Modular Group]
The \defn{modular group} is 
the group of
all matrices $\abcd{a}{b}{c}{d}$ with $a,b,c,d\in\Z$
and $ad-bc=1$.
\end{definition}
For example, the
matrices 
\begin{equation}\label{eqn:ST}
S=\mtwo{0}{-1}{1}{\hfill 0}
\qquad\text{and} \qquad
T=\mtwo{1}{1}{0}{1}
\end{equation} 
are both elements of $\SL_2(\Z)$;
the matrix $S$ induces the function $z\mapsto -1/z$ on~$\h$,
and~$T$ induces the function $z\mapsto z+1$.
\begin{theorem}\label{thm:sl2zgen}
The group $\SL_2(\Z)$ is generated by $S$ and $T$.
\end{theorem}
\begin{proof}
  See e.g. \cite[\S{}VII.1]{serre:arithmetic}.
\end{proof}

%begin_sage
\sageindex{$\SL_2(\Z)$}
In \sage{} we compute the group $\SL_2(\Z)$ and its generators as follows:
\begin{verbatim}
sage: G = SL(2,ZZ); G
Modular Group SL(2,Z)
sage: S, T = G.gens()
sage: S
[ 0 -1]
[ 1  0]
sage: T
[1 1]
[0 1]
\end{verbatim}
%end_sage

\begin{definition}[Holomorphic and Meromorphic]
  Let $R$ be an open subset of $\C$.  A function $f:R\to \C$ is
  \defn{holomorphic} if $f$ is complex differentiable at every point
  $z\in R$, i.e., for each $z\in R$ the limit 
$$
  f'(z) = \lim_{h\to 0}
  \frac{f(z+h)-f(z)}{h}
$$ 
  exists, where~$h$ may approach~$0$ along any path.
  A function $f: R\to\C\cup \{\infty\}$ is \defn{meromorphic}
    if it is holomorphic except (possibly) at a discrete set $S$ of
    points in~$ R$, and at each $\alpha\in S$ there is a positive
    integer $n$ such that $(z-\alpha)^n f(z)$ is holomorphic at
    $\alpha$.
\end{definition}

The function $f(z) = e^z$ is a holomorphic function on $\C$; in
contrast, $1/(z-i)$ is meromorphic on~$\C$ but not holomorphic
since it has a pole at~$i$.  The function $e^{-1/z}$ is not
even meromorphic on~$\C$. 

Modular forms are holomorphic functions on $\h$ that transform in a
particular way under a certain subgroup of $\SL_2(\Z)$.  Before
defining general modular forms, we define modular forms of level $1$.

\section{Modular Forms of Level $1$}\label{sec:modform1}

\begin{definition}[Weakly Modular Function]
  A \defn{weakly modular function} of \defn{weight} $k\in\Z$ is a
  meromorphic function $f$ on $\h$ such that for all
  $\gamma=\abcd{a}{b}{c}{d}\in\SL_2(\Z)$ and all $z\in\h$ we have
\begin{equation}\label{eqn:modfunc}
   f(z) = (cz+d)^{-k} f(\gamma(z)).
 \end{equation}
\end{definition}
The constant functions are weakly modular of weight $0$.
There are no nonzero weakly modular functions of odd weight
(see \exref{ch:modform}{ex:nomodformodd}), and it is 
not obvious that there are any weakly modular functions
of even weight $k\geq 2$ (but there are, as we will see!). 
The product of two weakly modular functions of weights~$k_1$
and~$k_2$ is a weakly modular function of weight $k_1+k_2$ 
(see \exref{ch:modform}{ex:wmfprod}).

When $k$ is even, \eqref{eqn:modfunc} has a possibly 
more conceptual interpretation; namely
 \eqref{eqn:modfunc} is the same as
$$
  f(\gamma(z))(d(\gamma(z)))^{k/2} = f(z)(dz)^{k/2}.
$$
Thus \eqref{eqn:modfunc} simply says that 
the weight~$k$ ``differential form'' $f(z)(dz)^{k/2}$
is fixed under the action of every element of
$\SL_2(\Z)$.  


By Theorem~\ref{thm:sl2zgen}, the group $\SL_2(\Z)$ is generated by
the matrices $S$ and $T$ of \eqref{eqn:ST}, so to show that a
meromorphic function~$f$ on $\h$ is a weakly modular function, 
all we have to do is show that for all $z\in\h$ we have
\begin{equation}\label{eqn:modfunc2}
  f(z+1) = f(z) \qquad\text{and}\qquad f(-1/z) = z^k f(z).
\end{equation}

Suppose~$f$ is a weakly modular function of weight~$k$.  A
\defn{Fourier expansion} of~$f$, if it exists, is a representation of
$f$ as $f(z) = \sum_{n=m}^{\infty} a_n e^{2\pi i nz}$, for all $z\in
\h$.  Let $q=q(z)=e^{2\pi i z}$, which we view as a holomorphic
function on $\C$.  Let $D'$ be the open unit disk with the origin
removed, and note that $q$ defines a map $\h \to D'$.  By \eqref{eqn:modfunc2} we have $f(z+1)=f(z)$, so there is a
function $F:D'\to \C$ such that $F(q(z)) = f(z)$.  This
function~$F$ is a complex-valued function on $D'$, but
it may or may not be well behaved at~$0$.


Suppose that $F$ is well behaved at $0$, in the sense that for some
$m\in\Z$ and all $q$ in a neighborhood of $0$ we have the equality
\begin{equation}\label{eqn:qexp}
 F(q) = \sum_{n=m}^{\infty} a_n q^n.
\end{equation}
If this is the case, we say that~$f$ is 
\defn{meromorphic at $\infty$}.  If, moreover,
$m\geq 0$, we say 
that~$f$ is \defn{holomorphic at  $\infty$}.   We also
call \eqref{eqn:qexp} the \defn{$q$-expansion} of~$f$ about~$\infty$.
 


\begin{definition}[Modular Function]
A \defn{modular function} of \defn{weight}~$k$ is a weakly modular function
of weight~$k$ that is meromorphic at $\infty$.
\end{definition}

\begin{definition}[Modular Form]
  A \defn{modular form} of \defn{weight}~$k$ (and \defn{level~$1$}) is a modular
  function of weight~$k$ that is holomorphic on~$\h$ and at~$\infty$.
\end{definition}

If $f$ is a modular form, then  there are numbers
$a_n$ such that for all $z\in \h$, 
\begin{equation}\label{eqn:qexp1}
f(z) = \sum_{n=0}^{\infty} a_n q^n.
\end{equation}
\begin{proposition}
The above series converges for all $z\in\h$.
\end{proposition}
\begin{proof}
The function $f(q)$ is holomorphic on $D$, so 
its Taylor series converges absolutely
in $D$.  
\end{proof}

Since $e^{2\pi i z}\to 0$ as  $z\to i \infty$, 
we set $f(\infty)=a_0$.


\begin{definition}[Cusp Form]
  A \defn{cusp form} of \defn{weight}~$k$ (and \defn{level~$1$}) is a
  modular form of weight~$k$ such that $f(\infty)=0$, i.e., $a_0=0$.
\end{definition}

Let \sym{\C[[q]]}
be the ring of all \defn{formal power series} 
in $q$. 
If $k=2$, then $dq = 2\pi i q dz$, 
so $dz = \frac{1}{2\pi i}\frac{dq}{q}$.
If $f(q)$ is a cusp form of weight~$2$, then 
$$
  2\pi i f(z) dz 
     = f(q) \frac{dq}{q}  = \frac{f(q)}{q} dq \in \C[[q]] dq.
$$
Thus the differential $2\pi i f(z) dz$ is holomorphic at $\infty$,
since $q$ is a local parameter at $\infty$.

\section{Modular Forms of Any Level}\label{sec:modformN}
In this section we define spaces of modular forms of arbitrary level.  

\begin{definition}[Congruence Subgroup]
A \defn{congruence subgroup} of
$\SL_2(\Z)$ is any subgroup of $\SL_2(\Z)$ that contains
$$\Gamma(N) = \Ker(\SL_2(\Z)\to \SL_2(\Z/N\Z))$$\index{$\Gamma(N)$}
for some positive integer $N$.  The smallest such~$N$ is the
\defn{level of $\Gamma$}.
\end{definition}
The most important congruence
subgroups in this book are\index{$\Gamma_1(N)$}
$$
  \Gamma_1(N) = \left\{ \mtwo{a}{b}{c}{d} \in \SL_2(\Z) : 
                    \mtwo{a}{b}{c}{d}\con \mtwo{1}{*}{0}{1}\pmod{N}
                \right\}
$$
and\index{$\Gamma_0(N)$}
$$
  \Gamma_0(N) = \left\{ \mtwo{a}{b}{c}{d} \in \SL_2(\Z) : 
                    \mtwo{a}{b}{c}{d}\con \mtwo{*}{*}{0}{*}\pmod{N}
                \right\},
$$
where $*$ means any element.  Both  groups have level~$N$ (see
\exref{ch:modform}{ex:conggamma1}).

Let $k$ be an integer.
Define the \defn{weight~$k$ right action}\index{right action of $\GL_2(\Q)$} 
of \sym{\GL_2(\Q)} on 
the set of all functions~$f:\h\to\C$  as follows.  
If $\gamma=\abcd{a}{b}{c}{d}\in\GL_2(\Q)$, let
\begin{equation}\label{eqn:bargamma}
  (f^{[\gamma]_k})(z) = \det(\gamma)^{k-1} (cz+d)^{-k} f(\gamma(z)).
\end{equation}
\begin{proposition}
Formula \eqref{eqn:bargamma} defines a right action
of $\GL_2(\Z)$ on the set of all functions $f:\h\to\C$; in
particular, 
$$
  f^{[\gamma_1\gamma_2]_k} = (f^{[\gamma_1]_k})^{[\gamma_2]_k}.
$$\index{$f^{[\gamma]_k}$}
\end{proposition}
\begin{proof}
See \exref{ch:modform}{ex:grpact}.
\end{proof}

\begin{definition}[Weakly Modular Function]
A \defn{weakly modular function} of weight $k$ for a
congruence subgroup $\Gamma$ is a meromorphic
function $f :\h \to \C$ such that 
$f^{[\gamma]_k} = f$ for all $\gamma \in \Gamma$.
\end{definition}

A central object in the theory of modular forms 
is the \defn{set of cusps}
$$
  \P^1(\Q) = \Q \union \{\infty\}.
$$
\index{cusps!action of $\SL_2(\Z)$ on}\index{$\SL_2(\Z)$}
An element $\gamma=\abcd{a}{b}{c}{d} \in \SL_2(\Z)$ 
acts on \sym{\P^1(\Q)} by 
$$
  \gamma(z) = \begin{cases} 
         \frac{az+b}{cz+d} & \text{if $z\neq \infty$},\\
         \frac{a}{c} & \text{if $z=\infty$.}
       \end{cases}
$$
Also, note that if the denominator $c$ or $cz+d$ is $0$ above, 
then $$\gamma(z)=\infty \in \P^1(\Q).$$

The set of \defn{cusps for a congruence subgroup $\Gamma$} is the set
\sym{C(\Gamma)} of $\Gamma$-orbits of $\P^1(\Q)$.  (We
will often identify elements of $C(\Gamma)$ with a representative
element from the orbit.)  For example, the lemma below
asserts that if $\Gamma=\SL_2(\Z)$, then there is
exactly one orbit, so $C(\SL_2(\Z)) = \{[\infty]\}$.
\begin{lemma}\label{lem:sl2ztrans}
  For any cusps $\alpha, \beta\in\P^1(\Q)$ there exists
  $\gamma\in\SL_2(\Z)$ such that $\gamma(\alpha) = \beta$.
\end{lemma}
\begin{proof}
This is \exref{ch:modform}{ex:sl2ztrans}.
\end{proof}

\begin{proposition}\label{prop:cuspfinite}
  For any congruence subgroup $\Gamma$, the set $C(\Gamma)$ of cusps
  is finite.
\end{proposition}
\begin{proof}
  This is \exref{ch:modform}{ex:cuspfinite}.
\end{proof}
See \cite[\S3.8]{diamond-shurman} and Algorithm~\ref{alg:cusplist} below for more discussion
of cusps and results relevant to their enumeration.

In order to define modular forms for general congruence subgroups, we
next explain what it means for a function to be holomorphic
on the \defn{extended upper half plane}
$$
  \h^* = \h\union \P^1(\Q).
$$\index{$\h^*$}

See \cite[\S1.3--1.5]{shimura:intro} for a detailed description of the
correct topology to consider on~$\h^*$.  In particular, a basis of
neighborhoods for $\alpha\in\Q$ is given by the sets $\{\alpha\}\union
D$, where~$D$ is an open disc in~$\h$ that is tangent to the real line
at~$\alpha$.

Recall from Section~\ref{sec:modform1} that a weakly modular
function~$f$ on $\SL_2(\Z)$ is holomorphic at~$\infty$ if its
$q$-expansion is of the form $\sum_{n=0}^{\infty} a_n q^n$.

In order to make sense of holomorphicity of a weakly modular
function~$f$ for an arbitrary congruence subgroup~$\Gamma$ at any
$\alpha\in\Q$, we first prove a lemma.
\begin{lemma}
If $f :\h\to\C$ is a weakly modular function of weight $k$
for a congruence subgroup $\Gamma$ and if $\delta \in \SL_2(\Z)$,
then $f^{[\delta]_k}$ is a weakly modular function for
$\delta^{-1}\Gamma \delta$.
\end{lemma}
\begin{proof}
If $s = \delta^{-1} \gamma\delta \in \delta^{-1}\Gamma \delta$,
then 
$$ 
   (f^{[\delta]_k})^{[s]_k} = f^{[\delta s]_k}
          = f^{[\delta \delta^{-1}\gamma\delta]_k}
          = f^{[\gamma\delta]_k} = f^{[\delta]_k}.
$$
\end{proof}

Fix a weakly modular function~$f$ of weight~$k$ for a congruence
subgroup $\Gamma$, and suppose $\alpha \in \Q$.  In
Section~\ref{sec:modform1} we constructed the $q$-expansion of $f$ by
using that $f(z) = f(z+1)$, which held since
$T=\mtwo{1}{1}{0}{1}\in\SL_2(\Z)$.  There are
congruence subgroups $\Gamma$ such that $T\not\in\Gamma$.  Moreover,
even if we are interested only in modular forms for $\Gamma_1(N)$,
where we have $T\in\Gamma_1(N)$ for all~$N$, we will still have to
consider $q$-expansions at infinity for modular forms on groups
$\delta^{-1} \Gamma_1(N)\delta$, and these need not contain~$T$.
Fortunately, $T^N = \abcd{1}{N}{0}{1} \in \Gamma(N)$, so a congruence
subgroup of level $N$ contains $T^N$.  Thus we have
$f(z+H) = f(H)$ for some positive integer $H$, e.g., $H=N$
always works, but there may be a smaller choice of $H$.
The minimal choice of $H>0$ such that
$\abcd{1}{H}{0}{1}\in \delta^{-1}\Gamma\delta$,
where $\delta(\infty)=\alpha$,
is called the \defn{width of the cusp}
$\alpha$ relative to the group $\Gamma$ (see Section~\ref{sec:cuspwidth}).
When~$f$ is meromorphic at infinity, we obtain 
a Fourier expansion
\begin{equation}\label{eqn:qser}
   f(z) = \sum_{n=m}^{\infty} a_n q^{n/H}
\end{equation}
in powers of the function $q^{1/H} = e^{2 \pi i z
  /H}$.  We say that $f$ is holomorphic at $\infty$ if in
\eqref{eqn:qser} we have $m\geq 0$.  

What about the other cusps $\alpha\in\P^1(\Q)$? By
Lemma~\ref{lem:sl2ztrans} there is a $\gamma\in\SL_2(\Z)$ such that
$\gamma(\infty)=\alpha$.  We declare~$f$ to be \defn{holomorphic at
  the cusp~$\alpha$} if the weakly modular function $f^{[\gamma]_k}$ is
holomorphic at $\infty$. 

\begin{definition}[Modular Form]
A \defn{modular form} of integer \defn{weight} $k$ for a congruence subgroup
$\Gamma$ is a weakly modular function $f:\h\to \C$ that
is holomorphic on $\h^*$.  We let $M_k(\Gamma)$\index{$M_k(\Gamma)$} 
denote the space
of weight~$k$ modular forms of weight~$k$ for $\Gamma$.
\end{definition}

\begin{proposition}
If a weakly modular function $f$ is holomorphic at a set of
representative elements for $C(\Gamma)$, then it is holomorphic at
every element of~$\P^1(\Q)$.
\end{proposition}
\begin{proof}
Let $c_1,\ldots, c_n \in \P^1(\Q)$ be representatives
for the set of cusps for $\Gamma$.
If $\alpha \in \P^1(\Q)$, then there is $\gamma\in\Gamma$
such that $\alpha = \gamma(c_i)$ for some~$i$.
By hypothesis~$f$ is holomorphic at $c_i$, so if $\delta\in\SL_2(\Z)$
is such that $\delta(\infty) = c_i$, then $f^{[\delta]_k}$
is holomorphic at~$\infty$.  Since~$f$ is a weakly modular
function for $\Gamma$,
\begin{equation}\label{eqn:prop:holoall}
  f^{[\delta]_k} = (f^{[\gamma]_k})^{[\delta]_k} 
   = f^{[\gamma\delta]_k}.
\end{equation}
But $\gamma(\delta(\infty)) = \gamma(c_i) = \alpha$, so 
\eqref{eqn:prop:holoall} implies that $f$ is holomorphic 
at~$\alpha$.
\end{proof}

\section{Remarks on Congruence Subgroups}\label{sec:congsgalg}

Recall that a congruence subgroup is a subgroup of $\SL_2(\Z)$
that contains $\Gamma(N)$ for some $N$.  Any congruence subgroup has
finite index in $\SL_2(\Z)$, since $\Gamma(N)$ does. What about the
converse: is every finite index subgroup of $\SL_2(\Z)$
a congruence subgroup? This is the \defn{congruence subgroup problem}.  One can ask about the congruence subgroup problem with $\SL_2(\Z)$
replaced by many similar groups.  If $p$ is a prime, then one can
prove that every finite index subgroup of $\SL_2(\Z[1/p])$ is a
congruence subgroup (i.e., contains the kernel of reduction modulo
some integer coprime to $p$), and for any $n>2$, all finite index
subgroups of $\SL_n(\Z)$ are congruence subgroups (see
\cite{humphreys:cong}).  However, there are numerous finite index
subgroups of $\SL_2(\Z)$ that are not congruence subgroups. The paper
\cite{hsu:congsub} contains an {\em algorithm} to decide if certain
finite index subgroups are congruence subgroups and gives an example
of a subgroup of index 12 that is not a congruence subgroup.

One can consider modular forms even for noncongruence subgroups.  See,
e.g., \cite{thompson:noncong} and the papers it references for work on
this topic.  We will not consider such modular forms further in this
book.
Note that modular symbols (which we define later in this book) {\em
  are} computable for noncongruence subgroups.

Finding coset representatives for $\Gamma_0(N)$, $\Gamma_1(N)$ and
$\Gamma(N)$ in $\SL_2(\Z)$ is straightforward and will be discussed
at length later in this book. To make the problem more explicit, note
that you can quotient out by $\Gamma(N)$ first. Then the question
amounts to finding coset representatives for a subgroup of
$\SL_2(\Z/N\Z)$ (and lifting), which is reasonably straightforward.

Given coset representatives for a finite index subgroup~$G$ of
$\SL_2(\Z)$, we can compute generators for~$G$ as follows.  Let~$R$ be
a set of coset representatives for $G$.  Let $\sigma, \tau \in
\SL_2(\Z)$ be the matrices denoted by $S$ and $T$ in \eqref{eqn:ST}.
Define maps $s, t: R \to G$ as follows. If $r \in R$, then there
exists a unique $\alpha_r \in R$ such that $Gr\sigma = G\alpha_r$.
Let $s(r) = r\sigma \alpha_r^{-1}$. Likewise, there is a unique
$\beta_r$ such that $Gr\tau = G\beta_r$ and we let $t(r) = r\tau
\beta_r^{-1}$. Note that $s(r)$ and $t(r)$ are in~$G$ for all~$r$.
Then $G$ is generated by $s(R)\cup t(R)$.

\begin{proposition}\label{prop:compgen}
The above procedure computes generators for~$G$.
\end{proposition}
\begin{proof}
  Without loss of generality, assume that $I=\abcd{1}{0}{0}{1}$
  represents the coset of~$G$.  Let~$g$ be an element of $G$.
  Since~$\sigma$ and~$\tau$ generate $\SL_2(\Z)$, it is possible to
  write~$g$ as a product of powers of~$\sigma$ and~$\tau$.  There is a
  procedure, which we explain below with an example in order to avoid
  cumbersome notation, which writes $g$ as a product of elements of
  $s(R) \union t(R)$ times a right coset representative~$r\in R$.  For
  example, if
$$g = \sigma \tau^2 \sigma \tau,$$
then $g=I\sigma \tau^2\sigma \tau = s(I)y\tau^2\sigma\tau$ 
for some $y \in R$.  Continuing, 
 $$s(I)y\tau^2\sigma\tau = s(I)(y\tau)\tau\sigma\tau 
= s(I)(t(y)z)\tau\sigma\tau$$ 
for some $z \in R$.  Again, 
$$s(I)(t(y)z)\tau\sigma\tau = s(I)t(y)(z\tau)\sigma\tau 
       = \cdots.
$$
The procedure illustrated above (with an example) makes sense for
arbitrary $g$ and, after carrying it out, writes~$g$ as a product of
elements of $s(R) \union t(R)$ times a right coset
representative~$r\in R$.  But $g\in G$ and~$I$ is the right coset
representative for~$G$, so this right coset representative must
be~$I$.
\end{proof}
\begin{remark}
We could also apply the proof of Proposition~\ref{prop:compgen} to
write any element of $G$ in terms of the given generators.
Moreover, we could use it to write any element $\gamma \in \SL_2(\Z)$
in the form $gr$, where $g\in G$ and $r\in R$, so we can decide
whether or not $\gamma \in G$.
\end{remark}

\subsection{Computing Widths of Cusps}\label{sec:cuspwidth}
Let $\Gamma$ be a congruence subgroup of level~$N$.  Suppose
$\alpha\in C(\Gamma)$ is a cusp, and choose $\gamma\in\SL_2(\Z)$ such
that $\gamma(\infty)=\alpha$.  Recall that the minimal $h$ such that
$\abcd{1}{h}{0}{1} \in \gamma^{-1}\Gamma\gamma$ is called the
\defn{width of the cusp} $\alpha$ for the group $\Gamma$.  In this
section we discuss how to compute~$h$.

\begin{algorithm}{Width of Cusp}\label{alg:cuspwidth}
  Given a congruence subgroup $\Gamma$ of level $N$ and a cusp
  $\alpha$ for $\Gamma$, this algorithm computes the width $h$ of
  $\alpha$.  We assume that $\Gamma$ is given by congruence conditions,
e.g., $\Gamma=\Gamma_0(N)$ or $\Gamma_1(N)$.
\begin{steps}
\item{}[Find $\gamma$] Use the extended Euclidean algorithm
to find $\gamma\in\SL_2(\Z)$ such that $\gamma(\infty)=\alpha$,
as follows.
If $\alpha=\infty$, set $\gamma\set 1$; otherwise, write 
$\alpha=a/b$, find $c,d$ such that $ad-bc=1$, and set
$\gamma \set \abcd{a}{b}{c}{d}$.
\item{}[Compute Conjugate Matrix] Compute the following element of
$\Mat_2(\Z[x])$:
$$
  \delta(x) \set \gamma \mtwo{1}{x}{0}{1} \gamma^{-1}.
$$
Note that the entries of $\delta(x)$ are constant or
linear in $x$.
\item{}[Solve] The congruence conditions that define $\Gamma$ give
  rise to four linear congruence conditions on $x$.  Use techniques
  from elementary number theory (or enumeration) to find the smallest
  simultaneous positive solution $h$ to these four equations.
\end{steps}
\end{algorithm}

\begin{example}
\begin{enumerate}
\item
Suppose $\alpha=0$ and $\Gamma=\Gamma_0(N)$ or $\Gamma_1(N)$.
Then $\gamma=\abcd{0}{-1}{1}{\hfill 0}$ has the property that $\gamma(\infty)=\alpha$.
Next, the congruence condition is 
$$\delta(x) = \gamma \mtwo{1}{x}{0}{1} \gamma^{-1} = \mtwo{1}{0}{-x}{1} 
  \con \mtwo{1}{*}{0}{1}\pmod{N}.
$$
Thus the smallest positive solution is $h=N$, so the width of $0$
is~$N$.

\item Suppose $N=pq$ where $p,q$ are distinct primes, and let $\alpha=1/p$.
Then $\gamma=\abcd{1}{0}{p}{1}$ sends $\infty$ to $\alpha$.
The congruence condition for $\Gamma_0(pq)$ is 
$$
\qquad \qquad\delta(x) = \gamma \mtwo{1}{x}{0}{1} \gamma^{-1} = \mtwo{1-px}{x}{-p^2 x}{px+1}
  \con \mtwo{*}{*}{0}{*}\pmod{pq}.
$$
Since $p^2x \con 0\pmod{pq}$, we see that $x=q$ is the smallest solution.
Thus $1/p$ has width $q$, and symmetrically $1/q$ has width~$p$.
\end{enumerate}
\end{example}
\begin{remark}
For $\Gamma_0(N)$, once we enforce that the bottom 
left entry is $0\pmod{N}$ and use that the determinant is 1, the 
coprimality from the other two congruences is automatic.
So there is one congruence to solve in the $\Gamma_0(N)$ case.
There are two congruences in the $\Gamma_1(N)$ case. 
\end{remark}

\section{Applications of Modular Forms}\label{sec:apps}
The above definition of modular forms might leave the impression that
modular forms occupy an obscure corner of complex analysis.  This is
{\em not} the case! Modular forms are highly geometric,  % typo found!
arithmetic, and topological objects that are of extreme interest all
over mathematics:

\begin{enumerate}
\item {\bf Fermat's last theorem:} Wiles' proof \cite{wiles:fermat}
  of Fermat's last theorem\index{Fermat's last theorem}
 uses modular forms extensively.  The work
  of Wiles et al. on modularity also massively extends computational
  methods for elliptic curves over~$\Q$, because many elliptic curve algorithms,
  e.g., for computing $L$-functions, modular degrees, Heegner points,
  etc., require that the elliptic curve be modular.


\item {\bf Diophantine equations:} Wiles'
  proof of Fermat's last theorem has made available a wide array of
  new techniques for solving certain diophantine equations.\index{Diophantine equations}  
Such work
  relies crucially on having access to tables or software for
  computing modular forms. See, e.g., \cite{darmon:eps,
    merel:diophantine, chen:dio, siksekcremona}.  (Wiles did not need
  a computer, because the relevant spaces of modular forms that arise
  in his proof have dimension $0$!)  Also, according to Siksek
  (personal communication) the paper \cite{siksek:dio} would ``have
  been entirely impossible to write without [the algorithms described
  in this book].''


\item {\bf Congruent number problem:} This ancient open
  problem is to determine which integers are the area of a right
  triangle with rational side lengths.  There is a potential solution
  that uses modular forms (of weight $3/2$) extensively (the solution
  is conditional on truth of the Birch and Swinnerton-Dyer
  conjecture, which is not yet known). See \cite{koblitz:cong}.
\index{congruent number problem}

\item {\bf Topology:} Topological modular forms are a major area
  of current research.

\item {\bf Construction of Ramanujan graphs:} Modular forms can be
  used to construct almost optimal expander graphs, which play
  a role in communications network theory.\index{Ramanujan graphs}

\item {\bf Cryptography and Coding Theory:} Point counting on elliptic
  curves over finite fields is crucial to the construction of elliptic
  curve cryptosystems, and modular forms are relevant to efficient
  algorithms for point counting (see \cite{elkies:ffield}).  Algebraic
  curves that are associated to modular forms are useful in constructing and
  studying certain error-correcting codes (see \cite{ebeling}).

\item {\bf The Birch and Swinnerton-Dyer conjecture:} This central
  open problem in arithmetic geometry relates arithmetic properties of
  elliptic curves (and abelian varieties) to special values of
  $L$-functions.  Most deep results toward this conjecture use modular
  forms extensively (e.g., work of Kolyvagin, Gross-Zagier, and Kato).
  Also, modular forms are used to compute and prove results about
  special values of these $L$-functions. See \cite{wiles:cmi}.
\index{Birch and Swinnerton-Dyer conjecture}

\item {\bf Serre's Conjecture on modularity of Galois representation:}
  Let $G_\Q=\Gal(\Qbar/\Q)$ be the Galois group of an algebraic
  closure of~$\Q$. Serre conjectured and many people have
  (nearly!)  proved that every continuous homomorphism $\rho:G_\Q \to
  \GL_2(\F_q)$, where $\F_q$ is a finite field and
  $\det(\rho(\text{complex conjugation}))=-1$, ``arises'' from a
  modular form.  More precisely, for almost all primes $p$ the
  coefficients $a_p$ of a modular (eigen-)form $\sum a_n q^n$ are
  congruent to the traces of elements $\rho(\Frob_p)$, where $\Frob_p$
  are certain special elements of $G_\Q$ called Frobenius elements.
  See \cite{ribet-stein:serre} and \cite[Ch.~9]{diamond-shurman}.
\index{Serre's conjecture}

\item {\bf Generating functions for partitions:} The generating
  functions for various kinds of partitions of an integer can often be
  related to modular forms. Deep theorems about modular forms
  then translate into results about partitions.  See work of Ramanujan,
  Gordon, Andrews, and Ahlgren and Ono (e.g., \cite{ahlgren-ono}).
\index{partitions}

\item {\bf Lattices:} If $L\subset \R^n$ is an even unimodular lattice
  (the basis matrix has determinant $\pm 1$ and $\lambda \cdot \lambda
  \in 2\Z$ for all $\lambda \in L$), then the theta series
  $$\theta_L(q) = \sum_{\lambda \in L} q^{\lambda \cdot \lambda}$$ is
  a modular form of weight $n/2$. The coefficient of $q^m$
  is the number of lattice vectors with squared length $m$.  Theorems
  and computational methods for modular forms translate into theorems
  and computational methods for lattices.  For example, the 290
  theorem of M.~Bharghava and J.~Hanke is a theorem about lattices,
which asserts that an integer-valued quadratic form represents all
  positive integers if and only if it represents the integers up to
  $290$; it is proved by doing many calculations with
  modular forms (both theoretical and with a computer).
\index{lattices}
\end{enumerate}


\begin{exercises}
\item\label{ex:upperhalfpres}
Suppose $\gamma = \abcd{a}{b}{c}{d}\in\GL_2(\R)$ has
positive determinant.  Prove that if $z\in\C$ is a complex
number with positive imaginary part, then the imaginary
part of $\gamma(z) = (az+b)/(cz+d)$ is also positive.

\item\label{ex:meromorphic}
Prove that every rational function (quotient of two polynomials)
is a meromorphic function on $\C$.

\item \label{ex:wmfprod}
Suppose $f$ and $g$ are weakly modular functions for a congruence
subgroup $\Gamma$ with $f\neq 0$.
\begin{enumerate}
\item Prove that the product $fg$ is a weakly modular function for $\Gamma$.
\item Prove that $1/f$ is a weakly modular function  for $\Gamma$.
\item If $f$ and $g$ are modular functions, show that $fg$
is a modular function for $\Gamma$.
\item If $f$ and $g$ are modular forms, show that $fg$
is a modular form  for $\Gamma$.
\end{enumerate}

\item \label{ex:nomodformodd}
Suppose $f$ is a weakly modular function of odd weight $k$ and
level $\Gamma_0(N)$ for some $N$.
Show that $f=0$.

\item \label{ex:allsame} 
Prove that $\SL_2(\Z)=\Gamma_0(1)=\Gamma_1(1)=\Gamma(1)$.

\item \label{ex:conggamma1}
\begin{enumerate}
\item Prove that $\Gamma_1(N)$ is a group.
\item Prove that $\Gamma_1(N)$ has finite index in $\SL_2(\Z)$
(Hint: It contains the kernel of the homomorphism
$\SL_2(\Z) \to \SL_2(\Z/N\Z)$.)
\item Prove that $\Gamma_0(N)$ has finite index in $\SL_2(\Z)$.
\item Prove that $\Gamma_0(N)$ and $\Gamma_1(N)$ have
level $N$.
\end{enumerate}

\item \label{ex:grpact}
Let $k$ be an integer, and for any function
$f:\h^*\to \C$ and $\gamma=\abcd{a}{b}{c}{d}\in\GL_2(\Q)$,
set $f^{[\gamma]_k}(z) = \det(\gamma)^{k-1} \cdot (cz+d)^{-k} \cdot f(\gamma(z))$. 
Prove that if $\gamma_1, \gamma_2 \in \GL_2(\Z)$,
then for all $z\in \h^*$ we have 
$$
 f^{[\gamma_1 \gamma_2]_k}(z)  = ((f^{[\gamma_1]_k})^{[\gamma_2]_k})(z).
$$


\item \label{ex:sl2ztrans}
Prove that for any $\alpha, \beta\in\P^1(\Q)$, there exists
  $\gamma\in\SL_2(\Z)$ such that $\gamma(\alpha)~=~\beta$.

\item \label{ex:cuspfinite}
Prove Proposition~\ref{prop:cuspfinite}, which
asserts that the set of cusps $C(\Gamma)$, for any
congruence subgroup $\Gamma$, is finite. 

\item \label{ex:cuspwidth}
Use Algorithm~\ref{alg:cuspwidth} to give an 
example of a group $\Gamma$ and cusp $\alpha$
with width $2$. 

\end{exercises}


\chapter{Modular Forms of Level $1$}\label{ch:level1}

In this chapter we study in detail the structure of level $1$ modular
forms, i.e., modular forms on $\SL_2(\Z)=\Gamma_0(1)=\Gamma_1(1)$.  
We assume some complex analysis (e.g., the residue theorem),
linear algebra, and that the reader has
read Chapter~\ref{ch:modform}.

\section{Examples of Modular Forms of Level $1$}\label{sec:level_one_eisen}

In this section we will finally see some examples of modular forms of
level $1$!  We first introduce the Eisenstein series and
then define $\Delta$, which is a cusp form of weight $12$.  In
Section~\ref{sec:struct1} we  prove the structure theorem,
which says that all modular forms of level $1$ are polynomials
in Eisenstein series. 

For an even integer $k\geq 4$, the \defn{nonnormalized weight~$k$
  Eisenstein series}\index{Eisenstein series}
is the function on the extended
upper half plane $\h^*=\h\cup\P^1(\Q)$ given
by\index{$G_k(z)$}
\begin{equation}\label{eqn:gk}
  G_k(z) = \sum_{m,n\in\Z}^{*} \frac{1}{(mz+n)^k}.
\end{equation}
The star on top of the sum symbol means that
for each $z$ the sum is over all $m,n\in\Z$ such that
$mz+n\neq 0$.

\begin{proposition}
The function $G_k(z)$ is a modular form of weight~$k$, i.e.,
 $ G_k \in M_k(\SL_2(\Z))$.
\end{proposition}
\begin{proof}
See \cite[\S{} VII.2.3]{serre:arithmetic} for a proof
that $G_k(z)$ defines a holomorphic function on $\h^*$.
To see that $G_k$ is modular, observe that
$$
  G_k(z+1) = \sum^* \frac{1}{(m(z+1)+n)^k} = \sum^* \frac{1}{(mz+(n+m))^k} = 
\sum^* \frac{1}{(mz+n)^k},
$$
where for the last equality we use that the map
$(m,n+m) \mapsto (m, n)$ on $\Z\times \Z$ is invertible.
Also,
\begin{align*}
G_k(-1/z) &= \sum^* \frac{1}{(-m/z+n)^k} \\
   &= \sum^* \frac{z^k}{(-m+nz)^k}\\
   &= z^k \sum^* \frac{1}{(mz+n)^k} = z^k G_k(z),
\end{align*}
where we use that $(n,-m) \mapsto (m,n)$ is invertible.
\end{proof}

\begin{proposition}
$G_k(\infty) = 2\zeta(k)$,
where $\zeta$ is the Riemann zeta function.
\end{proposition}
\begin{proof}
As $z\to \infty$ (along the imaginary axis) in \eqref{eqn:gk},
the terms that involve~$z$ with $m\neq 0$ 
go to~$0$.  
Thus 
$$
 G_k(\infty) = \sum^*_{n\in\Z} \frac{1}{n^k}.
$$
This sum is twice $\zeta(k) = \sum_{n\geq 1} \frac{1}{n^k}$, 
as claimed.
\end{proof}



\subsection{The Cusp Form $\Delta$}\index{cusp forms!$\Delta$}
Suppose $E=\C/\Lambda$ is an elliptic curve over $\C$, viewed as a
quotient of $\C$ by a lattice $\Lambda=\Z\omega_1+\Z\omega_2$, with
$\omega_1/\omega_2\in\h$ (see \cite[\S1.4]{diamond-shurman}).  
The \defn{Weierstrass $\wp$-function} of
the lattice $\Lambda$ is
$$
  \wp = \wp_{\Lambda}(u) = \frac{1}{u^2} + \sum_{k=4,6,8,\ldots} (k-1)
      G_k(\omega_1/\omega_2) u^{k-2},
$$
where the sum is over even integers $k\geq 4$.
It satisfies the differential equation
$$
  (\wp')^2 = 4\wp^3 - 60 G_4(\omega_1/\omega_2) \wp - 140 G_6(\omega_1/\omega_2).
$$
If we set $x=\wp$ and $y=\wp'$, the above is an (affine) equation of
the form $y^2 = ax^3 + bx + c$ for an elliptic curve that is complex
analytically isomorphic to $\C/\Lambda$ (see \cite[pg.~277]{ahlfors}
for why the cubic has distinct roots).

The discriminant of the cubic $$4x^3 - 60G_4(\omega_1/\omega_2) x - 
140 G_6(\omega_1/\omega_2)$$ is $16D(\omega_1/\omega_2)$, where
$$
 D(z) = (60 G_4(z))^3 - 27(140 G_6(z))^2.
$$
Since $D(z)$ is the difference of two modular forms of
weight $12$ it has weight $12$. 
 Moreover, 
\begin{align*}
 D(\infty) &=
\left(60 G_4(\infty)\right)^3 - 
27\left(140 G_6(\infty)\right)^2\\
 &= \left(\frac{60}{3^2\cdot 5} \pi^4\right)^3 - 
27\left(\frac{140\cdot 2}{3^3\cdot 5\cdot 7} \pi^6\right)^2\\
&=
0,
\end{align*}
so $D$ is a cusp form of weight $12$.
Let\indexp{$\Delta$}
$$
  \Delta = \frac{D}{(2\pi)^{12}}.
$$
  
\begin{lemma}\label{lem:delnz}
If $z\in\h$, then $\Delta(z)\neq 0$.
\end{lemma}
\begin{proof}
Let $\omega_1=z$ and $\omega_2=1$.
Since~$E=\C/(\Z\omega_1 + \Z\omega_2)$ 
is an elliptic curve,
it has nonzero discriminant
$\Delta(z) = \Delta(\omega_1/\omega_2)\neq 0$.
\end{proof}

\begin{proposition}
We have $\Delta = q\cdot \prod_{n=1}^{\infty}(1-q^n)^{24}$.
\end{proposition}
\begin{proof}
See \cite[Thm.~6, pg.~95]{serre:arithmetic}.
\end{proof}

\begin{remark}
\sage{} computes the $q$-expansion of $\Delta$ efficiently
to high precision using the command {\tt delta\_qexp}:
\sageindex{$q$-expansion of $\Delta$}
\begin{verbatim}
sage: delta_qexp(6)
q - 24*q^2 + 252*q^3 - 1472*q^4 + 4830*q^5 + O(q^6)
\end{verbatim}
\end{remark}

\subsection{Fourier Expansions of Eisenstein Series}\label{sec:fourexpeis}%
\index{Eisenstein series!Fourier expansion}
Recall from \eqref{eqn:qexp1} that elements $f$ of $M_k(\SL_2(\Z))$ can be
expressed as formal power series in terms of $q(z)=e^{2\pi i z}$ and
that this expansion is called the Fourier expansion of $f$.
The following proposition gives the Fourier expansion of the
Eisenstein series $G_k(z)$.

\begin{definition}[Sigma]
For any integer $t\geq 0$ and any positive integer $n$, 
the \defn{sigma function} 
$$
  \sigma_t(n) = \sum_{1\leq d \mid n} d^t
$$
is the sum of the $t$th powers of the positive divisors of~$n$.
Also, let $d(n) = \sigma_0(n)$, which is the number
of divisors of $n$, and let $\sigma(n) = \sigma_1(n)$.
For example, if $p$ is prime, then $\sigma_t(p) = 1 + p^t$.
\end{definition}

\begin{proposition}\label{prop:qexpGk}
For every even integer $k\geq 4$, we have
$$
  G_k(z) = 2\zeta(k) + 2\cdot \frac{(2\pi i)^{k}}{(k-1)!}\cdot
     \sum_{n=1}^{\infty} \sigma_{k-1}(n) q^n.
$$
\end{proposition}
\begin{proof}
  See \cite[Section~VII.4]{serre:arithmetic}, which uses clever
  manipulations of series, starting with the identity
$$
  \pi \cot(\pi z) = \frac{1}{z} + \sum_{m=1}^{\infty} 
     \left( \frac{1}{z+m} + \frac{1}{z-m}\right).
$$
\end{proof}

From a computational point of view, the $q$-expansion of
Proposition~\ref{prop:qexpGk} is unsatisfactory because it involves
transcendental numbers.  To understand these numbers, we introduce the
\defn{Bernoulli numbers} $B_n$ for $n\geq 0$ {\em defined} by the
following equality of formal power series:
\begin{equation}\label{eqn:def_bernoulli}
  \frac{x}{e^x - 1} = \sum_{n=0}^{\infty} B_n \frac{x^n}{n!}.
\end{equation}
Expanding the power series, we have 
$$
\frac{x}{e^x - 1} =
1 - \frac{x}{2} + \frac{x^2}{12} - \frac{x^4}{720} + \frac{x^6}{30240}
- \frac{x^8}{1209600} + \cdots. $$
As this expansion suggests, the
Bernoulli numbers $B_n$ with $n>1$ odd are~$0$ (see
\exref{ch:modform}{ex:odd_bernoulli}).
Expanding the series further, we obtain the following table:
\vspace{0.5ex}

\noindent{}$\ds B_{0}=1,\quad B_{1}=-\frac{1}{2},\quad B_{2}=\frac{1}{6},\quad
B_{4}=-\frac{1}{30},\quad B_{6}=\frac{1}{42},\quad
B_{8}=-\frac{1}{30},\quad $\vspace{1.5ex}

\noindent{}$\ds
B_{10}=\frac{5}{66},\quad
B_{12}=-\frac{691}{2730},\quad
B_{14}=\frac{7}{6},\quad
B_{16}=-\frac{3617}{510},\quad
B_{18}=\frac{43867}{798},\quad
$\vspace{1.5ex}

\noindent{}
$\ds{}\!\!B_{20}=-\frac{174611}{330},\quad
B_{22}=\frac{854513}{138},\quad
B_{24}=-\frac{236364091}{2730},\quad
B_{26}=\frac{8553103}{6}.
$\vspace{1ex}\label{tbl:bern}

See Section~\ref{sec:bercomp} for a discussion of fast
(analytic) methods for computing Bernoulli numbers. 

%begin_sage
\sageindex{Bernoulli numbers}
We compute some Bernoulli numbers in \sage:
\begin{verbatim}
sage: bernoulli(12)
-691/2730
sage: bernoulli(50)
495057205241079648212477525/66
sage: len(str(bernoulli(10000)))
27706
\end{verbatim}
%end_sage


A key fact is that Bernoulli numbers are
rational numbers and they are connected to values of~$\zeta$ at
positive even integers.
\begin{proposition}\label{prop:zeta_even}
If $k\geq 2$ is an even integer, then
$$
  \zeta(k) = -\frac{(2\pi i)^{k}}{2\cdot k!}\cdot B_k.
$$
\end{proposition}
\begin{proof}
This is proved 
by manipulating a series expansion of $z \cot(z)$
(see \cite[Section~VII.4]{serre:arithmetic}).
\end{proof}

\begin{definition}[Normalized Eisenstein Series]
The \defn{normalized Eisenstein series} of even weight~$k\geq 4$ is
$$
  E_k = \frac{(k-1)!}{2\cdot (2\pi i)^k}\cdot G_k.
$$
\end{definition} 
Combining Propositions~\ref{prop:qexpGk} and
\ref{prop:zeta_even}, we see that
\begin{equation}\label{eqn:ekexp}
 E_k = -\frac{B_k}{2k} + q + \sum_{n=2}^{\infty} \sigma_{k-1}(n) q^n.
\end{equation}


\begin{warning} Our series $E_k$ is normalized so that the coefficient
  of~$q$ is~$1$, but often in the literature $E_k$ is normalized so
  that the constant coefficient is~$1$.  We use the normalization with
  the coefficient of~$q$ equal to~$1$, because then the eigenvalue of
  the $n$th Hecke operator (see Section~\ref{sec:hecke_one}) is the
  coefficient of~$q^n$.  Our normalization is also 
  convenient when considering congruences between cusp forms and
  Eisenstein series.
\end{warning}

\section{Structure Theorem for Level $1$ Modular  Forms}\label{sec:struct1}
In this section we describe a structure theorem for modular
forms of level~$1$. 
If $f$ is a nonzero meromorphic function on $\h$ and $w\in\h$, let
$\ord_w(f)$ be the largest integer~$n$ such that $f(z)/(w-z)^n$ is
holomorphic at $w$.  If $f = \sum_{n=m}^{\infty} a_n q^n$ with
$a_m\neq 0$, we set $\ord_{\infty}(f)=m$.  We will use the following
theorem to give a presentation for the vector space of modular forms
of weight~$k$; this presentation yields an algorithm
to compute this space.

Let $M_k=M_k(\SL_2(\Z))$ \index{$M_k$}denote the complex vector space of modular
forms of weight~$k$ for $\SL_2(\Z)$.
The \defn{standard fundamental domain} \sym{\cF} for $\SL_2(\Z)$
is the set of $z\in \h$  with $|z|\geq 1$ and
$|\Re(z)|\leq 1/2$.  
Let $\rho = e^{2\pi i /3}$.

\begin{theorem}[Valence Formula]\label{thm:valence}\index{valence formula}
Let $k$ be any integer and
suppose $f \in M_k(\SL_2(\Z))$ is nonzero.  Then
$$
  \ord_{\infty}(f) + \frac{1}{2}\ord_i(f) + \frac{1}{3} \ord_{\rho}(f)
   + \sum_{w\in \cF}^* \ord_w(f) = \frac{k}{12},
$$
where $\ds \sum^*_{w \in \cF}$ is the sum over elements of $\cF$
other than~$i$ and~$\rho$.
\end{theorem}
\begin{proof}
The proof in \cite[\S{}VII.3]{serre:arithmetic} uses
the residue theorem.
\end{proof}

Let~$S_k=S_k(\SL_2(\Z))$\index{$S_k$}
denote the subspace of weight~$k$ cusp forms for $\SL_2(\Z)$.  We have
an exact sequence
$$
  0 \to S_k \to M_k \xrightarrow{\iota_\infty} \C 
$$
that sends $f\in M_k$ to $f(\infty)$.  When $k\geq 4$ is even, the
space $M_k$ contains the Eisenstein series $G_k$, and
$G_k(\infty)=2\zeta(k)\neq 0$, so the map $M_k\to \C$ is surjective.
This proves the following lemma.
\begin{lemma}
If $k\geq 4$ is even, then $
  M_k = S_k \oplus \C G_k
$
and the following sequence is exact:
$$
  0 \to S_k \to M_k \xrightarrow{\iota_\infty} \C  \to 0.
$$
\end{lemma}

\begin{proposition}\label{prop:mk_vanish}
For $k<0$ and $k=2$, we have $M_k=0$.  
\end{proposition}
\begin{proof}
Suppose $f\in M_k$ is nonzero yet $k=2$ or $k<0$.  By 
Theorem~\ref{thm:valence}, 
$$
  \ord_{\infty}(f) + \frac{1}{2}\ord_i(f) + \frac{1}{3} \ord_{\rho}(f)
   + \sum_{w\in D}^* \ord_w(f) = \frac{k}{12}\leq \frac{1}{6}.
$$
This is not possible because each quantity on the left is nonnegative so
whatever the sum is, it is too big (or $0$, in which case $k=0$).
\end{proof}

\begin{theorem}\label{thm:delta_iso}
Multiplication by $\Delta$ defines an isomorphism $M_{k-12}\to S_k$.
\end{theorem}
\begin{proof}
By Lemma~\ref{lem:delnz}, $\Delta$ is not identically~$0$, so
because $\Delta$ is holomorphic, multiplication 
by $\Delta$ defines an injective map $M_{k-12}\hra S_k$.
To see that this map is surjective, we show that if $f\in S_k$, then 
$f/\Delta \in M_{k-12}$.
Since $\Delta$ has weight~$12$ and $\ord_\infty(\Delta)\geq 1$,
Theorem~\ref{thm:valence} implies that~$\Delta$ has a simple
zero at $\infty$ and does not vanish on~$\h$.
Thus if $f\in S_k$ and if we let $g=f/\Delta$, then $g$ is holomorphic
and satisfies the appropriate transformation formula, so~$g \in M_{k-12}$.
\end{proof}

\begin{corollary}
  For $k=0,4,6,8,10,14$, the space $M_k$ has dimension~$1$, with
  basis $1$, $G_4$, $G_6$, $G_8$, $G_{10}$, and $G_{14}$, respectively, and
  $S_k=0$.
\end{corollary}
\begin{proof}
  Combining Proposition~\ref{prop:mk_vanish} with
  Theorem~\ref{thm:delta_iso}, we see that the spaces $M_k$ for $k\leq
  10$ cannot have dimension greater than~$1$, since otherwise $M_{k'}\neq
  0$ for some $k'<0$.  Also $M_{14}$ has dimension at most $1$, since
  $M_2$ has dimension $0$.  Each of the indicated spaces of weight
  $\geq 4$ contains the indicated Eisenstein series and so has
  dimension~$1$, as claimed.
\end{proof}

\begin{corollary}\label{cor:dim1}
$\dim M_k = \begin{cases}
   0 & \text{if }k\text{ is odd or negative}, \\ 
   \lfloor{} k/12 \rfloor{} & \text{if } k\con 2\pmod{12},\\
   \lfloor{} k/12 \rfloor{}+1 & \text{if } k\not\con 2\pmod{12}.
\end{cases}
$\\
Here $\lfloor{}x \rfloor{}$ is the biggest integer $\leq x$.
\end{corollary}
\begin{proof}
As we have already seen above, the formula is true when $k\leq 12$.
By Theorem~\ref{thm:delta_iso}, the dimension increases by $1$
when~$k$ is replaced by $k+12$.
\end{proof}

\begin{theorem}\label{thm:mk_one_basis}
The space $M_k$ has as basis the
modular forms $G_4^a G_6^b$, where $a,b$
run over all pairs of nonnegative integers such that 
$4a+6b=k$.
\end{theorem}
\begin{proof}
  Fix an even integer $k$.  We first prove by induction that the
  modular forms $G_4^a G_6^b$ generate $M_k$; the cases $k\leq 10$ and
  $k=14$ follow from the above arguments (e.g., when $k=0$, we have $a=b=0$ and
  basis~$1$).  Choose some pair of nonnegative integers $a,b$ 
  such that $4a+6b=k$.   The form
  $g=G_4^a G_6^b$ is not a cusp form, since it is nonzero at $\infty$.
  Now suppose $f\in M_k$ is arbitrary.  Since $g(\infty)\neq 0$, there
  exists $\alpha\in\C$ such that $f - \alpha g \in S_k$.  Then by
  Theorem~\ref{thm:delta_iso}, there is $h \in M_{k-12}$ such that $f
  - \alpha g = \Delta \cdot h$.  By induction, $h$ is a polynomial in
  $G_4$ and $G_6$ of the required type, and so is $\Delta$, so $f$ is
  as well.  Thus $$\{G_4^aG_6^b\ |\ a\geq 0,\ b\geq 0,\ 4a+6b=k\}$$
  spans $M_k$.

Suppose there is a nontrivial linear relation between the $G_4^a
G_6^b$ for a given~$k$.  By multiplying the linear relation by a
suitable power of $G_4$ and $G_6$, we may assume that we have
such a nontrivial relation with $k\con 0\pmod{12}$.  Now divide the
linear relation by the weight $k$ form $G_6^{k/6}$ to see that
$G_4^3/G_6^2$ satisfies a polynomial with coefficients in~$\C$
(see \exref{ch:level1}{ex:g4g6quo}).  Hence
$G_4^3/G_6^2$ is a root of a polynomial, hence a constant, which is a
contradiction since the $q$-expansion of $G_4^3/G_6^2$ is not
constant.  
\end{proof}

\begin{algorithm}{Basis for $M_k$}\label{alg:basis}
Given integers~$n$ and~$k$, this algorithm computes
a basis of $q$-expansions for the complex vector
space $M_k$ mod $q^n$.   The $q$-expansions output
by this algorithm have coefficients in $\Q$.
\begin{steps}
\item{}[Simple Case] If $k=0$, output the basis with
just $1$ in it and terminate; otherwise if $k<4$ or~$k$ 
is odd, output the empty basis and terminate.
\item{}[Power Series] Compute $E_4$ and $E_6$ mod $q^n$
using the formula from \eqref{eqn:ekexp} and
Section~\ref{sec:bercomp}.
\item{}[Initialize] Set $b\set 0$.   
\item{}[Enumerate Basis] For each integer~$b$ between~$0$ and
$\lfloor k/6\rfloor$, compute $a=(k-6b)/4$.
If~$a$ is an integer, compute and output the basis
element $E_4^a E_6^b \mod{q^n}$.  When computing
$E_4^a$, find
$E_4^m\pmod{q^n}$ for each $m\leq a$, and save these
intermediate powers, so they can be reused later, and
likewise for powers of $E_6$.
\end{steps}
\end{algorithm}
\begin{proof}
This is simply a translation of Theorem~\ref{thm:mk_one_basis}
into an algorithm, since $E_k$ is a nonzero scalar multiple
of $G_k$.  That the $q$-expansions have coefficients
in $\Q$ follows from \eqref{eqn:ekexp}.
\end{proof}

\begin{example}
We compute a basis for $M_{24}$, which is the space with
smallest weight whose dimension is greater than $1$.
It has as basis $E_4^6$, $E_4^3 E_6^2$, and $E_6^4$, whose explicit
expansions are
\begin{align*}
  E_4^6 &= \frac{1}{191102976000000} +
   \frac{1}{132710400000}q + \frac{203}{44236800000}q^2 + \cdots,\\
E_4^3 E_6^2 &= 
    \frac{1}{3511517184000} - \frac{1}{12192768000}q - \frac{377}{4064256000}q^2
  + \cdots, \\
  E_6^4 &= \frac{1}{64524128256} - \frac{1}{32006016}q + 
    \frac{241}{10668672}q^2 +\cdots. 
\end{align*}
\vspace{1ex}

%begin_sage
\sageindex{basis for $M_{24}$}
We compute this basis in \sage{} as follows:
\begin{verbatim}
sage: E4 = eisenstein_series_qexp(4, 3)
sage: E6 = eisenstein_series_qexp(6, 3)
sage: E4^6
1/191102976000000 + 1/132710400000*q 
                  + 203/44236800000*q^2 + O(q^3)
sage: E4^3*E6^2
1/3511517184000 - 1/12192768000*q 
                - 377/4064256000*q^2 + O(q^3)
sage: E6^4
1/64524128256 - 1/32006016*q + 241/10668672*q^2 + O(q^3)
\end{verbatim}
%end_sage
\end{example}

In Section~\ref{sec:vmthesis}, we will discuss the
reduced echelon form basis for $M_k$.


\section{The Miller Basis}\label{sec:vmthesis}
\begin{lemma}[V. Miller]\label{lem:vm}
The space $S_k$ has a basis $f_1,\ldots, f_{d}$ such
that if $a_i(f_j)$ is the $i$th coefficient of $f_j$, then
$a_i(f_j) = \delta_{i,j}$ for $i=1,\ldots, {d}$.  Moreover
the $f_j$ all lie in $\Z[[q]]$.  We call this basis the
\defn{Miller basis} for $S_k$.
\end{lemma}
This is a straightforward construction involving $E_4$, $E_6$ and
$\Delta$.  The following proof very closely follows \cite[Ch.~X,
Thm.~4.4]{lang:modular}, which in turn follows the
first lemma of V.~Miller's thesis.
\begin{proof}
Let $d=\dim S_k$.
Since $B_4 = -1/30$ and $B_6=1/42$, we note that
$$
 F_4 = -\frac{8}{B_4} \cdot E_4 = 1 + 240q + 2160q^2 + 6720q^3 + 17520q^4 + \cdots
$$ and 
$$
  F_6 = -\frac{12}{B_6}\cdot E_6 = 1 - 504q - 16632q^2 - 122976q^3 - 532728q^4
+\cdots 
$$
have $q$-expansions in $\Z[[q]]$ with leading
coefficient~$1$.
Choose integers $a,b\geq 0$ such that 
$$
  4a+6b\leq 14 \qquad\text{and}\qquad 4a + 6b \con k \pmod{12},
$$
with $a=b=0$ when $k\con 0\pmod{12}$, and
let 
$$
  g_j = \Delta^j F_6^{2(d-j)+b} F_4^a
      = \left(\frac{\Delta}{F_6^2}\right)^j F_6^{2d + b} F_4^a
,\qquad\qquad \text{for } j=1,\ldots,d.
$$
Then it is elementary to check that $g_j$ has weight $k$
$$
 a_j(g_j) = 1 \qquad\text{and}\qquad a_i(g_j)=0\qquad\text{when}\qquad i<j.
$$
Hence the $g_j$ are linearly independent over $\C$, so form
a basis for $S_k$.  Since $F_4, F_6$, and $\Delta$ are all in $\Z[[q]]$,
so are the $g_j$.  The $f_i$ may then be constructed from the $g_j$
by Gauss elimination.  The coefficients of the resulting power series lie in 
$\Z$ because each time we clear a column we use the power series
$g_j$ whose leading coefficient is $1$ (so no denominators are introduced).
\end{proof}
\begin{remark}
The basis coming from Miller's lemma is ``canonical'', 
since it is just the reduced row echelon form of any basis.
Also the set of all {\em integral} linear combinations of
the elements of the Miller basis are precisely
the modular forms of level $1$ with integral $q$-expansion.
\end{remark}

We extend the Miller basis to all $M_k$ by taking
a multiple of $G_k$ with constant term $1$ and subtracting
off the $f_i$ from the Miller basis so that the
coefficients of $q, q^2, \ldots q^d$ of 
the resulting expansion are $0$.  We call the extra basis
element $f_0$.


\begin{example}
If $k=24$, then $d=2$.  Choose $a=b=0$, since $k\con 0\pmod{12}$.
Then 
$$g_1 = \Delta F_6^2 = 
q - 1032q^2 + 245196q^3 + 10965568q^4 + 60177390q^5 - \cdots
$$
and 
$$
g_2 = \Delta^2 = q^2 - 48q^3 + 1080q^4 - 15040q^5 + \cdots.
$$
We let $f_2=g_2$ and 
$$ 
f_1 = g_1 + 1032 g_2
 = q + 195660q^3 + 12080128q^4 + 44656110q^5 - \cdots.
$$
\end{example}

\begin{example}\label{ex:vmb36}
When $k=36$, the Miller basis including $f_0$ is
\begin{align*}
 f_0 &= 1 + \quad\,\,\, 6218175600q^4 + 15281788354560q^5 + \cdots,\\
 f_1 &= \,\,\,\,\,\,  q + \quad\,\, 57093088q^4 + \,\,\,\,\,\,\,\,\,\,37927345230q^5 +\cdots,\\
 f_2 &=  \qquad q^2 +  \,\,\,\,\,\,194184q^4 + \quad\,\,\,\,\,\,\,\,\,\,\,\,\,\,\, 7442432q^5 + \cdots,\\
 f_3 &=  \qquad\quad\,\, q^3 -\,\,\,\,\,\,\,\,\,\, 72q^4 +\qquad\quad\,\,\,\,\,\,\,\,\,\,\,\,  2484q^5 + \cdots.
\end{align*}
\end{example}

%begin_sage
\sageindex{Miller basis}
\begin{example}
The \sage{} command {\tt victor\_miller\_basis} computes 
the Miller basis to any desired precision
for a given $k$.
\begin{verbatim}
sage: victor_miller_basis(28,5)
[
1 + 15590400*q^3 + 36957286800*q^4 + O(q^5),
q + 151740*q^3 + 61032448*q^4 + O(q^5),
q^2 + 192*q^3 - 8280*q^4 + O(q^5)
]
\end{verbatim}
\end{example}
%end_sage


\begin{remark}\label{rem:elkies}
  To write $f \in M_k$ as a polynomial in~$E_4$ and~$E_6$,
  it is wasteful to compute the Miller basis.  Instead,
  use the upper triangular (but not echelon!) basis $\Delta^j
  F_6^{2(d-j)+a} F_4^b$, and match coefficients from $q^0$ to $q^d$.

\end{remark}

\section{Hecke Operators}\label{sec:hecke_one}
In this section we define Hecke operators on level $1$ modular forms and
derive their basic properties.  We will not give proofs of the
analogous properties for Hecke operators on higher level modular
forms, since the proofs are clearest in the level $1$ case, and the
general case is similar (see, e.g., \cite{lang:modular}).

For any positive integer~$n$, let
$$
X_n = \left\{\mtwo{a}{b}{0}{d} \in \Mat_2(\Z) \,\,:\,\,
a\geq 1,\,\, ad=n, \text{ and } 0\leq b <d\right\}.
$$
Note that the set $X_n$ is in bijection with the set of subgroups of
$\Z^2$ of index~$n$, where $\abcd{a}{b}{c}{d}$ corresponds to
$L=\Z\cdot (a,b) + \Z\cdot (0,d)$, as one can see using
Hermite normal form, which is the analogue over $\Z$ of 
echelon form (see \exref{ch:linalg}{ex:hnf}).

Recall from \eqref{eqn:bargamma} that
if $\gamma=\abcd{a}{b}{c}{d}\in\GL_2(\Q)$, then
$$
   f^{[\gamma]_k} = \det(\gamma)^{k-1} (cz+d)^{-k} f(\gamma(z)). 
$$


\begin{definition}[Hecke Operator $T_{n,k}$]\label{defn:hecke_op}
The~$n$th Hecke operator $T_{n,k}$ of weight~$k$
is the operator on the set of functions on $\h$ defined by
$$
T_{n,k}(f) = \sum_{\gamma \in X_n} 
f^{[\gamma]_k}.
$$
\end{definition}
\begin{remark}
  It would make more sense to write $T_{n,k}$ on the right, e.g.,
  $f|T_{n,k}$, since $T_{n,k}$ is defined using
  a right group action.
  However, if $n,m$ are integers, then the action of $T_{n,k}$ and
  $T_{m,k}$ on weakly modular functions commutes (by
  Proposition~\ref{prop:hecke_com} below), so it makes no difference
  whether we view the Hecke operators of given weight~$k$ as
  acting on the right or left.
\end{remark}

\begin{proposition}\label{prop:tn_presweak}
  If~$f$ is a weakly modular function of weight~$k$, then so is
  $T_{n,k}(f)$; if~$f$ is a modular function, then so is $T_{n,k}(f)$.
\end{proposition}
\begin{proof}
Suppose $\gamma\in\SL_2(\Z)$.  Since~$\gamma$ 
induces an automorphism of $\Z^2$,
$$X_n \cdot \gamma = \{ \delta \gamma : \delta\in X_n\}$$
is also in bijection with the subgroups
of $\Z^2$ of index~$n$.  
For each element $\delta\gamma \in X_n\cdot \gamma$, there
is $\sigma\in\SL_2(\Z)$ such that $\sigma\delta\gamma\in X_n$ (the
element~$\sigma$ 
transforms $\delta\gamma$ to Hermite normal form),
and the set of elements $\sigma\delta\gamma$ is thus equal to $X_n$.
Thus 
$$
 T_{n,k}(f) = \sum_{\sigma\delta\gamma\in X_n} f^{[\sigma\delta\gamma]_k} 
     = \sum_{\delta\in X_n} f^{[\delta\gamma]_k}
     = T_{n,k}(f)^{[\gamma]_k}.
$$
A finite sum of meromorphic function is
meromorphic, so $T_{n,k}(f)$ is weakly modular. 
If $f$ is holomorphic on $\h$, then
each $f^{[\delta]_k}$ is holomorphic on $\h$ for $\delta \in X_n$.
A finite sum of holomorphic functions is holomorphic,
so $T_{n,k}(f)$ is holomorphic.

\end{proof}

We will frequently drop $k$ from the notation in $T_{n,k}$, since
the weight~$k$ is implicit in the modular function to which
we apply the Hecke operator.  Henceforth we make the
convention that if we write $T_{n}(f)$ and if $f$ is modular, 
then we mean $T_{n,k}(f)$, where~$k$ is the weight of~$f$.

\begin{proposition}\label{prop:hecke_com}
On weight $k$ modular functions we have
\begin{equation}\label{eqn:hecke_mul}
  T_{mn} = T_{m} T_{n} \qquad\qquad\qquad\qquad\quad
 \text{if }(m,n)=1,
\end{equation}
and
\begin{equation}\label{eqn:hecke_recur}
  T_{p^n} = T_{p^{n-1}} T_p\, -\, p^{k-1} T_{p^{n-2}} \qquad\!\!
\text{if $p$ is prime}.
\end{equation}
\end{proposition}
\begin{proof}
  Let $L$ be a subgroup of index $mn$.  The quotient $\Z^2/L$ is an
  abelian group of order $mn$, and $(m,n)=1$, so $\Z^2/L$ decomposes
  uniquely as a direct sum of a subgroup of order~$m$ with a subgroup of
  order~$n$.  Thus there exists a unique subgroup $L'$ such that
  $L\subset L'\subset \Z^2$, and $L'$ has index $m$ in $\Z^2$.  The
  subgroup $L'$ corresponds to an element of $X_m$, and the index~$n$
  subgroup $L\subset L'$ corresponds to multiplying that element on
  the right by some uniquely determined element of $X_n$.  We thus
  have $$ \SL_2(\Z) \cdot X_{m} \cdot X_n = \SL_2(\Z) \cdot X_{mn}, $$
  i.e., the set products of elements in $X_m$ with elements of $X_n$
  equal the elements of $X_{mn}$, up to $\SL_2(\Z)$-equivalence. 
  Thus for any~$f$, we have
  $T_{mn}(f) = T_{n}(T_m(f))$.  Applying this formula with~$m$ and~$n$
  swapped yields the equality $T_{mn} = T_{m} T_{n}$.

  We will show that $T_{p^n} + p^{k-1} T_{p^{n-2}} = T_p T_{p^{n-1}}$.
  Suppose~$f$ is a weight~$k$ weakly modular function.  
Using that $f^{[\abcd{p}{0}{0}{p}]_k} 
   = (p^2)^{k-1}p^{-k} f = p^{k-2} f$, we have
$$
 \sum_{x\in X_{p^n}} f^{[x]_k}\,\, +\,\, p^{k-1}\!\!\! \sum_{x\in
      X_{p^{n-2}}} f^{[x]_k}
   =  \sum_{x\in X_{p^n}} f^{[x]_k} \,\,\,+\, p\!\!\!\! \sum_{x\in pX_{p^{n-2}}} 
  f^{[x]_k}.
$$
Also $$T_p T_{p^{n-1}}(f) = \sum_{y\in X_p} \sum_{x\in X_{p^{n-1}}}
(f^{[x]_k})^{[y]_k} = \sum_{x \in X_{p^{n-1}}\cdot X_p} f^{[x]_k}. $$ Thus it
suffices to show that $X_{p^n}$ disjoint union~$p$ copies of $p X_{p^{n-2}}$ is
equal to $X_{p^{n-1}}\cdot X_p$, where we consider elements with
multiplicities and up to left $\SL_2(\Z)$-equivalence (i.e., the left
action of $\SL_2(\Z)$).

  Suppose~$L$ is a subgroup of $\Z^2$ of index $p^n$, so~$L$
  corresponds to an element of $X_{p^n}$.  First suppose~$L$ is not
  contained in $p\Z^2$.  Then the image of~$L$ in $\Z^2/p\Z^2 =
  (\Z/p\Z)^2$ is of order~$p$, so if $L'=p\Z^2 + L$, then
  $[\Z^2:L']=p$ and $[L:L']=p^{n-1}$, and $L'$ is the only subgroup
  with this property.  Second, suppose that $L\subset p\Z^2$ if of
  index $p^n$ and that $x\in X_{p^n}$ corresponds to $L$.  Then every
  one of the $p+1$ subgroups $L'\subset \Z^2$ of index~$p$
  contains~$L$. Thus there are $p+1$ chains $L\subset L'\subset \Z^2$
with $[\Z^2:L']=p$.

The chains $L\subset L'\subset\Z^2$ with $[\Z^2:L']=p$ and
$[\Z^2:L]=p^{n-1}$ are in bijection with the elements of
$X_{p^{n-1}}\cdot X_p$.  On the other hand the union of $X_{p^n}$ with
$p$ copies of $p X_{p^{n-2}}$ corresponds to the subgroups~$L$ of index
$p^{n}$, but with those that contain $p\Z^2$ counted $p+1$ times. The
structure of the set of chains $L\subset L'\subset\Z^2$ that we
derived in the previous paragraph gives the result.
\end{proof}

\begin{corollary}
The Hecke operator $T_{p^n}$, for prime~$p$, is a polynomial 
in $T_p$ with integer coefficients, i.e., $T_{p^n}\in\Z[T_p]$.
If $n,m$ are any integers, then $T_n T_m = T_m T_n$.
\end{corollary}
\begin{proof}
  The first statement follows from \eqref{eqn:hecke_recur}
  of Proposition~\ref{prop:hecke_com}. 
  It then follows that
  $T_n T_m = T_m T_n$ when $m$ and $n$ are both powers of a single prime~$p$.  Combining this with \eqref{eqn:hecke_mul} gives the second
  statement in general.
\end{proof}


\begin{proposition}\label{prop:qexpTn}
Let $f = \sum_{n\in\Z} a_n q^n$ be a modular function
of weight~$k$.
Then 
$$
   T_{n}(f)
     = \sum_{m\in\Z} 
\left(\sum_{1\leq d\,\mid\, \gcd(n,m)} d^{k-1} a_{mn/d^2}\right) q^m.
$$
In particular, if $n=p$ is prime, then
$$
  T_p(f) = \sum_{m\in \Z} \left(a_{mp} + p^{k-1} a_{m/p}\right) q^m,
$$
where $a_{m/p}=0$ if $m/p\not\in\Z$.
\end{proposition}
\begin{proof}
This is proved in \cite[\S
VII.5.3]{serre:arithmetic} by writing out $T_n(f)$ explicitly and
using that $\sum_{0\leq b<d} e^{2\pi i bm/d}$ is $d$ if $d\mid m$
and~$0$ otherwise. 
\end{proof}

\begin{corollary}\label{cor:tpres} 
 The Hecke operators preserve $M_k$ and $S_k$.
\end{corollary}
\begin{remark} Alternatively, for $M_k$ the above corollary is
  Proposition~\ref{prop:tn_presweak}, and for~$S_k$ we see from the
  definitions that if $f(\infty)=0$, then $T_n f$ also vanishes
  at~$\infty$.
\end{remark}

\begin{example}
Recall from \eqref{eqn:ekexp} that 
$$
E_4 = \frac{1}{240} + q + 9q^2 + 28q^3 + 73q^4 + 126q^5 + 252q^6 + 344q^7
+\cdots.
$$
Using the formula of Proposition~\ref{prop:qexpTn}, we see
that 
$$
T_2(E_4) = (1/240 + 2^3\cdot (1/240)) + 9q + (73 + 2^3\cdot 1) q^2 + \cdots.
$$
Since $M_4$ has dimension $1$ and since we have proved that $T_2$
preserves $M_4$, we know that $T_2$ acts as a scalar.  Thus
we know just from the constant coefficient of $T_2(E_4)$ that
$$T_2(E_4) = 9 E_4.$$
More generally, for $p$ prime we see by inspection of
the constant coefficient of $T_p(E_4)$ that 
$$
  T_p(E_4) = (1+p^3)E_4.
$$
In fact
$
 T_n(E_k) = \sigma_{k-1}(n) E_k,
$
for any integer $n\geq 1$ and even weight $k\geq 4$.
\end{example}

\begin{example}
  By Corollary~\ref{cor:tpres}, the Hecke operators $T_n$ also
  preserve the subspace $S_k$ of $M_k$.  Since $S_{12}$ has dimension
  $1$ (spanned by $\Delta$), we see that~$\Delta$ is an eigenvector
  for every $T_n$.  Since the coefficient of $q$ in the $q$-expansion
  of $\Delta$ is~$1$, the eigenvalue of $T_n$ on $\Delta$ is the $n$th
  coefficient of $\Delta$.  Since $T_{nm}=T_n T_m$ for $\gcd(n,m)=1$,
  we have proved the nonobvious fact that the \defn{Ramanujan
    function} $\tau(n)$ that gives the $n$th coefficient of $\Delta$
  is a multiplicative function, i.e., if $\gcd(n,m)=1$, then
  $\tau(nm)=\tau(n)\tau(m)$.
\end{example}

\begin{remark}
  The Hecke operators respect the decomposition $M_k = S_k \oplus \C E_k$,
  i.e., for all $k$ the series $E_k$ are eigenvectors for all $T_n$.
\end{remark}

\section{Computing Hecke Operators}
This section is about how to compute matrices
of Hecke operators on $M_k$.

\begin{algorithm}{Hecke Operator}
  This algorithm computes the matrix of the Hecke operator $T_n$ on the
 Miller basis for $M_k$.
\begin{steps}
\item{}[Dimension] Compute $d= \dim(M_k)-1$ using
  Corollary~\ref{cor:dim1}.
\item{}[Basis] Using 
  Lemma~\ref{lem:vm}, compute the echelon 
  basis $f_0,\ldots, f_d$ for $M_k$ $\pmod{q^{dn+1}}$.
\item{}[Hecke operator] Using 
  Proposition~\ref{prop:qexpTn}, compute for
  each $i$ the image $T_n(f_i) \pmod{q^{d+1}}$ .
\item{}[Write in terms of basis]\label{usevm} The elements 
  $T_n(f_i) \pmod{q^{d+1}}$
  determine linear combinations of $$f_0,f_1,\ldots, f_d
  \pmod{q^d}.$$  These linear combinations are easy to find
once we compute $T_n(f_i) \pmod{q^{d+1}}$,
since our basis of~$f_i$ is in echelon form.
The linear combinations are just the coefficients 
of the power series $T_n(f_i)$ up to and including $q^d$.
\item{}[Write down matrix] The matrix of $T_n$ acting from the right
  relative to the basis $f_0, \ldots, f_d$ is the matrix whose rows
  are the linear combinations found in the previous step, i.e., whose
  rows are the coefficients of $T_n(f_i)$.
\end{steps}
\end{algorithm}
\begin{proof}
  By Proposition~\ref{prop:qexpTn}, the $d$th coefficient
  of $T_n(f)$ involves only $a_{dn}$ and smaller-indexed coefficients
  of~$f$. We need only compute a modular form~$f$ modulo $q^{dn+1}$ in
  order to compute $T_n(f)$ modulo $q^{d+1}$.
  Uniqueness in step~(\ref{usevm}) follows from Lemma~\ref{lem:vm}
  above.
\end{proof}


\begin{example}
We compute the Hecke operator $T_2$ on $M_{12}$ using
the above algorithm.
\begin{steps}
\item{}[Compute dimension] We have $d=2 - 1 = 1$.
\item{}[Compute basis]  Compute up to (but not including)
the coefficient of $q^{dn+1} = q^{1\cdot 2 + 1} = q^3$.
As given in the proof of Lemma~\ref{lem:vm}, we have
$$
 F_4 = 1 + 240q + 2160q^2 + \cdots \quad\text{and}\quad
 F_6 = 1 - 504q - 16632q^2 + \cdots.
$$
Thus $M_{12}$ has basis
$$
 F_4^3 = 1 + 720q + 179280q^{2} + \cdots
\quad\text{and}\quad
 \Delta = (F_4^3 - F_6^2)/1728 = q - 24q^{2} + \cdots.
$$
Subtracting $720\Delta$ from $F_4^3$ yields the 
echelon basis, which is
$$
f_0 = 1 + 196560q^{2} + \cdots
\quad\text{and}\quad
f_1 = q - 24q^{2} + \cdots.
$$
%begin_sage
\sageindex{Eisenstein arithmetic}%
\sage{} does the arithmetic in the above calculation
as follows:
\begin{verbatim}
sage: R.<q> = QQ[['q']]    
sage: F4 =  240 * eisenstein_series_qexp(4,3)
sage: F6 = -504 * eisenstein_series_qexp(6,3)
sage: F4^3
1 + 720*q + 179280*q^2 + O(q^3)
sage: Delta = (F4^3 - F6^2)/1728; Delta
q - 24*q^2 + O(q^3)
sage: F4^3 - 720*Delta
1 + 196560*q^2 + O(q^3)
\end{verbatim}
%end_sage

\item{}[Compute Hecke operator] In each case letting $a_n$ denote
the $n$th coefficient of $f_0$ or $f_1$, respectively, we have
\begin{align*}
  T_2(f_0) &= T_2 (1 + 196560q^{2} + \cdots) \\
           &= 
           (a_0 + 2^{11} a_0) q^0 + (a_2 + 2^{11} a_{1/2})q^1 + \cdots\\
           &= 2049  + 196560q + \cdots,
\end{align*}
and
\begin{align*}
  T_2(f_1) &= T_2 (q - 24q^{2} + \cdots) \\
           &= (a_0 + 2^{11} a_0) q^0 + (a_2 + 2^{11} a_{1/2})q^1 + \cdots\\
           &= 0 -24 q + \cdots.
\end{align*}
(Note that $a_{1/2} = 0$.)
         
\item{}[Write in terms of basis]
We read off at once that
$$
  T_2(f_0) = 2049 f_0 + 196560 f_1
\quad\text{and}\quad
 T_2(f_1) = 0 f_0 + (-24)f_1.
 $$

\item{}[Write down matrix] Thus the matrix of $T_2$, acting
from the right on the basis $f_0$, $f_1$, is 
$$
  T_2 = \mtwo{2049}{196560}{0}{-24}.
$$
\end{steps}
As a check note that the characteristic polynomial
of $T_2$ is $(x-2049)(x+24)$ and that 
$2049 = 1 + 2^{11}$ is the sum of the $11$th powers of the
divisors of~$2$.
\end{example}

\begin{example}
The Hecke operator $T_2$ on $M_{36}$ with respect to the
echelon basis is
$$
\left(
\begin{matrix}
\hfill 34359738369&\hfill0&\hfill6218175600&\hfill9026867482214400\\
0&\hfill0&\hfill34416831456&\hfill5681332472832\\
0&\hfill1&\hfill194184&\hfill-197264484\\
0&\hfill0&\hfill-72&\hfill-54528\\
\end{matrix}\right).
$$
It has characteristic polynomial 
$$
(x - 34359738369) \cdot
(x^3 - 139656x^2 - 59208339456x - 1467625047588864),
$$
where the cubic factor is irreducible.

%begin_sage
\sageindex{$M_{36}$}
The {\tt echelon\_form()} command creates the space of modular forms
but with basis in echelon form (which is not the default).
\begin{verbatim}
sage: M = ModularForms(1,36, prec=6).echelon_form()
sage: M.basis()
[
1 + 6218175600*q^4 + 15281788354560*q^5 + O(q^6),
q + 57093088*q^4 + 37927345230*q^5 + O(q^6),
q^2 + 194184*q^4 + 7442432*q^5 + O(q^6),
q^3 - 72*q^4 + 2484*q^5 + O(q^6)
]
\end{verbatim}%link

Next we compute the matrix of the Hecke operator $T_2$.
%link
\begin{verbatim}
sage: T2 = M.hecke_matrix(2); T2
[34359738369   0       6218175600 9026867482214400]
[          0   0      34416831456    5681332472832]
[          0   1           194184       -197264484]
[          0   0              -72           -54528]
\end{verbatim}%link

Finally we compute and factor its characteristic polynomial.

%link
\begin{verbatim}
sage: T2.charpoly().factor()
(x - 34359738369) * 
   (x^3 - 139656*x^2 - 59208339456*x - 1467625047588864)
\end{verbatim}
%end_sage
\end{example}


The following is a famous  open problem about
Hecke operators on modular forms of level $1$.  It generalizes
our above observation that the characteristic polynomial of $T_2$
on $M_k$, for $k=12, 36$, factors as a product of a linear factor
and an irreducible factor.
\begin{conjecture}[Maeda]\index{Maeda's conjecture}\index{conjecture!Maeda}
The characteristic polynomial of $T_2$ on $S_k$ is irreducible for any~$k$.
\end{conjecture}

Kevin Buzzard observed that in several specific cases the Galois group of
the characteristic polynomial of $T_2$ is the full symmetric group
(see \cite{buzzard:t2}).  See also \cite{farmer-james:maeda} for more
evidence for the following conjecture:
\begin{conjecture}
For all primes $p$ and all even $k\geq 2$
the characteristic polynomial of $T_{p,k}$ acting
on $S_k$ is irreducible.
\end{conjecture}



\section{Fast Computation of Fourier Coefficients}\label{sec:fasttau}

How difficult is it to compute prime-indexed coefficients of 
$$
 \Delta = \sum_{n=1}^{\infty} \tau(n) q^n 
   = q - 24q^2 + 252q^3 - 1472q^4 + 4830q^5 - 6048q^6 + \cdots?
$$

\begin{theorem}[Bosman, Couveignes, Edixhoven, de Jong, Merkl]\label{conj:edixhoven}
  Let~$p$ be a prime.  There is a probabilistic algorithm to compute
  $\tau(p)$, for prime $p$, that has expected running time polynomial
  in $\log(p)$.
\end{theorem}
\begin{proof}
See \cite{edix:tau}.
\end{proof}

More generally, if $f=\sum a_n q^n$ is an eigenform in some space
$M_k(\Gamma_1(N))$, where $k\geq 2$, then one expects that there is an
algorithm to compute $a_p$ in time polynomial in $\log(p)$.  Bas
Edixhoven, Jean-Marc Couveignes and Robin de Jong have proved that
$\tau(p)$ can be computed in polynomial time; their approach involves
sophisticated techniques from arithmetic geometry (e.g., \'etale
cohomology, motives, Arakelov theory). The ideas they use are inspired
by the ones introduced by Schoof, Elkies and Atkin for quickly
counting points on elliptic curves over finite fields (see
\cite{MR1413578}).

Edixhoven describes (in an email to the author) the strategy as follows:
\begin{enumerate}
\item We compute the mod $\ell$ Galois representation $\rho$
  associated to $\Delta$.  In particular, we produce a polynomial $f$
  such that $\Q[x]/(f)$ is the fixed field of $\ker(\rho)$. This is
  then used to obtain $\tau(p) \pmod{\ell}$ and to do a Schoof-like
  algorithm for computing $\tau(p)$.

\item We compute the field of definition of suitable points of
  order~$\ell$ on the modular Jacobian $J_1(\ell)$ to do part (1) (see
  \cite[Ch.~6]{diamond-shurman} for the definition of $J_1(\ell)$).

\item The method is to approximate the polynomial $f$ in some sense
  (e.g., over the complex numbers or modulo many small primes $r$)
  and to use an estimate from Arakelov theory to determine a
  precision that will suffice.
\end{enumerate}

\section{Fast Computation of Bernoulli Numbers}
\label{sec:bercomp}
This section, which was written jointly with Kevin McGown,
is about computing the 
Bernoulli numbers $B_n$, for $n\geq 0$, defined 
in Section~\ref{sec:fourexpeis} by
\begin{equation}\label{eqn:def_bernoulli2}
  \frac{x}{e^x - 1} = \sum_{n=0}^{\infty} B_n \frac{x^n}{n!}.
\end{equation}

One way to compute $B_n$ is to
multiply both sides of \eqref{eqn:def_bernoulli2}
by $e^x-1$ and equate
coefficients of $x^{n+1}$ to obtain the
recurrence
$$
  B_0=1\text{,}
  \qquad
  B_n=-\frac{1}{n+1}\cdot \sum_{k=0}^{n-1}\binom{n+1}{k}B_{k}.
$$
This recurrence provides a straightforward and easy-to-implement
method for calculating $B_n$ if one is interested in
computing $B_n$ for all~$n$ up to some bound.
For example,
$$
 B_1 = - \frac{1}{2} \cdot \left( \binom{2}{0} B_0\right) = -\frac{1}{2}
$$
and
$$
 B_2 = -\frac{1}{3} \cdot \left( \binom{3}{0} B_0 + \binom{3}{1} B_1  \right)
   = -\frac{1}{3} \cdot \left( 1 - \frac{3}{2}\right) = \frac{1}{6}.
$$
However, computing $B_n$ via the recurrence is slow; it requires
summing over many large terms, it requires storing the numbers
$B_0,\dots,B_{n-1}$, and it takes only limited advantage of
asymptotically fast arithmetic algorithms.  There is also an inductive
procedure to compute Bernoulli numbers that resembles Pascal's
triangle called the Akiyama-Tanigawa algorithm (see
\cite{akiyama-tanigawa}).

Another approach to computing $B_n$ is to use Newton
iteration and asymptotically fast polynomial arithmetic to approximate
$1/(e^x -1)$.  This method yields a very fast algorithm to compute
$B_0, B_2,\ldots, B_{p-3}$ modulo $p$.  
See \cite{buhler:million} for an application of this
method modulo a prime~$p$ to the verification of
Fermat's last theorem for irregular primes up to one million.  
\begin{example}
David Harvey implemented the algorithm of \cite{buhler:million}  
in \sage{} as the command {\tt bernoulli\_mod\_p}:
\sageindex{Bernoulli numbers modulo $p$}
\begin{verbatim}
sage: bernoulli_mod_p(23)
[1, 4, 13, 17, 13, 6, 10, 5, 10, 9, 15]
\end{verbatim}
\end{example}

A third way to compute $B_n$ uses
an algorithm based on Proposition~\ref{prop:zeta_even},
which we explain below (Algorithm~\ref{alg:ber}).
This algorithm appears to
have been independently invented by several people: by Bernd C.
Kellner (see \cite{bernoulli}); by Bill Dayl; and by H. Cohen and
K.~Belabas.

We compute $B_n$ as an exact rational number
by approximating $\zeta(n)$ to very high precision using 
Proposition~\ref{prop:zeta_even}, the
Euler product
$$
  \zeta(s)=\sum_{m=1}^\infty m^{-s} =\prod_{p\text{ prime}} (1-p^{-s})^{-1},
$$
and the following theorem:
\begin{theorem}[Clausen, von Staudt] \label{thm:clausenvs}
For even $n\geq 2$,
$$
  \denom(B_n)=\prod_{p-1 \,\mid\, n} p.
$$
\end{theorem}
\begin{proof}
See \cite[Ch.~X, Thm.~2.1]{lang:modular}.
\end{proof}

\subsection{The Number of Digits of $B_n$}\label{sec:bernnum}
The following is a new quick way to
compute the number of digits of the numerator of $B_n$.  For example,
using it we can compute the number of digits of $B_{10^{50}}$ in
less than a second.

By Theorem~\ref{thm:clausenvs} we have $d_n = \denom(B_n) =
  \prod_{p-1\mid n} p$.  The number of digits of the numerator is thus
$$
  \lceil \log_{10}(d_n \cdot |B_n|) \rceil.
$$
But 
\begin{align*}
\log(|B_n|) &= \log\left(\frac{2\cdot n!}{(2\pi)^n}\,\zeta(n)\right)\\
     &= \log(2) + \log(n!) - n\log(2) - n\log(\pi) + \log(\zeta(n)),
   \end{align*}
and $\zeta(n)\sim 1$ so $\log(\zeta(n))\sim 0$. 
 Finally, Stirling's formula (see \cite[pg.~198--206]{ahlfors}) 
gives a fast way to compute $\log(n!) = \log(\Gamma(n+1))$:
\begin{equation}\label{eqn:loggamma}
  \log(\Gamma(z)) ``=\text{''}
  \frac{\log(2\pi)}{2} + \left(z - \frac{1}{2}\right)\log(z)
            - z + \sum_{m=1}^{\infty} 
            \frac{B_{2m}}{2m(2m-1)z^{2m-1}}.
\end{equation}
We put quotes around the equality
sign because $\log(\Gamma(z))$ does
not converge to its Laurent series.   Indeed, note that
for any fixed value of~$z$ the summands
on the right side go to $\infty$ as $m\to \infty$!
Nonetheless, we can use this formula to very
efficiently compute $\log(\Gamma(z))$,
since if we truncate the
sum, then the error is smaller than the next
term in the infinite sum.
% (this is explained 
%in \cite[pg.~198--206]{ahlfors}).




\subsection{Computing $B_n$ Exactly}
We return to the problem of computing $B_n$.  Let
$$
  K=\frac{2\cdot n!}{(2\pi)^n}
$$
so that $|B_n|=K\zeta(n)$.
Write
$$
  B_n=\frac{a}{d},
$$
with $a,d\in\Z$, $d\geq 1$, and $\gcd(a,d)=1$.  
It is elementary
to show that $a=(-1)^{n/2+1}\;|a|$ for even $n\geq 2$. 
Suppose that using the Euler product we approximate
$\zeta(n)$ from below by a number~$z$ such that
$$
  0\leq \zeta(m)-z<\frac{1}{Kd}.
$$
Then 
$
  0\leq |B_n| - zK < d^{-1};
$
hence
$
  0\leq |a| - zKd < 1.
$
It follows that $|a|=\lceil zKd\rceil$
and hence $a=(-1)^{n/2+1}\;\lceil zKd\rceil$.

It remains to compute~$z$.
Consider the following problem:
given $s>1$ and
$\varepsilon>0$, find $M\in\Z_+$ so that
$$
  z=\prod_{p\leq M} (1-p^{-s})^{-1}
$$
satisfies $0\leq\zeta(s)-z<\varepsilon$.
We always have $0\leq\zeta(s)-z$.
Also,
$$
  \sum_{n\leq M} n^{-s}\leq\prod_{p\leq M} (1-p^{-s})^{-1},
$$
so
$$
  \zeta(s)-z \leq
  \sum_{n=M+1}^\infty n^{-s}
\leq
  \int_M^\infty x^{-s}\;dx
=
  \frac{1}{(s-1)M^{s-1}}.
$$
Thus if $M>\varepsilon^{-1/(s-1)}$, then
$$
  \frac{1}{(s-1)M^{s-1}}
  \leq
  \frac{1}{M^{s-1}}
  <
  \varepsilon
  \,,
$$
so
$\zeta(s)-z<\varepsilon$, as required.
For our purposes, we have $s=n$ and $\varepsilon=(Kd)^{-1}$,
so it suffices to take $M>(Kd)^{1/(n-1)}$.

\begin{algorithm}{Bernoulli Number $B_n$}\label{alg:ber}
Given an integer $n\geq 0$, this algorithm computes
the Bernoulli number $B_n$ as an exact rational number. 
\begin{steps}
\item{}[Special cases] If $n=0$, return $1$; if $n=1$, return $-1/2$;
if $n\geq 3$ is odd, return $0$.
  \item{}[Factorial factor]\label{alg:ber:compk} 
Compute $\ds K=\frac{2 \cdot n!}{(2\pi)^n}$ to sufficiently many
digits of precision so the ceiling in step~(\ref{alg:ber:num}) is
uniquely determined (this precision can be determined using
Section~\ref{sec:bernnum}). 
  \item{}[Denominator] Compute $\ds d=\prod_{p-1 | n} p$.
  \item{}[Bound]\label{alg:ber:compM} 
                  Compute $\ds M = \left\lceil (K d)^{1/(n-1)} \right\rceil$.
  \item{}[Approximate $\zeta(n)$]\label{alg:ber:approx} 
              Compute $\ds z = \prod_{p\leq M} (1-p^{-n})^{-1}$.
  \item{}[Numerator]\label{alg:ber:num} 
    Compute $\ds a = (-1)^{n/2+1}\,\lceil d K z \rceil$.
  \item{}[Output $B_n$] Return $\ds \frac{a}{d}$.
\end{steps}
\end{algorithm}

In step~(\ref{alg:ber:approx}) use a sieve to compute all primes $p\leq
M$ efficiently (which is fast, since $M$ is so small).  In
step~(\ref{alg:ber:compM}) we may replace $M$ by any integer greater
than the one specified by the formula, so we do not have to compute
$(Kd)^{1/(n-1)}$ to very high precision.

In Section~\ref{sec:genber} below we will
generalize the above algorithm.

\begin{example}
We illustrate Algorithm~\ref{alg:ber} by computing $B_{50}$.
Using 135 binary digits of precision, we compute
\begin{eqnarray*}
  K=7500866746076957704747736.71552473164563479.
\end{eqnarray*}
The divisors of $n$ are $1, 2, 5, 10, 25, 50$, so
$$
  d = 2\cdot 3\cdot 11=66.
$$
We find $M=4$ and compute
$$
  z = 1.00000000000000088817842109308159029835012.
$$
Finally we compute
$$
  dKz =
 495057205241079648212477524.999999994425778,
$$
so
$$
  B_{50}=\frac{495057205241079648212477525}{66}.
$$
\end{example}








\begin{exercises}

\item \label{ex:zetaval} Using Proposition~\ref{prop:zeta_even} and the
table on page~\pageref{tbl:bern}, compute
$\sum_{n=1}^{\infty} \frac{1}{n^{26}}$ explicitly. 

\item\label{ex:odd_bernoulli} 
Prove that if $n>1$ is odd, then the Bernoulli number $B_n$ is $0$.


\item \label{ex:expeis} Use \eqref{eqn:ekexp} to write down the
  coefficients of $1$, $q$, $q^2$, and $q^3$ of the Eisenstein series
  $E_8$.

\item \label{ex:g4g6quo}
Suppose $k$ is a positive integer with $k\equiv 0 \pmod{12}$.
Suppose $a,b\geq 0$ are integers with $4a + 6b = k$.
\begin{enumerate}
\item Prove $3 \mid a$.
\item Show that 
  $\ds G_4^a \cdot G_6^b \,/\, G_6^{\frac{k}{6}} = 
     \left(G_4^3 / G_6^2 \right)^{\frac{a}{3}}.$
\end{enumerate}

\item \label{ex:vm} Compute the Miller basis for
$M_{28}(\SL_2(\Z))$ with precision $O(q^8)$.  Your answer will look 
like  Example~\ref{ex:vmb36}.

\item \label{ex:elkies28}
Consider the cusp form $f = q^2 + 192 q^3 - 8280 q^4 + \cdots $
in $S_{28}(\SL_2(\Z))$.  Write~$f$ as a polynomial in $E_4$
and $E_6$ (see Remark~\ref{rem:elkies}).

\item \label{ex:fkint} Let $G_k$ be the weight~$k$ Eisenstein series from
  equation~\eqref{eqn:gk}.  Let~$c$ be the complex number so that the
  constant coefficient of the $q$-expansion of $g = c \cdot G_k$
  is~$1$.  Is it always the case that the $q$-expansion of~$g$ lies in
  $\Z[[q]]$?

\item \label{ex:vm32t2} Compute the matrix of the Hecke operator $T_2$
  on the Miller basis for $M_{32}(\SL_2(\Z))$.  Then compute
  its characteristic polynomial and verify it factors as a product of
  two irreducible polynomials.


\end{exercises}


\par\vspace{2ex}\noindent{\bf What Next?} Much of the rest 
of this book is about methods for computing subspaces of
$M_k(\Gamma_1(N))$ for general $N$ and $k$.  These general methods are
more complicated than the methods presented in this chapter, since
there are many more modular forms of small weight and it can be
difficult to obtain them.  Forms of level $N>1$ have subtle
connections with elliptic curves, abelian varieties, and motives.
Read on for more!





\chapter{Modular Forms of Weight $2$}\label{ch:modform2}

We saw in Chapter~\ref{ch:level1} (especially
Section~\ref{sec:struct1}) that we can compute each space
$M_k(\SL_2(\Z))$ explicitly.  This involves computing Eisenstein
series $E_4$ and $E_6$ to some precision, then forming the basis
$\{E_4^a E_6^b : 4a+6b = k, 0\leq a,b \in \Z\}$ for $M_k(\SL_2(\Z))$.
In this chapter we consider the more general problem of computing
$S_2(\Gamma_0(N))$, for any positive integer $N$.  Again we have a
decomposition
$$
   M_2(\Gamma_0(N)) = S_2(\Gamma_0(N)) \oplus \Eis_2(\Gamma_0(N)),
$$
where $\Eis_2(\Gamma_0(N))$ is spanned by generalized
Eisenstein
 series and $S_2(\Gamma_0(N))$ is the space of cusp forms,
i.e., elements of $M_2(\Gamma_0(N))$ that vanish at {\em all} cusps.

In Chapter~\ref{ch:eisen} we compute the space $\Eis_2(\Gamma_0(N))$
in a similar way to how we computed $M_k(\SL_2(\Z))$.  On the other
hand, elements of $S_2(\Gamma_0(N))$ often cannot be written as sums
or products of generalized Eisenstein series.  In fact, the structure
of $M_2(\Gamma_0(N))$ is, in general, much more complicated than that
of $M_k(\SL_2(\Z))$.  For example, when $p$ is a prime,
$\Eis_2(\Gamma_0(p))$ has dimension $1$, whereas $S_2(\Gamma_0(p))$
has dimension about $p/12$.

Fortunately an idea of Birch, which he called modular symbols, provides a
method for computing $S_2(\Gamma_0(N))$ and indeed for much more
that is relevant to understanding special values of $L$-functions.
Modular symbols are also a powerful theoretical tool.  In this
chapter, we explain how $S_2(\Gamma_0(N))$ is related to modular
symbols and how to use this relationship to explicitly compute a
basis for $S_2(\Gamma_0(N))$. In Chapter~\ref{ch:modsym} we will
introduce more general modular symbols and explain how to use them to
compute $S_k(\Gamma_0(N))$, $S_k(\Gamma_1(N))$ and $S_k(N,\eps)$ for
any integers $k\geq 2$ and $N$ and character~$\eps$.

Section~\ref{sec:modformhecke} contains a very brief summary of basic facts
about modular forms of weight $2$, modular curves, Hecke operators,
and integral homology.  Section~\ref{sec:modsym} introduces modular
symbols and describes how to compute with them.  In
Section~\ref{sec:maninboundary} we talk about how to cut out the
subspace of modular symbols corresponding to cusp forms using the
boundary map.  Section~\ref{sec:basiss2one} is about a straightforward
method to compute a basis for $S_2(\Gamma_0(N))$ using modular
symbols, and Section~\ref{sec:computingmodform} outlines a more
sophisticated algorithm for computing newforms that uses
Atkin-Lehner theory.

Before reading this chapter, you should have read
Chapter~\ref{ch:modform} and Chapter~\ref{ch:level1}.
We also assume familiarity with algebraic curves, Riemann surfaces,
and homology groups of compact Riemann surfaces.

\section{Hecke Operators}
\label{sec:modformhecke}

Recall from Chapter~\ref{ch:modform} that the group $\Gamma_0(N)$ acts
on $\h^*=\h\cup \P^1(\Q)$ by linear fractional transformations.
The quotient $\Gamma_0(N)\backslash \h^*$ is a Riemann surface, which
we denote by $X_0(N)$.  See \cite[Ch.~2]{diamond-shurman} for 
a detailed description of the topology on $X_0(N)$.
The Riemann surface $X_0(N)$
also has a canonical structure of algebraic curve over~$\Q$,
as is explained in \cite[Ch.~7]{diamond-shurman}
(see also \cite[\S6.7]{shimura:intro}).

Recall from Section~\ref{sec:modformN} that a cusp form of weight~$2$
for $\Gamma_0(N)$ is a function~$f$ on~$\h$ such that $f(z)dz$ defines
a holomorphic differential on $X_0(N)$.
Equivalently, a cusp form is a holomorphic function~$f$ on $\h$ such
that
\begin{itemize} 
\item[(a)] the expression $f(z)dz$ is invariant
under replacing~$z$ by $\gamma(z)$ for each $\gamma\in \Gamma_0(N)$ and 
\item[(b)] $f(z)$ vanishes at every cusp for $\Gamma_0(N)$.
\end{itemize}
The space $S_2(\Gamma_0(N))$ of weight $2$ cusp forms on $\Gamma_0(N)$ 
is a finite-dimensional complex
vector space, of dimension equal
to the genus~$g$ of $X_0(N)$. 
The space $X_0(N)(\C)$ is a compact oriented Riemann surface,
so it is a $2$-dimensional oriented real manifold, i.e., 
$X_0(N)(\C)$ is a $g$-holed torus (see Figure~\ref{diagram:homology}
on page~\pageref{diagram:homology}).

Condition (b) in the definition of~$f$ means that~$f$ has a Fourier
expansion about each element of $\P^1(\Q)$.  Thus, at~$\infty$ we have
\begin{align*}
f(z) &= a_1 e^{2\pi i z} + a_2 e^{2 \pi i 2z}
       + a_3 e^{2 \pi i 3z} + \cdots\\
     &= a_1 q + a_2 q^2 + a_3 q^3 + \cdots,
\end{align*}
where, for brevity, we write $q=q(z)=e^{2\pi i z}$.


%\begin{figure}[ht]
%\begin{center}
%\includegraphics[width=.8\textwidth]{diagrams/plot39a}
%\caption{Plot of $q\mapsto \sum a_n q^n$ for $q\in(-1,1)$\label{fig:39amf}}
%\end{center}
%\end{figure}
\begin{example}
Let~$E$ be the elliptic curve defined by the equation
$y^2 + xy = x^3 + x^2 - 4x - 5$.
Let $a_p = p+1-\#\tilde{E}(\F_p)$, where~$\tilde{E}$ 
is the reduction of~$E$ mod~$p$ (note that for the 
primes that divide the conductor of $E$ we have $a_3=-1$, $a_{13}=1$).
For~$n$ composite, define $a_n$ using the relations at the
end of Section~\ref{sec:computingmodform}. 
Then the Shimura-Taniyama conjecture asserts that
\begin{align*}
  f&=q + a_2 q^2 + a_3 q^3 + a_4 q^4 + a_5 q^5 + \cdots\\
   &= q + q^2 - q^3 -q^4 + 2q^5 + \cdots
\end{align*}
is the $q$-expansion of an element of $S_2(\Gamma_0(39))$.
%(see Figure~\ref{fig:39amf}).
This conjecture,\index{conjecture!Shimura-Taniyama}
\index{Shimura-Taniyama conjecture}
which is now a theorem (see
\cite{breuil-conrad-diamond-taylor}), asserts that any
$q$-expansion constructed as above from an elliptic 
curve over~$\Q$ is a modular form.    This conjecture was
mostly proved first by Wiles \cite{wiles:fermat} 
as a key step in the proof of Fermat's last theorem. 
\end{example}


Just as is the case for level $1$ modular forms (see
Section~\ref{sec:hecke_one}) there are commuting Hecke operators $T_1,
T_2, T_3, \ldots$ that act on $S_2(\Gamma_0(N))$.  To define them
conceptually, we introduce an interpretation of the modular curve
$X_0(N)$ as an object whose points {\em parameterize} elliptic curves
with extra structure.

\begin{proposition}\label{prop:modulix0n}
  The complex points of 
  $Y_0(N)=\Gamma_0(N)\backslash \h$ are in natural bijection with
  isomorphism classes of pairs $(E,C)$, where~$E$ is an elliptic
  curve over~$\C$ and~$C$ is a cyclic subgroup of $E(\C)$ of
  order~$N$.  The class of the point $\lambda \in \h$ corresponds to
  the pair
$$
  \left(\C/(\Z + \Z  \lambda),
    \quad \left(\frac{1}{N}\Z + \Z\lambda\right)/(\Z +  \Z \lambda)\right).
$$
\end{proposition}
\begin{proof}
See \exref{ch:modform2}{ex:x0n}.
\end{proof}

Suppose $n$ and $N$ are coprime positive integers.
There
are two natural maps $\pi_1$ and $\pi_2$ from $Y_0(n\cdot N)$ to $Y_0(N)$;
the first, $\pi_1$, sends $(E,C)\in Y_0(n\cdot N)(\C)$ to $(E,C')$, where~$C'$ is
the unique cyclic subgroup of~$C$ of order~$N$, and the second,
$\pi_2$, sends $(E,C)$ to $(E/D, C/D)$,
where~$D$ is the unique cyclic subgroup of~$C$ of order~$n$.  These
maps extend in a unique way to algebraic maps
from $X_0(n\cdot N)$ to $X_0(N)$:
\begin{equation}\label{eqn:corrx0n}
\xymatrix{ & X_0(n\cdot N) \ar[dl]_{\pi_2}\ar[dr]^{\pi_1}\\
  X_0(N) & & X_0(N).}
\end{equation}
The $n$th \defn{Hecke operator} $T_n$ is ${\pi_1}_* \circ \pi_2^*$,
where $\pi_2^*$ and ${\pi_1}_*$ denote pullback and pushforward
of differentials, respectively. (There is a similar definition
of $T_n$ when $\gcd(n,N)\neq 1$.)
Using our interpretation of $S_2(\Gamma_0(N))$ as differentials
on $X_0(N)$, this gives an action of Hecke operators on $S_2(\Gamma_0(N))$.
One can show that these induce the maps of Proposition~\ref{prop:qexpTn}
on $q$-expansions.

\begin{example}
There is a basis of $S_2(39)$ so that 
$$T_2 = \mthree{\hfill 1}{\hfill 1}{\hfill 0}{-2}{-3}{-2}{\hfill 0}{\hfill 0}{\hfill 1} \quad\text{and}\quad
  T_5 = \mthree{-4}{-2}{-6}{\hfill4}{\hfill4}{\hfill4}{\hfill0}{\hfill0}{\hfill2}.$$
Notice that these matrices commute.  Also, the characteristic
polynomial of $T_2$ is $(x - 1) \cdot (x^{2} + 2x - 1)$.
\end{example}

\subsection{Homology}
The first homology group $\H_1(X_0(N),\Z)$ is the group of closed
$1$-cycles modulo boundaries of $2$-cycles (formal sums of images of
$2$-simplexes).  Topologically $X_0(N)$ is a $g$-holed torus,
where~$g$ is the genus of $X_0(N)$.  Thus $\H_1(X_0(N),\Z)$ is a free
abelian group of rank $2g$ (see, e.g.,
\cite[Ex.~19.30]{greenberg-harper} and \cite[\S6.1]{diamond-shurman}),
with two generators corresponding to each hole, as illustrated in the
case $N=39$ in Figure~\ref{diagram:homology}.
\begin{figure}[ht]
\begin{center}
\includegraphics[width=4in]{diagrams/torus39}
$$\H_1(X_0(39),\Z) \isom \Z\cross\Z\cross\Z\cross\Z\cross\Z\cross\Z$$
\end{center}
\caption{The homology of $X_0(39)$.\label{diagram:homology}}
\end{figure}

The homology of $X_0(N)$ is closely related to modular forms, since
the Hecke operators~$T_n$ also act on $\H_1(X_0(N),\Z)$.  The action
is by pullback of homology classes by~$\pi_2$ followed by taking the
image under~$\pi_1$, where $\pi_1$ and $\pi_2$ are as in
\eqref{eqn:corrx0n}.

Integration defines a pairing
\begin{equation}\label{eqn:pairing}
  \langle\,,\,\rangle: S_2(\Gamma_0(N)) \cross \H_1(X_0(N),\Z) \ra \C.
\end{equation}
Explicitly, for a path~$x$,
$$
  \langle f , x \rangle = 2\pi i \cdot \int_{x} f(z) dz.
$$ 
\begin{theorem}\label{thm:pairing2}
The pairing \eqref{eqn:pairing} is nondegenerate and
Hecke equivariant in the sense that for every Hecke operator $T_n$, we have 
$\langle f T_n, x\rangle = \langle f, T_n x\rangle$.
Moreover, it induces a perfect pairing
\begin{equation}
  \langle\,,\,\rangle: S_2(\Gamma_0(N)) \cross \H_1(X_0(N),\R) \ra \C.
\end{equation}
\end{theorem}
This is a special case of the results in Section~\ref{sec:pairing}.
%(See also Theorem~\ref{thm:shokoruvpairing} below.)

As we will see, modular symbols allow us to make explicit the action
of the Hecke operators on $\H_1(X_0(N),\Z)$; the above pairing then
translates this into a wealth of information about cusp forms.

We will also consider the relative homology group
$\H_1(X_0(N),\Z;\{\text{cusps}\})$ of $X_0(N)$ \defn{relative to the
  cusps}; it is the same as usual homology, but in addition we
allow paths with endpoints in the cusps instead of restricting to
closed loops.
Modular symbols provide a ``combinatorial'' presentation of
$\H_1(X_0(N),\Z)$ in terms of paths between elements of $\P^1(\Q)$.

\section{Modular Symbols}\label{sec:modsym}
Let~$\sM_2$ be the free abelian group with basis the set of 
symbols $\{\alpha,\beta\}$ with $\alpha,\beta\in\P^1(\Q)$
modulo the 3-term relations
$$
  \{\alpha,\beta\} + \{\beta,\gamma\} + \{\gamma,\alpha\} = 0
$$  
above and modulo any torsion.
Since $\sM_2$ is torsion-free, we have
$$
  \{\alpha,\alpha\} = 0\qquad\text{and}\qquad
  \{\alpha,\beta\} = -\{\beta,\alpha\}.
$$
\begin{remark}[Warning]
  The symbols $\{\alpha,\beta\}$ satisfy the relations
  $\{\alpha,\beta\} = -\{\beta,\alpha\}$, so order matters.  The notation
  $\{\alpha,\beta\}$ looks like the set containing two elements, which
  strongly (and incorrectly) suggests that the order does not matter.
  This is the standard notation in the literature.
\end{remark}
\begin{figure}[ht]
\psfrag{infinity}{$\infty$}
\psfrag{alpha}{$\alpha$}
\psfrag{beta}{$\beta$}
\psfrag{0}{$0$}
\psfrag{Q}{$\Q$}
\begin{center}
\includegraphics[height=2.5in]{diagrams/alpha-beta}
\end{center}
\caption{The modular symbols $\{\alpha,\beta\}$ and $\{0,\infty\}$.
\label{figure:modsym}}
\end{figure}
As illustrated in Figure~\ref{figure:modsym}, we ``think of'' this
modular symbol as the homology class, relative to the cusps, of a
path from~$\alpha$ to~$\beta$ in $\h^*$.

Define a \defn{left action of $\GL_2(\Q)$} on $\sM_2$ by
letting $g\in\GL_2(\Q)$ act by
$$
  g\{\alpha, \beta\} = \{g(\alpha), g(\beta)\},
$$
and~$g$ acts on~$\alpha$ and~$\beta$ via the
corresponding linear fractional transformation.
The space $\sM_2(\Gamma_0(N))$ of 
\defn {modular symbols for $\Gamma_0(N)$}
is the quotient of $\sM_2$ by the submodule 
generated by the infinitely many elements
of the form $x - g(x)$, for~$x$ in ~$\sM_2$
and~$g$ in $\Gamma_0(N)$, and modulo any torsion.  
A modular symbol for $\Gamma_0(N)$ is an element of 
this space.   We frequently denote the equivalence
class of a modular symbol by giving a 
representative element. 
\begin{example}\label{ex:01}
  Some modular symbols are $0$ no matter what the level $N$ is!  For
  example, since $\gamma=\abcd{1}{1}{0}{1}\in \Gamma_0(N)$, we have
$$\{\infty,0\} = \{\gamma(\infty),\gamma(0)\} = \{\infty,1\},$$
so
$$0 = \{\infty,1\}-\{\infty,0\} = \{\infty,1\}  + \{0,\infty\} 
  =  \{0,\infty\} + \{\infty,1\} = \{0,1\}.$$
  See \exref{ch:modform2}{ex:modsymzero} for a generalization of this
  observation.
\end{example}

There is a natural homomorphism 
\begin{equation}\label{eqn:homtoms}
  \varphi: \sM_2(\Gamma_0(N)) \to \H_1(X_0(N),\{\text{cusps}\},\Z)
\end{equation}
that sends a formal linear combination of geodesic paths in the upper
half plane to their image as paths on $X_0(N)$.  In
\cite{manin:parabolic} Manin proved that \eqref{eqn:homtoms} is an
isomorphism (this is a fairly involved topological argument).  

Manin identified the subspace of $\sM_2(\Gamma_0(N))$ that is
sent isomorphically onto $\H_1(X_0(N),\Z)$.  Let $\sB_2(\Gamma_0(N))$
denote the free abelian group whose basis is the finite set
$C(\Gamma_0(N)) = \Gamma_0(N)\backslash \P^1(\Q)$ of cusps for
$\Gamma_0(N)$.  The \defn{boundary map}
$$\delta: \sM_2(\Gamma_0(N))\ra \sB_2(\Gamma_0(N))$$ sends
$\{\alpha,\beta\}$ to $\{\beta\}-\{\alpha\}$, where $\{\beta\}$ denotes the
basis element of $\sB_2(\Gamma_0(N))$ corresponding to
$\beta\in\P^1(\Q)$.  The kernel $\sS_2(\Gamma_0(N))$ of~$\delta$ is
the subspace of \defn{cuspidal modular symbols}.  Thus an element of
$\sS_2(\Gamma_0(N))$ can be thought of as a linear combination of
paths in $\h^*$ whose endpoints are cusps and whose images in
$X_0(N)$ are homologous to a $\Z$-linear combination of closed
paths.
\begin{theorem}[Manin]
The map~$\varphi$ above induces a canonical isomorphism 
   $$\sS_2(\Gamma_0(N)) \isom \H_1(X_0(N),\Z).$$
\end{theorem}
\begin{proof}
This is \cite[Thm.~1.9]{manin:parabolic}.
\end{proof}
For any (commutative) ring $R$ let 
$$
  \sM_2(\Gamma_0(N), R) = \sM_2(\Gamma_0(N)) \tensor_\Z R
$$
and 
$$
  \sS_2(\Gamma_0(N), R) = \sS_2(\Gamma_0(N)) \tensor_\Z R.
$$
\begin{proposition}\label{prop:dimsym2}
We have 
$$
  \dim_\C\sS_2(\Gamma_0(N),\C) = 2\dim_\C S_2(\Gamma_0(N)).
$$
\end{proposition}
\begin{proof}
We have
 $$\dim_\C\sS_2(\Gamma_0(N),\C) = \rank_\Z \sS_2(\Gamma_0(N)) = 
 \rank_\Z \H_1(X_0(N),\Z) = 2g.$$
\end{proof}

%begin_sage
\begin{example}  
  We illustrate modular symbols in the case when $N=11$.  
  Using \sage{} (below),
  which implements the algorithm that we describe below
  over $\Q$, we find that $\sM_2(\Gamma_0(11);\Q)$ has basis
  $\{\infty,0\}$, $\{-1/8,0\}$, $\{-1/9,0\}$.  
A basis for the integral homology $\H_1(X_0(11),\Z)$ is the
subgroup generated by $\{-1/8,0\}$ and $\{-1/9,0\}$.
\sageindex{modular symbols of level $11$}
\begin{verbatim}
sage: set_modsym_print_mode ('modular')
sage: M = ModularSymbols(11, 2)
sage: M.basis()
({Infinity,0}, {-1/8,0}, {-1/9,0})
sage: S = M.cuspidal_submodule()
sage: S.integral_basis()     # basis over ZZ.
({-1/8,0}, {-1/9,0})
sage: set_modsym_print_mode ('manin')    # set it back
\end{verbatim}

\end{example}
%end_sage

\section{Computing with Modular Symbols}
\subsection{Manin's Trick}\label{sec:manintrick}
In this section, we describe a trick of Manin that we will use to
prove that spaces of modular symbols are computable.

By \exref{ch:modform}{ex:conggamma1} the 
group $\Gamma_0(N)$ has finite index in $\SL_2(\Z)$.
Fix right coset representatives 
$r_0, r_1, \ldots, r_m$ for
$\Gamma_0(N)$ in $\SL_2(\Z)$, so that
$$\SL_2(\Z) = \Gamma_0(N)r_0 \, \union\, \Gamma_0(N)r_1 \, \union\, 
   \cdots \, \union\, \Gamma_0(N)r_m,$$
where the union is disjoint.
For example, when~$N$ is prime, a list of coset representatives is
 $$\mtwo{1}{0}{0}{1}, 
   \mtwo{1}{0}{1}{1},
   \mtwo{1}{0}{2}{1},
   \mtwo{1}{0}{3}{1},
\ldots,
   \mtwo{\hfill 1}{0}{N-1}{1},
   \mtwo{0}{-1}{1}{\hfill 0}.$$

Let 
\begin{equation}\label{eqn:p1znz}
   \P^1(\Z/N\Z)= \{(a:b) \,: \, a,b\in\Z/N\Z,  \, \,\gcd(a,b,N)=1  \,\}/\sim
 \end{equation}
where $(a:b)\sim (a':b')$ if there is $u\in(\Z/N\Z)^*$ such
that $a=ua', b=ub'$.

\begin{proposition}\label{prop:p1bij}
There is a bijection between $\P^1(\Z/N\Z)$ and the 
right cosets of $\Gamma_0(N)$ in $\SL_2(\Z)$, which
sends a coset representative $\abcd{a}{b}{c}{d}$ 
to the class of $(c:d)$ in $\P^1(\Z/N\Z)$.
\end{proposition}
\begin{proof}
See \exref{ch:modform2}{ex:bij}.
\end{proof}
See Proposition~\ref{prop:gamma1cosets} for the
analogous statement for $\Gamma_1(N)$.


We now describe an observation of Manin (see
\cite[\S1.5]{manin:parabolic}) that is crucial to making
$\sM_2(\Gamma_0(N))$ computable.  It allows us to write any modular
symbol $\{\alpha, \beta\}$ as a $\Z$-linear combination of symbols of
the form $r_i\{0,\infty\}$, where the $r_i\in \SL_2(\Z)$ are coset
representatives as above.  In particular, the finitely many symbols
$r_0\{0,\infty\}, \ldots, r_m\{0,\infty\}$ generate
$\sM_2(\Gamma_0(N))$.

\begin{proposition}[Manin]\label{prop:manintrick}
Let $N$ be a positive integer and $r_0,\ldots, r_m$
a set of right coset representatives for $\Gamma_0(N)$ in 
$\SL_2(\Z)$.
Every $\{\alpha, \beta\} \in \sM_2(\Gamma_0(N))$
is a $\Z$-linear combination of 
$r_0\{0,\infty\}, \ldots, r_m\{0,\infty\}$.
\end{proposition}
We give two proofs of the proposition.  The first is useful for
computation (see \cite[\S2.1.6]{cremona:algs}); the second (see
\cite[\S2]{mtt}) is easier to understand conceptually since it does
not require any knowledge of continued fractions.
\begin{proof}[Continued Fractions Proof of Proposition~\ref{prop:manintrick}]
Since $$\{\alpha,\beta\}=\{0,\beta\}-\{0,\alpha\},$$
it suffices to consider modular symbols of the form $\{0,b/a\}$, where
the rational number $b/a$ is in lowest terms.  Expand $b/a$ as a
continued fraction and consider the successive convergents in lowest
terms:
$$
\frac{b_{-2}}{a_{-2}} = \frac{0}{1},\quad
\frac{b_{-1}}{a_{-1}} = \frac{1}{0},\quad
\frac{b_0}{a_0} = \frac{b_0}{1},\,
\ldots,\quad
\frac{b_{n-1}}{a_{n-1}},\quad
\frac{b_n}{a_n}= \frac{b}{a}
$$
where the first two are included formally.
Then
$$b_k a_{k-1} - b_{k-1} a_k = (-1)^{k-1},$$
so that
$$g_k = \mtwo{b_k}{\hfill (-1)^{k-1}b_{k-1}}{a_k}{(-1)^{k-1}a_{k-1}}\in \SL_2(\Z).$$
Hence
$$\left\{\frac{b_{k-1}}{a_{k-1}}, \frac{b_k}{a_k}\right\} 
  = g_k \{ 0, \infty\} = r_i \{0,\infty\},$$
for some~$i$, is of the required special form. 
Since
$$
\{0,b/a\} = \{0,\infty\} + \{\infty,b_0\} + \left\{\frac{b_0}{1}, \frac{b_1}{a_1}\right\}
   + \cdots + \left\{\frac{b_{n-1}}{a_{n-1}}, \frac{b_n}{a_n}\right\},
$$
this completes the proof.
\end{proof}
\begin{proof}[Inductive Proof of Proposition~\ref{prop:manintrick}]
As in the first proof it suffices to prove the proposition
for any symbol $\{0,b/a\}$, where $b/a$ is in lowest terms.
We will induct on $a\in\Z_{\geq 0}$.  If $a=0$, then the symbol
is $\{0,\infty\}$, which corresponds to the identity coset,
so assume that $a>0$.  Find $a'\in\Z$ such that 
$$
  ba'\con 1\pmod{a};
$$
then
$
  b' = (b a' - 1)/a \in\Z
$
so the matrix
$$ 
  \delta = \mtwo{b}{b'}{a}{a'}
$$
is an element of $\SL_2(\Z)$.  Thus
 $\delta = \gamma \cdot r_j$ for some right coset representative $r_j$
and $\gamma \in \Gamma_0(N)$.  Then
$$
   \{0, b/a\} - \{0, b'/a'\} = 
    \{b'/a', b/a\} = \mtwo{b}{b'}{a}{a'}\cdot \{0,\infty\} = r_j \{0,\infty\},
$$
as elements of $\M_2(\Gamma_0(N))$.
By induction, $\{0, b'/a'\}$ is a linear combination of symbols of the
form $r_k\{0,\infty\}$, which completes the proof.
\end{proof}

\begin{example}
Let $N=11$, and consider the modular symbol $\{0,4/7\}$. 
We have
$$\frac{4}{7} 
  = 0 + \frac{1}{1+\frac{1}{1+\frac{1}{3}}},$$
so the partial convergents are
$$\frac{b_{-2}}{a_{-2}} = \frac{0}{1},\quad
\frac{b_{-1}}{a_{-1}} = \frac{1}{0},\quad
\frac{b_0}{a_0} = \frac{0}{1},\quad
\frac{b_1}{a_1} = \frac{1}{1},\quad
\frac{b_2}{a_2} = \frac{1}{2},\quad
\frac{b_3}{a_3} = \frac{4}{7}.
$$

Thus, noting as in Example~\ref{ex:01} that $\{0,1\}=0$, we have
\begin{eqnarray*}
\{0,4/7\}
 &=& \{0,\infty\} +\{\infty,0\} + \{0,1\}+ \{1,1/2\} + \{1/2,4/7\}\\
 &=&
    \mtwo{1}{-1}{2}{-1}\{0,\infty\}
    + \mtwo{4}{1}{7}{2}\{0,\infty\}\\
 &=&
    \mtwo{1}{0}{9}{1}\{0,\infty\}
    + \mtwo{1}{0}{9}{1}\{0,\infty\}\\
 &=&
    2 \cdot \left[\mtwo{1}{0}{9}{1}\{0,\infty\}\right].
\end{eqnarray*}
We compute the convergents of $4/7$ in \sage{}
as follows (note that $0$ and $\infty$ are excluded):
\sageindex{continued fraction convergents}
\begin{verbatim}
sage: convergents(4/7)
[0, 1, 1/2, 4/7]
\end{verbatim}
\end{example}

\subsection{Manin Symbols}\label{sec:maninsyms}\index{Manin symbols}
As above, fix coset representatives $r_0,\ldots,r_m$ for $\Gamma_0(N)$
in $\SL_2(\Z)$.  Consider formal symbols $[r_i]'$ for $i=0,\ldots, m$.
Let $[r_i]$ be the modular symbol $r_i\{0,\infty\} = \{r_i(0),
r_i(\infty)\}$.  We equip
the symbols $[r_0]',\ldots,[r_m]'$ with a \defn{right action of
$\SL_2(\Z)$}, which is given by $[r_i]'.g = [r_j]'$, where $\Gamma_0(N)
r_j = \Gamma_0(N) r_i g$.   We extend the notation by writing  $[\gamma]' = 
[\Gamma_0(N)\gamma]' = [r_i]'$, where $\gamma \in \Gamma_0(N) r_i$.
Then the right action of $\Gamma_0(N)$ is simply $[\gamma]'.g = [\gamma g]'$.

Theorem~\ref{thm:sl2zgen} implies that $\SL_2(\Z)$ 
is generated by the two matrices
$\sigma = \abcd{0}{-1}{1}{\hfill0}$ and 
$\tau=\abcd{1}{-1}{1}{\hfill0}$.
Note that $\sigma=S$ from Theorem~\ref{thm:sl2zgen} and 
$\tau = TS$, so $T = \tau\sigma \in \langle \sigma, \tau\rangle.$

The following theorem provides us with a finite presentation for the
space $\M_2(\Gamma_0(N))$ of modular symbols.\index{modular symbols!finite presentation}
\begin{theorem}[Manin]\label{thm:manin_present}
Consider the quotient $M$ of the free abelian group
on Manin symbols $[r_0]',\ldots,[r_m]'$ by the subgroup
generated by the elements (for all $i$):
$$
{[r_i]'} + [r_i]'\sigma \qquad\text{ and }\qquad
{[r_i]'} + [r_i]'\tau + [r_i]'\tau^2 ,
$$
and modulo any torsion.
Then there is an isomorphism 
$$\Psi: M \xrightarrow{\,\sim\,} \sM_2(\Gamma_0(N))$$
given by $[r_i]' \mapsto [r_i] = r_i \{0,\infty\}$.
\end{theorem}
\begin{proof}
We will only prove that $\Psi$ is surjective;
the proof that $\Psi$ is injective requires much more work and
will be omitted from this book (see \cite[\S1.7]{manin:parabolic}
for a complete proof). 

Proposition~\ref{prop:manintrick} implies that 
$\Psi$ is surjective, assuming that $\Psi$ is
well defined. 
We next verify that $\Psi$ is well defined, i.e.,
that the listed 2-term and 3-term relations hold in
the image. 
To see that the first relation holds, note that
\begin{align*}
  {[r_i]} + [r_i]\sigma 
  &= \{r_i(0),r_i(\infty)\} +  \{r_i\sigma(0), r_i\sigma(\infty)\}\\
&= \{r_i(0),r_i(\infty)\} +  \{r_i(\infty), r_i(0)\} \\
 & = 0.
\end{align*}
For the second relation we have
\begin{align*}
{[r_i]} + [r_i]\tau + [r_i]\tau^2 &= 
    \{r_i(0),r_i(\infty)\} +  \{r_i\tau(0), r_i\tau(\infty)\}
                  + \{r_i\tau^2(0), r_i\tau^2(\infty)\} \\
  &= \{r_i(0),r_i(\infty)\} +  \{r_i(\infty), r_i(1)\}
                  + \{r_i(1), r_i(0)\} \\
  &= 0.
\end{align*}

\end{proof}

%begin_sage
\sageindex{modular symbols}
\begin{example}\label{ex:sagemodsym1}
By default \sage{} computes modular symbols spaces over $\Q$, i.e., 
$\sM_2(\Gamma_0(N);\Q) \isom \sM_2(\Gamma_0(N))\tensor\Q$.
\sage{} represents (weight $2$) Manin symbols as pairs $(c,d)$.
Here $c,d$ are integers that satisfy $0\leq c,d < N$; they define a point
$(c:d) \in \P^1(\Z/N\Z)$, hence a right coset of $\Gamma_0(N)$
in $\SL_2(\Z)$ (see Proposition~\ref{prop:p1bij}).  

Create $\sM_2(\Gamma_0(N);\Q)$ in \sage{} by typing
{\tt ModularSymbols(N, 2)}.  
We then use the \sage{} command {\tt manin\_generators} to enumerate a list
of generators $[r_0], \ldots, [r_n]$ as in
Theorem~\ref{thm:manin_present} for several spaces of modular symbols.
\sageindex{Manin symbols}
\begin{verbatim}
sage: M = ModularSymbols(2,2)
sage: M
Modular Symbols space of dimension 1 for Gamma_0(2) 
of weight 2 with sign 0 over Rational Field
sage: M.manin_generators()
[(0,1), (1,0), (1,1)]

sage: M = ModularSymbols(3,2)
sage: M.manin_generators()
[(0,1), (1,0), (1,1), (1,2)]

sage: M = ModularSymbols(6,2)
sage: M.manin_generators()
[(0,1), (1,0), (1,1), (1,2), (1,3), (1,4), (1,5), (2,1), 
 (2,3), (2,5), (3,1), (3,2)]
\end{verbatim}

Given {\tt x=(c,d)}, the command {\tt x.lift\_to\_sl2z(N)}
computes an element of $\SL_2(\Z)$ whose lower two entries are
congruent to $(c,d)$ modulo~$N$.
\begin{verbatim}
sage: M = ModularSymbols(2,2)
sage: [x.lift_to_sl2z(2) for x in M.manin_generators()]
[[1, 0, 0, 1], [0, -1, 1, 0], [0, -1, 1, 1]]
sage: M = ModularSymbols(6,2)
sage: x = M.manin_generators()[9]
sage: x
(2,5)
sage: x.lift_to_sl2z(6)
[1, 2, 2, 5]
\end{verbatim}

The {\tt manin\_basis} command returns a list of indices into the
Manin generator list such that the corresponding symbols form a basis
for the quotient of the $\Q$-vector space spanned by Manin symbols
modulo the $2$-term and $3$-term relations of
Theorem~\ref{thm:manin_present}.
\begin{verbatim}
sage: M = ModularSymbols(2,2)
sage: M.manin_basis()
[1]
sage: [M.manin_generators()[i] for i in M.manin_basis()]
[(1,0)]
sage: M = ModularSymbols(6,2)
sage: M.manin_basis()
[1, 10, 11]
sage: [M.manin_generators()[i] for i in M.manin_basis()]
[(1,0), (3,1), (3,2)]
\end{verbatim}%link

Thus, e.g., every element of $\sM_2(\Gamma_0(6))$ is a $\Q$-linear
combination of the three symbols $[(1,0)]$, $[(3,1)]$, and $[(3,2)]$.
We can write each of these as a modular symbol using the 
{\tt modular\_symbol\_rep} function.

\sageindex{modular symbols printing}
%link
\begin{verbatim}
sage: M.basis()
((1,0), (3,1), (3,2))
sage: [x.modular_symbol_rep() for x in M.basis()]
[{Infinity,0}, {0,1/3}, {-1/2,-1/3}]
\end{verbatim}

The {\tt manin\_gens\_to\_basis} function returns a matrix
whose rows express each Manin symbol generator in terms
of the subset of Manin symbols that forms a basis (as
returned by {\tt manin\_basis}). 
\begin{verbatim}
sage: M = ModularSymbols(2,2)
sage: M.manin_gens_to_basis()
[-1]
[ 1]
[ 0]
\end{verbatim}
Since the basis is $(1,0)$, this 
means that in $\sM_2(\Gamma_0(2);\Q)$, we have
$[(0,1)] = -[(1,0)]$ and $[(1,1)] = 0$.
(Since no denominators are involved, we have in
fact computed a presentation of $\sM_2(\Gamma_0(2);\Z)$.)

To convert a Manin symbol $x=(c,d)$ to an element of a modular
symbols space $M$, use {\tt M(x)}:
\begin{verbatim}
sage: M = ModularSymbols(2,2)
sage: x = (1,0); M(x)
(1,0)
\end{verbatim}
%sage: M( (3,1) )    # entries are reduced modulo 2 first
%0
%sage: M( (10,19) )    
%-(1,0)

Next consider $\sM_2(\Gamma_0(6);\Q)$:
\begin{verbatim}
sage: M = ModularSymbols(6,2)
sage: M.manin_gens_to_basis()
[-1  0  0]
[ 1  0  0]
[ 0  0  0]
[ 0 -1  1]
[ 0 -1  0]
[ 0 -1  1]
[ 0  0  0]
[ 0  1 -1]
[ 0  0 -1]
[ 0  1 -1]
[ 0  1  0]
[ 0  0  1]
\end{verbatim}
Recall that our choice of basis for 
$\sM_2(\Gamma_0(6);\Q)$ is $[(1,0)], [(3,1)], [(3,2)]$.
Thus, e.g., the first row of this matrix says that
$[(0,1)] = -[(1,0)]$, and the fourth row asserts
that 
$[(1,2)] = -[(3,1)] + [(3,2)]$. 
\begin{verbatim}
sage: M = ModularSymbols(6,2)
sage: M((0,1))
-(1,0)
sage: M((1,2))
-(3,1) + (3,2)
\end{verbatim}

\end{example}
%end_sage


\section{Hecke Operators}




\subsection{Hecke Operators on Modular Symbols}\label{sec:heckemodsym}
When~$p$ is a prime not dividing~$N$, define
$$T_p(\{\alpha,\beta\}) = \mtwo{p}{0}{0}{1}\{\alpha,\beta\}
                  + \sum_{r \hspace{-.6em}\mod p}
                    \mtwo{1}{r}{0}{p} \{\alpha,\beta\}.$$
The Hecke operators are compatible with 
the integration pairing $\langle\, , \, \rangle$
of Section~\ref{sec:modformhecke}, in the
sense that $\langle f T_p, x\rangle = \langle f, T_p x \rangle$.
When $p\mid N$, the definition is the same, except that the matrix
$\abcd{p}{0}{0}{1}$ is not included in the sum (see Theorem~\ref{thm:modsymformpairing}).
There is a similar definition of $T_n$ for $n$ composite
(see Section~\ref{sec:gendefn}).


\begin{example}For example, when $N=11$, we have
\begin{eqnarray*}
T_2\{0,1/5\} &=& \{0,2/5\} + \{0,1/10\} + \{1/2,3/5\}\\
&=& -2\{0,1/5\}. 
\end{eqnarray*}
See Figure~\ref{diagram:t2}.
\begin{figure}[ht]
\begin{center}
\includegraphics[width=0.8\textwidth]{diagrams/t2}
\end{center}
\caption{Image of $\{0,1/5\}$ under $T_2$\label{diagram:t2}}
\end{figure}
\end{example}


\subsection{Hecke Operators on Manin Symbols}\label{sec:heckemanin}
In \cite{merel:1585}, L.~Merel gives a description of the action of
$T_p$ directly on Manin symbols $[r_i]$ (see
Section~\ref{sec:hecke_on_manin} for details).  For example, when
$p=2$ and~$N$ is odd, we have
\begin{equation}\label{eqn:t2merel}
T_2([r_i]) = [r_i]\mtwo{1}{0}{0}{2} + [r_i]\mtwo{2}{0}{0}{1}
              + [r_i]\mtwo{2}{1}{0}{1} + [r_i]\mtwo{1}{0}{1}{2}.
\end{equation}


For any prime, let $C_p$ be the set of matrices
constructed using the following algorithm (see \cite[\S2.4]{cremona:algs}):
\begin{algorithm}{Cremona's Heilbronn Matrices}\index{Heilbronn matrices}\label{alg:crheil}
  Given a prime $p$, this algorithm outputs a list of $2\times 2$
  matrices of determinant $p$ that can be used to compute the Hecke
  operator $T_p$.
\begin{steps}
\item Output $\mtwo{1}{0}{0}{p}$.
\item For $\ds r = \left\lceil -\frac{p}{2} \right\rceil, \ldots, 
                   \left\lfloor \frac{p}{2} \right\rfloor$:
\begin{steps}
\item Let $x_1 = p$, $x_2 = -r$, $y_1 = 0$, $y_2 = 1$, $a=-p$, $b=r$.
\item Output $\mtwo{x_1}{x_2}{y_1}{y_2}$.
\item As long as $b\neq 0$, do the following:
\begin{steps}
\item Let $q$ be the integer closest to $a/b$ (if $a/b$ is a half
  integer, round away from $0$).
\item Let $c=a-bq$, $a=-b$, $b=c$.
\item Set
  $x_3=qx_2 - x_1,\,  x_1 = x_2,\,  x_2 = x_3$, and\\
 $y_3=qy_2 - y_1,\,  y_1 = y_2,\,  y_2 = y_3$.
\item Output $\mtwo{x_1}{x_2}{y_1}{y_2}$.
\end{steps}
\end{steps}
\end{steps}
\end{algorithm}


\begin{proposition}[Cremona, Merel]
Let $C_p$ be as above. Then
for $p\nmid N$ and $[x]\in \M_2(\Gamma_0(N))$ a Manin 
symbol, we have
$$
 T_p([x]) = \sum_{g \in C_p} [xg].
$$
\end{proposition}
\begin{proof}
See Proposition~2.4.1 of \cite{cremona:algs}.
\end{proof}

There are other lists of matrices, due to Merel, that work even when
$p\mid N$ (see Section~\ref{sec:hecke_on_manin}).

%begin_sage
\sageindex{Heilbronn matrices}\index{Heilbronn matrices!\sage}
The command {\tt HeilbronnCremonaList(p)}, for~$p$ prime, 
outputs the list of matrices from Algorithm~\ref{alg:crheil}.
\begin{verbatim}
sage: HeilbronnCremonaList(2)
[[1, 0, 0, 2], [2, 0, 0, 1], [2, 1, 0, 1], [1, 0, 1, 2]]
sage: HeilbronnCremonaList(3)
[[1, 0, 0, 3], [3, 1, 0, 1], [1, 0, 1, 3], [3, 0, 0, 1], 
 [3, -1, 0, 1], [-1, 0, 1, -3]]
sage: HeilbronnCremonaList(5)
[[1, 0, 0, 5], [5, 2, 0, 1], [2, 1, 1, 3], [1, 0, 3, 5], 
 [5, 1, 0, 1], [1, 0, 1, 5], [5, 0, 0, 1], [5, -1, 0, 1], 
 [-1, 0, 1, -5], [5, -2, 0, 1], [-2, 1, 1, -3], 
 [1, 0, -3, 5]]
sage: len(HeilbronnCremonaList(37))
128
sage: len(HeilbronnCremonaList(389))
1892
sage: len(HeilbronnCremonaList(2003))
11662
\end{verbatim}

%end_sage

%begin_sage
\sageindex{Hecke operator $T_2$}
\begin{example}
We compute the matrix of $T_2$ on $\sM_2(\Gamma_0(2))$:
\begin{verbatim}
sage: M = ModularSymbols(2,2)
sage: M.T(2).matrix()
[1]
\end{verbatim}
\end{example}
%end_sage


%begin_sage
\sageindex{Hecke operators $\sM_2(\Gamma_0(6))$}
\begin{example}
We compute some Hecke operators on $\sM_2(\Gamma_0(6))$:
\begin{verbatim}
sage: M = ModularSymbols(6, 2)
sage: M.T(2).matrix()
[ 2  1 -1]
[-1  0  1]
[-1 -1  2]
sage: M.T(3).matrix()
[3 2 0]
[0 1 0]
[2 2 1]
sage: M.T(3).fcp()  # factored characteristic polynomial
(x - 3) * (x - 1)^2
\end{verbatim}
For $p\geq 5$ we have $T_p = p+1$, since $M_2(\Gamma_0(6))$ is
spanned by generalized Eisenstein series (see Chapter~\ref{ch:eisen}).
\end{example}
%end_sage

%begin_sage
\sageindex{Hecke operators $\sM_2(\Gamma_0(39))$}
\begin{example}
We compute the Hecke operators on $\sM_2(\Gamma_0(39))$:
\begin{verbatim}
sage: M = ModularSymbols(39, 2)
sage: T2 = M.T(2)
sage: T2.matrix()
[ 3  0 -1  0  0  1  1 -1  0]
[ 0  0  2  0 -1  1  0  1 -1]
[ 0  1  0 -1  1  1  0  1 -1]
[ 0  0  1  0  0  1  0  1 -1]
[ 0 -1  2  0  0  1  0  1 -1]
[ 0  0  1  1  0  1  1 -1  0]
[ 0  0  0 -1  0  1  1  2  0]
[ 0  0  0  1  0  0  2  0  1]
[ 0  0 -1  0  0  0  1  0  2]
sage: T2.fcp()     # factored characteristic polynomial
(x - 3)^3 * (x - 1)^2 * (x^2 + 2*x - 1)^2
\end{verbatim}%link

The Hecke operators commute, so their
eigenspace structures are related. 

%link
\begin{verbatim}
sage: T2 = M.T(2).matrix()
sage: T5 = M.T(5).matrix()
sage: T2*T5 - T5*T2 == 0
True
sage: T5.charpoly().factor()
(x^2 - 8)^2 * (x - 6)^3 * (x - 2)^2
\end{verbatim}

The decomposition of $T_2$ is a list of the kernels of
$(f^e)(T_2)$, where~$f$ runs through the irreducible factors of the
characteristic polynomial of $T_2$ and $f^e$ exactly divides this
characteristic polynomial.  Using \sage{}, we find them:
\begin{verbatim}
sage: M = ModularSymbols(39, 2)
sage: M.T(2).decomposition()
[
Modular Symbols subspace of dimension 3 of Modular 
Symbols space of dimension 9 for Gamma_0(39) of weight 
2 with sign 0 over Rational Field,
Modular Symbols subspace of dimension 2 of Modular 
Symbols space of dimension 9 for Gamma_0(39) of weight 
2 with sign 0 over Rational Field,
Modular Symbols subspace of dimension 4 of Modular 
Symbols space of dimension 9 for Gamma_0(39) of weight 
2 with sign 0 over Rational Field
]
\end{verbatim}
\end{example}
%end_sage


\section{Computing the Boundary Map}\label{sec:maninboundary}
In Section~\ref{sec:modsym} we defined a map $\delta: \sM_2(\Gamma_0(N)) \to
\sB_2(\Gamma_0(N))$.  The kernel of this map is the space
$\sS_2(\Gamma_0(N))$ of cuspidal modular symbols.  This kernel will be
important in computing cusp forms in Section~\ref{sec:computingmodform}
below.

To compute the boundary map\index{boundary map!computing} on $[\gamma]$, note
that $[\gamma] = \{\gamma(0),\gamma(\infty)\}$, so 
if $\gamma=\abcd{a}{b}{c}{d}$, then 
$$
  \delta([\gamma]) = \{\gamma(\infty)\} - \{\gamma(0)\}
                   = \{a/c\} - \{b/d\}.
$$

Computing this boundary map would appear to first require an algorithm
to compute the set $C(\Gamma_0(N)) = \Gamma_0(N) \backslash \P^1(\Q)$
of cusps for $\Gamma_0(N)$. In fact, there is a trick that computes the
set of cusps in the course of running the algorithm.  First, give an
algorithm for deciding whether or not two elements of $\P^1(\Q)$ are
equivalent modulo the action of $\Gamma_0(N)$.  Then simply construct
$C(\Gamma_0(N))$ in the course of computing the boundary map, i.e.,
keep a list of cusps found so far, and whenever a new cusp class is
discovered, add it to the list.  The following proposition,
which is proved in \cite[Prop.~2.2.3]{cremona:algs}, explains how to determine
whether two cusps are equivalent.

\begin{proposition}[Cremona]\label{prop:cusp1g0}
Let $(c_i,d_i)$, $i=1,2$, be pairs of 
integers with $\gcd(c_i,d_i)=1$ and possibly
$d_i=0$.
There is $g\in\Gamma_0(N)$ such that 
    $g (c_1/d_1)=c_2/d_2$ in $\P^1(\Q)$ if and only if
 $$s_1 d_2 \con s_2 d_1 \pmod{\gcd(d_1 d_2,N)}$$
where $s_j$ satisfies $c_j s_j\con 1\pmod{d_j}$.
\end{proposition}

%begin_sage
\sageindex{boundary map}
\sageindex{cuspidal submodule}
In \sage{} the command {\tt boundary\_map()} computes the boundary map from
$\sM_2(\Gamma_0(N))$ to $\sB_2(\Gamma_0(N))$, and the {\tt
  cuspidal\_submodule()} command computes its kernel.  For example,
for level $2$ the boundary map is given by the matrix $[1 \,\,\, -1]$,
and its kernel is the $0$ space:
\begin{verbatim}
sage: M = ModularSymbols(2, 2)
sage: M.boundary_map()
Hecke module morphism boundary map defined by the matrix
[ 1 -1]
Domain: Modular Symbols space of dimension 1 for 
Gamma_0(2) of weight ...
Codomain: Space of Boundary Modular Symbols for 
Congruence Subgroup Gamma0(2) ...
sage: M.cuspidal_submodule()
Modular Symbols subspace of dimension 0 of Modular 
Symbols space of dimension 1 for Gamma_0(2) of weight 
2 with sign 0 over Rational Field
\end{verbatim}

The smallest level for which the boundary map has nontrivial kernel,
i.e., for which $\sS_2(\Gamma_0(N))\neq 0$, is $N=11$.
\begin{verbatim}
sage: M = ModularSymbols(11, 2)
sage: M.boundary_map().matrix()
[ 1 -1]
[ 0  0]
[ 0  0]
sage: M.cuspidal_submodule()
Modular Symbols subspace of dimension 2 of Modular 
Symbols space of dimension 3 for Gamma_0(11) of weight 
2 with sign 0 over Rational Field
sage: S = M.cuspidal_submodule(); S
Modular Symbols subspace of dimension 2 of Modular 
Symbols space of dimension 3 for Gamma_0(11) of weight 
2 with sign 0 over Rational Field
sage: S.basis()
((1,8), (1,9))
\end{verbatim}%link

The following illustrates that the
Hecke operators preserve $\sS_2(\Gamma_0(N))$:
%link
\begin{verbatim}
sage: S.T(2).matrix()
[-2  0]
[ 0 -2]
sage: S.T(3).matrix()
[-1  0]
[ 0 -1]
sage: S.T(5).matrix()
[1 0]
[0 1]
\end{verbatim}%link

A nontrivial fact is that for $p$ prime the eigenvalue of each
of these matrices is $p+1 - \#E(\F_p)$, where $E$ is the
elliptic curve $X_0(11)$ defined by the (affine) equation
$y^2 + y = x^3 - x^2 - 10x - 20.$ 
For example, we have
%link
\begin{verbatim}
sage: E = EllipticCurve([0,-1,1,-10,-20])
sage: 2 + 1 - E.Np(2)
-2
sage: 3 + 1 - E.Np(3)
-1
sage: 5 + 1 - E.Np(5)
1
sage: 7 + 1 - E.Np(7)
-2
\end{verbatim}%link

The same numbers appear as the eigenvalues
of Hecke operators:
%link
\begin{verbatim}
sage: [S.T(p).matrix()[0,0] for p in [2,3,5,7]]
[-2, -1, 1, -2]
\end{verbatim}
In fact, something similar happens for every elliptic curve over $\Q$.
The book \cite{diamond-shurman} (especially Chapter~8) is about this
striking numerical relationship between the number of 
points on elliptic curves
modulo $p$ and coefficients of modular forms.
%end_sage


\section{Computing a  Basis for $S_2(\Gamma_0(N))$}\label{sec:basiss2one}
This section is about a method for using modular symbols to compute a
basis for $S_2(\Gamma_0(N))$.  It is not the most efficient for
certain applications, but it is  easy to explain and understand.
See Section~\ref{sec:computingmodform} for a method that takes
advantage of additional structure of $S_2(\Gamma_0(N))$.

Let $\sM_2(\Gamma_0(N);\Q)$ and $\sS_2(\Gamma_0(N);\Q)$ be the spaces
of modular symbols and cuspidal modular symbols over $\Q$.  Before we begin, we
describe a simple but crucial fact about the relation between cusp
forms and Hecke operators.\index{Hecke operator}

If $f = \sum b_n q^n \in \C[[q]]$ is a power series,
let $a_n(f) = b_n$ be the $n$ coefficient of $f$.
Notice that $a_n$ is a $\C$-linear map $\C[[q]] \to \C$.

As explained in \cite[Prop.~5.3.1]{diamond-shurman} and
\cite[\S{}VII.3]{lang:modular} (recall also Proposition~\ref{prop:qexpTn}), the
Hecke operators\index{Hecke operator} $T_n$ act on elements of
$M_2(\Gamma_0(N))$ as follows (where $k=2$ below):
\begin{equation}\label{eqn:hecke2}
  T_n\left(\sum_{m=0}^{\infty} a_m q^m\right) = 
\sum_{m=0}^{\infty}
   \left(\sum_{1\leq d\,\mid\, \gcd(n,m)} \eps(d) \cdot d^{k-1} \cdot a_{mn/d^2}\right) q^m,
 \end{equation}
 where $\eps(d)=1$ if $\gcd(d,N)=1$ and $\eps(d)=0$ if $\gcd(d,N)\neq 1$.
(Note: More generally, if $f \in M_k(\Gamma_1(N))$ is a modular form
with Dirichlet character $\eps$, then the above formula holds; above
we are considering this formula in the special case when $\eps$
is the trivial character and $k=2$.)

\begin{lemma}\label{lem:hecketn}
  Suppose $f\in \C[[q]]$ and $n$ is a positive integer.  Let $T_n$
be the operator on $q$-expansions (formal power series) defined
by \eqref{eqn:hecke2}.  Then
$$ 
  a_1(T_n(f)) = a_n(f).
$$
\end{lemma}
\begin{proof}
The coefficient of $q$ in \eqref{eqn:hecke2} is 
$\eps(1)\cdot 1 \cdot a_{1 \cdot n/1^2} = a_n$.
\end{proof}

The \defn{Hecke algebra} $\T$ is the ring generated by all
Hecke operators $T_n$ acting on $M_k(\Gamma_1(N))$.
Let $\T'$ denote the image of the Hecke algebra in
$\End(S_2(\Gamma_0(N)))$, and let $\T'_\C=\T'\tensor_\Z\C$ be the
$\C$-span of the Hecke operators.  Let $\tT_\C$ denote the subring of
$\End(\C[[q]])$ generated over $\C$ by all Hecke operators acting on
formal power series via definition \eqref{eqn:hecke2}.

\begin{proposition}\label{prop:tildeduality}
There is a  bilinear pairing of complex vector spaces
$$
  \C[[q]] \cross \tT_\C \to \C
$$
given by 
$$
  \langle f, t \rangle = a_1(t(f)).
$$
If $f$ is such that $\langle f, t\rangle = 0$ for
all $t\in \tT_\C$, then $f=0$.
\end{proposition}
\begin{proof}
The pairing is bilinear since both $t$ and $a_1$ are linear.

Suppose $f\in \C[[q]]$ is such that 
$\langle f, t\rangle = 0$ for all $t\in \tT_\C$.
Then $\langle f, T_n\rangle = 0$
for each positive integer $n$.
But by Lemma~\ref{lem:hecketn} we have 
$$a_n(f) = a_1(T_n(f)) = 0$$
for all $n$; thus $f=0$. 
\end{proof}

\begin{proposition}\label{prop:duality}
There is a perfect bilinear pairing of complex vector spaces
$$
  S_2(\Gamma_0(N)) \cross \T'_\C \to \C
$$
given by 
$$
  \langle f, t \rangle = a_1(t(f)).
$$
\end{proposition}
\begin{proof}
The pairing has $0$ kernel on the left by
Proposition~\ref{prop:tildeduality}.
Suppose that $t\in\T'_\C$ is such that 
$\langle f, t\rangle =0$ for all $f \in S_2(\Gamma_0(N))$.
Then $a_1(t(f)) = 0$ for all $f$.  For any $n$, the image
$T_n(f)$ is also a cusp form, so $a_1(t(T_n(f))) = 0$
for all $n$ and $f$.  Finally the fact that $\T'$ is commutative
and Lemma~\ref{lem:hecketn} together imply that for all $n$
and $f$,
$$
  0 = a_1(t(T_n(f))) = a_1(T_n(t(f))) = a_n(t(f)),
$$
so $t(f)=0$ for all $f$.  Thus $t$ is the $0$ operator.

Since $S_2(\Gamma_0(N))$ has finite dimension and the 
kernel on each side of the pairing is $0$, it follows
that the pairing is perfect, i.e., defines an {\em isomorphism}
$$
  \T'_\C \isom \Hom_{\C}(S_2(\Gamma_0(N));\C).
$$
\end{proof}


By Proposition~\ref{prop:duality} there is an isomorphism of vector spaces
\begin{equation}\label{eqn:s2isomtdual}
 \Psi: S_2(\Gamma_0(N)) \xrightarrow{\,\,\isom\,\,} \Hom(\T'_\C,\C)
\end{equation}
that sends $f \in S_2(\Gamma_0(N))$ to the homomorphism
$$t \mapsto a_1(t(f)).$$

For any $\C$-linear map $\vphi:\T'_\C\to \C$, let 
$$
 f_{\vphi} = \sum_{n=1}^{\infty} \vphi(T_n)q^n \in\C[[q]].
$$
\begin{lemma}
The series $f_{\vphi}$ is the $q$-expansion of 
$\Psi^{-1}(\vphi) \in S_2(\Gamma_0(N))$.
\end{lemma}
\begin{proof}
Note that it is not even {\em a priori} obvious that $f_{\vphi}$
is the $q$-expansion of a modular form. 
Let $g = \Psi^{-1}(\vphi)$, which is by definition the unique element of 
$S_2(\Gamma_0(N))$ such that $\langle g, T_n \rangle = \vphi(T_n)$
for all $n$.
By Lemma~\ref{lem:hecketn}, we have 
$$
 \langle f_{\vphi}, T_n \rangle = a_1(T_n(f_\vphi)) = a_n(f_{\vphi}) = \vphi(T_n),
$$
so $\langle f_{\vphi} - g, T_n \rangle = 0$ for all $n$.
Proposition~\ref{prop:tildeduality} implies that
$f_\vphi - g = 0$, so $f_{\vphi} = g = \Psi^{-1}(\vphi)$,
as claimed.
\end{proof}
{\bf \noindent{}Conclusion:}  
The cusp forms $f_{\vphi}$, as $\vphi$ varies through a 
basis of $\Hom(\T'_\C,\C)$, form a basis
for $S_2(\Gamma_0(N))$.  
In particular, we can compute $S_2(\Gamma_0(N))$
by computing $\Hom(\T'_{\C},\C)$, where we compute $\T'$ in any way we want,
e.g., using a space that contains an isomorphic copy of $S_2(\Gamma_0(N))$.

\begin{algorithm}{Basis of Cusp Forms}\label{alg:s2basis}
Given positive integers $N$ and $B$, this algorithm computes
a basis for $S_2(\Gamma_0(N))$ to precision $O(q^B)$.
\begin{steps}
\item Compute $\sM_2(\Gamma_0(N);\Q)$
via the presentation of Section~\ref{sec:maninsyms}.
\item Compute the subspace $\sS_2(\Gamma_0(N);\Q)$ 
of cuspidal modular symbols as in Section~\ref{sec:maninboundary}.
\item Let $d=\frac{1}{2}\cdot \dim\sS_2(\Gamma_0(N);\Q)$.  
By Proposition~\ref{prop:dimsym2}, 
$d$ is the dimension of $S_2(\Gamma_0(N))$.
\item Let $[T_n]$ denote the matrix of $T_n$ acting
 on a basis of $\sS_2(\Gamma_0(N);\Q)$.  For a matrix $A$, let 
$a_{ij}(A)$ denote the $ij$th entry of $A$.  
For various integers $i,j$ with $0\leq i,j \leq d-1$,
compute formal $q$-expansions
$$
  f_{ij}(q) = \sum_{n=1}^{B-1} a_{ij}([T_n])q^n + O(q^{B}) \in \Q[[q]]
$$
until we find enough to span a space of dimension $d$ (or exhaust
all of them).  These $f_{ij}$ are
a basis for $S_2(\Gamma_0(N))$ to precision $O(q^{B})$.
\end{steps}
\end{algorithm}

%begin_sage
\sageindex{basis for $S_2(\Gamma_0(N))$}
\subsection{Examples}
We use \sage{} to demonstrate
Algorithm~\ref{alg:s2basis}.

\begin{example}
The smallest $N$ with $S_2(\Gamma_0(N))\neq 0$ is $N=11$.
\begin{verbatim}
sage: M = ModularSymbols(11); M.basis()
((1,0), (1,8), (1,9))
sage: S = M.cuspidal_submodule(); S
Modular Symbols subspace of dimension 2 of Modular 
Symbols space of dimension 3 for Gamma_0(11) of weight 
2 with sign 0 over Rational Field
\end{verbatim}%link

We compute a few Hecke operators, and then read off a 
nonzero cusp form, which forms a basis for $S_2(\Gamma_0(11))$:
%link
\begin{verbatim}
sage: S.T(2).matrix()
[-2  0]
[ 0 -2]
sage: S.T(3).matrix()
[-1  0]
[ 0 -1]
\end{verbatim}
Thus 
$$
 f_{0,0} = q - 2q^2 - q^3 + \cdots \in S_2(\Gamma_0(11))
$$
forms a basis for $S_2(\Gamma_0(11))$.
\end{example}

\begin{example}
We compute a basis for $S_2(\Gamma_0(33))$ to precision $O(q^6)$.
\begin{verbatim}
sage: M = ModularSymbols(33)
sage: S = M.cuspidal_submodule(); S
Modular Symbols subspace of dimension 6 of Modular 
Symbols space of dimension 9 for Gamma_0(33) of weight 
2 with sign 0 over Rational Field
\end{verbatim}%link

Thus $\dim S_2(\Gamma_0(33)) = 3$.
%link
\begin{verbatim}
sage: R.<q> = PowerSeriesRing(QQ)  
sage: v = [S.T(n).matrix()[0,0] for n in range(1,6)]
sage: f00 = sum(v[n-1]*q^n for n in range(1,6)) + O(q^6)
sage: f00
q - q^2 - q^3 + q^4 + O(q^6)
\end{verbatim}%link

This gives us one basis element of $S_2(\Gamma_0(33))$.  It remains
to find two others.   We find
%link
\begin{verbatim}
sage: v = [S.T(n).matrix()[0,1] for n in range(1,6)]
sage: f01 = sum(v[n-1]*q^n for n in range(1,6)) + O(q^6)
sage: f01
-2*q^3 + O(q^6)
\end{verbatim}%link

and 
%link
\begin{verbatim}
sage: v = [S.T(n).matrix()[1,0] for n in range(1,6)]
sage: f10 = sum(v[n-1]*q^n for n in range(1,6)) + O(q^6)
sage: f10
q^3 + O(q^6)
\end{verbatim}%link

This third one is (to our precision) a scalar multiple of the second,
so we look further. 
%link
\begin{verbatim}
sage: v = [S.T(n).matrix()[1,1] for n in range(1,6)]
sage: f11 = sum(v[n-1]*q^n for n in range(1,6)) + O(q^6)
sage: f11
q - 2*q^2 + 2*q^4 + q^5 + O(q^6)
\end{verbatim}
This latter form is clearly not in the span of the first two.
Thus we have the following basis for $S_2(\Gamma_0(33))$ (to
precision $O(q^6)$):
\begin{align*}
f_{00} &= q - q^{2} - q^{3} + q^{4} + \cdots,\\
f_{11} &= q - 2q^{2} + 2q^{4} + q^{5} + \cdots,\\
f_{10} &= q^{3} + \cdots.
\end{align*}
\end{example}

\begin{example}
Next consider $N=23$, where we have 
$$d = \dim S_2(\Gamma_0(23)) = 2.$$
The command {\tt q\_expansion\_cuspforms} computes 
matrices $T_n$ and returns a function
$f$ such that $f(i,j)$ is the $q$-expansion of 
$f_{i,j}$ to some precision. 
(For efficiency reasons, $f(i,j)$ in \sage{} actually
computes matrices of $T_n$ acting
on a basis for the linear dual of $\sS_2(\Gamma_0(N))$.)
\begin{verbatim}
sage: M = ModularSymbols(23)
sage: S = M.cuspidal_submodule()
sage: S
Modular Symbols subspace of dimension 4 of Modular 
Symbols space of dimension 5 for Gamma_0(23) of weight 
2 with sign 0 over Rational Field
sage: f = S.q_expansion_cuspforms(6)
sage: f(0,0)
q - 2/3*q^2 + 1/3*q^3 - 1/3*q^4 - 4/3*q^5 + O(q^6)
sage: f(0,1)
O(q^6)
sage: f(1,0)
-1/3*q^2 + 2/3*q^3 + 1/3*q^4 - 2/3*q^5 + O(q^6)
\end{verbatim}%link

Thus a basis for $S_2(\Gamma_0(23))$ is
\begin{align*}
f_{0,0} &= q - \frac{2}{3}q^{2} + \frac{1}{3}q^{3} - \frac{1}{3}q^{4} - \frac{4}{3}q^{5} + \cdots, \\
f_{1,0} &= -\frac{1}{3}q^{2} + \frac{2}{3}q^{3} + \frac{1}{3}q^{4} - \frac{2}{3}q^{5} + \cdots.
\end{align*}
Or, in echelon form, 
\begin{align*}
 & q - q^{3} - q^{4} + \cdots\\
 & \,\,\,\,\,q^{2} - 2q^{3} - q^{4} + 2q^{5} + \cdots
\end{align*}
which we computed using
%link
\begin{verbatim}
sage: S.q_expansion_basis(6)
[
q - q^3 - q^4 + O(q^6),
q^2 - 2*q^3 - q^4 + 2*q^5 + O(q^6)
]
\end{verbatim}

\end{example}

%end_sage


\section{Computing $S_2(\Gamma_0(N))$ Using Eigenvectors}
\label{sec:computingmodform}
In this section we describe how to use modular symbols to construct a
basis of $S_2(\Gamma_0(N))$ consisting of modular forms that are
eigenvectors for every element of the ring $\T^{(N)}$ generated by the
Hecke operator $T_p$, with $p\nmid N$.  Such eigenvectors are called
\defn{eigenforms}.

Suppose~$M$ is a positive integer that divides~$N$.  As explained in
\cite[VIII.1--2]{lang:modular}, for each divisor~$d$ of $N/M$ there is
a natural \defn{degeneracy map} $\alpha_{M,d} : S_2(\Gamma_0(M))\ra
S_2(\Gamma_0(N))$ given by $\alpha_{M,d}(f(q)) = f(q^d)$.  The
\defn{new subspace} of $S_2(\Gamma_0(N))$, denoted
$S_2(\Gamma_0(N))_{\new}$, is the complementary $\T$-submodule of the
$\T$-module generated by the images of all maps $\alpha_{M,d}$,
with~$M$ and~$d$ as above.  It is a nontrivial fact that this
complement is well defined; one possible proof uses the Petersson
inner product (see \cite[\S{}VII.5]{lang:modular}).\index{Petersson inner
product}

The theory of Atkin and Lehner \cite{atkin-lehner} (see
Theorem~\ref{thm:ALL} below) asserts that, as a $\T^{(N)}$-module,
$S_2(\Gamma_0(N))$ decomposes as follows:
$$S_2(\Gamma_0(N)) \quad = 
\bigoplus_{M | N, \,\, d | N/M}
\beta_{M,d}(S_2(\Gamma_0(M))_{\new}).$$ 
To compute $S_2(\Gamma_0(N))$ it suffices to compute
$S_2(\Gamma_0(M))_{\new}$ for each $M\mid N$.

We now turn to the problem of computing $S_2(\Gamma_0(N))_{\new}$.
Atkin and Lehner \cite{atkin-lehner} proved that
$S_2(\Gamma_0(N))_{\new}$ is spanned by eigenforms for all $T_p$ with
$p\nmid N$ and that the common eigenspaces of all the $T_p$ with
$p\nmid N$ each have dimension~$1$.  Moreover, if $f\in
S_2(\Gamma_0(N))_{\new}$ is an eigenform then the coefficient of~$q$
in the $q$-expansion of~$f$ is nonzero, so it is possible to
normalize~$f$ so the coefficient of~$q$ is~$1$ (such a {\em
  normalized} eigenform in the new subspace is called a
\defn{newform}).  With~$f$ so normalized, if $T_p(f) = a_p f$, then
the $p$th Fourier coefficient of~$f$ is~$a_p$.  If
$f=\sum_{n=1}^{\infty} a_n q^n$ is a normalized eigenvector for all
$T_p$, then the $a_n$, with~$n$ composite, are determined by the
$a_p$, with~$p$ prime, by the following formulas: $a_{nm}=a_n a_m$
when $n$ and $m$ are relatively prime and $a_{p^r} = a_{p^{r-1}} a_p
- p a_{p^{r-2}}$ for $p\nmid N$ prime.  When $p \mid N$,
$a_{p^r}=a_p^r$.  We conclude that in order to compute
$S_2(\Gamma_0(N))_{\new}$, it suffices to compute all systems of
eigenvalues $\{a_2, a_3, a_5, \ldots\}$ of the prime-indexed Hecke
operators $T_2, T_3, T_5, \ldots$ acting on $S_2(\Gamma_0(N))_{\new}$.
Given a system of eigenvalues, the corresponding eigenform is
$f=\sum_{n=1}^{\infty} a_n q^n$, where the $a_n$, for~$n$ composite,
are determined by the recurrence given above.

In light of the pairing $\langle\, ,\, \rangle$ introduced in 
Section~\ref{sec:modformhecke}, 
computing the above systems of eigenvalues $\{a_2,a_3, a_5, \ldots \}$
amounts to computing the systems of eigenvalues of the Hecke operators
$T_p$ on the subspace~$V$ of $\sS_2(\Gamma_0(N))$ that corresponds
to the new subspace of $S_2(\Gamma_0(N)).$
For each proper divisor~$M$ of~$N$ and each divisor~$d$ of
$N/M$, let $\phi_{M,d}:\sS_2(\Gamma_0(N)) \ra \sS_2(\Gamma_0(M))$ be the map
sending~$x$ to $\abcd{d}{0}{0}{1}x$. 
Then~$V$ is the intersection of the kernels of all
maps $\phi_{M,d}$.

Computing the systems of eigenvalues of a collection of
commuting diagonalizable endomorphisms is a problem in 
linear algebra (see Chapter~\ref{ch:linalg}).

\begin{example}
All forms in $S_2(\Gamma_0(39))$ are new. 
Up to Galois conjugacy, the eigenvalues of the Hecke operators
$T_2$, $T_3$, $T_5$, and $T_7$ on $\sS_2(\Gamma_0(39))$ are
$\{1, -1, 2, -4 \}$ and $\{a, 1, -2a-2, 2a+2\}$, where
$a^2+2a-1=0$.  Each of these eigenvalues occur
in $\sS_2(\Gamma_0(39))$ with multiplicity two; for example,
the characteristic polynomial of $T_2$ on $\sS_2(\Gamma_0(39))$ is
$(x - 1)^2\cdot(x^2+2x-1)^2$.
Thus $S_2(\Gamma_0(39))$ is spanned by 
\begin{align*}
f_1&=q + q^2 - q^3 - q^4 + 2q^5 - q^6 - 4q^7 + \cdots, \\
f_2&= q + aq^2 + q^3 + (-2a - 1)q^4 + (-2a - 2)q^5 + aq^6 + (2a + 2)q^7 
  + \cdots, \\
f_3&= q + \sigma(a)q^2 + q^3 + (-2\sigma(a) - 1)q^4 + (-2\sigma(a) - 2)q^5 + \sigma(a)q^6 
  + \cdots, 
\end{align*}
where $\sigma(a)$ is the other $\Gal(\Qbar/\Q)$-conjugate of $a$.
\end{example}

\subsection{Summary} 
To compute the $q$-expansion of a basis for $S_2(\Gamma_0(N))$, we use
the degeneracy maps so that we only have to solve the problem for
$S_2(\Gamma_0(M))_{\new}$, for all integers $M\mid N$.  Using modular
symbols, we compute all systems of eigenvalues $\{a_2, a_3,
a_5,\ldots\}$, and then write down the corresponding 
eigenforms $\sum a_n q^n$.

\begin{exercises}
\item \label{ex:x0n} Suppose that $\lambda, \lambda' \in \h$ are in
the same orbit for the action of $\Gamma_0(N)$, i.e., that there
exists $g\in \Gamma_0(N)$ such that $g(\lambda) = \lambda'$.
Let $\Lambda = \Z + \Z\lambda$ and $\Lambda' = \Z + \Z\lambda'$.
Prove that the pairs $(\C/\Lambda, (\frac{1}{N}\Z+\Lambda)/\Lambda)$
and  $(\C/\Lambda', (\frac{1}{N}\Z + \Lambda')/\Lambda')$ are isomorphic. 
(By an isomorphism $(E,C)\to (F,D)$ of pairs, we mean
an isomorphism $\phi: E\to F$ of elliptic curves that sends
$C$ to $D$.  You may use the fact that an isomorphism of
elliptic curves over $\C$ is a $\C$-linear map $\C\to\C$ that 
sends the lattice corresponding to one curve onto the lattice
corresponding to the other.)

\item \label{ex:modsymzero} 
Let $n,m$ be integers and $N$ a positive integer.
Prove that the modular symbol $\{n,m\}$ is $0$ as an element of 
$\sM_2(\Gamma_0(N))$.  [Hint: See Example~\ref{ex:01}.]

\item \label{ex:bij} Let $p$ be a prime. 
\begin{enumerate}
\item List representative elements of $\P^1(\Z/p\Z)$.
\item What is the cardinality of $\P^1(\Z/p\Z)$ as a function of $p$?
\item Prove that there is a bijection
between the right cosets of $\Gamma_0(p)$ in $\SL_2(\Z)$
and the elements of $\P^1(\Z/p\Z)$ that sends
$\abcd{a}{b}{c}{d}$ to $(c:d)$.
(As mentioned in this chapter, the analogous statement
is also true when the level is composite; 
see \cite[\S2.2]{cremona:algs} for complete details.)
\end{enumerate}

\item \label{ex:intermsofmanin}
Use the inductive proof of Proposition~\ref{prop:manintrick} to
  write $\{0,4/7\}$ in terms of Manin symbols for $\Gamma_0(7)$.

\item \label{ex:heckes2g03}
Show that the Hecke operator $T_2$ acts as multiplication
by $3$ on the space $\sM_2(\Gamma_0(3))$  as follows:
\begin{enumerate}
\item Write down right coset representatives for $\Gamma_0(3)$ in $\SL_2(\Z)$.
\item List all eight relations coming from Theorem~\ref{thm:manin_present}.
\item Find a single Manin symbols $[r_i]$ so that the three other
Manin symbols are a nonzero multiple of $[r_i]$ modulo the relations found
in the previous step.
\item Use formula \eqref{eqn:t2merel} to compute $T_2([r_i])$.  You
  will obtain a sum of four symbols.  Using the relations above, write
  this sum as a multiple of $[r_i]$.  (The multiple must be $3$ or you
  made a mistake.)
\end{enumerate}

\end{exercises}





\chapter{Dirichlet Characters}\label{ch:dirichlet}

In this chapter we develop a theory for computing with Dirichlet
characters, which are extremely important to computations with modular
forms for (at least) two reasons:
\begin{enumerate}
\item To compute the Eisenstein subspace $\Eis_k(\Gamma_1(N))$
of $M_k(\Gamma_1(N))$, we write down 
Eisenstein series\index{Eisenstein series!compute} attached to 
pairs of Dirichlet characters (the space 
$\Eis_k(\Gamma_1(N))$ will be defined in Chapter~\ref{ch:eisen}).

\item To compute $S_k(\Gamma_1(N))$, we instead compute
a decomposition
$$
    M_k(\Gamma_1(N)) = \bigoplus M_k(\Gamma_1(N), \eps)
$$
and then compute each factor (see Section~\ref{sec:decompdirchar}).  
Here the sum is over all Dirichlet characters $\eps$ of modulus~$N$.
\end{enumerate}

Dirichlet characters appear frequently in many other areas of
number theory.  For example, by the Kronecker-Weber theorem, Dirichlet
characters correspond to the $1$-dimensional representations of
$\Gal(\Qbar/\Q)$.

After defining Dirichlet characters in Section~\ref{sec:defndir}, in
Section~\ref{sec:repdir} we describe a good way to represent Dirichlet
characters using a computer.  Section~\ref{sec:evaldir} is about how
to evaluate Dirichlet characters and leads naturally to a discussion
of the baby-step giant-step algorithm for solving the discrete log
problem and methods for efficiently computing the Kronecker symbol.
In Section~\ref{sec:conddir} we explain how to factor Dirichlet
characters into their prime power constituents and apply this to the
computations of conductors.   We describe how to carry out a number
of standard operations with Dirichlet characters in 
Section~\ref{sec:dirresext} and discuss alternative ways
to represent them in Section~\ref{sec:alter_rep}. 
%begin_sage
Finally, in Section~\ref{sec:sagedir} we give a very 
short tutorial about how to compute with Dirichlet characters using \sage.
%end_sage


%%begin_sage
%%end_sage


\section{The Definition}\label{sec:defndir}
Fix an integral domain~$R$ and a root~$\zeta$ of unity in~$R$.
\begin{definition}[Dirichlet Character]
  A \defn{Dirichlet character} of modulus~$N$ over~$R$ is a map
  $\eps:\Z\to R$ such that there is a homomorphism $f:(\Z/N\Z)^* \to
  \langle \zeta \rangle$ for which
$$
 \eps(a) = \begin{cases} 
       0 &\text{if }\gcd(a,N)>1,\\
     f\,(a\!\!\!\!\mod{N}) &\text{if }\gcd(a,N)=1.
   \end{cases}
$$
\end{definition}   
We denote the group of such Dirichlet characters by $D(N,R)$.  Note
that elements of $D(N,R)$ are in bijection with homomorphisms
$(\Z/N\Z)^* \to \langle \zeta \rangle$.

A familiar Dirichlet character is the Legendre symbol
$\kr{a}{p}$, with $p$ an odd prime, that appears in quadratic
reciprocity theory. It is a Dirichlet character of modulus~$p$ that
takes the value $1$ on integers that are congruent to a nonzero square
modulo~$p$, the value $-1$ on integers that are congruent to a nonzero
nonsquare modulo~$p$, and~$0$ on integers divisible by $p$.




\section{Representing Dirichlet Characters}\label{sec:repdir}
\begin{lemma}\label{lem:dual}
  The groups $(\Z/N\Z)^*$ and $D(N,\C)$
  are isomorphic.
\end{lemma}
\begin{proof}
  We prove the more general fact that for any finite abelian
  group~$G$, we have that $G\ncisom \Hom(G,\C^*)$.  To deduce this
  latter  isomorphism, first reduce to the case
  when~$G$ is cyclic by writing $G$ as a product of
cyclic groups. The cyclic case follows
  because if $G$ is cyclic of order $n$, then 
$\C^*$ contains an $n$th root of unity,
 so $\Hom(G,\C^*)$ is also cyclic of order~$n$.  
Any two  cyclic groups of the same order are isomorphic,
so $G$ and $\Hom(G,\C^*)$ are isomorphic.
\end{proof}

\begin{corollary}\label{cor:dir_ord}
We have $\#D(N,R) \mid \vphi(N)$, with equality
if and only if the order of our choice of $\zeta\in R$
is a multiple of the exponent of the group $(\Z/N\Z)^*$.
\end{corollary}
\begin{proof}
This is because $\#(\Z/N\Z)^* = \vphi(N)$.
\end{proof}

Fix a positive integer~$N$.  To find the set of ``canonical''
generators for the group $(Z/NZ)^*$, write $N=\prod_{i=0}^{n}
{p_i^{e_i}}$ where $p_0<p_1<\cdots < p_n$ are the prime divisors
of~$N$.  By \exref{ch:dirichlet}{ex:cyclic}, each factor
$(\Z/p_i^{e_i}\Z)^*$ is a cyclic group $C_i=\langle g_i \rangle$,
except if $p_0=2$ and $e_0\geq 3$, in which case $(\Z/p_0^{e_0}\Z)^*$
is a product of the cyclic subgroup $C_0 = \langle -1 \rangle$ of
order $2$ with the cyclic subgroup $C_1 = \langle 5\rangle$.  In all
cases we have 
$$ 
(\Z/N\Z)^* \isom \prod_{0\leq i\leq n} C_i =
\prod_{0\leq i \leq n} \langle g_i \rangle.
$$
For~$i$ such that $p_i>2$, choose the generator $g_i$ of $C_i$ to be
the element of $\{2,3,\ldots, p_i^{e_i}-1\}$ that is smallest and
generates.  Finally, use the Chinese Remainder Theorem (see
\cite[\S1.3.3]{cohen:course_ant}) to lift each $g_i$ to an element in
$(\Z/N\Z)^*$, also denoted $g_i$, that is~$1$ modulo each $p_j^{e_j}$
for $j\neq i$.

\begin{algorithm}{Minimal Generator for $(\Z/p^r\Z)^*$}\label{alg:mingens}
Given a prime power~$p^r$ with $p$ odd, this algorithm computes
the minimal generator of $(\Z/p^r\Z)^*$.
\begin{steps}
\item{}[Factor Group Order] Factor $n=\phi(p^r) = p^{r-1}\cdot 2\cdot
  ((p-1)/2)$ as a product $\prod p_i^{e_i}$ of primes.  This is
  equivalent in difficulty to factoring $(p-1)/2$.  (See, e.g.,
  \cite[Ch.8, Ch.~10]{cohen:course_ant} for an excellent discussion of 
  factorization algorithms, though of course much progress has been
  made since then.)
\item{}[Initialize]\label{step:gen_init} Set $g\set 2$.  
\item{}[Generator?] Using the binary powering algorithm (see
  \cite[\S1.2]{cohen:course_ant}), compute $g^{n/p_i}\pmod{p^r}$, for
  each prime divisor $p_i$ of~$n$.  If any of these powers are $1$,
then $g$ is not a generator, so
  set $g\set g+1$ and go to step~(\ref{step:gen_init}). 
If no powers are $1$, output $g$ and terminate.
\end{steps}
\end{algorithm}
See \exref{ch:dirichlet}{ex:orderalg} for a proof that this
algorithm is correct.

\begin{example}
A minimal generator for $(\Z/49\Z)^*$ is $3$.  
We have $n=\vphi(49)=42=2\cdot 3\cdot 7$ and 
$$
  2^{n/2} \con 1, \qquad 2^{n/3} \con 18, \qquad 2^{n/7} \con 15 \pmod{49},
$$
so $2$ is not a generator for $(\Z/49\Z)^*$.  (We see
this just from $2^{n/2}\con 1\pmod{49}$.)  However $3$ is 
a generator since
$$
  3^{n/2} \con 48, \qquad 3^{n/3} \con 30, \qquad 3^{n/7} \con 43 \pmod{49}.
$$
\end{example}

\begin{example}\label{ex:mingens}
In this example
we compute minimal generators for $N=25$, $100$, and $200$:
\begin{enumerate}
\item
The minimal generator for $(\Z/25\Z)^*$ is $2$.

\item
The minimal generators for $(\Z/100\Z)^*$, lifted
to numbers modulo $100$, are $g_0=51$ and $g_1=77$.
Notice that $g_0\con -1\pmod{4}$ and $g_0\con 1\pmod{25}$
and that $g_1\con 2\pmod{25}$ is the minimal generator modulo~$25$.

\item
The minimal generators for $(\Z/200\Z)^*$, lifted
to numbers modulo $200$, are 
$g_0 = 151$, $g_1 = 101$, and $g_2=177$.  Note
that $g_0\con -1\pmod{4}$, that $g_1\con 5\pmod{8}$
and $g_2\con 2\pmod{25}$.
\end{enumerate}
\end{example}

%begin_sage
\sageindex{$\Z/N\Z$}
In \sage, the command {\tt Integers(N)} creates $\Z/N\Z$. 
\begin{verbatim}
sage: R = Integers(49)
sage: R
Ring of integers modulo 49
\end{verbatim}%link

The {\tt unit\_gens} command computes the minimal generators 
for $(\Z/N\Z)^*$, as defined above.
%link
\begin{verbatim}
sage: R.unit_gens()
[3]
sage: Integers(25).unit_gens()
[2]
sage: Integers(100).unit_gens()
[51, 77]
sage: Integers(200).unit_gens()
[151, 101, 177]
sage: Integers(2005).unit_gens()
[402, 1206]
sage: Integers(200000000).unit_gens()
[174218751, 51562501, 187109377]
\end{verbatim}
%end_sage

Fix an element~$\zeta$ of finite multiplicative order in a ring~$R$,
and let $D(N,R)$ denote the group of Dirichlet characters of modulus~$N$
over~$R$, with image in $\langle \zeta\rangle \union \{0\}$. 
In most of this chapter, we
specify an element $\eps\in D(N,R)$ by giving the list
\begin{equation}\label{eqn:epslist}
[\eps(g_0), \eps(g_1), \ldots, \eps(g_n)]
\end{equation}
of images of the generators of $(\Z/N\Z)^*$.  (Note that if $N$ is
even, the number of elements of the list \eqref{eqn:epslist} {\em
  does} depend on whether or not $8\mid N$---there are two factors
corresponding to $2$ if $8\mid N$, but only one if $8\nmid N$.)  This
representation completely determines~$\eps$ and is convenient for
arithmetic operations.  It is analogous to representing a linear
transformation by a matrix. 
\begin{remark}
In any actual implementation (e.g., the one in \sage), it
is better to represent the $\eps(g_i)$ by recording an integer
$j$ such that $\eps(g_i) = \zeta^j$, where $\zeta \in R$
is a fixed root of unity.  Then \eqref{eqn:epslist} is
internally represented as an element of $(\Z/m\Z)^{n+1}$, where
$m$ is the multiplicative order of $\zeta$.  When
the representation of \eqref{eqn:epslist} is needed for an algorithm,
it can be quickly computed on the fly using a table of the powers
of $\zeta$.
 See Section~\ref{sec:alter_rep} for further discussion about
 ways to represent characters.
\end{remark}

\begin{example}
The group $D(5,\C)$ has elements 
$\{[1], [i], [-1], [-i]\}$,
so it is cyclic of order $\vphi(5)=4$.
In contrast,  the group $D(5,\Q)$ has only the
two elements $[1]$ and $[-1]$ and order $2$.
%begin_sage
\sageindex{Dirichlet group}
The command {\tt DirichletGroup(N)}
with no second argument creates the group of
Dirichlet characters with values in the
cyclotomic field $\Q(\zeta_n)$, where $n$
is the exponent of the group $(\Z/N\Z)^*$.
Every element in $D(N,\C)$ takes values
in $\Q(\zeta_n)$, so $D(N,\Q(\zeta_n)) \ncisom D(N,\C)$.
\begin{verbatim}
sage: list(DirichletGroup(5))
[[1], [zeta4], [-1], [-zeta4]]
sage: list(DirichletGroup(5, QQ))
[[1], [-1]]
\end{verbatim}
%end_sage
\end{example}

\section{Evaluation of Dirichlet Characters}\label{sec:evaldir}
This section is about how to compute $\eps(n)$, where $\eps$
is a Dirichlet character and~$n$ is an integer.
We begin with an example.
\begin{example}\label{ex:char}
If $N=200$, then $g_0=151$, $g_1=101$ and $g_2=177$, as
we saw in Example~\ref{ex:mingens}.  The
exponent of $(\Z/200\Z)^*$ is $20$, since that
is the least common multiple of the exponents of
$4=\#(\Z/8\Z)^*$ and $20=\#(\Z/25\Z)^*$.
The orders of $g_0$, $g_1$, and $g_2$ are $2$, $2$, and $20$.
Let $\zeta=\zeta_{20}$ be a primitive $20$th root of
unity in $\C$.  Then the following are generators
for $D(200,\C)$:
$$
  \eps_0 = [-1,1,1],\qquad
  \eps_1 = [1,-1,1],\qquad
  \eps_2 = [1,1,\zeta],
$$
and $\eps=[1,-1,\zeta^5]$ is an example element of order~$4$.
To evaluate $\eps(3)$, we write~$3$ in
terms of $g_0$, $g_1$, and $g_2$.
First, reducing $3$ modulo~$8$, we see that 
$3 \con g_0 \cdot g_1\pmod{8}$.  Next 
reducing $3$ modulo $25$ and trying powers of 
$g_2=2$, we find that $e\con g_2^7\pmod{25}$.  
Thus 
\begin{align*}
  \eps(3) &= \eps(g_0 \cdot g_1 \cdot g_2^7)\\
          &= \eps(g_0) \eps(g_1) \eps(g_2)^7\\
          &= 1 \cdot (-1) \cdot (\zeta^5)^7\\
          &= -\zeta^{35} = -\zeta^{15}.
\end{align*}

%begin_sage
\sageindex{evaluation of character}
We next illustrate the above computation of $\eps(3)$ in \sage.
First we make the group $D(200,\Q(\zeta_8))$ and list its
generators.
\begin{verbatim}
sage: G = DirichletGroup(200)
sage: G
Group of Dirichlet characters of modulus 200 over 
Cyclotomic Field of order 20 and degree 8
sage: G.exponent()
20
sage: G.gens()
([-1, 1, 1], [1, -1, 1], [1, 1, zeta20])
\end{verbatim}%link

We construct $\eps$.
%link
\begin{verbatim}
sage: K = G.base_ring()
sage: zeta = K.0
sage: eps = G([1,-1,zeta^5])
sage: eps
[1, -1, zeta20^5]
\end{verbatim}%link

Finally, we evaluate $\eps$ at $3$.
%link
\begin{verbatim}
sage: eps(3)
zeta20^5
sage: -zeta^15
zeta20^5
\end{verbatim}

%end_sage

\end{example}


Example~\ref{ex:char} illustrates that if~$\eps$ is represented using
a list as described above, evaluation of~$\eps$ is inefficient without
extra information; it requires solving the discrete log problem in
$(\Z/N\Z)^*$. 

\begin{remark}
For a general character~$\eps$, is calculation
of~$\eps$ at least as hard as finding discrete
logarithms? Quadratic characters
are easier---see Algorithm~\ref{alg:jacobi}.
\end{remark}

\begin{algorithm}{Evaluate $\eps$}\label{alg:eval_eps}
Given a Dirichlet character $\eps$ of modulus~$N$, represented
by a list $[\eps(g_0), \eps(g_1), \ldots, \eps(g_n)]$,
and an integer~$a$, this algorithm computes~$\eps(a)$.
\begin{steps}
\item{}[GCD] Compute $g=\gcd(a,N)$.  If $g>1$, output~$0$ and terminate.
\item{}[Discrete Log]\label{step:eval_fb} For each $i$, write $a\pmod{p_i^{e_i}}$ as a
  power $m_i$ of $g_i$ using some algorithm for solving the discrete
  log problem (see below).  If $p_i=2$, write
  $a\pmod{p_i^{e_i}}$ as $(-1)^{m_0}\cdot 5^{m_1}$.  (This step is
  analogous to writing a vector in terms of a basis.)
\item{}[Multiply] Output $\prod \eps(g_i)^{m_i}$ as an
  element of~$R$, and terminate.  (This is analogous to multiplying a
  matrix times a vector.)
\end{steps}
\end{algorithm}

\subsection{The Discrete Log Problem}
\exref{ch:dirichlet}{ex:dlogadd} gives  an isomorphism
of groups
$$
(1+p^{n-1}(\Z/p^n\Z),\,\times) \isom (\Z/p\Z,\,+),
$$
so one sees by induction that step~(\ref{step:eval_fb}) is ``about as
difficult'' as finding a discrete log in $(\Z/p\Z)^*$.  There is an
algorithm called ``baby-step giant-step'', which solves the discrete
log problem in $(\Z/p\Z)^*$ in time $O(\sqrt{\ell})$, where $\ell$ is
the largest prime factor of $p-1=\#(\Z/p\Z)^*$ (note that the discrete
log problem in $(\Z/p\Z)^*$ reduces to a series of discrete log
problems in each prime-order cyclic factor).  This is unfortunately
still exponential in the number of digits of~$\ell$; it also uses
$O(\sqrt{\ell})$ memory.  We now describe this algorithm without
any specific optimizations.

\begin{algorithm}{Baby-step Giant-step Discrete Log}
\label{alg:baby_giant_dlog}
Given a prime $p$, a generator $g$ of $(\Z/p\Z)^*$,
and an element $a\in (\Z/p\Z)^*$, this algorithm
finds an $n$ such that $g^n=a$.  
(Note that this algorithm works in any cyclic group,
not just $(\Z/p\Z)^*$.)
\begin{steps}
\item{}[Make Lists] 
Let $m=\lceil{} \sqrt{p}\rceil{}$ be the ceiling of $\sqrt{p}$,
and
construct two lists $$
1,\, g^m,\, \ldots, \,g^{(m-1)m}
\qquad\qquad\text{(giant steps)} $$
and 
$$
a, ag,\, ag^2,\, \ldots, \,
ag^{m-1}\qquad\qquad\text{(baby steps)}.
$$
\item{}[Find Match]\label{step:find_match}
Sort the two lists and
find a match $g^{im} = a g^j$. Then $a = g^{im-j}$.
\end{steps}
\end{algorithm}
\begin{proof}
  We prove that there will always be a match. Since we know that $a=g^k$ for some $k$
  with $0\leq k\leq p-1$ and any such $k$ can be written in the form
  $im-j$ for $0\leq i,j\leq m-1$, we will find such a match.
\end{proof}

Algorithm~\ref{alg:baby_giant_dlog} uses nothing special about
$(\Z/p\Z)^*$, so it works in a generic group.  It is a theorem that
there is no faster algorithm to find discrete logs in a ``generic
group'' (see \cite{shoup:lower,nechaev:lower}).  There are
much better subexponential algorithms for solving the discrete log
problem in $(\Z/p\Z)^*$, which use the special structure of this
group.  They use the number field sieve (see, e.g.,\index{number field sieve}
\cite{gordon:dlog}), which is also the best-known algorithm for
factoring integers.  This class of algorithms has been very well
studied by cryptographers; though sub-exponential, solving discrete
log problems when $p$ is large is still extremely difficult.  For a
more in-depth survey see \cite{gordon:dlp}. 
%These algorithms are
%particularly relevant when~$p$ is large, and their development is
%motivated mainly by their application to breaking cryptosystems.
For computing Dirichlet characters in our context,~$p$ is not 
too large, so Algorithm~\ref{alg:baby_giant_dlog} works well.

\subsection{Enumeration of All Values}
For many applications of Dirichlet characters to computing modular
forms, $N$ is fairly small, e.g., $N<10^6$, and we 
evaluate~$\eps$ on a {\em huge} number of random elements, inside
inner loops of algorithms.  Thus for such purposes it will often be
better to make a table of all values of~$\eps$, so that evaluation
of~$\eps$ is extremely fast.  The following algorithm computes a table
of all values of $\eps$, and it does not require computing any
discrete logs since we are computing {\em all} values.

\begin{algorithm}{Values of $\eps$}
Given a Dirichlet character $\eps$ represented by the list
of values of $\eps$ on the minimal generators $g_i$ of $(\Z/N\Z)^*$, this algorithm
creates a list of all the values of $\eps$.
\begin{steps}
\item{}[Initialize] For each minimal generator $g_i$, set $a_i\set 0$.
  Let $n = \prod g_i^{a_i}$, and set $z\set 1$.  Create a list~$v$
  of~$N$ values, all initially set equal to~$0$.  When this algorithm
  terminates, the list~$v$ will have the property that 
  $$v\,\,[x
  \!\!\!\!\!\pmod{N}] = \eps(x).$$
  Notice that we index $v$ starting
  at $0$.
\item{}[Add Value to Table]\label{step:add_value} Set $v[n] \set z$.
\item{}[Finished?] If each $a_i$ is one less than the order of $g_i$, output~$v$ 
and terminate.
\item{}[Increment] Set $a_0\set a_0 + 1$, $n\set n\cdot g_0\pmod{N}$,
and $z\set z\cdot \eps(g_0)$.  If $a_0\geq \ord(g_0)$, set $a_0\to 0$,
and then set $a_1\set a_1 + 1$, $n\set n\cdot g_1\pmod{N}$, and 
$z\set z\cdot \eps(g_1)$.  If $a_1\geq \ord(g_1)$, do what you just
did with $a_0$ but with all subscripts replaced by $1$.  Etc.
(Imagine a car odometer.)  Go to step~(\ref{step:add_value}).
\end{steps}
\end{algorithm}





\section{Conductors of Dirichlet Characters}\label{sec:conddir}

The following algorithm for computing the order of~$\eps$ reduces
the problem to computing the orders of powers of $\zeta$ in $R$.
\begin{algorithm}{Order of Character}\label{alg:dir_order}
  This algorithm computes the order of a Dirichlet
  character $\eps\in D(N,R)$.
\begin{steps}
\item Compute the order $r_i$ of each $\eps(g_i)$, for each minimal
  generator $g_i$ of $(\Z/N\Z)^*$.  The order of $\eps(g_i)$ is
  a divisor of $n=\#(\Z/p_i^{e_i}\Z)^*$ so we can compute its order by
  considering the divisors of~$n$.
\item Compute and output the least common multiple of the integers
  $r_i$.
\end{steps}
\end{algorithm}
\begin{remark}
  Computing the order of $\eps(g_i) \in R$ is potentially difficult.
  Simultaneously using a different representation of Dirichlet
  characters avoids having to compute the order of elements of $R$
  (see Section~\ref{sec:alter_rep}).
\end{remark}


The next algorithm factors $\eps$ as a product of ``local'' characters,
one for each prime divisor of $N$.  It is useful for other algorithms,
e.g., for explicit computations with trace formulas
(see \cite{hijikata:trace}). This factorization is easy to compute
because of how we represent $\eps$.

\begin{algorithm}{Factorization of Character}\label{alg:dirfac}
Given a Dirichlet character $\eps\in D(N,R)$, with $N=\prod p_i^{e_i}$, 
this algorithm finds Dirichlet characters $\eps_i$ modulo $p_i^{e_i}$, such
that for all $a\in(\Z/N\Z)^*$, we have 
$\eps(a) = \prod \eps_i(a\!\!\pmod{p_i^{e_i}})$.
If $2\mid N$, the steps are as follows:
\begin{steps}
\item Let $g_i$ be the minimal generators of $(\Z/N\Z)^*$, so $\eps$
is given by a list $$[\eps(g_0),\ldots, \eps(g_n)].$$
\item\label{step:singletons} For $i=2,\ldots, n$, let $\eps_i$ be the element of $D(p_i^{e_i},R)$
defined by the singleton list $[\eps(g_i)]$.
\item\label{step:extra2} Let $\eps_1$ be the element of $D(2^{e_1},R)$ defined
by the list $[\eps(g_0), \eps(g_1)]$ of length~$2$.   Output the
$\eps_i$ and terminate.
\end{steps}
If $2\nmid N$, then omit step~(\ref{step:extra2}), and include
all $i$ in step~(\ref{step:singletons}).
\end{algorithm}
The factorization of Algorithm~\ref{alg:dirfac} is unique since each
$\eps_i$ is determined by the image of the canonical map
$(\Z/p_i^{e_i}\Z)^*$ in $(\Z/N\Z)^*$, which sends $a\pmod
{p_i^{e_i}}$ to the element of $(\Z/N\Z)^*$ that is
$a\pmod{p_i^{e_i}}$ and $1 \pmod{p_j^{e_j}}$ for $j\neq i$.

\begin{example}\label{ex:prodlocal}
If $\eps = [1,-1,\zeta^5] \in D(200,\C)$, then 
$\eps_1 = [1,-1]\in D(8,\C)$ and $\eps_2 = [\zeta^5]\in D(25,\C)$.
\end{example}



\begin{definition}[Conductor]\label{defn:conductordir}
The \defn{conductor} of a Dirichlet character $\eps\in D(N,R)$
is the smallest positive divisor $c\mid N$ such that 
there is a character $\eps' \in D(c,R)$ for which 
$\eps(a) = \eps'(a)$ for all $a\in\Z$ with $(a,N)=1$.
A Dirichlet character is \defn{primitive} if its modulus equals
its conductor.  The character $\eps'$ associated to~$\eps$
with modulus equal to the conductor of~$\eps$ is called
the \defn{primitive character associated to}~$\eps$.
\end{definition}
We will be interested in conductors later, when computing new
subspaces of spaces of modular forms with character.  Also certain
formulas for special values of~$L$ functions are only valid for
primitive characters.

\begin{algorithm}{Conductor}\label{alg:conductor}
  This algorithm computes the conductor of a Dirichlet
  character~$\eps \in D(N,R)$.
\begin{steps}
\item{}[Factor Character]\label{step:factor_dir} Using Algorithm~\ref{alg:dirfac}, find
  characters~$\eps_i$ whose product is~$\eps$.
\item{}[Compute Orders] Using Algorithm~\ref{alg:dir_order}, compute
  the orders~$r_i$ of each~$\eps_i$.
\item{}[Conductors of Factors]\label{step:cond_fac}
For each $i$, either set $c_i\to 1$ if $\eps_i$
is the trivial character (i.e., of order $1$) or set $c_i \set
  p_i^{\ord_{p_i}(r_i)+1}$, where $\ord_{p}(n)$ is the largest power
  of $p$ that divides~$n$.
\item{}[Adjust at $2$?] If $p_1=2$ and $\eps_1(5)\neq 1$, set
$c_1\set 2c_1$.  
\item{}[Finished] Output $c = \prod c_i$ and terminate.
\end{steps}
\end{algorithm}
\begin{proof}
Let $\eps_i$ be the local factors of~$\eps$, as in step~(\ref{step:factor_dir}).
We first show that the product of the conductors $f_i$ of the $\eps_i$ is
the conductor~$f$ of~$\eps$.  Since $\eps_i$ factors through $(\Z/f_i\Z)^*$,
the product~$\eps$ of the $\eps_i$ factors through $(\Z/\prod f_i\Z)^*$,
so the conductor of~$\eps$ divides $\prod f_i$.  Conversely, if 
$\ord_{p_i}(f) < \ord_{p_i}(f_i)$ for some~$i$, then we could factor $\eps$
as a product of local (prime power) characters differently, which contradicts
that this factorization is unique.

It remains to prove that if~$\eps$ is a nontrivial character of modulus
$p^n$, where~$p$ is a prime, and if~$r$ is the order of~$\eps$, then the
conductor of~$\eps$ is $p^{\ord_p(r)+1}$, except possibly if $8\mid
p^n$.  Since the order and conductor of $\eps$ and of the associated
primitive character $\eps'$ are the same, we may assume $\eps$ is
primitive, i.e., that $p^n$ is the conductor of~$\eps$; note that 
$n>0$, since~$\eps$ is nontrivial.  


First suppose $p$ is odd.
Then the abelian group $D(p^n,R)$ splits as a direct sum $D(p,R)\oplus D(p^n,R)'$,
where $D(p^n,R)'$ is the $p$-power torsion subgroup of $D(p^n,R)$.
Also $\eps$ has order $u\cdot p^m$, where $u$, which
is coprime to~$p$, is the order of the image of $\eps$
in $D(p,R)$ and $p^m$ is the order of the image in $D(p^n,R)'$.
If $m=0$, then the order of $\eps$ is coprime to $p$, so 
$\eps$ is in $D(p,R)$, which means that $n=1$,
so $n=m+1$, as required.  If $m>0$, then $\zeta\in R$ 
must have order divisible by $p$, so~$R$ has characteristic not
equal to~$p$.  The conductor of $\eps$ does not change if we
adjoin roots of unity to~$R$, so in light of Lemma~\ref{lem:dual}
we may assume that $D(N,R) \ncisom (\Z/N\Z)^*$.
It follows that for each $n'\leq n$, the $p$-power subgroup
$D(p^{n'},R)'$ of $D(p^{n'},R)$ is the $p^{n'-1}$-torsion
subgroup of $D(p^n,R)'$.   Thus $m=n-1$, since
$D(p^n,R)'$ is by assumption the smallest such group that 
contains the projection of~$\eps$.  This proves the 
formula of step~(\ref{step:cond_fac}).   We leave the argument
when $p=2$ as an exercise (see \exref{ch:dirichlet}{ex:cond2}).
\end{proof}

\begin{example}\label{ex:char_cond}
If $\eps = [1,-1,\zeta^5] \in D(200,\C)$, then
as in Example~\ref{ex:prodlocal},~$\eps$ 
is the product of $\eps_1=[1,-1]$ and $\eps_2 = [\zeta^5]$.
Because $\eps_1(5)=-1$, the conductor of $\eps_1$ is $8$.
The order of $\eps_2$ is $4$ (since $\zeta$ is a $20$th
root of unity), so the conductor of $\eps_2$ is $5$.
Thus the conductor of~$\eps$ is $40=8\cdot 5$.
\end{example}

\section{The Kronecker Symbol}
In this section all characters have values in $\C$.

Frequently quadratic characters are described in terms of the
Kronecker symbol $\kr{a}{n}$, which we define for any integer~$a$
and positive integer~$n$ as follows.
First, if $n=p$ is an odd prime, then for any integer~$a$,
$$
\kr{a}{p}
= \begin{cases} \hfill 0 & \text{if $\gcd(a,p)\neq 1$,}\\
  \hfill  1 & \text{if $a$ is a square mod $p$,}\\
  -1 & \text{if $a$ is not a square mod $p$.}
\end{cases}
$$
If $p=2$, then
$$
\kr{a}{2} = 
  \begin{cases} \hfill 0 & \text{if $a$ is even,}\\
  \hfill 1 & \text{if $a \con \pm 1\pmod{8}$,}\\
 -1 & \text{if $a \con \pm 3\pmod{8}$.}
\end{cases}
$$
More generally, if $n=\prod p_i^{e_i}$ with the $p_i$ prime,
then 
$$
  \kr{a}{n} = \prod \kr{a}{p_i}^{e_i}.
$$
\begin{remark}
  One can also extend $\kr{a}{n}$ to $n<0$, but we will not need this.
  The extension is to set $\kr{a}{-1} = -1$ and $\kr{a}{1} = 1$, for
  $a\neq 0$, and to extend multiplicatively (in the denominator).
  Note that the map $\kr{\bullet}{-1}$ is not a Dirichlet
character (see \exref{ch:dirichlet}{ex:notdir}).
\end{remark}

Let $M$ be the product of the primes $p$ such
that $\ord_p(n)$ is odd.
If $M$ is odd, let $N=M$; otherwise, let $N=8M$.
\begin{lemma}\label{lem:jacdir}
The function 
$$
  \eps(a) = \begin{cases}
         \kr{a}{n} & \text{if $\gcd(a,N)=1$,}\\
         \hfill 0\hfill          & \text{otherwise}
       \end{cases}
$$ 
is a Dirichlet character of modulus~$N$.
The function 
$$
  \eps(a) = \begin{cases}
     \kr{-1}{a} &\text{if $a$ is odd,}\\
    \hfill 0 \hfill & \text{if $a$ is even}
   \end{cases}
$$
is a Dirichlet character of modulus~$N$.
\end{lemma}
\begin{proof}
  When restricted to $(\Z/N\Z)^*$, each map $\kr{\bullet}{p}$, for
  $p$ prime, is a homomorphism, so $\eps$ a product of homomorphisms.
  The second statement follows from the definition and the
  fact that $-1$ is a
  square modulo an odd prime $p$ if and only if $p\con 1\pmod{4}$.
\end{proof}

This section is about going between representing quadratic
characters as row matrices and via Kronecker symbols.  This is
valuable because the algorithms in
\cite[\S1.1.4]{cohen:course_ant} for computing Kronecker symbols
run in time quadratic in the number of digits of the input.
They do
not require computing discrete logarithms; instead, 
they use, e.g., that
$\kr{a}{p} \con a^{(p-1)/2}\pmod{p}$, when~$p$ is an odd prime.

\begin{algorithm}{Kronecker Symbol as Dirichlet Character}\label{alg:jacobi}
  Given $n>0$, this algorithm computes a representation of
  the Kronecker symbol $\kr{\bullet}{n}$ as a Dirichlet character.
\begin{steps}
\item{}[Modulus] Compute $N$ as in Lemma~\ref{lem:jacdir}.
\item{}[Minimal Generators] Compute minimal generators $g_i$ of
  $(\Z/N\Z)^*$ using Algorithm~\ref{alg:mingens}.
\item{}[Images] Compute $\kr{g_i}{N}$ for each $g_i$ using one
of the algorithms of \cite[\S1.1.4]{cohen:course_ant}.
\end{steps}
\end{algorithm}

\begin{example}
We compute the Dirichlet character associated to $\kr{\bullet}{200}$.
Using \sage, we compute the $\kr{g_i}{200}$, for $i=0,1,2$, where the
$g_i$ are as in Example~\ref{ex:char}:
\begin{verbatim}
sage: kronecker(151,200)
1
sage: kronecker(101,200)
-1
sage: kronecker(177,200)
1
\end{verbatim}
Thus the corresponding character is defined by $[1,-1,1]$.
\end{example}




\begin{example}
We compute the character associated to $\kr{\bullet}{420}$.
We have $420=4\cdot 3\cdot 5\cdot 7$, and minimal generators
are 
$$
  g_0 = 211,\quad g_1=1, \quad g_2 = 281, \quad g_3=337, \quad g_4=241.
$$
We have $g_0\con -1\pmod{4}$, $g_2\con 2\pmod{3}$, $g_3\con 2\pmod{5}$
and $g_4\con 3\pmod{7}$.
We find $\kr{g_0}{420}=\kr{g_1}{420}=1$
and $\kr{g_2}{420}=\kr{g_3}{420}=\kr{g_4}{420}=-1$.
The corresponding character is $[1,1,-1,-1,-1]$.
\end{example}

Using the following algorithm, we can go in the other direction, i.e.,
write any quadratic Dirichlet character as a Kronecker symbol.

\begin{algorithm}{Dirichlet Character as Kronecker Symbol}
  Given~$\eps$ of order $2$ with modulus~$N$, this
  algorithm writes~$\eps$ as a Kronecker symbol.
\begin{steps}
\item{}[Conductor] 
Use Algorithm~\ref{alg:conductor} to compute the conductor~$f$ of~$\eps$.
\item{}[Odd] If $f$ is odd, output $\kr{\bullet}{f}$.
\item{}[Even] If $\eps(-1)=1$, output $\kr{\bullet}{f}$;
if $\eps(-1)=-1$, output 
$\kr{\bullet}{f} \cdot \kr{-1}{\bullet}$.
\end{steps}
\end{algorithm}
\begin{proof}
Since~$f$ is the conductor of a quadratic Dirichlet character,
it is a square-free product $g$ of odd primes times either
$4$ or $8$, so the
group $(\Z/f\Z)^*$ does not inject into $(\Z/g\Z)^*$ for
any proper divisor~$g$ of~$f$ (see this by 
reducing to the prime power case).  
Since $g$ is odd and square-free, the character $\kr{\bullet }{g}$ 
has conductor~$g$.  
For each odd prime $p$, 
by step~(\ref{step:cond_fac}) of Algorithm~\ref{alg:conductor}
the factor at~$p$ of both $\eps$ and $\kr{\bullet}{g}$ 
is a quadratic character with modulus~$p$.  
By \exref{ch:dirichlet}{ex:cyclic} and Lemma~\ref{lem:dual} 
the group $D(p,\C)$ is cyclic, so it has a unique element
of order~$2$, so 
the factors of~$\eps$ and $\kr{\bullet}{g}$ at~$p$ are equal.

The quadratic characters with conductor a power of~$2$ are $[-1]$,
$[1,-1]$, and $[-1,-1]$.  The character $[1,-1]$ is $\kr{\bullet}{2}$
and the character $[-1]$ is $\kr{-1}{\bullet}$.
\end{proof}

\begin{example}
  Consider $\eps=[-1,-1,-1,-1,-1]$ with modulus
  $840=8\cdot 3\cdot 5\cdot 7$.  It has conductor $840$, and
  $\eps(-1)=-1$, so for all $a$ with $\gcd(a,840)=1$, we have $\eps(a)
  = \kr{a}{840}\cdot \kr{-1}{a}$.
\end{example}


\section{Restriction, Extension, and Galois Orbits}\label{sec:dirresext}
The following two algorithms restrict and extend characters to a
compatible modulus.  Using them, it is easy to define multiplication of
two characters $\eps\in D(N,R)$ and $\eps'\in D(N',R')$, as long as
$R$ and $R'$ are subrings of a common ring.  To carry out the
multiplication, extend both characters to a common base ring,
and then extend them to characters modulo $\lcm(N,N')$ and multiply.

\begin{algorithm}{Restriction of Character}\label{alg:restrict}
Given a Dirichlet character $\eps\in D(N,R)$ and a divisor
$N'$ of $N$ that is a multiple of the conductor of $\eps$,
this algorithm finds a characters $\eps' \in D(N',R)$, such
that $\eps'(a) = \eps(a)$, for all $a\in\Z$ with $(a,N)=1$.
\begin{steps}
\item{}[Conductor] Compute the conductor of~$\eps$ using
  Algorithm~\ref{alg:conductor}, and verify that $N'$ is
  divisible by the conductor and divides $N$.
\item{}[Minimal Generators] Compute minimal generators $g_i$ for
  $(\Z/N'\Z)^*$.
\item{}[Values of Restriction]\label{step:addmod} For each $i$,
  compute $\eps'(g_i)$ as follows.  Find a multiple $aN'$ of $N'$ such
  that $(g_i+aN',\,N)=1$; then $\eps'(g_i) = \eps(g_i+aN')$.
\item{}[Output Character] Output the Dirichlet character of modulus~$N'$
  defined by $[\eps'(g_0),\ldots, \eps'(g_n)]$.
\end{steps}
\end{algorithm}
\begin{proof}
The only part that is not clear is that in step~(\ref{step:addmod})
there is an $a$ such that $(g_i+aN', N)=1$.  If
we write $N=N_1\cdot N_2$, with $(N_1,N_2)=1$ and $N_1$ divisible
by all primes that divide $N'$, then $(g_i,N_1)=1$ since
$(g_i,N')=1$.  By the Chinese Remainder Theorem, 
there is an $x\in\Z$ such that $x\con g_i\pmod{N_1}$ and
$x\con 1\pmod{N_2}$. Then $x = g_i + b N_1 = g_i + (bN_1/N')\cdot N'$
and $(x,N)=1$, which completes the proof.
\end{proof}


\begin{algorithm}{Extension of Character}\label{alg:extend}
Given a Dirichlet character $\eps\in D(N,R)$ and a multiple
$N'$ of $N$,
this algorithm finds a character $\eps' \in D(N',R)$, such
that $\eps'(a) = \eps(a)$, for all $a\in\Z$ with $(a,N')=1$.
\begin{steps}
\item{}[Minimal Generators] Compute  minimal generators $g_i$ for
  $(\Z/N'\Z)^*$.
\item{}[Evaluate] Compute $\eps(g_i)$ for each~$i$.  Since $(g_i,N')=1$,
we also have $(g_i,N)=1$.
\item{}[Output Character] Output the character
$[\eps(g_0),\ldots, \eps(g_n)]$.
\end{steps}
\end{algorithm}


Let $F$ be the prime subfield of $R$, and assume
that $R\subset \overline{F}$, where $\overline{F}$
is a separable closure of~$F$.  If $\sigma \in \Gal(\overline{F}/F)$
and $\eps \in D(N,R)$, let $(\sigma\eps)(n) = \sigma(\eps(n))$;
this defines an action of $\Gal(\overline{F}/F)$ on $D(N,R)$.
Our next algorithm computes the orbits for
the action of $\Gal(\overline{F}/F)$ on $D(N,R)$.  This
algorithm can provide huge savings for modular
forms computations because the spaces $M_k(N,\eps)$
and $M_k(N,\eps')$ are canonically isomorphic if $\eps$
and $\eps'$ are conjugate.
\begin{algorithm}{Galois Orbit}
Given a Dirichlet character $\eps\in D(N,R)$, this algorithm
computes the orbit of~$\eps$ under the action of $G=\Gal(\overline{F}/F)$,
where~$F$ is the prime subfield of $\Frac(R)$, so $F=\Fp$ or $\Q$.
\begin{steps}
\item{}[Order of $\zeta$] Let $n$ be the order of the chosen root $\zeta\in R$.
\item{}[Nontrivial Automorphisms]\label{step:nontriv_aut}
If $\Char(R)=0$, let
$$
 A = \{a : 2\leq a <n \text{ and }(a,n) = 1\}.
$$
If $\Char(R)=p>0$, compute the multiplicative order~$r$ of~$p\!\!\pmod{n}$,
and let
$$
  A = \{p^m : 1\leq m < r\}.
$$
\item{}[Compute Orbit] Compute and output the {\em set} of unique elements $\eps^a$ 
for each $a\in A$ (there could be repeats, so we output unique elements only).
\end{steps}
\end{algorithm}
\begin{proof}
  We prove that the nontrivial automorphisms of $\langle \zeta\rangle$
  in characteristic~$p$ are as in step~(\ref{step:nontriv_aut}).  It is
  well known that every automorphism in characteristic~$p$ on
  $\zeta\in\Fpbar$ is of the form $x\mapsto x^{p^s}$, for some~$s$.
  The images of $\zeta$ under such automorphisms are
$$
  \zeta, \zeta^p, \zeta^{p^2}, \ldots.
$$
Suppose $r>0$ is minimal such that $\zeta=\zeta^{p^r}$.  Then the
orbit of $\zeta$ is $\zeta,\ldots, \zeta^{p^{r-1}}$.
Also $p^r\con 1\pmod{n}$, where $n$ is the multiplicative order of~$\zeta$,
so~$r$ is the multiplicative order of~$p$ modulo~$n$, which completes
the proof.
\end{proof}

\begin{example}
The Galois orbits of characters in $D(20,\C^*)$ are as follows:
\begin{align*}
  G_0 &= \{ [1,1,1]\},\\
  G_1 &= \{[-1,1,1]\}, \\
  G_2 &= \{[1,1,\zeta_4],\,\, [1,1,-\zeta_4]\}\\
  G_3 &= \{[-1,1,\zeta_4],\,\, [-1,1,-\zeta_4]\}\\  
  G_4 &= \{[1,1,-1]\}, \\
  G_5 &= \{[-1,1,-1]\}.
\end{align*}
The conductors of the characters in orbit $G_0$ are $1$, in
orbit $G_1$ they are $4$, in orbit $G_2$ they are $5$, in
$G_3$ they are $20$, in $G_4$ the conductor is $5$, 
and in $G_5$ the conductor is $20$. (You should verify this.)

\sage{} computes Galois orbits as follows:
\begin{verbatim}
sage: G = DirichletGroup(20)
sage: G.galois_orbits()
[
[[1, 1]],
[[1, zeta4], [1, -zeta4]],
[[1, -1]],
[[-1, 1]],
[[-1, zeta4], [-1, -zeta4]],
[[-1, -1]]
]
\end{verbatim}
\end{example}

\section{Alternative Representations of Characters}\label{sec:alter_rep}
Let~$N$ be a positive integer and~$R$ an integral domain, with fixed
root of unity~$\zeta$ of 
order~$n$, and let $D(N,R)=D(N,R,\zeta)$.  As in
the rest of this chapter, write $N=\prod p_i^{e_i}$, and let
$C_i=\langle g_i\rangle $ be the corresponding cyclic factors of
$(\Z/N\Z)^*$.  In this section we discuss other ways to represent
elements $\eps \in D(N,R)$.  Each representation has advantages and
disadvantages, and no single representation is best. 
%It emerged while writing this chapter that simultaneously using
%more than one representation of elements of $D(N,R)$ would be best.
It is easy to convert between them, and some algorithms are much
easier using one representation than when using another.  
In this section we present two other representations, each 
having advantages and disadvantages.  There
is no reason to restrict to only one representation; for example,
\sage{} internally uses both.


We could represent $\eps$ by giving a list $[b_0,\ldots, b_r]$, where
each $b_i\in\Z/n\Z$ and $\eps(g_i) = \zeta^{b_i}$.  Then arithmetic in
$D(N,R)$ is arithmetic in $(\Z/n\Z)^{r+1}$, which is very efficient.
A drawback to this approach (in practice) is that it is easy to
accidentally consider sequences that do not actually correspond to
elements of $D(N,R)$.  Also the choice of $\zeta$ is less clear, which
can cause confusion.  Finally, the orders of the local factors is more
opaque, e.g., compare $[-1,\zeta_{40}]$ with $[20,1]$.  Overall this
representation is not too bad and is more like representing a linear
transformation by a matrix.  It has the {\em advantage} over the
representation discussed earlier in this chapter that arithmetic in
$D(N,R)$ is very efficient and does not require any operations in the
ring $R$.

Another way to represent~$\eps$ would be to give a list $[b_0,\ldots,
b_r]$ of integers, but this time with $b_i\in \Z/\gcd(s_i,n)\Z$, where
$s_i$ is the order of $g_i$.  Then 
$$\eps(g_i) = \zeta^{b_i \cdot
  n/(\gcd(s_i,n))},$$
which is already pretty complicated.  With this representation we set up an identification
$$
D(N,R)\isom \bigoplus_i \Z/\gcd(s_i,n)\Z, $$
and arithmetic is
efficient.  This approach is seductive because every sequence of
integers determines a character, and the sizes of the integers in the
sequence nicely indicate the local orders of the character.  However,
giving analogues of many of the algorithms discussed in this chapter
that operate on characters represented this way is tricky.  For
example, the representation depends very much on the order of $\zeta$,
so it is difficult to correctly compute natural maps $D(N,R) \to
D(N,S)$, for $R\subset S$ rings. 



%begin_sage
\section{Dirichlet Characters in \sage}\label{sec:sagedir}
\sageindex{Dirichlet character tutorial}
To create a Dirichlet character in \sage{}, first create
the group $D(N,R)$ of Dirichlet characters then construct
elements of that group.  First we make $D(11,\Q)$:
\begin{verbatim}
sage: G = DirichletGroup(11, QQ); G
Group of Dirichlet characters of modulus 11 over 
Rational Field
\end{verbatim}%link

A Dirichlet character prints as a matrix that gives
the values of the character on canonical generators
of $(\Z/N\Z)^*$ (as discussed below).
%link
\begin{verbatim}
sage: list(G)
[[1], [-1]]
sage: eps = G.0      # 0th generator for Dirichlet group
sage: eps
[-1]
\end{verbatim}%link

The character $\varepsilon$ takes the value $-1$ on the
unit generator.
%link
\begin{verbatim}
sage: G.unit_gens()  
[2]
sage: eps(2)
-1
sage: eps(3)
1
\end{verbatim}%link

It is $0$ on any integer not coprime to $11$:
%link
\begin{verbatim}
sage: [eps(11*n) for n in range(10)]  
[0, 0, 0, 0, 0, 0, 0, 0, 0, 0]
\end{verbatim}

We can also create groups of Dirichlet characters 
taking values in other rings or fields. 
For example, we create the cyclotomic field $\Q(\zeta_4)$.
\begin{verbatim}
sage: R = CyclotomicField(4)
sage: CyclotomicField(4)
Cyclotomic Field of order 4 and degree 2
\end{verbatim}%link

Then we define $G = D(15,\Q({\zeta_4}))$.
%link
\begin{verbatim}
sage: G = DirichletGroup(15, R)
sage: G
Group of Dirichlet characters of modulus 15 over 
Cyclotomic Field of order 4 and degree 2
\end{verbatim}%link

Next we list each of its elements. 
%link
\begin{verbatim}
sage: list(G)
[[1, 1], [-1, 1], [1, zeta4], [-1, zeta4], [1, -1], 
[-1, -1], [1, -zeta4], [-1, -zeta4]]
\end{verbatim}%link

Now we evaluate the second generator of $G$ on various
integers:
%link
\begin{verbatim}
sage: e = G.1
sage: e(4)
-1
sage: e(-1)
-1
sage: e(5)
0
\end{verbatim}%link

Finally we list all the values of $e$.
%link
\begin{verbatim}
sage: [e(n) for n in range(15)]
[0, 1, zeta4, 0, -1, 0, 0, zeta4, -zeta4, 
    0, 0, 1, 0, -zeta4, -1]
\end{verbatim}

We can also compute with groups of Dirichlet characters
with values in a finite field.
\begin{verbatim}
sage: G = DirichletGroup(15, GF(5)); G
Group of Dirichlet characters of modulus 15 
      over Finite Field of size 5
\end{verbatim}%link

We list all the elements of $G$, again represented by lists
that give the images of each unit generator, as an element of $\F_5$.
%link
\begin{verbatim}
sage: list(G)
  [[1, 1], [4, 1], [1, 2], [4, 2], [1, 4], [4, 4], 
   [1, 3], [4, 3]]
\end{verbatim}%link

We evaluate the second generator of $G$ on several integers. 
%link
\begin{verbatim}
sage: e = G.1
sage: e(-1)
4
sage: e(2)
2
sage: e(5)
0
sage: print [e(n) for n in range(15)]
[0, 1, 2, 0, 4, 0, 0, 2, 3, 0, 0, 1, 0, 3, 4]
\end{verbatim}

%end_sage



\begin{exercises}
\item \label{ex:notdir}
Let $f:\Z\to\C$ be the map given by 
$$
  f(a) = \begin{cases} \hfill 0 & \text{if }a=0,\\
        -1 & \text{if } a <0,\\
        \hfill 1  &\text{if } a > 0.
      \end{cases}
$$      
Prove that $f$ is not a Dirichlet character
of any modulus $N$.
\item\label{ex:cyclic}
This exercise is about the structure of the units of $\Z/p^n\Z$.
\begin{enumerate}
\item If~$p$ is odd and~$n$ is a positive integer,
prove that $(\Z/p^n\Z)^*$  is cyclic.
\item For $n\geq 3$, prove that $(\Z/2^n\Z)^*$ is a direct
sum of the cyclic subgroups $\langle -1 \rangle$
and $\langle 5 \rangle$, of orders 2 and $2^{n-2}$, 
respectively.
\end{enumerate}

\item \label{ex:orderalg}
Prove that Algorithm~\ref{alg:mingens} works, i.e., that 
if $g\in(\Z/p^r\Z)^*$ and $g^{n/p_i}\neq 1$ for all $p_i\mid n=\vphi(p^r)$,
then $g$ is a generator of $(\Z/p^r\Z)^*$.

\item \label{ex:dlogadd}
\begin{enumerate}
\item
Let $p$ be an odd prime and $n\geq 2$ an integer, and prove
that
$$
  \bigl((1+p^{n-1}\Z/p^n\Z),\,\times\bigr) \isom (\Z/p\Z,\,+).
$$
\item
Use the first part to show that solving the discrete log problem
in $(\Z/p^n\Z)^*$ is ``not much harder'' than
solving the discrete log problem in $(\Z/p\Z)^*$.
\end{enumerate}

\item \label{ex:cond2}
Suppose $\eps$ is a nontrivial Dirichlet character of modulus $2^n$
of order~$r$ over the complex numbers $\C$.
Prove that the conductor of $\eps$ is
$$
  c = \begin{cases} 
       2^{\ord_2(r) + 1} & \text{if $\eps(5)=1$},\\
       2^{\ord_2(r) + 2} & \text{if $\eps(5)\neq 1$.}
     \end{cases}
$$

\item\label{ex:charmod5}
\begin{enumerate}
\item Find an irreducible quadratic polynomial~$f$ over $\F_5$.
\item Then $\F_{25} = \F_5[x]/(f)$.  Find an element with multiplicative
order $4$ in $\F_{25}$.
\item Make a list of all Dirichlet characters in $D(25,\F_{25},\zeta)$.
\item Divide these characters into orbits for the action
of $\Gal(\Fbar_5/\F_5)$.
\end{enumerate}


\end{exercises}






















\chapter{Eisenstein Series and Bernoulli Numbers}\label{ch:eisen}
We introduce generalized Bernoulli numbers attached to Dirichlet
characters and give an algorithm to enumerate the Eisenstein series
in $M_k(N,\eps)$.\index{Eisenstein series!and Bernoulli numbers}

\section{The Eisenstein Subspace}\index{Eisenstein subspace}

Let $M_k(\Gamma_1(N))$ be the space of modular forms of weight
$k$ for $\Gamma_1(N)$, and let $\T$ be the \defn{Hecke algebra}
acting on $M_k(\Gamma_1(N))$, which is the subring of 
$\End(M_k(\Gamma_1(N)))$ generated by all Hecke operators.   
Then there is a $\T$-module decomposition
$$
  M_k(\Gamma_1(N)) = \Eis_k(\Gamma_1(N)) \oplus S_k(\Gamma_1(N)),
$$
where $S_k(\Gamma_1(N))$ is the subspace of modular forms that vanish
at all cusps and $\Eis_k(\Gamma_1(N))$ is the \defn{Eisenstein
  subspace}, which is uniquely determined by this decomposition.  
The above decomposition
induces a decomposition of $M_k(\Gamma_0(N))$ and of
$M_k(N,\eps)$, for any Dirichlet character~$\eps$ of modulus~$N$.


\section{Generalized Bernoulli Numbers}\index{Bernoulli numbers!generalized}
Suppose $\eps$ is a Dirichlet character of modulus~$N$ over~$\C$.
Leopoldt \cite{leopoldt} defined generalized Bernoulli 
numbers attached to $\eps$.
\begin{definition}[Generalized Bernoulli Number]
We define the \defn{generalized Bernoulli numbers} $B_{k,\eps}$
attached to~$\eps$ by the following 
identity of infinite series:
$$
\sum_{a=1}^{N} \frac{\eps(a) \cdot x \cdot e^{ax}}{e^{Nx}-1}
                \,\, = \,\, 
\sum_{k=0}^{\infty} B_{k,\eps} \cdot \frac{x^k}{k!}.
$$
\end{definition}
If $\eps$ is the trivial character of modulus $1$ and $B_k$
are as in Section~\ref{sec:level_one_eisen}, then
$B_{k,\eps} = B_k$, except when $k=1$, in which case
$B_{1,\eps} = -B_1 = 1/2$ (see \exref{ch:eisen}{ex:bern_triv}).

\subsection{Algebraically Computing Generalized Bernoulli Numbers}
Let $\Q(\eps)$ denote the field generated by the image
of the character $\eps$; thus $\Q(\eps)$ is the cyclotomic
extension $\Q(\zeta_n)$, where $n$ is the order of $\eps$.
\begin{algorithm}{Generalized Bernoulli Numbers}\label{alg:gen_bernoulli}
Given an integer $k\geq 0$ and any
Dirichlet character  $\eps$ with modulus~$N$, this algorithm computes
the generalized Bernoulli numbers $B_{j,\eps}$, for $j\leq k$.
\begin{steps}
\item{}\label{step:ber1} Compute $g \set x/(e^{Nx}-1) \in \Q[[x]]$ to
  precision $O(x^{k+1})$
by computing $e^{Nx}-1 = \sum_{n\geq 1} N^n x^n/n!$ 
to precision $O(x^{k+2})$ and computing the inverse
$1/(e^{Nx}-1)$, then multiplying by $x$.  
%Note that if 
%$f = a_0 + a_1 x + a_2x^2 + \cdots$, then we have
%the following recursive formula for the coefficients $b_n$ of the
%expansion of $1/f$:
%$$
% b_n \set - \frac{b_0}{a_0} \cdot (b_{n-1}a_1 + b_{n-2}a_2 + \cdots + b_0 a_n).
%$$
\item{}\label{step:ber2} For each $a=1,\ldots, N$, compute $f_a \set g
  \cdot e^{ax}\in\Q[[x]]$, to precision $O(x^{k+1})$.  This requires
computing $e^{ax}=\sum_{n\geq 0} a^n x^n/n!$ to precision $O(x^{k+1})$.
(Omit computation of $e^{Nx}$ if $N>1$ since then $\eps(N)=0$.)

\item{}\label{step:ber4} Then for $j\leq k$, we have 
$$B_{j,\eps} \set j!\cdot 
  \sum_{a=1}^{N} \eps(a) \cdot c_j(f_a),
$$ where $c_j(f_a)$ is the
  coefficient of $x^j$ in~$f_a$.
\end{steps}
\end{algorithm}
Note that in steps~(\ref{step:ber1}) and (\ref{step:ber2}) we compute the
power series doing arithmetic only in $\Q[[x]]$, not in
$\Q(\eps)[[x]]$, which could be much less efficient if~$\eps$ has
large order.   In step~(\ref{step:ber1}) if $k$ is huge,
we could compute the inverse $1/(e^{Nx}-1)$ using
asymptotically fast arithmetic and
Newton iteration.

\begin{example}
The nontrivial character~$\eps$  with modulus~$4$
has order~$2$ and takes values in~$\Q$.  The Bernoulli
numbers $B_{k,\eps}$ for $k$ even are all $0$ and for $k$ odd
they are
\begin{align*}
B_{1,\eps} &= -1/2,\\
B_{3,\eps} &= 3/2,\\
B_{5,\eps} &= -25/2,\\
B_{7,\eps} &= 427/2,\\
B_{9,\eps} &= -12465/2,\\
B_{11,\eps} &= 555731/2,\\
B_{13,\eps} &= -35135945/2,\\
B_{15,\eps} &= 2990414715/2,\\
B_{17,\eps} &= -329655706465/2,\\
B_{19,\eps} &= 45692713833379/2.
\end{align*}
\end{example}

\begin{example}
  The generalized Bernoulli numbers
  need not be in $\Q$.  Suppose~$\eps$ is the mod~$5$ character
such that $\eps(2) = i = \sqrt{-1}$.  Then  $B_{k,\eps}=0$ for $k$ even
and
\begin{align*}
B_{1,\eps} &= \frac{-i - 3}{5},\\
B_{3,\eps} &= \frac{6i + 12}{5},\\
B_{5,\eps} &= \frac{-86i - 148}{5},\\
B_{7,\eps} &= \frac{2366i + 3892}{5},\\
B_{9,\eps} &= \frac{-108846i - 176868}{5},\\
B_{11,\eps} &= \frac{7599526i + 12309572}{5},\\
B_{13,\eps} &= \frac{-751182406i - 1215768788}{5},\\
B_{15,\eps} &= \frac{99909993486i + 161668772052}{5},\\
B_{17,\eps} &= \frac{-17209733596766i - 27846408467908}{5}.\\
\end{align*}


\end{example}

%begin_sage
\begin{example}\sageindex{generalized Bernoulli numbers}
We use \sage{} to compute some of the above generalized Bernoulli numbers.
First we define the character and verify that $\eps(2)=i$ (note that
in \sage{} {\tt zeta4} is $\sqrt{-1}$).
\begin{verbatim}
sage: G = DirichletGroup(5)
sage: e = G.0
sage: e(2)
zeta4
\end{verbatim}%link

We compute the Bernoulli number $B_{1,\eps}$.
%link
\begin{verbatim}
sage: e.bernoulli(1)
-1/5*zeta4 - 3/5
\end{verbatim}%link

We compute $B_{9,\eps}$.
%link
\begin{verbatim}
sage: e.bernoulli(9)
-108846/5*zeta4 - 176868/5
\end{verbatim}
\end{example}
%end_sage

\begin{proposition}
If $\eps(-1) \neq (-1)^k$ and $k\geq 2$, then $B_{k,\eps}=0$.
\end{proposition}
\begin{proof}
See \exref{ch:eisen}{ex:gbparity}.
\end{proof}

\subsection{Computing Generalized Bernoulli Numbers Analytically}\label{sec:genber}
\mbox{}\newline

This section, which was written jointly with Kevin McGown,
is about a way to compute generalized
Bernoulli numbers, which is similar to the algorithm
in Section~\ref{sec:bercomp}.

Let $\chi$ be a primitive Dirichlet character
modulo its conductor~$f$. 
Note from the definition of Bernoulli numbers that
if $\sigma\in\Gal(\Qbar/\Q)$, then 
\begin{equation}\label{eqn:galois_act_chi}
\sigma(B_{n,\chi}) =  B_{n,\sigma(\chi)}.
\end{equation}

For any character $\chi$, we define the Gauss sum $\tau(\chi)$ as
$$
  \tau(\chi)=\sum_{r=1}^{f-1} \chi(r)\,\zeta^r
  \,,
$$
where $\zeta=\exp(2\pi i/f)$ is the principal $f$th root of unity.
The Dirichlet $L$-function for $\chi$ for $\Re(s)>1$ is
$$
  L(s,\chi)=\sum_{n=1}^\infty \chi(n)\,n^{-s}
  \,.
$$
In the right half plane $\{s\in\C\mid \Re(s)>1\}$
this function is analytic, and because~$\chi$ is
multiplicative, we have the Euler product representation
\begin{align}\label{eqn:direp}
  L(s,\chi)=\prod_{p \text{ prime}} \left(1-\chi(p) p^{-s}\right)^{-1}
  \,.
\end{align}
We note (but will not use) that through analytic continuation $L(s,\chi)$
can be extended to a meromorphic function on the entire complex plane.

If $\chi$ is a nonprincipal primitive Dirichlet character of
conductor $f$ such that $\chi(-1)=(-1)^n$, then (see, e.g.,
\cite{wangkai})
$$
  L(n,\chi)=(-1)^{n-1}\, \frac{\tau(\chi)}{2} \left(\frac{2\pi i}{f}\right)^n
  \frac{B_{n,\overline{\chi}}}{n!}
  \,.
$$
Solving for the Bernoulli number yields
$$
  B_{n,\chi}=(-1)^{n-1}\frac{2 n!}{\tau(\overline{\chi})}\left(\frac{f}{2\pi i}\right)^n
  L(n,\overline{\chi})
  \,.
$$
This allows us to give decimal approximations for $B_{n,\chi}$.
It remains to compute $B_{n,\chi}$ exactly (i.e., as an algebraic integer).
To simplify the above expression, we define
$$
  K_{n,\chi} =(-1)^{n-1}\, 2 n!\left(\frac{f}{2 i}\right)^n
$$
and write
\begin{equation}\label{eqn:bnchi}
  B_{n,\chi}=\frac{K_{n,\chi}}{\pi^n\,\tau(\overline{\chi})}\;
  L(n,\overline{\chi})
  \,.
\end{equation}
Note that we can compute $K_{n,\chi}$ exactly in the field
$\mathbb{Q}(i)$.

The following result
identifies the denominator of $B_{n,\chi}$.
\begin{theorem}
Let $n$ and $\chi$ be as above, and
  define an integer $d$ as follows:
  $$
  d =
  \begin{cases}
    1    & \text{if $f$ is divisible by two distinct primes,}\\
    2 & \text{if $f=4$,}\\
    1 & \text{if $f=2^\mu$, $\mu>2$,}\\
    np   & \text{if $f=p$, $p>2$,}\\
    (1-\chi(1+p)) & \text{if $f=p^\mu$, $p>2$, $\mu>1$.}\\    
  \end{cases}
  $$
Then $dn^{-1}\,B_{n,\chi}$ is integral.
\end{theorem}
\begin{proof}
See \cite{carlitz1} for the proof and \cite{carlitz2}
for further details.
\end{proof}


To compute the algebraic integer $d n^{-1} B_{n,\chi}$, and we compute
$L(n,\overline{\chi})$ to very high precision using the Euler product
\eqref{eqn:direp} and the formula \eqref{eqn:bnchi}.  We carry out the
same computation for each of the $\Gal(\Qbar/\Q)$ conjugates of
$\chi$, which by \eqref{eqn:galois_act_chi} yields the conjugates of
$d n^{-1} B_{n,\chi}$.  We can then write down the characteristic
polynomial of $d n^{-1} B_{n,\chi}$ to very high precision and
recognize the coefficients as rational integers.  Finally, we determine
which of the roots of the characteristic polynomial is $d n^{-1}
B_{n,\chi}$ by approximating them all numerically to high precision
and seeing which is closest to our numerical approximation to $d
n^{-1} B_{n,\chi}$.  The details are similar
to what is explained in Section~\ref{sec:bercomp}.


\section{Explicit Basis for the Eisenstein Subspace}

Suppose $\chi$ and $\psi$ are primitive Dirichlet characters with
conductors $L$ and $R$, respectively.  
Let
\begin{equation}\label{eqn:eisen}
 E_{k,\chi,\psi}(q) = c_0 + \sum_{m \geq 1} \left(
                  \sum_{n|m} \psi(n) \cdot \chi(m/n) \cdot n^{k-1}\right) q^{m}
\in \Q(\chi, \psi)[[q]],
\end{equation}
where 
$$
c_0 = \begin{cases} 0 & \text{ if } L>1, \\
\ds- \frac{B_{k,\psi}}{2k} & \text{ if } L=1.
\end{cases}
$$
Note that when $\chi=\psi=1$ and $k\geq 4$, then $E_{k,\chi,\psi} = E_k$,
where $E_k$ is from Chapter~\ref{ch:modform}.

Miyake proves statements that imply the following in
\cite[Ch.~7]{miyake}. 
\begin{theorem}\label{thm:eisser}
Suppose $t$ is a positive integer and $\chi$, $\psi$ are as above
and that~$k$ is a positive integer such that $\chi(-1)\psi(-1) = (-1)^k$.
Except when $k=2$ and $\chi=\psi=1$, the power series 
$E_{k,\chi,\psi}(q^t)$ defines an element of 
$M_k(RLt,\chi \psi)$.
If $\chi=\psi=1$, $k=2$, $t>1$,  and $E_2(q) = E_{k,\chi,\psi}(q)$, then 
$E_2(q) - t E_2(q^t)$ is a modular form in $M_2(\Gamma_0(t))$.
\end{theorem}

\index{Eisenstein series!basis of}
\begin{theorem}\label{thm:eisgen}
  The Eisenstein series in $M_k(N, \eps)$ coming from
  Theorem~\ref{thm:eisser} with $RLt \mid N$ and $\chi \psi = \eps$
  form a basis for the
  Eisenstein subspace $E_k(N,\eps)$.
\end{theorem}

\index{Eisenstein series!are eigenforms}
\begin{theorem}\label{thm:eiseigen}
The Eisenstein series 
$E_{k,\chi,\psi}(q) \in M_k(RL)$ defined above
are eigenforms (i.e., eigenvectors for all Hecke operators $T_n$).  
Also $E_2(q) - t E_2(q^t)$, for $t>1$, is an eigenform.
\end{theorem}
Since $E_{k,\chi,\psi}(q)$ is normalized so the coefficient
of~$q$ is $1$, the eigenvalue of $T_m$  is  the coefficient
$$\sum_{n|m} \psi(n) \cdot \chi(m/n) \cdot n^{k-1}$$
of $q^m$ (see Proposition~\ref{prop:eigncoeffs}).
Also for $f = E_2(q) - t E_2(q^t)$ with $t>1$ prime, the coefficient of~$q$ 
is~$1$, $T_m(f) = \sigma_{1}(m)\cdot f$ for $(m,t)=1$, and $T_t(f) = 
((t+1)-t)f = f$.

\begin{algorithm}{Enumerating Eisenstein Series}\label{alg:enum_eisen}
\index{Eisenstein series!algorithm to enumerate}
  Given a weight~$k$ and a Dirichlet character~$\eps$ of modulus~$N$,
  this algorithm computes a basis for the Eisenstein
  subspace $E_k(N,\eps)$ of $M_k(N,\eps)$ to precision $O(q^r)$.
\begin{steps}
\item{}[Weight $2$ Trivial Character?] If $k=2$ and $\eps=1$, output the
  Eisenstein series $E_2(q) - tE_2(q^t)$, for each divisor $t\mid N$
  with $t\neq 1$, and then terminate.
\item{}[Empty Space?] If $\eps(-1)\neq (-1)^k$, output the empty list.
\item{}[Compute Dirichlet Group] Let $G\set D(N,\Q(\zeta_n))$ be the
  group of Dirichlet characters with values in $\Q(\zeta_n)$,
  where~$n$ is the exponent of $(\Z/N\Z)^*$.
\item{}[Compute Conductors]\label{step:enum_eisen3} Compute the
  conductor of every element of $G$ using
  Algorithm~\ref{alg:conductor}.
\item{}[List Characters $\chi$]\label{step:enum_eisen4} Form a list
  $V$ of all Dirichlet characters $\chi \in G$ such that
  $\cond(\chi)\cdot \cond(\chi/\eps)$ divides~$N$.
\item{}[Compute Eisenstein Series] For each character $\chi$ in $V$,
  let $\psi = \chi/\eps$ and compute $E_{k,\chi,\psi}(q^t)\pmod{q^r}$
  for each divisor~$t$ of $N/(\cond(\chi)\cdot \cond(\psi))$.  Here we
  compute $E_{k,\chi,\psi}(q^t) \pmod{q^r}$ using \eqref{eqn:eisen}
  and Algorithm~\ref{alg:gen_bernoulli}.

\end{steps}
\end{algorithm}

\begin{remark}
Algorithm~\ref{alg:enum_eisen} is what is currently used in \sage.
It might be better to first reduce to the prime power case by writing
all characters as a product of local characters and combine
steps~(\ref{step:enum_eisen3}) and (\ref{step:enum_eisen4}) into a single
step that involves orders.  However, this might make things more
obscure.
\end{remark}
\begin{example}
The following is a basis of Eisenstein series 
for $E_2(\Gamma_1(13))$.
\begin{align*}
f_{1} &= \frac{1}{2} + q + 3q^{2} + 4q^{3} + \cdots,\\
f_{2} &= -\frac{7}{13}\zeta_{12}^{2} - \frac{11}{13} + q + \left(2\zeta_{12}^{2} + 1\right)q^{2} + \left(-3\zeta_{12}^{2} + 1\right)q^{3} + \cdots,\\
f_{3} &= q + \left(\zeta_{12}^{2} + 2\right)q^{2} + \left(-\zeta_{12}^{2} + 3\right)q^{3} + \cdots,\\
f_{4} &= -\zeta_{12}^{2} + q + \left(2\zeta_{12}^{2} - 1\right)q^{2} + \left(3\zeta_{12}^{2} - 2\right)q^{3} + \cdots,\\
f_{5} &= q + \left(\zeta_{12}^{2} + 1\right)q^{2} + \left(\zeta_{12}^{2} + 2\right)q^{3} + \cdots,\\
f_{6} &= -1 + q + -q^{2} + 4q^{3} + \cdots,\\
f_{7} &= q + q^{2} + 4q^{3} + \cdots,\\
f_{8} &= \zeta_{12}^{2} - 1 + q + \left(-2\zeta_{12}^{2} + 1\right)q^{2} + \left(-3\zeta_{12}^{2} + 1\right)q^{3} + \cdots,\\
f_{9} &= q + \left(-\zeta_{12}^{2} + 2\right)q^{2} + \left(-\zeta_{12}^{2} + 3\right)q^{3} + \cdots,\\
f_{10} &= \frac{7}{13}\zeta_{12}^{2} - \frac{18}{13} + q + \left(-2\zeta_{12}^{2} + 3\right)q^{2} + \left(3\zeta_{12}^{2} - 2\right)q^{3} + \cdots,\\
f_{11} &= q + \left(-\zeta_{12}^{2} + 3\right)q^{2} + \left(\zeta_{12}^{2} + 2\right)q^{3} + \cdots.\\
\end{align*}

%begin_sage
\sageindex{Eisenstein series}
We computed it as follows:\index{Eisenstein series!compute using \sage}
\begin{verbatim}
sage: E = EisensteinForms(Gamma1(13),2)
sage: E.eisenstein_series()
\end{verbatim}%link

We can also compute the parameters $\chi,\psi,t$ that
define each series:
%link
\begin{verbatim}
sage: e = E.eisenstein_series()
sage: for e in E.eisenstein_series():
...       print e.parameters()
...
([1], [1], 13)
([1], [zeta6], 1)
([zeta6], [1], 1)
([1], [zeta6 - 1], 1)
([zeta6 - 1], [1], 1)
([1], [-1], 1)
([-1], [1], 1)
([1], [-zeta6], 1)
([-zeta6], [1], 1)
([1], [-zeta6 + 1], 1)
([-zeta6 + 1], [1], 1)
\end{verbatim}
%end_sage

\end{example}

\begin{exercises}


\item\label{ex:diag} Suppose $A$ and $B$ are diagonalizable linear
  transformations of a finite-dimensional vector space~$V$ over an
  algebraically closed field $K$ and that $AB=BA$.  Prove there is a
  basis for $V$ so that the matrices of $A$ and $B$ with respect to
  that basis are both simultaneously diagonal.

\item\label{ex:bern_triv} If $\eps$ is the trivial character of
  modulus $1$ and $B_k$ are as in Section~\ref{sec:level_one_eisen},
  then $B_{k,\eps} = B_k$, except when $k=1$, in which case
  $B_{1,\eps} = -B_1 = 1/2$.

\item \label{ex:gbparity} Prove that for $k\geq 2$ if $\eps(-1) \neq (-1)^k$,
then $B_{k,\eps} = 0$.

\item\label{ex:dime3} Show that the dimension of the Eisenstein subspace
  $E_3(\Gamma_1(13))$ is $12$ by finding a basis of series
  $E_{k,\chi,\psi}$.  You do not have to write down the $q$-expansions
  of the series, but you do have to figure out which $\chi, \psi$ to
  use.

\end{exercises}




















\chapter{Dimension Formulas}
\label{ch:dim}

When computing with spaces of modular forms, it is helpful to have
easy-to-compute formulas for dimensions of these spaces.  Such
formulas provide a check on the output of the algorithms from
Chapter~\ref{ch:modsym} that compute explicit bases for spaces of
modular forms.  We can also use dimension formulas to improve the
efficiency of some of the algorithms in Chapter~\ref{ch:modsym}, since
we can use them to determine the ranks of certain matrices without
having to explicitly compute those matrices.  Dimension formulas can also be
used in generating bases of $q$-expansions; if we know the dimension
of $M_k(N,\eps)$ and if we have a process for computing $q$-expansions
of elements of $M_k(N,\eps)$, e.g., multiplying together
$q$-expansions of certain forms of smaller weight, then we can tell
when we are done generating $M_k(N,\eps)$.

This chapter contains formulas for dimensions of spaces of modular
forms, along with some remarks about how to evaluate these formulas.
In some cases we give dimension formulas for spaces that we
will define in later chapters.  
We also give
many examples, some of which were computed using the modular symbols
algorithms from Chapter~\ref{ch:modsym}.

Many of the dimension formulas and algorithms we give below grew out
of Shimura's book \cite{shimura:intro} and a program that Bruce Kaskel
wrote (around 1996) in PARI, which Kevin Buzzard extended.  That
program codified dimension formulas that Buzzard and Kaskel found or
extracted from the literature (mainly \cite[\S2.6]{shimura:intro}).
The algorithms for dimensions of spaces with nontrivial character are
from~\cite{cohen-oesterle:dimensions}, with some refinements
suggested by Kevin Buzzard.

For the rest of this chapter,~$N$ denotes a positive integer and
$k\geq 2$ is an integer.  We will give {\em no simple formulas} for
dimensions of spaces of weight $1$ modular forms; in fact, it 
might not be possible to give such formulas since the 
methods used to derive the formulas below do not apply in the case
$k=1$.  If $k=0$, the only modular forms are the constants, and for
$k<0$ the dimension of $M_k(N,\eps)$ is $0$.

For a nonzero integer $N$ and a prime~$p$, let $v_p(N)$ be the largest
integer~$e$ such that $p^e \mid N$.  In the formulas in this
chapter,~$p$ always denotes a prime number.  Let $M_k(N,\eps)$ be the
space of modular forms of level~$N$ weight~$k$ and character~$\eps$,
and let $S_k(N,\eps)$ and $E_k(N,\eps)$ be the cuspidal and Eisenstein
subspaces, respectively.

The dimension formulas below for $S_k(\Gamma_0(N))$,
$S_k(\Gamma_1(N))$, $E_k(\Gamma_0(N))$ and $E_k(\Gamma_1(N))$ can be
found in \cite[Ch.~3]{diamond-shurman},
\cite[\S2.6]{shimura:intro}\footnote{The formulas in
  \cite[\S2.6]{shimura:intro} contain some minor mistakes.} and
\cite[\S2.5]{miyake}. They are derived using the Riemann-Roch Theorem
applied to the covering $X_0(N)\to X_0(1)$ or $X_1(N)\to X_1(1)$ and
appropriately chosen divisors.  It would be natural to give a sample
argument along these lines at this point, but we will not since it
easy to find such arguments in other books and survey papers (see,
e.g., \cite{diamond-im}).  So you will not learn much about how to
derive dimension formulas from this chapter.  What you will learn is
precisely what the dimension formulas are, which is something that is
often hard to extract from obscure references.

In addition to reading this chapter, the reader may wish to consult
\cite{martin:dims} for proofs of similar dimension formulas,
asymptotic results, and a nonrecursive formula for dimensions of
certain new subspaces.

\section{Modular Forms for $\Gamma_0(N)$}\label{sec:dimg0}
For any prime $p$ and any positive integer $N$, let $v_p(N)$ be 
the power of $p$ that divides $N$. Also, let
\begin{align*}
\mu_0(N) &= \prod_{p\mid N} \left( p^{v_p(N)} + p^{v_p(N)-1} \right),\\
\mu_{0,2}(N) &= \begin{cases}
    0 & \text{if $4\mid N$,}\\
  \prod_{p\mid N} \left(1 + \kr{-4}{p}\right) & \text{otherwise,}
\end{cases}\\
\mu_{0,3}(N) &= \begin{cases}
    0 & \text{if $2\mid N$ or $9\mid N$,}\\
  \prod_{p\mid N} \left(1 + \kr{-3}{p}\right) & \text{otherwise,}
\end{cases}\\
c_0(N) &= \sum_{d\mid N} \vphi(\gcd(d,N/d)),\\
g_0(N) &= 1 + \frac{\mu_0(N)}{12} - \frac{\mu_{0,2}(N)}{4} - \
               \frac{\mu_{0,3}(N)}{3} - \frac{c_0(N)}{2}.\\
\end{align*}
Note that $\mu_0(N)$ is the index of $\Gamma_0(N)$ in $\SL_2(\Z)$
(see \exref{ch:dim}{ex:mu0n}).

\index{cusp forms!higher level dimension}
\index{dimension!cusp forms of higher level}
\begin{proposition}\label{prop:dimg0}
We have $\dim S_2(\Gamma_0(N)) = g_0(N)$, and
for $k\geq 4$ even, 
\begin{align*}
\dim S_k(\Gamma_0(N)) &= 
 (k-1) \cdot (g_0(N) - 1) \,\,+ \,\,
\left( \frac{k}{2} - 1 \right) \cdot c_0(N) \,\,   \\
  &\qquad\qquad   + \mu_{0,2}(N)\cdot \left\lfloor\frac{k}{4}\right\rfloor \,\,+ \,\,
     \mu_{0,3}(N)\cdot \left\lfloor \frac{k}{3}\right\rfloor.
\end{align*}
The dimension of the Eisenstein subspace is
$$
  \dim E_k(\Gamma_0(N)) = 
\begin{cases}
c_0(N) & \text{if $k\neq 2$,}\\
c_0(N)-1 & \text{if $k=2$.}
\end{cases}
$$
\end{proposition}

The following is a table of $\dim S_k(\Gamma_0(N))$
for some values of $N$ and $k$:
\begin{center}
\begin{tabular}{|l|llll|}\hline
$N$ & $S_2(\Gamma_0(N))$ & $S_4(\Gamma_0(N))$ & 
      $S_6(\Gamma_0(N))$ & $S_{24}(\Gamma_0(N))$\\\hline
1 & 0 & 0 & 0 & 2\\
10 & 0 & 3 & 5 & 33\\
11 & 1 & 2 & 4 & 22\\
100 & 7 & 36 & 66 & 336\\
389 & 32 & 97 & 161 & 747\\
1000 & 131 & 430 & 730 & 3430\\
2007 & 221 & 806 & 1346 & 6206\\
100000 & 14801 & 44800 & 74800 & 344800\\\hline
\end{tabular}
\end{center}

%begin_sage
\begin{example}\sageindex{dimension formulas}
Use the commands {\tt dimension\_cusp\_forms}, 
{\tt dimension\_eis}, and {\tt dimension\_modular\_forms}
to compute the dimensions of the three spaces $S_k(\Gamma_0(N))$,
$E_k(\Gamma_0(N))$ and $M_k(\Gamma_0(N))$, respectively.
For example, 
\begin{verbatim}
sage: dimension_cusp_forms(Gamma0(2007),2)
221
sage: dimension_eis(Gamma0(2007),2)
7
sage: dimension_modular_forms(Gamma0(2007),2)
228
\end{verbatim}
\end{example}
%end_sage

\begin{figure}[ht]
\begin{center}
\includegraphics[width=0.8\textwidth]{diagrams/dims2g0}
\caption{Dimension of $S_2(\Gamma_0(N))$ as a function of $N$.\label{fig:dims2g0}}
\end{center}
\end{figure}

\begin{figure}[ht]
\begin{center}
\includegraphics[width=0.8\textwidth]{diagrams/dimskg012}
\caption{Dimension of $S_{2k}(\Gamma_0(12))$ as a function of $k$.\label{fig:dimskg012}}
\end{center}
\end{figure}

\begin{remark}
  Csirik, Wetherell, and Zieve prove in
  \cite{csirik-wetherell-zieve:g0} that a random positive integer has
  probability~$0$ of being a value of $$g_0(N)=\dim S_2(\Gamma_0(N)),$$
  and they give bounds on the size of the set of values of $g_0(N)$ below
  some given~$x$.  For example, they show that $150, 180, 210, 286,
  304, 312, \ldots$ are the first few integers that are not of the
  form $g_0(N)$ for any~$N$.  See Figure~\ref{fig:dims2g0}
  for a plot of the very erratic function $g_0(N)$.  In
  contrast, the function $k\mapsto \dim S_{2k}(\Gamma_0(12))$ is very
  well behaved (see Figure~\ref{fig:dimskg012}). 
\end{remark}


\subsection{New and Old Subspaces}\label{sec:newold}





In this section we assume the reader is either
familiar with newforms or has read Section~\ref{sec:ALL}.
 
For any integer $R$, let
$$
\overline{\mu}(R) = 
\begin{cases} 
 0 & \text{if $p^3\mid R$ for some $p$,}\\
\ds \prod_{p\mid\mid R} -2 & \text{otherwise},
\end{cases}
$$
where the product is over primes that exactly divide~$R$.
Note that $\overline{\mu}$ is {\em not} the Moebius function, but it
has a similar flavor. 

\begin{proposition}\label{prop:newg0}
The dimension of the new subspace is
$$
 \dim S_k(\Gamma_0(N))_{\new} = \sum_{M\mid N} \overline{\mu}(N/M) 
   \cdot \dim S_k(\Gamma_0(M)),
$$
where the sum is over the positive divisors $M$ of $N$.
As a consequence of Theorem~\ref{thm:ALL},
we also have
$$
  \dim S_k(\Gamma_0(N)) = \sum_{M\mid N} \sigma_0(N/M)
                             \dim S_k(\Gamma_0(M))_{\new},
$$
where $\sigma_0(N/M)$ is the number of divisors of $N/M$.
\end{proposition}

%begin_sage
\begin{example}\sageindex{dimension $S_k(\Gamma_0(N))$}
We compute the dimension of the new subspace of $S_k(\Gamma_0(N))$
using the \sage{} command {\tt dimension\_new\_cusp\_forms}
as follows:
\begin{verbatim}
sage: dimension_new_cusp_forms(Gamma0(11),12)
8
sage: dimension_cusp_forms(Gamma0(11),12)
10
sage: dimension_new_cusp_forms(Gamma0(2007),12)
1017
sage: dimension_cusp_forms(Gamma0(2007),12)
2460
\end{verbatim}
\end{example}
%end_sage








\section{Modular Forms for $\Gamma_1(N)$}
This section follows Section~\ref{sec:dimg0} closely, but with
suitable modifications with $\Gamma_0(N)$ replaced by $\Gamma_1(N)$.

Define functions of a positive integer~$N$ by the following formulas:
\begin{align*}
\mu_1(N) &= 
\begin{cases}
\mu_0(N) & \text{if $N=1,2$,}\\
\ds\frac{\phi(N)\cdot \mu_0(N)}{2} & \text{otherwise,}
\end{cases} \\
\mu_{1,2}(N) &= \begin{cases}
    0 & \text{if $N\geq 4$,}\\
  \mu_{0,2}(N) & \text{otherwise,}
\end{cases}\\
\mu_{1,3}(N) &= \begin{cases}
    0 & \text{if $N\geq 4$,}\\
    \mu_{0,3}(N) & \text{otherwise,}
\end{cases}\\
c_1(N) &= 
\begin{cases}
  c_0(N) & \text{if $N=1,2$,}\\
  3      & \text{if $N=4$,}\\
  \ds \sum_{d \mid N} \frac{\phi(d)\phi(N/d)}{2} & \text{otherwise},
\end{cases} \\
g_1(N) &= 1 + \frac{\mu_1(N)}{12} - \frac{\mu_{1,2}(N)}{4} - \frac{\mu_{1,3}(N)}{3}
            - \frac{c_1(N)}{2}.
\end{align*}
Note that $g_1(N)$ is the genus of the modular curve $X_1(N)$ (associated
to $\Gamma_1(N)$) and $c_1(N)$ is the number of cusps of $X_1(N)$.

\index{cusp forms!higher level dimension}
\index{dimension!cusp forms of higher level}
\begin{proposition}
  We have $\dim S_2(\Gamma_1(N)) = g_1(N)$.  If $N\leq 2$, then 
$\Gamma_0(N)=\Gamma_1(N)$ so 
$$\dim
  S_k(\Gamma_1(N)) = \dim S_k(\Gamma_0(N)),$$
  where $\dim
  S_k(\Gamma_0(N))$ is given by the formula of
  Proposition~\ref{prop:dimg0}.  If $k\geq 3$, let $$
  a(N,k) = (k-1)(g_1(N)
  - 1) + \left(\frac{k}{2} -1\right)\cdot c_1(N).  $$
  Then for $N\geq
  3$, $$
  \dim S_k(\Gamma_1(N)) =
\begin{cases}
 a + 1/2 & \text{if $N=4$ and $2\nmid k$,}\\
 a + \lfloor{} k/3\rfloor{} & \text{if $N=3$,}\\
 a & \text{otherwise.}
\end{cases}
$$
The dimension of the Eisenstein subspace is as follows:
$$
  \dim E_k(\Gamma_1(N)) = 
\begin{cases}
c_1(N) & \text{if $k\neq 2$,}\\
c_1(N)-1 & \text{if $k=2$.}
\end{cases}
$$
The dimension of the new subspace of $M_k(\Gamma_1(N))$ is
$$
\dim S_k(\Gamma_1(N))_{\new} = 
     \sum_{M \mid N} \overline{\mu}(N/M) \cdot \dim S_k(\Gamma_1(M)),
$$
where $\overline{\mu}$ is as in the statement
of Proposition~\ref{prop:newg0}.
\end{proposition}
\begin{remark}
Since $M_k = S_k \oplus E_k$, the formulas above for
$\dim S_k$ and $\dim E_k$ also yield
a formula for the dimension of $M_k$.  
\end{remark}

\begin{figure}[ht]
\begin{center}
\includegraphics[width=0.8\textwidth]{diagrams/dims2g1}
\caption{Dimension of $S_2(\Gamma_1(N))$ as a function of $N$.\label{fig:dims2g1}}
\end{center}
\end{figure}


The following table contains the dimension of $S_k(\Gamma_1(N))$
for some sample values of $N$ and $k$:
\begin{center}
\begin{tabular}{|l|llll|}\hline
$N$ & $S_2(\Gamma_1(N))$ & $S_3(\Gamma_1(N))$ & 
      $S_4(\Gamma_1(N))$ & $S_{24}(\Gamma_1(N))$\\\hline
1 & 0 & 0 & 0 & 2\\
10 & 0 & 2 & 5 & 65\\
11 & 1 & 5 & 10 & 110\\
100 & 231 & 530 & 830 & 6830\\
389 & 6112 & 12416 & 18721 & 144821\\
1000 & 28921 & 58920 & 88920 & 688920\\
2007 & 147409 & 296592 & 445776 & 3429456\\
 100000 & 299792001 & 599792000 & 899792000 & 6899792000\\\hline
\end{tabular}
\end{center}

%begin_sage
\begin{example}\sageindex{dimension $S_k(\Gamma_1(N))$}
We compute dimensions of
spaces of modular forms for $\Gamma_1(N)$:
\begin{verbatim}
sage: dimension_cusp_forms(Gamma1(2007),2)
147409
sage: dimension_eis(Gamma1(2007),2)
3551
sage: dimension_modular_forms(Gamma1(2007),2)
150960
\end{verbatim}
\end{example}
%end_sage

\section{Modular Forms with Character}
Fix a Dirichlet character~$\eps$ of modulus~$N$,
and let $c$ be the conductor of~$\eps$ (we do {\em not}
assume that $\eps$ is primitive).
 Assume that $\eps \neq 1$, since otherwise
$M_k(N,\eps) = M_k(\Gamma_0(N))$ and the formulas of
Section~\ref{sec:dimg0} apply.  Also, assume that $\eps(-1)=(-1)^k$,
since otherwise $\dim M_k(\Gamma_0(N)) = 0$.  In this section we
discuss formulas for computing each of $M_k(N,\eps)$, $S_k(N,\eps)$
and $E_k(N,\eps)$.

In \cite{cohen-oesterle:dimensions}, Cohen and 
Oesterl\'e assert (without {\em published} proof; see Remark~\ref{rem:co} below) 
that for any $k\in\Z$ and $N$, $\eps$
as above,
\begin{align*}
\dim S_k(N,\eps) &- \dim M_{2-k}(N,\eps)\\
  &= \frac{k-1}{12} \cdot \mu_0(N) \,\,-\,\,
         \frac{1}{2} \cdot \prod_{p\mid N} \lambda(p, N, v_p(c))\\
  &\qquad + \gamma_4(k) \cdot\!\! 
\sum_{x \in A_4(N)} \eps(x) \,\,\,+\,\,\, 
   \gamma_3(k) \cdot \!\! \sum_{x \in A_3(N)} \eps(x)
\end{align*}
where $\mu_0(N)$ is as in Section~\ref{sec:dimg0},
$A_4(N) = \{x \in \Z/N\Z : x^2 +1 = 0\}$ and
$A_3(N) = \{x \in \Z/N\Z : x^2 + x + 1 = 0 \}$,
and $\gamma_3, \gamma_4$ are
\begin{align*}
\gamma_4(k) &= 
    \begin{cases}
     -1/4 & \text{if $k\con 2\pmod{4}$,}\\
      1/4 & \text{if $k\con 0\pmod{4}$,}\\
        0 & \text{if $k$ is odd.}
     \end{cases}\\
\gamma_3(k) &= 
    \begin{cases}
     -1/3 & \text{if $k\con 2\pmod{3}$,}\\
      1/3 & \text{if $k\con 0\pmod{3}$,}\\
        0 & \text{if $k\con 1\pmod{3}$.}
     \end{cases}
\end{align*}
It remains to define $\lambda$.  
Fix a prime divisor $p\mid N$ and let $r=v_p(N)$.
Then 
$$
\lambda(p,N,v_p(c)) = 
    \begin{cases}
      p^{\frac{r}{2}} + p^{\frac{r}{2} - 1} & \text{if $2\cdot v_p(c)\leq r$ and $2\mid r$,}\\
     2\cdot p^{\frac{r-1}{2}} & \text{if $2\cdot v_p(c)\leq r$ and $2\nmid r$,}\\
     2\cdot p^{r-v_p(c)} & \text{if $2\cdot v_p(c) > r$.}
   \end{cases}
$$   

This flexible formula can be used to compute the dimension of
$M_k(N,\eps)$, 
$S_k(N,\eps)$, and $E_k(N,\eps)$ for any $N$, $\eps$, $k\neq 1$,
by using that
\begin{align*}
  \dim S_k(N,\eps) &= 0 \qquad \text{if $k\leq 0$,}\\
  \dim M_k(N,\eps) &= 0 \qquad \text{if $k<0$,}\\
  \dim M_0(N,\eps) &= 1 \qquad \text{if $k=0$.}
\end{align*}

One thing that is not straightforward when implementing an algorithm
to compute the above dimension formulas is how to efficiently compute
the sets $A_4(N)$ and $A_6(N)$.  Kevin Buzzard suggested the
following two algorithms.  Note that if~$k$
is odd, then $\gamma_4(k)=0$, so the sum over
$A_4(N)$ is only needed when $k$ is even.  
\begin{algorithm}{Sum over $A_4(N)$}\label{alg:suma4}
Given a positive integer $N$ and an even 
Dirichlet character $\eps$ of modulus~$N$,
this algorithm computes  $\sum_{x \in A_4(N)} \eps(x)$.
\begin{steps}
\item{}[Factor $N$]  Compute the prime factorization 
$p_1^{e_1} \cdots p_n^{e_n}$ of $N$.
\item{}[Initialize] Set $t\set 1$ and $i\set 0$.
\item{}[Loop Over Prime Divisors]\label{alg:sum1go} Set $i\set i+1$.
If $i>n$, return $t$.  Otherwise 
set $p\set p_i$ and $e\set e_i$.
\begin{steps}
\item\label{alg:a4_3a} If $p\con 3\pmod{4}$, return $0$.
\item\label{alg:a4_3b} If $p=2$ and $e>1$, return $0$.
\item If $p=2$ and $e=1$, go to step~(\ref{alg:sum1go}).
\item\label{alg:a4_3d} Compute a generator $a\in (\Z/p\Z)^*$ using
Algorithm~\ref{alg:mingens}.
\item Compute $\omega = a^{(p-1)/4}$.
\item Use the Chinese Remainder Theorem to find $x\in\Z/N\Z$
such that $x\con a\pmod{p}$ and $x\con 1\pmod{N/p^e}$.
\item\label{alg:a4_3g} Set $x \set x^{p^{r-1}}$.
\item\label{alg:a4_3h} Set $s\set \eps(x)$.
\item\label{alg:a4_3i} If $s=1$, set $t\set 2t$ and go to 
step~(\ref{alg:sum1go}).
\item\label{alg:a4_3j} If $s=-1$, set $t\set -2t$ and go to 
step~(\ref{alg:sum1go}).
\end{steps}
\end{steps}
\end{algorithm}
\begin{proof}
  Note that $\eps(-x)=\eps(x)$, since $\eps$ is even.  By the Chinese
  Remainder Theorem, the set $A_4(N)$ is empty if and only if there is
  no square root of $-1$ modulo some prime power divisor of $p$.  If
  $A_4(N)$ is empty, the algorithm correctly detects this fact in
  steps \eqref{alg:a4_3a}--\eqref{alg:a4_3b}.  
  Thus assume $A_4(N)$ is
  nonempty.  For each prime power $p_i^{e_i}$ that exactly divides
  $N$, let $x_i \in Z/N\Z$ be such that $x_i^2 = -1$ and $x_i\con
  1\pmod{p_j^{e_j}}$ for $i \neq j$.  This is the value of $x$
  computed in steps \eqref{alg:a4_3d}--\eqref{alg:a4_3g} (as one sees
  using elementary number theory).

  The next key observation is that 
\begin{equation}\label{eqn:a4prodsum}
  \prod_i (\eps(x_i) + \eps(-x_i)) = 
    \sum_{x \in A_4(N)} \eps(x),
 \end{equation}
 since by the Chinese Remainder Theorem the elements of $A_4(N)$ are
 in bijection with the choices for a square root of $-1$ modulo each
 prime power divisors of~$N$.  The observation \eqref{eqn:a4prodsum}
 is a huge gain from an efficiency point of view---if $N$ had $r$
 prime factors, then $A_4(N)$ would have size $2^r$, which could be
 prohibitive, where the product involves only $r$ factors.  To finish
 the proof, just note that steps~(\ref{alg:a4_3h})--(\ref{alg:a4_3j})
 compute the local factors $\eps(x_i) + \eps(-x_i) = 2\eps(x_i)$,
 where again we use that $\eps$ is even.  Note that a solution
 of $x^2+1 \con 0\pmod{p}$ lifts uniquely to a solution mod $p^n$ for
 any~$n$, because the kernel of the natural homomorphism $(\Z/p^n\Z)^*
 \to (\Z/p\Z)^*$ is a group of $p$-power order.
\end{proof}
The algorithm for computing the sum over $A_3(N)$ is similar.

For $k\geq 2$, to compute $\dim S_k(N,\eps)$, use the formula directly
and the fact that $\dim M_{2-k}(N,\eps) = 0$, unless $\eps=1$ and $k=2$.
To compute $\dim M_k(N,\eps)$ for $k\geq 2$, use the fact that the big
formula at the beginning of this section 
is valid for any integer $k$ to replace $k$
by $2-k$ and 
that $\dim S_k(N,\eps)=0$ for $k\leq 0$ to rewrite the formula as 
\begin{align*}
 \dim M_k(N,\eps) &= - (\dim S_{2-k}(N,\eps) - \dim M_{k}(N,\eps))\\
  &=-\Bigl( \frac{1-k}{12} \cdot \mu_0(N) \,\,-\,\,
         \frac{1}{2} \cdot \prod_{p\mid N} \lambda(p, N, v_p(c))\\
  &\qquad + \gamma_4(2-k) \cdot\!\! 
\sum_{x \in A_4(N)} \eps(x) \,\,\,+\,\,\, 
   \gamma_3(2-k) \cdot \!\! \sum_{x \in A_3(N)} \eps(x)\Bigr).
\end{align*}
Note also that for $k=0$, $\dim E_k(N,\eps) = 1$ if 
and only if~$\eps$ is trivial and it equals~$0$ otherwise.
We then also obtain
$$
 \dim E_k(N,\eps) = \dim M_k(N,\eps) - \dim S_k(N,\eps).
$$
We can also compute $\dim E_k(N,\eps)$ when $k=1$
directly, since $$\dim S_{2-1}(N,\eps) = \dim S_{1}(N,\eps).$$


The following table contains the dimension of $S_k(N,\eps)$
for some sample values of $N$ and $k$.  In each case, $\eps$
is the product of characters $\eps_p$ of maximal order
corresponding to the prime power factors of~$N$ (i.e.,
the product of the generators of the group
$D(N,\C^*)$ of Dirichlet characters of modulus~$N$).
\begin{center}
\begin{tabular}{|l|llll|}\hline
$N$ & $\dim S_2(N,\eps)$ & $\dim S_3(N,\eps)$ & 
      $\dim S_4(N,\eps)$ & $\dim S_{24}(N,\eps)$\\\hline
1 & 0 & 0 & 0 & 2\\
10 & 0 & 1 & 0 & 0\\
11 & 0 & 1 & 0 & 0\\
100 & 13 & 0 & 43 & 343\\
389 & 0 & 64 & 0 & 0\\
1000 & 148 & 0 & 448 & 3448\\
2007 & 222 & 0 & 670 & 5150\\\hline
\end{tabular}
\end{center}

\begin{example}
We compute the last line of the above table.  First 
we create the character $\eps$.
\sageindex{dimension with character}
\begin{verbatim}
sage: G = DirichletGroup(2007)
sage: e = prod(G.gens(), G(1))
\end{verbatim}%link

Next we compute the dimension of the four spaces.
%link
\begin{verbatim}
sage: dimension_cusp_forms(e,2)
222
sage: dimension_cusp_forms(e,3)
0
sage: dimension_cusp_forms(e,4)
670
sage: dimension_cusp_forms(e,24)
5150
\end{verbatim}%link

We can also compute dimensions of the corresponding
spaces of Eisenstein series. 
%link
\begin{verbatim}
sage: dimension_eis(e,2)
4
sage: dimension_eis(e,3)
0
sage: dimension_eis(e,4)
4
sage: dimension_eis(e,24)
4
\end{verbatim}
\end{example}



\begin{remark}\label{rem:co}
  Cohen and Oesterl\'e also give dimension formulas for spaces of
  half-integral weight modular forms, which we do not give in this
  chapter.  Note that \cite{cohen-oesterle:dimensions} does not
  contain any {\em proofs} that their claimed formulas are correct,
  but instead they say only that ``Les formules qui les donnent sont
  connues de beaucoup de gens et il existe plusieurs m\'ethodes
  permettant de les obtenir (th\'eor\`eme de Riemann-Roch, application
  des formules de trace donn\'ees par Shimura).''\footnote{The formulas that
  we give here are well known and there exist many methods to prove
  them, e.g., the Riemann-Roch theorem and applications of the trace
  formula of Shimura.} Fortunately, in \cite{quer:dim},
  Jordi Quer derives the 
  (integral weight) formulas of \cite{cohen-oesterle:dimensions}
  along with formulas for dimensions of spaces $S_k(G)$ and
  $M_k(G)$ for more general congruence subgroups. 
\end{remark}

Let $f$ be the conductor of a Dirichlet character $\eps$ of modulus~$N$.  
Then the dimension of the new subspace of $M_k(N,\eps)$ is
$$
\dim S_k(N,\eps)_{\new} = 
     \sum_{M\text{ such that }f \mid M \mid N} \overline{\mu}(N/M) \cdot \dim S_k(M,\eps'),
$$
where $\overline{\mu}$ is as in the statement
of Proposition~\ref{prop:newg0}, and $\eps'$ is the restriction
of~$\eps$ mod~$M$.

\begin{example}
We compute the dimension of $S_2(2007,\eps)_{\new}$
for $\eps$ a quadratic character of modulus $2007$.
\begin{verbatim}
sage: G = DirichletGroup(2007, QQ)
sage: e = prod(G.gens(), G(1))
sage: dimension_new_cusp_forms(e,2)
76
\end{verbatim}
\end{example}


 \begin{exercises}
\item\label{ex:mu0n} 
Let $\mu_0$ and $\mu_1$ be as in this chapter. 
\begin{enumerate}
\item Prove that $\mu_0(N) = [\SL_2(\Z):\Gamma_0(N)]$.
\item Prove that for $N\geq 3$, $\mu_1(N) = [\SL_2(\Z):\Gamma_1(N)]/2$,
so $\mu_1(N)$ is the index of $\Gamma_1(N) \cdot \{\pm 1\}$
in $\PSL_2(\Z) = \SL_2(\Z)/\{\pm 1\}$.
\end{enumerate}

\item \label{ex:formsl2z}
Use Proposition~\ref{prop:newg0} to find a formula
for $\dim S_k(\SL_2(\Z))$.  Verify that this formula is
the same as the one in Corollary~\ref{cor:dim1}.

\item \label{ex:newmatch}
Suppose either that $N=1$ or that $N$ is prime and $k=2$.
Prove that $M_k(\Gamma_0(N))_{\new} = M_k(\Gamma_0(N))$.

\item \label{ex:suma4} Fill in the details of the 
proof of Algorithm~\ref{alg:suma4}.

\item \label{ex:sagedim} Implement a computer program
to compute $\dim S_k(\Gamma_0(N))$ as a function of $k$ and $N$.


\end{exercises}











\chapter{Linear Algebra}\label{ch:linalg}
This chapter is about several algorithms for matrix algebra over the
rational numbers and cyclotomic fields.  Algorithms for linear algebra
over exact fields are necessary in order to implement the modular
symbols algorithms that we will describe in Chapter~\ref{ch:linalg}.
This chapter partly overlaps with \cite[Sections 2.1--2.4]{cohen:course_ant}.

\vspace{2ex}
\noindent{\bf Note:} We view all matrices as defining
linear transformations by acting on row vectors from the right.

\section{Echelon Forms of Matrices}\label{sec:echelon_form}

\begin{definition}[Reduced Row Echelon Form]
  A matrix is in (reduced row) \defn{echelon form} if each row in the matrix
  has more zeros at the beginning than the row above it, the first
  nonzero entry of every row is $1$, and the first nonzero entry of
  any row is the only nonzero entry in its column.
\end{definition}

Given a matrix~$A$, there is another matrix~$B$ such that~$B$ is
obtained from~$A$ by left multiplication by an invertible matrix
and~$B$ is in reduced row echelon form.  This matrix~$B$ is called the
echelon form of~$A$.  It is unique.

A \defn{pivot column} of~$A$ is a column of $A$ such that the reduced
row echelon form of~$A$ contains a leading~$1$.

\begin{example}\label{ex:rref1}
The following matrix is not in reduced row echelon form:
$$
\left(\begin{array}{rrrrr}
14&2&7&228&-224\\
0&0&3&78&-70\\
0&0&0&-405&381
\end{array}\right).
$$
The reduced row echelon form of the above matrix is 
$$
\left(\begin{array}{rrrrr}
1&\frac{1}{7}&0&0&-\frac{1174}{945}\vspace{2pt}\\
0&0&1&0&\frac{152}{135}\vspace{2pt}\\
0&0&0&1&-\frac{127}{135}
\end{array}\right).
$$
Notice that the entries of the reduced row echelon form can be
rationals with large denominators even though the entries of the
original matrix $A$ are integers.  Another example is the simple
looking matrix
$$
\left(\begin{array}{rrrrrrrr}
-9&6&7&3&1&0&0&0\\
-10&3&8&2&0&1&0&0\\
3&-6&2&8&0&0&1&0\\
-8&-6&-8&6&0&0&0&1
\end{array}\right)
$$
whose echelon form is 
$$
\left(\begin{array}{rrrrrrrr}
1&0&0&0&\frac{42}{1025}&-\frac{92}{1025}&\frac{1}{25}&-\frac{9}{205}\vspace{2pt}\\
0&1&0&0&\frac{716}{3075}&-\frac{641}{3075}&-\frac{2}{75}&-\frac{7}{615}\vspace{2pt}\\
0&0&1&0&-\frac{83}{1025}&\frac{133}{1025}&\frac{1}{25}&-\frac{23}{410}\vspace{2pt}\\
0&0&0&1&\frac{184}{1025}&-\frac{159}{1025}&\frac{2}{25}&\frac{9}{410}
\end{array}\right).
$$
\end{example}

A basic fact is that two matrices~$A$ and~$B$ have the same reduced
row echelon form if and only if there is an invertible matrix~$E$ such
that $EA = B$.  Also, many standard operations in linear algebra,
e.g., computation of the kernel of a linear map, intersection of
subspaces, membership checking, etc., can be encoded as a question
about computing the echelon form of a matrix.


The following standard algorithm computes the echelon form of a
matrix.
\begin{algorithm}{Gauss Elimination}\label{alg:gauss}%
Given an  $m\times n$ matrix $A$ over a field,
the algorithm outputs the reduced row echelon form of~$A$.
Write $a_{i,j}$ for the $i,j$ entry of $A$,
where $0\leq i \leq m-1$ and $0\leq j \leq n-1$.
\begin{steps}
  \item{} [Initialize] Set $k=0$.
  \item{} [Clear Each Column] For each column $c=0,1,\ldots, n-1$,
     clear the $c$th column as follows:
     \begin{steps}
       \item{} [First Nonzero] Find the smallest $r$ such that
       $a_{r,c}\neq 0$, or if there is no such $r$, go to the next column. 
       \item{} [Rescale] Replace row $r$ of $A$ by $\frac{1}{a_{r,c}}$ times row $r$.
       \item{} [Swap] Swap row $r$ with row $k$.
       \item{} [Clear] For each $i=0,\ldots, m-1$ with $i\neq k$,
         if $a_{i,c}\neq 0$, add $-a_{i,c}$ times row $k$ of $A$ to 
         row $i$ to clear the leading entry of the $i$th row. 
       \item{} [Increment] Set $k=k+1$.
    \end{steps}
\end{steps}
\end{algorithm}
This algorithm takes $O(mn^2)$ arithmetic operations in the base field,
where~$A$ is an $m\times n$ matrix.  If the base field is $\Q$, the
entries can become huge and arithmetic operations are then
very expensive.  See Section~\ref{sec:modularmethod} for ways
to mitigate this problem.

To conclude this section, we mention how to convert a few standard
problems into questions about reduced row echelon forms of matrices.
Note that one can also phrase some of these answers in terms of the
echelon form, which might be easier to compute, or an LUP
decomposition (lower triangular times upper triangular times
permutation matrix), which the numerical analysts use.
\begin{enumerate}
\item{} {\bf Kernel of $A$:}\label{alg:ker} We explain how to compute
  the kernel of $A$ acting on column vectors from the right (first
  transpose to obtain the kernel of $A$ acting on row vectors).  Since
  passing to the reduced row echelon form of $A$ is the same as
  multiplying on the left by an invertible matrix, the kernel of the
  reduce row echelon form $E$ of $A$ is the same as the kernel of $A$.
  There is a basis vector of $\ker(E)$ that corresponds to each
  nonpivot column of $E$.  That vector has a~$1$ at the nonpivot
  column, $0$'s at all other nonpivot columns, and for each pivot
  column, the negative of the entry of~$A$ at the nonpivot column in
  the row with that pivot element.

\item{} {\bf Intersection of Subspaces:} Suppose $W_1$ and $W_2$ are
  subspace of a finite-dimensional vector space $V$.  Let $A_1$ and
  $A_2$ be matrices whose columns form a basis for $W_1$ and $W_2$,
  respectively.  Let $A=[A_1 | A_2]$ be the augmented matrix formed
  from $A_1$ and $A_2$.  Let $K$ be the kernel of the linear
  transformation defined by $A$.  Then $K$ is isomorphic to the desired
  intersection.  To write down the intersection explicitly, suppose
  that $\dim(W_1) \leq \dim(W_2)$ and do the following: For each $b$ in a
  basis for~$K$, write down the linear combination of a basis for $W_1$
  obtained by taking the first $\dim(W_1)$ entries of the vector $b$.  The
  fact that $b$ is in $\Ker(A)$ implies that the vector we just
  wrote down is also in $W_2$.  
This is because a linear relation
  $$\sum a_i w_{1,i}  + \sum b_j w_{2,j} = 0,$$
i.e., an element of that kernel, is the same
as 
  $$\sum a_i w_{1,i} = \sum -b_j w_{2,j}.$$
For more details, see \cite[Alg.~2.3.9]{cohen:course_ant}.
\end{enumerate}

\section{Rational Reconstruction}
Rational reconstruction is a process that allows one to sometimes lift
an integer modulo~$m$ uniquely to a bounded rational number.

\begin{algorithm}{Rational Reconstruction}\label{alg:ratrecon}%
  Given an integer $a\geq 0$ and an integer $m> 1$, this algorithm
  computes the numerator $n$ and denominator $d$ of the unique
  rational number $n/d$, if it exists, with
\begin{equation}\label{eqn:rr}
  |n|, d \leq \sqrt{\frac{m}{2}}
\qquad\text{and}\qquad n \con a d \pmod{m},
\end{equation}
or it reports that there is no such number.
\begin{steps}
\item{}[Reduce mod $m$] Replace $a$ with the least integer between $0$
  and $m-1$ that is congruent to $a$ modulo $m$.
\item{} [Trivial Cases] If $a=0$ or $a=1$, return $a$.
\item{} [Initialize] Let $b=\sqrt{m/2}$, $u=m$, $v=a$, and set
   $U=(1,0,u)$ and $V=(0,1,v)$.  Use the notation $U_i$ and $V_i$
to refer to the $i$th entries of $U, V$, for $i=0,1,2$.
\item{} [Iterate] Do the following as long as $|V_2|>b$:
Set $q = \lfloor U_2 / V_2 \rfloor$, 
set $T = U - q V$, and set $U=V$ and $V=T$.
\item{} [Numerator and Denominator] Set $d=|V_1|$ and $n=V_2$.
\item{} [Good?] If $d\leq b$ and $\gcd(n,d)=1$, return $n/d$;
otherwise report that there is no rational number as 
in \eqref{eqn:rr}.
\end{steps}
\end{algorithm}


Algorithm~\ref{alg:ratrecon} for rational reconstruction is described
(with proof) in \cite[pgs. 656--657]{knuth2} as
the solution to Exercise~51 on page 379 in that book.  See, in particular, the
paragraph right in the middle of page 657, which describes the
algorithm.  Knuth attributes this rational reconstruction algorithm to Wang, Kornerup, and Gregory from around 1983.

We now give an indication of why Algorithm~\ref{alg:ratrecon} computes
the rational reconstruction of $a \pmod{m}$, leaving the precise
details and uniqueness to \cite[pgs. 656--657]{knuth2}.  At each
step in Algorithm~\ref{alg:ratrecon}, the $3$-tuple $V=(v_0, v_1,
v_2)$ satisfies
\begin{equation}\label{eqn:xgcdeqn}
   m \cdot v_0 + a\cdot v_1 = v_2,
 \end{equation}
and similarly for $U$.
When computing the usual extended $\gcd$, at the end $v_2=\gcd(a,m)$
and~$v_0$,~$v_1$ give a representation of the~$v_2$ as a $\Z$-linear
combination of~$m$ and~$a$.
In Algorithm~\ref{alg:ratrecon}, we are instead interested
in finding a rational number $n/d$ such that 
$n \con a\cdot d\pmod{m}$.  If we set $n=v_2$ and $d=v_1$
in \eqref{eqn:xgcdeqn} and rearrange, we obtain
$$
  n  = a \cdot d + m \cdot v_0.
$$
Thus at {\em every} step of the algorithm we find a rational number
$n/d$ such that $n\con ad\pmod{m}$.  The problem at intermediate
steps is that, e.g., $v_0$ could be~$0$, or~$n$ or~$d$
could be too large.  



%begin_sage
\begin{example}\sageindex{rational reconstruction}
We compute an example using \sage.
\begin{verbatim} 
sage: p = 389
sage: k = GF(p)
sage: a = k(7/13); a
210
sage: a.rational_reconstruction()
7/13
\end{verbatim}
\end{example}
%end_sage


\section{Echelon Forms over $\Q$}\label{sec:modularmethod}
A difficulty with computation of the echelon form of a matrix over the
rational numbers is that arithmetic with large rational numbers is
time-consuming; each addition potentially requires a $\gcd$
and numerous additions and multiplications of integers.  Moreover, the
entries of $A$ during intermediate steps of Algorithm~\ref{alg:gauss}
can be huge even though the entries of $A$ and the answer are small.
For example, suppose $A$ is an invertible square matrix.  Then the
echelon form of $A$ is the identity matrix, but during intermediate
steps the numbers involved could be quite large.  One technique for
mitigating this is to compute the echelon form using a
multimodular method.  



If~$A$ is a matrix with rational entries, let $H(A)$ be the
\defn{height} of~$A$, which is the maximum of the absolute values of
the numerators and denominators of all entries of~$A$.  If $x,y$ are
rational numbers and~$p$ is a prime, we write $x\con y\pmod{p}$ to
mean that the denominators of~$x$ and~$y$ are not divisible by~$p$ but
the numerator of the rational number $x-y$ (in reduced form) is
divisible by~$p$.  For example, if $x=5/7$ and $y=2/11$, then
$x-y=41/77$, so $x\con y\pmod{41}$.


\begin{algorithm}{Multimodular Echelon Form}\label{alg:modech}%
Given an $m\times n$ matrix $A$ with entries in $\Q$,
this algorithm computes the reduced row echelon form of~$A$.
\begin{steps}
\item Rescale the input matrix $A$ to have integer entries.  This does
  not change the echelon form and makes reduction modulo many primes
  easier.  We may thus assume $A$ has integer entries.

\item Let $c$ be a guess for the height of the echelon form. 

\item List successive primes $p_1, p_2, \ldots$
 such that the product of the $p_i$ is greater than 
$n \cdot c \cdot H(A) + 1$, where $n$ is the number of columns of~$A$.

\item Compute the echelon forms $B_i$ of the reduction $A\pmod{p_i}$
  using, e.g., Algorithm~\ref{alg:gauss} or any other echelon
  algorithm. % (This is where most of the work takes place.)

\item\label{step:discard} Discard any $B_i$ whose pivot column list is
  not maximal among pivot lists of all $B_j$ found so far.  (The pivot
  list associated to $B_i$ is the ordered list of integers $k$ such
  that the $k$th column of $B_j$ is a pivot column.  We mean maximal
  with respect to the following ordering on integer sequences: shorter
  integer sequences are smaller, and if two sequences have the same
  length, then order in reverse lexicographic order.  Thus $[1,2]$ is
  smaller than $[1,2,3]$, and $[1,2,7]$ is smaller than $[1,2,5]$.
  Think of maximal as ``optimal'', i.e., best possible pivot columns.)

\item Use the Chinese Remainder Theorem to find a matrix $B$ with
  integer entries such that $B \con B_i \pmod{p_i}$ for all $p_i$.

\item Use Algorithm~\ref{alg:ratrecon}
to try to find a matrix~$C$ whose coefficients are rational 
numbers $n/r$ such that
$|n|, r \leq \sqrt{M/2}$, where $M=\prod p_i$, 
and $C \con B_i \pmod{p_i}$ for each prime~$p$.
If rational reconstruction fails,
compute a few more echelon forms mod the next few primes (using
the above steps)  and attempt rational reconstruction again.
Let $E$ be the matrix over~$\Q$ so obtained.
(A trick here is to keep track of denominators found
so far to avoid doing very many rational reconstructions.)

\item Compute the denominator $d$ of $E$, i.e., the smallest positive
  integer such that $dE$ has integer entries.  If 
\begin{equation}\label{eqn:modalgbound}
  H(dE) \cdot H(A)
  \cdot n < \prod p_i, 
\end{equation}
then $E$ is the reduced row echelon form of $A$.  If not,
repeat the above steps with a few more primes.  
\end{steps}
\end{algorithm}
\begin{proof}
  We prove that if \eqref{eqn:modalgbound} is satisfied, then the
  matrix~$E$ computed by the algorithm really is the reduced row
  echelon form~$R$ of~$A$.  First note that $E$ is in reduced row
  echelon form since the set of pivot columns of all matrices $B_i$
  used to construct $E$ are the same, so the pivot columns of $E$ are
  the same as those of any $B_i$ and all other entries in the $B_i$
  pivot columns are $0$, so the other entries of $E$ in the pivot
  columns are also $0$.

  Recall from the end of Section~\ref{sec:echelon_form} that a matrix
  whose columns are a basis for the kernel of~$A$ can be obtained from
  the reduced row echelon form $R$.  Let~$K$ be the matrix whose
  columns are the vectors in the kernel algorithm applied to $E$, 
so $EK=0$.  Since the reduced row echelon form is obtained by left
multiplying by an invertible matrix, for each $i$, there is 
an invertible matrix $V_i$ mod $p_i$ such that $A \con V_i B_i\pmod{p_i}$
so
$$
  A\cdot d K \con V_i B_i \cdot d K \con V_i \cdot dE \cdot K \con 0\pmod{p_i}.
$$
Since $d K$ and $A$ are integer matrices, the Chinese
remainder theorem implies that
$$
 A \cdot dK \con 0\,\,\,\left(\text{mod } \prod{p_i}\right).
$$
The integer entries $a$ of $A\cdot dK$ all satisfy
 $|a| \leq H(A)\cdot{} H(dK)\cdot{} n$,
where $n$ is the number of columns of $A$.  Since $H(K) \leq H(E)$,
the bound \eqref{eqn:modalgbound} implies that $A \cdot dK = 0$.  Thus
$A K = 0$, so $\Ker(E) \subset \Ker(A)$.  On the other hand,
the rank of~$E$ equals the rank of each $B_i$ (since the pivot columns
are the same), so $$\rank(E) = \rank(B_i) = \rank(A\!\!\!\pmod{p_i}) \leq
\rank(A).$$  Thus $\dim(\Ker(A)) \leq \dim(\Ker(E))$, and combining
this with the bound obtained above, we see that 
$\Ker(E) = \Ker(A)$.   This implies that $E$ is the reduced row
echelon form of~$A$, since two matrices have the same kernel if
and only if they have the same reduced row echelon form 
(the echelon form is an invariant of the row space, and the
kernel is the orthogonal complement of the row space).

The reason for step~(\ref{step:discard}) is that the matrices $B_i$ need
{\em not} be the reduction of $R$ modulo~$p_i$, and indeed this
reduction might not even be defined, e.g., if~$p_i$ divides the
denominator of some element of $R$, then this reduction makes no
sense.  For example, set $p=p_i$ and suppose $A=\abcd{p}{1}{0}{0}$.
Then $R=\abcd{1}{1/p}{0}{0}$, which has no reduction modulo~$p$; also,
the reduction of $A$ modulo~$B_i$ is $B_i = \abcd{0}{1}{0}{0}
\pmod{p}$, which is already in reduced row echelon form.  However if
we were to combine $B_i$ with the echelon form of $A$ modulo another
prime, the result could never be lifted using rational reconstruction.
Thus the reason we exclude all $B_i$ with nonmaximal pivot column
sequence is so that a rational reconstruction will exist.  There are
only finitely many primes that divide denominators of entries of $R$,
so eventually all $B_i$ will have maximal pivot column sequences,
i.e., they are the reduction of the true reduced row echelon form $R$, so
the algorithm terminates.
\end{proof}

\begin{remark}




 Algorithm~\ref{alg:modech}, with {\em sparse} matrices seems to
  work very well in practice.  A simple but helpful
  modification to Algorithm~\ref{alg:gauss} in the sparse case is to
  clear each column using a row with a minimal number of nonzero
  entries, so as to reduce the amount of ``fill in'' (denseness) of
  the matrix.  There are much more sophisticated methods along these
  lines called ``intelligent Gauss elimination''.  (Cryptographers are
  interested in linear algebra mod $p$ with huge sparse matrices,
  since they come up in attacks on the discrete log problem and
  integer factorization.)
\end{remark}


One can adapt Algorithm~\ref{alg:modech} to computation of 
echelon forms of matrices $A$ over cyclotomic fields $\Q(\zeta_n)$.
Assume~$A$ has denominator~$1$. Let~$p$ be a prime that splits
completely in $\Q(\zeta_n)$.  Compute the homomorphisms $f_i :
\Z_p[\zeta_n] \to \F_p$ by finding the elements of order~$n$ in
$\F_p^*$.  Then compute the mod~$p$ matrix $f_i(A)$ for each~$i$, and
find its reduced row echelon form.  Taken together, the maps $f_i$
together induce an isomorphism $\Psi:\F_p[X]/\Phi_n(X) \isom \F_p^d$,
where $\Phi_n(X)$ is the $n$th cyclotomic polynomial and $d$ is its
degree.  It is easy to compute $\Psi(f(x))$ by evaluating $f(x)$ at
each element of order~$n$ in $\F_p$.  To compute $\Psi^{-1}$, simply
use linear algebra over $\F_p$ to invert a matrix that represents
$\Psi$.  Use $\Psi^{-1}$ to compute the reduced row echelon form of
$A\pmod{p}$, where~$(p)$ is the nonprime ideal in $\Z[\zeta_n]$
generated by~$p$.  Do this for several primes $p$, and use rational
reconstruction on each coefficient of each power of $\zeta_n$, to
recover the echelon form of~$A$. 

\section{Echelon Forms via Matrix Multiplication}
\label{sec:echmatmul}

In this section we explain
how to compute echelon forms using matrix multiplication.  This is
valuable because there are asymptotically fast, i.e., better than
$O(n^3)$ field operations, algorithms for matrix multiplication, and
implementations of linear algebra libraries often include
highly optimized matrix multiplication algorithms.
We only sketch the basic ideas 
behind these asymptotically fast algorithms (following \cite{steel:talk}), since more 
detail would take us too far from modular forms.  

The naive algorithm for multiplying two $m\times m$ matrices requires
$O(m^3)$ arithmetic operations in the base ring.  In
\cite{strassen:gauss}, Strassen described a clever algorithm that
computes the product of two $m\times m$ matrices in $O(m^{\log_2(7)})
= O(m^{2.807\ldots})$ arithmetic operations in the base ring.  Because
of numerical stability issues, Strassen's algorithm is rarely used in
numerical analysis.  But for matrix arithmetic over exact base rings
(e.g., the rational numbers, finite fields, etc.) it is of extreme
importance.

In \cite{strassen:gauss}, Strassen also sketched a new algorithm for
computing the inverse of a square matrix using matrix multiplication.
Using this algorithm, the number of operations to invert an $m\times m$ 
matrix is (roughly) the same as the number needed to multiply two
$m\times m$ matrices.
Suppose the input matrix is $2^n\times 2^n$ and we
write it in block form as $\abcd{A}{B}{C}{D}$ where $A,B,C,D$ are all
$2^{n-1}\times 2^{n-1}$ matrices.  Assume that any intermediate
matrices below that we invert are invertible.  Consider the augmented
matrix
$$
\left(\begin{array}{cc|cc}
A&B&I&0\\
C&D&0&I
\end{array}\right).
$$
Multiply the top row by $A^{-1}$ to obtain
$$
\left(\begin{array}{cc|cc}
I&A^{-1}B&A^{-1}&0\\
C&D&0&I
\end{array}\right),
$$
and write $E=A^{-1}B$.
Subtract $C$ times the first row from the second row to get
$$
\left(\begin{array}{cc|cc}
I&E&A^{-1}&0\\
0&D - CE &-CA^{-1}&I
\end{array}\right).
$$
Set $F=D-CE$ and multiply
the bottom row by $F^{-1}$ on the left to obtain
$$
\left(\begin{array}{cc|cc}
I&E&A^{-1}&0\\
0&I &-F^{-1}CA^{-1}&F^{-1}
\end{array}\right).
$$
Set $G=-F^{-1} C A^{-1}$, and subtract~$E$ 
times the second from the first row to
arrive at 
$$
\left(\begin{array}{cc|cc}
I&0&A^{-1}-EG&-EF^{-1}\\
0&I &G & F^{-1}
\end{array}\right).
$$

The idea listed above can, with significant work,
be extended to a general algorithm (as is done in \cite{sage}). 



%This algorithm can be extended to matrices $A$ of any size by padding
%the matrices by $1$ in the lower right to make $A$ a $2^n\times 2^n$
%matrix.  If one of the intermediate matrices is not invertible, either
%attempt to apply
%a random permutation of the rows of $A$ then retry the algorithm 
%or fall back to a standard matrix inverse algorithm. 
%Applying a random permutation of the rows is the same as multiplying
%$A$ on the left by an invertible matrix $E$.  If we compute
%$B=(EA)^{-1}$ we recover $A^{-1}$ by noting that $B=A^{-1}E^{-1}$, so
%$A^{-1} = BE$.  I.e., to obtain $A^{-1}$ we permute the {\em columns}
%of $B$ using the permutation that we applied to the rows of $A$.

%\begin{algorithm}{Inverse Using Matrix Multiplication}\label{alg:matinv}
%  Given an $m \times m$ matrix $M$ with $m=2^n$
%  this algorithm computes the inverse of $M$.
%\begin{steps}
%\item{}[Submatrices]\label{alg:invmatmul}
% Write $M=\abcd{A}{B}{C}{D}$ with each of $A,B,C,D$
%a $2^{n-1}\times 2^{n-1}$ matrix. 
%\item{}[Multiply and Invert] Compute each of $E=A^{-1}B$, $F=D-CE$
%  and $G=-F^{-1}CA^{-1}$.  The inverses involve $2^{n-1}\times
%  2^{n-1}$ and we do them recursively using this algorithm, except if
%  $n=2$. If either $A$ or $F$ is not invertible choose a random
%  permutation $\sigma$ of $2^n$, apply this permutation to the rows of
%  $M$, and go to step~(\ref{alg:invmatmul}) (or fall back to a standard
%matrix inverse algorithm).
%\item{}[Compute Inverse] Compute  
%    $$M' = \mtwo{A^{-1} - EG}{-EF^{-1}}{G}{F^{-1}}.$$
%\item{}[Apply Permutation] Apply the permutation $\sigma$ to 
%the columns of $M'$.  Output the resulting matrix, which is 
%$M^{-1}$.  (If $m\neq 2^n$ output only the upper left $m\times m$
%submatrix of this matrix.)
%\end{steps}
%\end{algorithm}

Next we very briefly sketch how to compute echelon forms of matrices using matrix
multiplication and inversion.  Its complexity is comparable
to the complexity of matrix multiplication. 

As motivation, recall the standard algorithm from undergraduate linear
algebra for inverting an invertible square matrix $A$: form the
augmented matrix $[A|I]$, and then compute the echelon form of this
matrix, which is $[I|A^{-1}]$. If $T$ is the transformation matrix to
echelon form, then $T[A|I] = [I|T]$, so $T=A^{-1}$.  In particular, we
could find the echelon form of $[A|I]$ by multiplying on the left by
$A^{-1}$.  Likewise, for any matrix $B$ with the same number of rows
as $A$, we could find the echelon form of $[A|B]$ by multiplying on
the left by $A^{-1}$. Next we extend this idea to give
an algorithm to compute echelon forms using only matrix multiplication (and
echelon form modulo one prime).

\begin{algorithm}{Asymptotically Fast Echelon Form}\label{alg:fe}
  Given a matrix $A$ over the rational numbers (or a number field),
  this algorithm computes the echelon form of $A$.
\begin{steps}
\item{}[Find Pivots]\label{alg:fepivots} Choose a random prime $p$ (coprime to the
  denominator of any entry of $A$) and compute the echelon form of
  $A\pmod{p}$, e.g., using Algorithm~\ref{alg:gauss}.  Let $c_0, \ldots,
  c_{n-1}$ be the pivot columns of $A\pmod{p}$.  When computing the
  echelon form, save the positions $r_0,\ldots, r_{n-1}$ of the rows
  used to clear each column. 
\item{}[Extract Submatrix] Extract the $n\times n$ submatrix
  $B$ of $A$ whose entries are $A_{r_i,c_j}$ for $0\leq i,j \leq {n-1}$.
\item{}[Compute Inverse] 
  Compute the inverse $B^{-1}$ of $B$.  Note that $B$ must be
  invertible since its reduction modulo $p$ is invertible.
\item{}[Multiply] Let $C$ be the matrix whose rows are the
rows $r_0,\ldots, r_{n-1}$ of $A$.   Compute $E=B^{-1}C$.
If $E$ is not in echelon form, go to step~(\ref{alg:fepivots}).
\item{}[Done?]\label{alg:fedone} 
 Write down a matrix $D$ whose columns are a basis for
  $\ker(E)$ as explained on page~\pageref{alg:ker}.  Let $F$ be the
  matrix whose rows are the rows of $A$ other than rows $r_0,\ldots,
  r_{n-1}$.  Compute the product $FD$.  If $FD=0$, output $E$, which is
  the echelon form of $A$.  If $FD\neq 0$, go to
  step~(\ref{alg:fepivots}) and run the whole algorithm again.
\end{steps}
\end{algorithm}
\begin{proof}
  We prove both that the algorithm terminates and that when it
  terminates, the matrix~$E$ is the echelon form of~$A$.  

  First we prove that the algorithm terminates.  Let~$E$ be the
  echelon form of~$A$. 
  By \exref{ch:linalg}{ex:rrefsmodp}, 
  for all but finitely many primes~$p$ (i.e., any
  prime where $A\pmod{p}$ has the same rank as~$A$) the echelon form of
  $A\pmod{p}$ equals $E\pmod{p}$.  For any such prime~$p$ the pivot
  columns of $E\pmod{p}$ are the pivot columns of~$E$, so the
  algorithm will terminate for that choice of~$p$.

  We next prove that when the algorithm terminates, $E$ is the
  echelon form of $A$.  By assumption, $E$ is in echelon form and is
  obtained by multiplying $C$ on the left by an invertible matrix,
  so~$E$ must be {\em the} echelon form of~$C$.  The rows of~$C$ are a
  subset of those of~$A$, so the rows of~$E$ are a subset of the rows
  of the echelon form of~$A$.  Thus $\ker(A) \subset \ker(E)$.  To
  show that~$E$ equals the echelon form of~$A$, we just need to verify
  that $\ker(E)\subset \ker(A)$, i.e., that $AD=0$, where~$D$ is as in
  step~(\ref{alg:fedone}).  Since~$E$ is the echelon form of~$C$, we
  know that $CD=0$.  By step~(\ref{alg:fedone}) we also know that
  $FD=0$.  Thus $AD=0$, since the rows of~$A$ are the union of the
  rows of~$F$ and~$C$.

\end{proof}



\begin{example}
Let $A$ be the $4\times 8$ matrix 
$$
A =\left(\begin{array}{rrrrrrrr}
-9&6&7&3&1&0&0&0\\
-10&3&8&2&0&1&0&0\\
3&-6&2&8&0&0&1&0\\
-8&-6&-8&6&0&0&0&1
\end{array}\right)
$$
from Example~\ref{ex:rref1}.
%begin_sage   % this fails the doctest, but that's ok.
\begin{verbatim}
sage: M = MatrixSpace(QQ,4,8)
sage: A = M([[-9,6,7,3,1,0,0,0],[-10,3,8,2,0,1,0,0],
             [3,-6,2,8,0,0,1,0],[-8,-6,-8,6,0,0,0,1]])
\end{verbatim}%link

%end_sage
First choose the ``random'' prime $p=41$, which does not
divide any of the entries of $A$, and compute the
echelon form of the reduction of $A$ modulo $41$.\sageindex{echelon form}
%begin_sage
%link
\begin{verbatim}
sage: A41 = MatrixSpace(GF(41),4,8)(A)
sage: E41 = A41.echelon_form()
\end{verbatim}%link

%end_sage
The echelon form of $A\pmod{41}$ is 
$$
\left(\begin{array}{rrrrrrrr}
1&0&0&2&0&20&33&18\\
0&1&0&40&0&30&7&1\\
0&0&1&39&0&19&13&17\\
0&0&0&0&1&31&0&37
\end{array}\right).
$$
Thus we take $c_0=0$, $c_1=1$, $c_2=2$, and $c_3 = 4$.
Also $r_i=i$ for $i=0,1,2,3$.  
%begin_sage
Next extract the submatrix $B$.
%link
\begin{verbatim}
sage: B = A.matrix_from_columns([0,1,2,4])
\end{verbatim}%link

%end_sage
The submatrix $B$ is 
$$
  B = \left(\begin{array}{rrrr}
-9&6&7&1\\
-10&3&8&0\\
3&-6&2&0\\
-8&-6&-8&0
\end{array}\right).
$$
The inverse of $B$ is 
$$
 B^{-1} = \left(\begin{array}{rrrr}
0&-\frac{5}{92}&\frac{1}{46}&-\frac{9}{184}\vspace{2pt}\\
0&-\frac{1}{138}&-\frac{3}{23}&-\frac{11}{276}\vspace{2pt}\\
0&\frac{11}{184}&\frac{7}{92}&-\frac{17}{368}\vspace{2pt}\\
1&-\frac{159}{184}&\frac{41}{92}&\frac{45}{368}
\end{array}\right).
$$
Multiplying by $A$ yields
$$
E = B^{-1} A = 
\left(\begin{array}{rrrrrrrr}
1&0&0&-\frac{21}{92}&0&-\frac{5}{92}&\frac{1}{46}&-\frac{9}{184}\vspace{2pt}\\
0&1&0&-\frac{179}{138}&0&-\frac{1}{138}&-\frac{3}{23}&-\frac{11}{276}\vspace{2pt}\\
0&0&1&\frac{83}{184}&0&\frac{11}{184}&\frac{7}{92}&-\frac{17}{368}\vspace{2pt}\\
0&0&0&\frac{1025}{184}&1&-\frac{159}{184}&\frac{41}{92}&\frac{45}{368}
\end{array}\right).
$$
%begin_sage
%link
\begin{verbatim}
sage: E = B^(-1)*A
\end{verbatim}%link

%end_sage
This is {\em not} the echelon form of $A$.  Indeed, it is not
even in echelon form, since the last row is not
normalized so the leftmost nonzero entry is $1$.
We thus choose another random prime, say $p=43$.
The echelon form mod $43$ has columns $0,1,2,3$ as pivot columns.
We thus extract the matrix
$$
 B = \left(\begin{array}{rrrr}
-9&6&7&3\\
-10&3&8&2\\
3&-6&2&8\\
-8&-6&-8&6
\end{array}\right).
$$
%begin_sage
%link
\begin{verbatim}
sage: B = A.matrix_from_columns([0,1,2,3])
\end{verbatim}
%end_sage

This matrix has inverse
$$
B^{-1} = 
\left(\begin{array}{rrrr}
\frac{42}{1025}&-\frac{92}{1025}&\frac{1}{25}&-\frac{9}{205}\vspace{2pt}\\
\frac{716}{3075}&-\frac{641}{3075}&-\frac{2}{75}&-\frac{7}{615}\vspace{2pt}\\
-\frac{83}{1025}&\frac{133}{1025}&\frac{1}{25}&-\frac{23}{410}\vspace{2pt}\\
\frac{184}{1025}&-\frac{159}{1025}&\frac{2}{25}&\frac{9}{410}
\end{array}\right).
$$
Finally, the echelon form of $A$ is 
$
 E = B^{-1}A.
$
No further checking is needed since the product so
obtained is in echelon form, and the matrix $F$
of the last step of the algorithm has $0$ rows.

\end{example}

\begin{remark}
Above we have given only the barest sketch of asymptotically fast
``block'' algorithms for linear algebra.  For optimized algorithms
that work in the general case, please see the source code of \cite{sage}.
\end{remark}


\section{Decomposing Spaces under the Action of Matrix}
\label{sec:decompmodsym}

Efficiently solving the following problem is a crucial step in
computing a basis of eigenforms for a space of modular forms 
(see Sections~\ref{sec:computingmodform} and \ref{sec:syseigen}).

\begin{problem}\label{prob:decomp}
  Suppose $T$ is an $n\times n$ matrix with entries in a field~$K$
  (typically a number field or finite field) and that the minimal
  polynomial of~$T$ is square-free and has degree~$n$.  View $T$ as
  acting on $V=K^n$.  Find a simple module decomposition $W_0 \oplus
  \cdots \oplus W_m$ of~$V$ as a direct sum of simple $K[T]$-modules.
  Equivalently, find an invertible matrix~$A$ such that $A^{-1} T A$
  is a block direct sum of matrices $T_0, \ldots, T_m$ such that the
  minimal polynomial of each $T_i$ is irreducible.
\end{problem}

\begin{remark}\label{rem:rjf}
A generalization of Problem~\ref{prob:decomp} 
to arbitrary matrices with entries in $\Q$ is
finding the \defn{rational Jordan form} (or rational canonical
form, or Frobenius form) of $T$.
This is like the usual Jordan form, but the resulting matrix is
over~$\Q$ and the summands of the matrix corresponding to
eigenvalues are replaced by matrices whose minimal polynomials
are the minimal polynomials (over~$\Q$) of the eigenvalues.  The
rational Jordan form was extensively studied by Giesbrecht in his
Ph.D. thesis and many successive papers, where he analyzes the
complexity of his algorithms and observes that the limiting step is
factoring polynomials over~$K$.  The reason is that given a polynomial
$f\in K[x]$, one can easily write down a matrix $T$ such that one can
read off the factorization of~$f$ from the rational Jordan form of
$T$ (see also \cite{steel:canonical}).
\end{remark}

\subsection{Characteristic Polynomials}
The computation of characteristic polynomials of matrices is crucial to
modular forms computations.  There are many approaches to this
problems: compute $\det(xI-A)$ symbolically (bad), compute the traces
of the powers of $A$ (bad), or compute the Hessenberg form modulo many
primes and use CRT (bad; see for \cite[\S2.2.4]{cohen:course_ant} the
definition of Hessenberg form and the algorithm).  A more
sophisticated method is to compute the rational canonical form of $A$
using Giesbrecht's algorithm\footnote{Allan Steel also
  invented a similar algorithm.} (see \cite{giesbrecht-storjohann}),
which involves computing Krylov subspaces\index{Krylov subspace} 
(see Remark~\ref{rem:krylov} below), and building up the whole
space on which $A$ acts.  This latter method is a generalization of
Wiedemann's algorithm for computing minimal polynomials (see
Section~\ref{sec:wiedemann}), but with more structure to handle the
case when the characteristic polynomial is not equal to the minimal
polynomial.

\subsection{Polynomial Factorization}\label{sec:polyfac}
Factorization of polynomials in $\Q[X]$ (or over number fields) is an
important step in computing an explicit basis of Hecke eigenforms for
spaces of modular forms.  The best algorithm is the van Hoeij method
\cite{hoeij:nf},
which uses the LLL lattice basis reduction algorithm \cite{lll} in a
novel way to solve the optimization problems that come up in trying to
lift factorizations mod~$p$ to~$\Z$. 
It has been generalized by Belebas, van Hoeij, Kl\"uners, and Steel
to number fields.




\subsection{Wiedemann's Minimal Polynomial Algorithm}
\label{sec:wiedemann}
In this section we describe 
an algorithm due to Wiedemann for computing the
minimal polynomial of an $n\times n$ matrix $A$ over a field. 

Choose a random vector $v$ and compute the iterates
\begin{equation}\label{eqn:iterates}
   v_0=v, \quad v_1 = A(v),\quad v_2 = A^2(v),\quad \ldots, \quad v_{2n-1} = A^{2n-1}(v).
\end{equation}
If $f=x^m + c_{m-1}x^{m-1} + \cdots + c_1 x + c_0 $ 
is the minimal polynomial of $A$, then 
$$
   A^m + c_{m-1} A^{m-1} + \cdots + c_0 I_n = 0,
$$
where $I_n$ is the $n\times n$ identity matrix. 
For any $k\geq 0$, by
multiplying both sides on the right by the vector $A^kv$, we see that
$$
   A^{m+k} v + c_{m-1} A^{m-1+k}v + \cdots + c_0 A^k v = 0;
$$
hence
$$
   v_{m+k} + c_{m-1} v_{m-1+k} + \cdots + c_0 v_k = 0, \qquad \text{all $k\geq 0$}.
$$

Wiedemann's idea is to observe that any single component of the
vectors $v_0, \ldots, v_{2n-1}$ satisfies the linear recurrence with
coefficients $1, c_{m-1}, \ldots, c_0$.  The Berlekamp-Massey
algorithm (see Algorithm~\ref{alg:bm} below) was introduced in the
1960s in the context of coding theory to find the minimal polynomial
of a linear recurrence sequence $\{a_r\}$.  The minimal polynomial of
this linear recurrence is by definition the unique monic
polynomial~$g$, such that if $\{a_r\}$ satisfies a linear recurrence
$a_{j+k} + b_{j-1} a_{j-1+k} + \cdots + b_0 a_k=0$ (for all $k\geq
0$), then $g$ divides the polynomial $x^j + \sum_{i=0}^{j-1} b_i x^i$.
If we apply Berlekamp-Massey to the top coordinates of the $v_i$, we
obtain a polynomial $g_0$, which divides $f$.  We then apply it to the
second to the top coordinates and find a polynomial $g_1$ that 
divides~$f$, etc. Taking the least common multiple of the first few $g_i$, we
find a divisor of the minimal polynomial of~$f$.  One can show that
with ``high probability'' one quickly finds~$f$, instead of just a
proper divisor of~$f$.

\begin{remark}\label{rem:krylov}
In the literature, techniques that involve iterating a vector 
as in \eqref{eqn:iterates} are 
often called \defn{Krylov methods}.  The subspace generated by the iterates
of a vector under a matrix is called a \defn{Krylov subspace}.
\end{remark}


\begin{algorithm}{Berlekamp-Massey}\label{alg:bm}%
Suppose $a_{0},\ldots , a_{2n-1}$ are the first $2n$ terms of a
sequence that satisfies a linear recurrence of degree at most $n$.
This algorithm computes the minimal 
polynomial $f$ of the sequence.
\begin{steps}
\item Let $R_0 = x^{2n}$, $R_1 = \sum_{i=0}^{2n-1} a_i x^i$, 
$V_0 = 0$, $V_1 = 1$.
\item While $\deg(R_1) \geq n$, do the following:
\begin{enumerate}
\item Compute $Q$ and $R$ such that 
$R_0 = Q R_1 + R$.
\item Let $(V_0, V_1, R_0, R_1) = (V_1, V_0 - QV_1, R_1, R)$.
\end{enumerate}
\item Let $d = \max(\deg(V_1), 1+\deg(R_1))$ and set
$P = x^d V_1(1/x)$.  
\item Let $c$ be the leading coefficient of $P$ and output
      $f = P/c$.
\end{steps}
\end{algorithm}
The above description of Berlekamp-Massey is taken
from \cite{attitoca:berlekamp}, which contains 
some additional ideas for improvements.


Now suppose $T$ is an $n\times n$ matrix as in
Problem~\ref{prob:decomp}.  We find the minimal polynomial of~$T$ by
computing the minimal polynomial of $T\!\!\pmod{p}$ using Wiedemann's
algorithm, for many primes~$p$, and using the Chinese Remainder
Theorem.  (One has to bound the number of primes that must be
considered; see, e.g., \cite{cohen:course_ant}.)

One can also compute the characteristic polynomial of~$T$ directly
from the Hessenberg form of~$T$, which can be computed in $O(n^4)$
field operations, as described in \cite{cohen:course_ant}.  This is
simple but slow.  Also, the~$T$ we consider will often be sparse,
and Wiedemann is particularly good when~$T$ is sparse.


\begin{example}
We compute the minimal polynomial of 
$$
A= \left(\begin{array}{ccc}
3&0&0\\
0&0&2\\
-1&1/2&-1
\end{array}\right)
$$
using Wiedemann's algorithm.
Let $v=(1,0,0)^t$.  Then 
\begin{align*}
  v&=(1,0,0)^t,\quad Av = (3,0,-1)^t,\quad A^2v = (9,-2,-2)^t, \\
  A^3v&=(27,-4,-8)^t, \quad A^4v=(81,-16,-21)^t, \quad A^5v=(243,-42,-68)^t.
\end{align*}
The linear recurrence sequence coming from the first entries is
$$
   1, 3, 9, 27, 81, 243.
$$
This sequence satisfies the linear recurrence
$$
    a_{k+1} -3 a_k = 0, \qquad\text{all $k>0$},
$$
so its minimal polynomial is $x-3$.  This implies that
$x-3$ divides the minimal polynomial of the matrix $A$.
Next we use the sequence of second coordinates of the 
iterates of $v$, which is
$$
  0, 0, -2, -4, -16, -42.
$$
The recurrence that this sequence satisfies is slightly 
less obvious, so we apply the Berlekamp-Massey algorithm
to find it, with $n=3$.
\begin{enumerate}
\item We have $R_0 = x^6$, 
$
 R_1 = -42x^5 - 16x^4 - 4x^3 - 2x^2,
$
$V_0 = 0, V_1 = 1$.

\item 
\begin{enumerate}
\item Dividing $R_0$ by $R_1$, we find
$$
  R_0 = R_1 \left(  -\frac{1}{42}x + \frac{4}{441}  \right) + 
    \left(\frac{22}{441}x^4 - \frac{5}{441}x^3 + \frac{8}{441}x^2 \right).
$$
\item The new $V_0, V_1, R_0, R_1$ are 
\begin{align*}
  V_0 &= 1,\\
  V_1 &= \frac{1}{42}x - \frac{4}{441},\\
  R_0 &= -42x^5 - 16x^4 - 4x^3 - 2x^2,\\
  R_1 &= \frac{22}{441}x^4 - \frac{5}{441}x^3 + \frac{8}{441}x^2.\\
\end{align*}
Since $\deg(R_1)\geq n=3$, we do the above three steps again.
\end{enumerate}

\item We repeat the above three steps.
\begin{enumerate}
\item Dividing $R_0$ by $R_1$, we find
$$
  R_0 = R_1 \left(-\frac{9261}{11}x - \frac{123921}{242}\right) + 
    \left( \frac{1323}{242}x^3 + \frac{882}{121}x^2\right).
$$
\item The new $V_0, V_1, R_0, R_1$ are:
\begin{align*}
  V_0 &= \frac{1}{42}x - \frac{4}{441},\\
  V_1 &= \frac{441}{22}x^{2} + \frac{2205}{484}x + \frac{441}{121},\\
  R_0 &= \frac{22}{441}x^{4} - \frac{5}{441}x^{3} + \frac{8}{441}x^{2},\\
  R_1 &= \frac{1323}{242}x^{3} + \frac{882}{121}x^{2}.\\
\end{align*}
\end{enumerate}

\item We have to repeat the steps yet again:
\begin{align*}
 V_0 &= \frac{441}{22}x^{2} + \frac{2205}{484}x + \frac{441}{121},\\
 V_1 &= -\frac{242}{1323}x^{3} + \frac{968}{3969}x^{2} + 
             \frac{484}{3969}x - \frac{242}{3969},\\
 R_0 &= \frac{1323}{242}x^{3} + \frac{882}{121}x^{2},\\
 R_1 &= \frac{484}{3969}x^{2}.
\end{align*}

\item We have $d=3$, 
 so $P = -\frac{242}{3969}x^{3} + \frac{484}{3969}x^{2} + \frac{968}{3969}x - \frac{242}{1323}$.
\item Multiply through by $-3969/242$ and
output $$x^3 - 2x^2 - 4x + 3 = (x - 3)(x^2 + x - 1).$$
\end{enumerate}
The minimal polynomial of $T_2$ is $(x - 3)(x^2 + x - 1)$,
since the minimal polynomial has degree at most $3$
and is divisible by $(x - 3)(x^2 + x - 1)$.
\end{example}


\subsection{$p$-adic Nullspace}
 
We will use the following algorithm of Dixon \cite{dixon:exact} to
compute $p$-adic approximations to solutions of linear equations
over~$\Q$.  Rational reconstruction modulo $p^n$ then allows us
to recover the corresponding solutions over~$\Q$.
\begin{algorithm}{$p$-adic Nullspace}\label{alg:dixon}
Given a matrix $A$ with integer entries and nonzero kernel, this
algorithm computes a nonzero element of $\ker(A)$
using successive $p$-adic approximations.
\begin{steps}
\item{}[Prime] Choose a random prime $p$.
\item{}[Echelon] Compute the echelon form of $A$ modulo $p$.
\item{}[Done?] If $A$ has full rank modulo $p$, it has full rank,
so we terminate the algorithm.
\item{}[Setup] Let $b_0 = 0$.
\item{}[Iterate] For each $m=0,1,2,\ldots, k$, 
use the echelon form of $A$ modulo $p$ to find a vector $y_m$
with integer entries such that $A y_m \con b_m\pmod{p}$,
and then set 
  $$
    b_{m+1} = \frac{b_m - A y_m}{p}.
  $$
(If $m=0$, choose $y_m\neq 0$.)
\item{}[$p$-adic Solution]
Let $\ds x=y_0 + y_1 p + y_2 p^2 + y_3 p^3 + \cdots + y_k p^k$.
\item{}[Lift] Use rational reconstruction (Algorithm~\ref{alg:ratrecon})
to find a vector $z$ with rational entries such that
$z \con x \pmod{p^{k+1}}$, if such a vector exists.  If the
vector does not exist, increase $k$ or use a different $p$.
Otherwise, output $z$.
\end{steps}
\end{algorithm}
\begin{proof}
We verify the case $k=2$ only.
We have $A y_0 = 0\pmod{p}$ and $Ay_1 = -\frac{Ay_0}{p}\pmod{p}$.
Thus
$$
  A y_0 + p A y_1 \con A y_0 + (-Ay_0) \pmod{p^2}.
$$
\end{proof}


\subsection{Decomposition Using Kernels}
We now know enough to give an algorithm to solve 
Problem~\ref{prob:decomp}.
\begin{algorithm}{Decomposition Using Kernels}\label{alg:decomp}%
Given an $n\times n$ matrix~$T$ over a field~$K$ as
in Problem~\ref{prob:decomp},
this algorithm computes the 
decomposition of $V$ as a direct sum of 
simple $K[T]$ modules.
\begin{steps}
\item{}[Minimal Polynomial] Compute the minimal polynomial~$f$ of~$T$,
e.g., using the multimodular Wiedemann algorithm.
\item{}[Factorization]\label{alg:decomp:factor} Factor~$f$ using the 
algorithm in Section~\ref{sec:polyfac}.
\item{}[Compute Kernels]\label{alg:decomp:ker}
For each irreducible factor $g_i$ of $f$, compute
the following.
\begin{enumerate}
\item Compute the matrix $A_i = g_i(T)$.  
\item Compute $W_i = \ker(A_i)$, e.g., using Algorithm~\ref{alg:dixon}.
\end{enumerate}
\item{}[Output Answer] Then $V=\bigoplus W_i$.
\end{steps}
\end{algorithm}

\begin{remark}
As mentioned in Remark~\ref{rem:rjf}, 
if one can compute such decompositions $V=\bigoplus W_i$, then
one can easily factor polynomials~$f$; hence the difficulty of
polynomial factorization
is a lower bound on the complexity of writing~$V$
as a direct sum of simples. 
\end{remark}






\begin{exercises}
\item \label{ex:matrixwithkernel}  % written by D. Joyner
Given a subspace $W$ of $k^n$, where $k$ is a field and
$n\geq 0$ is an integer, give an algorithm to find
a matrix~$A$ such that $W = \Ker(A)$.

\item \label{ex:rrefmod} 
If $\rref(A)$ denotes the row reduced echelon form of $A$
and $p$ is a prime not dividing any denominator
of any entry of $A$ or of $\rref(A)$, is 
$\rref(A \pmod{p}) = \rref(A) \pmod{p}$?

\item \label{ex:rrefsmodp}
Let $A$ be a matrix with entries in $\Q$.  Prove that
for all but finitely many primes~$p$ we have
$\rref(A \pmod{p}) = \rref(A) \pmod{p}$.

\item \label{ex:linalg1to9} Let 
$$A = \left(\begin{array}{rrr}
1&2&3\\
4&5&6\\
7&8&9
\end{array}\right).$$
\begin{enumerate}
\item Compute the echelon form of $A$
using each of Algorithm~\ref{alg:gauss}, Algorithm~\ref{alg:modech},
and Algorithm~\ref{alg:fe}. 
\item Compute the kernel of $A$.
\item Find the characteristic polynomial of $A$
using the algorithm of Section~\ref{sec:wiedemann}.
\end{enumerate}


\item \label{ex:hnf} The notion of echelon form extends to
  matrices whose entries come from certain rings other than fields,
  e.g., Euclidean domains.  In the case of 
  matrices over $\Z$ we define a matrix to be in echelon
form (or \defn{Hermite normal form})  if it satisfies
\begin{itemize}
  \item $a_{ij}=0$, for $i>j$,
  \item $a_{ii}\geq 0$,
  \item $a_{ij}<a_{ii}$ for all $j<i$ (unless $a_{ii}=0$, in 
which case all $a_{ij}=0$).
\end{itemize}
There are algorithms for computing with finitely generated
modules over $\Z$ that are analogous to the ones in this chapter
for vector spaces, which depend on computation of Hermite 
forms. 
\begin{enumerate}
  \item Show that the Hermite form of 
$\left(\begin{matrix}1&2&3\\4&5&6\\7&8&9\end{matrix}\right)$ is 
$\left(\begin{matrix}1&2&3\\0&3&6\\0&0&0\end{matrix}\right)$.
\item Describe an algorithm for transforming
  an $n\times n$ matrix $A$ with integer entries 
into Hermite form using row
  operations and the Euclidean algorithm.
\end{enumerate}


\end{exercises}















\chapter{General Modular Symbols}\label{ch:modsym}

In this chapter we explain how to generalize the notion of modular
symbols given in Chapter~\ref{ch:modform2} to higher weight and more
general level.  We define Hecke operators on them and their relation to
modular forms via the integration pairing.  We omit many difficult 
proofs that modular symbols have certain properties and instead
focus on how to compute with modular symbols.  For more details
see the references given in this section (especially \cite{merel:1585})
and  \cite{gabor:phd}.

Modular symbols are a formalism that make it 
elementary to compute with homology or cohomology related to certain
Kuga-Sato varieties (these are $\cE\times_X \cdots \times_X \cE$,
where $X$ is a modular curve and $\cE$ is the universal elliptic curve
over it).  It is not necessary to know anything about these Kuga-Sato
varieties in order to compute with modular symbols.

This chapter is about spaces of modular symbols and how to compute
with them.  It is by far the most important chapter in this book.  The
algorithms that build on the theory in this chapter are central to all
the computations we will do later in the book.

This chapter closely follows Lo\"{i}c Merel's paper \cite{merel:1585}.
First we define modular symbols of weight $k\geq 2$.  Then we define
the corresponding Manin symbols and state a theorem of
Merel-Shokurov, which gives all relations between Manin symbols.  (The
proof of the Merel-Shokurov theorem is beyond the scope of this book
but is presented nicely in \cite{gabor:phd}.)
Next we describe how the Hecke operators act on both modular and Manin
symbols and how to compute trace and inclusion maps between spaces of
modular symbols of different levels.

Not only are modular symbols useful for computation, but they have
been used to prove theoretical results about modular forms.  For
example, certain technical calculations with modular symbols are used
in Lo\"{i}c Merel's proof of the uniform boundedness conjecture for
torsion points on elliptic curves over number fields (modular symbols
are used to understand linear independence of Hecke
operators).  Another example is \cite{grigor:phd}, which
distills hypotheses about Kato's Euler system in~$K_2$ of modular
curves to a simple formula involving modular symbols (when the
hypotheses are satisfied, one obtains a lower bound on the
Shafarevich-Tate group of an elliptic curve).

\section{Modular Symbols}\label{sec:modsymgen}
We recall from Chapter~\ref{ch:modform2}
the free abelian group $\sM_2$ of modular symbols.
We view these as elements of the relative 
homology of the extended upper half plane
$\h^* = \h\union \P^1(\Q)$ relative to the cusps.  
The group $\sM_2$
is the free abelian group on symbols $\{\alpha,\beta\}$ with 
$$\alpha, \beta \in \P^1(\Q) = \Q \union \{\infty\}$$
modulo the relations
$$
  \{\alpha,\beta\} + \{\beta,\gamma\} + \{\gamma,\alpha\} = 0,
$$
for all $\alpha,\beta,\gamma\in\P^1(\Q)$, and all torsion.
More precisely, $$\sM_2 =
(F/R)/(F/R)_{\tor},$$ where $F$ is the free abelian group on all pairs
$(\alpha,\beta)$ and $R$ is the subgroup generated by all elements of
the form $(\alpha,\beta) + (\beta,\gamma) + (\gamma,\alpha)$.
Note that $\sM_2$ is a huge free abelian group of countable rank.


For any integer $n\geq 0$, let
$\Z[X,Y]_n$ be the abelian group of homogeneous polynomials of
degree $n$ in two variables $X,Y$.

\begin{remark}
Note that $\Z[X,Y]_n$ is isomorphic
to $\Sym^{n}(\Z\cross \Z)$ as a group, but certain natural actions are
different. In \cite{merel:1585}, Merel
uses the notation $\Z_{n}[X,Y]$ for what we
denote by $\Z[X,Y]_{n}$.
\end{remark}
Now fix an integer $k\geq 2$.
Set
$$
     \sM_k = \Z[X,Y]_{k-2} \tensor_\Z \sM_2,
$$
which is a torsion-free abelian group whose elements are
sums of expressions of the form $X^{i}Y^{k-2-i}\tensor \{\alpha,\beta\}$.
For example, 
$$
    X^3 \tensor \{0,1/2\} -17 XY^2 \tensor \{\infty,1/7\} \in \sM_{5}.
$$

Fix a finite index subgroup $G$ of $\SL_2(\Z)$. 
Define a \defn{left action of $G$} on $\Z[X,Y]_{k-2}$ as follows.  If
$g=\abcd{a}{b}{c}{d}\in G$ and $P(X,Y)\in \Z[X,Y]_{k-2}$, let
$$(g P)(X,Y) = P(dX-bY, -cX+aY).$$
Note that if we think of $z=(X,Y)$
as a column vector, then $$(g P)(z) = P(g^{-1}z),$$
since $g^{-1} =
\abcd{\hfill{}d}{-b}{-c}{\hfill{}a}$.  The reason
for the inverse is so that this is a left action instead of a right
action, e.g., if $g,h\in G$, then 
$$
((gh)P)(z) = P((gh)^{-1}z) = P(h^{-1}g^{-1}z) =
(hP)(g^{-1}z) = (g(hP))(z).  
$$

Recall that we let $G$ act on the left on $\sM_2$ by 
$$
 g\{\alpha, \beta\} = \{g(\alpha), g(\beta)\},
$$
where $G$ acts via linear fractional transformations,
so if $g=\abcd{a}{b}{c}{d}$, then 
$$g(\alpha)=\frac{a\alpha + b}{c\alpha + d}.$$
For example, useful special cases to remember are that if $g=\abcd{a}{b}{c}{d}$, then
$$
g(0) = \frac{b}{d} \qquad\text{and}\qquad g(\infty) = \frac{a}{c}.
$$
(Here we view $\infty$ as $1/0$ in order to describe the
action.)

We now combine these two actions to obtain a left action of $G$ on $\sM_{k}$, 
which is given by
$$
  g(P \tensor \{\alpha, \beta\}) = (gP) \tensor\{g(\alpha), g(\beta)\}.
$$
For example, 
\begin{align*}
  \mtwo{\hfill 1}{\hfill 2}{-2}{-3} (X^3 \tensor \{0,1/2\}) &=  
     (-3X - 2Y)^3 \tensor\left\{-\frac{2}{3},  -\frac{5}{8}\right\} \\
  &\hspace{-4em}= (-27X^3 - 54X^2Y - 36XY^2 - 8Y^3)\tensor \left\{-\frac{2}{3},  -\frac{5}{8}\right\}.
\end{align*}

We will often write 
$P(X,Y)\{\alpha,\beta\}$ for $P(X,Y)\tensor\{\alpha,\beta\}$.

\begin{definition}[Modular Symbols]\label{defn:modsym}
Let $k\geq 2$ be an integer and let $G$ be a finite index subgroup of $\SL_2(\Z)$.
The space \sym{\sM_k(G)} 
of \defn{weight~$k$ modular symbols for $G$} is the quotient of
$\sM_k$ by all relations $gx - x$ for $x \in \sM_k$, $g\in G$, and by any torsion.
\end{definition}
Note that $\sM_k$ is a torsion-free abelian group, and it is a
nontrivial fact that $\sM_k$ has finite rank.  We denote modular
symbols for $G$ in exactly the same way we denote elements of $\sM_k$;
the group~$G$ will be clear from context.

The space of \defn{modular symbols over a ring $R$} is\index{$\sM_k(G;R)$}
$$
  \sM_k(G;R) = \sM_k(G)\tensor_\Z R.
$$


\section{Manin Symbols}\label{sec:manin}

Let $G$ be a finite index subgroup of $\SL_2(\Z)$ and $k\geq 2$ 
an integer.
Just as in Chapter~\ref{ch:modform2} it is possible to compute
$\sM_k(G)$ using a computer, despite that, as defined above,
$\sM_k(G)$ is the quotient of one infinitely generated abelian group
by another one.  This section is about Manin symbols, which are a
distinguished subset of $\sM_k(G)$ that lead to a finite presentation
for $\sM_k(G)$.  Formulas written in terms of Manin symbols are
frequently much easier to compute using a computer than formulas in
terms of modular symbols.

Suppose $P\in \Z[X,Y]_{k-2}$ and  $g\in \SL_2(\Z)$.
Then the \defn{Manin symbol} associated to this
pair of elements is
$$
  [P,g] = g(P\{0,\infty\}) \in \sM_k(G).
$$
Notice that if $Gg = Gh$, then $[P,g] = [P,h]$, since 
the symbol $g(P\{0,\infty\})$ is invariant
by the action of $G$ on the left (by definition,
since it is a modular symbol for $G$).
Thus for a right coset $Gg$ it makes sense to 
write $[P, Gg]$ for the symbol $[P,h]$ for any $h\in Gg$.
Since $G$ has finite index in $\SL_2(\Z)$,
the abelian group generated by Manin symbols is of finite rank, generated by
$$
 \bigl\{ [X^{k-2-i}Y^i, \, Gg_j] \, : \, i = 0,\ldots, k-2 \quad\text{and}\quad j = 0,\ldots, r \bigr\},
$$
where $g_0, \ldots, g_r$ run through representatives for the 
right cosets $G\backslash \SL_2(\Z)$.

We next show that every modular symbol can be written as
a $\Z$-linear combination of Manin symbols, 
so they generate $\sM_k(G)$. 
\begin{proposition}\label{prop:mangen}
The Manin symbols generate $\sM_k(G)$.
\end{proposition}
\begin{proof}
The proof if very similar to that
of Proposition~\ref{prop:manintrick} except we introduce an
extra twist to deal with the polynomial part.
Suppose that we are given a modular
symbol $P\{\alpha,\beta\}$ and wish to represent it as a 
sum of Manin symbols.   Because 
    $$P\{a/b,c/d\} = P\{a/b,0\}+P\{0,c/d\},$$
it suffices to write $P\{0,a/b\}$ in
terms of Manin symbols. 
Let
$$0=\frac{p_{-2}}{q_{-2}} = \frac{0}{1},\,\,
\frac{p_{-1}}{q_{-1}}=\frac{1}{0},\,\,
\frac{p_0}{1}=\frac{p_0}{q_0},\,\,
\frac{p_1}{q_1},\,\,
\frac{p_2}{q_2},\,\ldots,\,\frac{p_r}{q_r}=\frac{a}{b}$$
denote the continued fraction convergents of the 
rational number $a/b$. 
Then
$$p_j q_{j-1} 
  - p_{j-1} q_j = (-1)^{j-1}\qquad \text{for }-1\leq j\leq r.$$
If we let
$g_j = \mtwo{(-1)^{j-1}p_j}{p_{j-1}}{(-1)^{j-1}q_j}{q_{j-1}}$,
then $g_j\in\sltwoz$ and 
\begin{align*}
  P\{0,a/b\}
 &=P\sum_{j=-1}^{r}\left\{\frac{p_{j-1}}{q_{j-1}},\frac{p_j}{q_j}\right\}\\
 &=\sum_{j=-1}^{r} g_j((g_j^{-1}P)\{0,\infty\})\\
 &=\sum_{j=-1}^{r} [g_j^{-1}P,\,\,g_j].
\end{align*}
Since $g_j\in\sltwoz$ and $P$ has integer coefficients, 
the polynomial $g_j^{-1}P$ also has integer coefficients,
so we introduce no denominators.
\end{proof}

Now that we know the Manin symbols generate $\sM_k(G)$, we
next consider the relations between Manin symbols.  Fortunately,
the answer is fairly simple (though the proof is not).  Let
$$
\sigma = \mtwo{0}{-1}{1}{\hfill 0},\qquad
\tau = \mtwo{0}{-1}{1}{-1}, \qquad
J = \mtwo{-1}{\hfill 0}{\hfill 0}{-1}.
$$
Define a \defn{right action of $\SL_2(\Z)$} on Manin symbols
as follows.  If $h \in \SL_2(\Z)$, let
$$
   [P, g]h = [h^{-1} P, gh].
$$
This is a right action because both $P\mapsto h^{-1}P$ 
and $g\mapsto gh$ are right actions.
\begin{theorem}\label{thm:mansym}
If $x$ is a Manin symbol, then
\begin{align}
x + x\sigma &= 0,\\
x + x\tau + x\tau^{2} &= 0,\\
x - xJ &= 0.
\end{align}
Moreover, these are all the relations between Manin symbols, in the
sense that the space $\sM_k(G)$ of modular symbols is isomorphic to
the quotient of the free abelian group on the finitely many symbols
$[X^{i}Y^{k-2-i},Gg]$ (for $i=0,\ldots, k-2$ and 
$Gg \in G\backslash \SL_2(\Z)$) by the above relations and any torsion.
\end{theorem}
\begin{proof}
First we prove that the Manin symbols
satisfy the above relations.
We follow Merel's proof (see \cite[\S1.2]{merel:1585}).
Note that 
$$\sigma(0) = \sigma^2(\infty) = \infty
\qquad\text{and}\qquad
  \tau(1) = \tau^2(0) = \infty.
$$
Writing $x=[P,g]$, we have 
\begin{align*}
[P,g] + [P,g]\sigma &= 
  [P,g]  + [\sigma^{-1}P, g\sigma]\\
 &= g(P\{0,\infty\}) + g\sigma (\sigma^{-1}P \{0,\infty\})\\
 &=
  (gP) \{ g(0), g(\infty) \} + 
  (g\sigma)(\sigma^{-1} P)\{ g\sigma(0), g\sigma(\infty)\} \\
  &=   (gP) \{ g(0), g(\infty) \} + 
  (gP)\{ g(\infty), g(0)\} \\ 
  &= (gP)\{ g(0), g(\infty) \} + \{ g(\infty), g(0)\})\\
  &= 0.
\end{align*}
Also, 
\begin{align*}
[P,g] &+ [P,g]\tau + [P,g]\tau^2 = 
  [P,g]  + [\tau^{-1}P, g\tau] + [\tau^{-2}P, g\tau^2]\\
  &= 
g(P\{0,\infty\}) + 
   g\tau(\tau^{-1}P \{0,\infty\}) 
   + g\tau^2(\tau^{-2}P \{0,\infty\}) \\
&= (gP) \{ g(0), g(\infty) \} + 
   (gP) \{ g \tau(0),g\tau(\infty)\}
   + (gP) \{ g\tau^2(0),\tau^2(\infty)\} \\
&= (gP) \{ g(0), g(\infty) \} + 
   (gP) \{ g(1),g(0)\}) 
   + (gP) \{ g(\infty), g(1)\} \\
 &= (gP)\bigl(\{ g(0), g(\infty)\} + 
\{ g(\infty), g(1)\} + \{ g(1),g(0)\}\bigr) \\
&= 0.
\end{align*}
Finally, 
\begin{align*}
[P,g] + [P,g]J &= 
  g(P\{0,\infty\}) - gJ(J^{-1} P \{gJ(0), gJ(\infty)\}\\
  &= (gP)\{g(0),g(\infty)\} - (gP) \{g(0),g(\infty)\}\\
  &= 0,
\end{align*}
where we use that $J$ acts trivially via linear fractional
transformations.  This proves that the
listed relations are all satisfied. 

That the listed relations are all relations is more difficult to
prove.  One
approach is to show (as in \cite[\S1.3]{merel:1585}) that the quotient
of Manin symbols by the above relations and torsion is isomorphic to a
space of Shokurov symbols, which is in turn isomorphic to
$\sM_k(G)$.  A much better approach is to apply some results
from group cohomology, as in 
\cite{gabor:phd}.
\end{proof}


If $G$ is a finite index subgroup and we have an algorithm to
enumerate the right cosets $G\backslash \SL_2(\Z)$ and to decide
which coset an arbitrary element of $\SL_2(\Z)$ belongs to, then
Theorem~\ref{thm:mansym} and the algorithms of Chapter~\ref{ch:linalg}
yield an algorithm to compute $\sM_k(G;\Q)$.  
Note that  if $J\in G$, then the relation $x-xJ=0$ is automatic.  

\begin{remark}
The matrices $\sigma$ and $\tau$ {\em do not commute},
so in computing $\sM_k(G;\Q)$, 
one cannot first quotient out by the two-term $\sigma$
relations, then quotient out only the remaining free generators
by the $\tau$ relations, and get the right answer in general.
\end{remark}


\subsection{Coset Representatives and Manin Symbols}

The following is analogous to Proposition~\ref{prop:p1bij}:

\begin{proposition}\label{prop:gamma1cosets}
The right cosets $\Gamma_1(N)\backslash \SL_2(\Z)$ are in bijection
with pairs $(c,d)$ where $c,d\in\Z/N\Z$ and $\gcd(c,d,N)=1$.
The coset containing
a matrix $\abcd{a}{b}{c}{d}$ corresponds to $(c,d)$.
\end{proposition}
\begin{proof}
This proof is copied from \cite[pg.~203]{cremona:gammaone},
except in that paper Cremona works with
the analogue of $\Gamma_1(N)$ in $\PSL_2(\Z)$, so
his result is slightly different.
Suppose $\gamma_i = \abcd{a_i}{b_i}{c_i}{d_i} \in \SL_2(\Z)$,
for $i=1,2$.
We have
$$
 \gamma_1 \gamma_2^{-1} = 
  \mtwo{a_1}{b_1}{c_1}{d_1}
  \mtwo{\hfill{}d_2}{-b_2}{-c_2}{\hfill{}a_2}
  =\mtwo{a_1 d_2 - b_1 c_2}{*}{c_1 d_2 - d_1 c_2}{a_2 d_1 - b_2 c_1},$$
which is in $\Gamma_1(N)$ if and only if 
\begin{equation}\label{eqn:cos1a}
  c_1 d_2 - d_1 c_2 \con 0 \pmod{N}
\end{equation}
and 
\begin{equation}\label{eqn:cos1b}
  a_2 d_1 - b_2 c_1 \con a_1 d_2 - b_1 c_2 \con 1\pmod{N}.
\end{equation}
Since the $\gamma_i$ have determinant~$1$, if $(c_1,d_1)=(c_2,d_2)\pmod{N}$,
then the congruences \eqref{eqn:cos1a}--\eqref{eqn:cos1b} hold.
Conversely, if \eqref{eqn:cos1a}--\eqref{eqn:cos1b} hold, then 
\begin{align*}
 c_2 &\con a_2 d_1 c_2 - b_2  c_1 c_2 \\
    &\con a_2 d_2 c_1 - b_2 c_2 c_1 \quad\text{ since }d_1 c_2 \con d_2 c_1 \pmod{N}\\
    &\con c_1 \qquad\qquad\qquad\quad\,\text{since }a_2 d_2 - b_2 c_2 = 1,
  \end{align*}
and likewise 
$$
d_2 \con a_2 d_1 d_2 - b_2 c_1 d_2
   \con a_2 d_1 d_2 - b_2 d_1 c_2
   \con d_1\pmod{N}.
$$
\end{proof}

Thus we may view weight~$k$ Manin symbols for $\Gamma_1(N)$ as triples
of integers $(i,c,d)$, where $0\leq i\leq k-2$ and $c, d\in\Z/N\Z$
with $\gcd(c,d,N)=1$.  Here $(i,c,d)$ corresponds to the Manin symbol
$[X^iY^{k-2-i}, \abcd{a}{b}{c'}{d'}]$, where $c'$ and $d'$ are
congruent to $c,d \pmod{N}$.  The relations of
Theorem~\ref{thm:mansym} become
\begin{align*}
(i,\,c,d) + (-1)^i (k-2-i, \,d, -c) &= 0,\\
(i,\,c,d) \quad + \quad\,        (-1)^{k-2} \sum_{j=0}^{k-2-i}  (-1)^{j} \binom{k-2-i}{j} (j,\,d,-c-d)&\\
       \qquad\qquad + \quad
(-1)^{k-2-i} \sum_{j=0}^i (-1)^j \binom{i}{j} (k-2-i+j,\,-c-d,c)
 &=0,\\
(i,\,c,d) - (-1)^{k-2}(i,\,-c,-d) &= 0.
\end{align*}

Recall that Proposition~\ref{prop:p1bij}
gives a similar description for $\Gamma_0(N)$.

\subsection{Modular Symbols with Character}
Suppose $G=\Gamma_1(N).$
Define an action of \defn{diamond-bracket operators} $\dbd{d}$o,
with $\gcd(d,N)=1$ on
Manin symbols as follows:
$$
\langle d \rangle ([P,(u,v)]) = \left[P,(du, dv)\right].
$$

Let 
$$
  \eps : (\Z/N\Z)^* \to \Q(\zeta)^*
$$
be a Dirichlet character, where $\zeta$ is
an $n$th root of unity and~$n$ is the order of~$\eps$.  
Let \sym{\sM_k(N,\eps)} be the quotient of 
$\sM_k(\Gamma_1(N);\Z[\zeta])$ by the relations (given
in terms of Manin symbols)
$$
  \dbd{d}x - \eps(d)x = 0,
$$
for all $x\in \sM_k(\Gamma_1(N);\Z[\zeta])$,
$d\in (\Z/N\Z)^*$, and by any $\Z$-torsion.  
Thus $\sM_k(N,\eps)$ is a $\Z[\eps]$-module
that has no torsion when viewed as a $\Z$-module.



\section{Hecke Operators}

Suppose $\Gamma$ is a subgroup of $\SL_2(\Z)$ of level $N$
that contains $\Gamma_1(N)$. 
Just as for modular forms, there is a commutative \defn{Hecke algebra}
$\T=\Z[T_1,T_2, \ldots]$, which is the subring of $\End(\sM_k(\Gamma))$
generated by all Hecke operators. 
Let 
$$
R_p = \left\{\mtwo{1}{r}{0}{p} : r = 0,1,\ldots, p-1 \right\} 
         \union \left\{ \mtwo{p}{0}{0}{1}\right\},
$$
where we omit $\abcd{p}{0}{0}{1}$ if $p\mid N$.  Then 
the \defn{Hecke operator} $T_p$ on $\sM_k(\Gamma)$
is given by
$$
  T_p(x) = \sum_{g \in R_p} gx.
$$
Notice when $p\nmid N$ that $T_p$ is defined by summing over 
$p+1$ matrices that correspond to the $p+1$ subgroups of $\Z\times \Z$
of index~$p$.  This is exactly how we defined $T_p$ on modular forms
in Definition~\ref{defn:hecke_op}.


\subsection{General Definition of Hecke Operators}\label{sec:gendefn}
Let $\Gamma$ be a finite index subgroup of $\SL_2(\Z)$ and suppose
$$
  \Delta\subset \GL_2(\Q)
$$
is a set such that $\Gamma\Delta = \Delta\Gamma = \Delta$ 
and $\Gamma \backslash \Delta$
is finite.  For example, $\Delta=\Gamma$
satisfies this condition.  Also,
if $\Gamma=\Gamma_1(N)$, then for any positive integer $n$, the set 
$$
  \Delta_n = \left\{
\!\mtwo{a}{b}{c}{d}\in \Mat_2(\Z) \,:\, ad-bc=n, \quad
\mtwo{a}{b}{c}{d}\con \mtwo{1}{*}{0}{n}\!\!\!\!\!\!\pmod{N}
   \right\}
$$
also satisfies this condition, as we will now prove.
\begin{lemma}\label{lem:deltan}
We have 
$$
  \Gamma_1(N)\cdot \Delta_n = \Delta_n \cdot \Gamma_1(N) = \Delta_n
$$
and
$$
 \Delta_n = \bigcup_{a,b} \Gamma_1(N) \cdot \sigma_a \mtwo{a}{b}{0}{n/a},
$$
where $\sigma_a \con \abcd{\hfill 1/a}{\,\, 0}{0}{ \hfill a}\pmod{N}$,
the union is disjoint and $1\leq a \leq n$ with $a\mid n$, $\gcd(a,N)=1$,
and $0\leq b < n/a$.  In particular, the set of cosets
$\Gamma_1(N)\backslash \Delta_n$ is finite.
\end{lemma}
\begin{proof}
(This is Lemma 1 of \cite[\S2.3]{merel:1585}.)
If $\gamma\in\Gamma_1(N)$ and $\delta\in\Delta_n$, then
$$
\mtwo{1}{*}{0}{1} \cdot \mtwo{1}{*}{0}{n} \con 
\mtwo{1}{*}{0}{n} \cdot \mtwo{1}{*}{0}{1}
\con \mtwo{1}{*}{0}{n}\pmod{N}.
$$ 
Thus $\Gamma_1(N)\Delta_n \subset \Delta_n$, and since $\Gamma_1(N)$
is a group, $\Gamma_1(N)\Delta_n = \Delta_n$; likewise 
$\Delta_n\Gamma_1(N) = \Delta_n$.

For the coset decomposition, we first prove the statement for $N=1$, i.e., for 
$\Gamma_1(N)=\SL_2(\Z)$.   If $A$ is an arbitrary element
of $\Mat_2(\Z)$ with determinant $n$, then using row operators on the left
with determinant $1$, i.e., left multiplication by elements of $\SL_2(\Z)$,
we can transform $A$ into the form $\abcd{a}{b}{0}{n/a}$, with
$1\leq a\leq n$ and $0\leq b < n$.  (Just imagine applying the Euclidean
algorithm to the two entries in the first column of $A$.  Then $a$
is the $\gcd$ of the two entries in the first column, and the lower left
entry is $0$.  Next subtract $n/a$ from $b$ until $0\leq b < n/a$.)

Next suppose $N$ is arbitrary.  Let $g_1, \ldots, g_r$ be such
that $$g_1\Gamma_1(N) \union \cdots \union g_r \Gamma_1(N) = \SL_2(\Z)$$
is a disjoint union.  If $A \in \Delta_n$ is arbitrary, then 
as we showed above, there is some $\gamma\in\SL_2(\Z)$, so that
$\gamma\cdot A = \abcd{a}{b}{0}{n/a}$, with $1\leq a \leq n$ and 
$0\leq b < n/a$, and $a\mid n$.   Write $\gamma = g_i \cdot \alpha$, with $\alpha\in\Gamma_1(N)$.
Then 
$$
  \alpha\cdot A = g_i^{-1} \cdot \mtwo{a}{b}{0}{n/a} \con \mtwo{1}{*}{0}{n}\pmod{N}.
$$
It follows that 
$$
  g_i^{-1} \con \mtwo{1}{*}{0}{n} \cdot \mtwo{a}{b}{0}{n/a}^{-1} 
    \con \mtwo{1/a}{*}{0}{a}\pmod{N}.
$$
Since $\abcd{1}{1}{0}{1}\in \Gamma_1(N)$ and $\gcd(a,N)=1$,
there is $\gamma'\in\Gamma_1(N)$ such
that $$\gamma' g_i^{-1} \con \mtwo{1/a}{0}{0}{a}\pmod{N}.$$
We may then choose $\sigma_a = \gamma' g_i^{-1}$.
Thus every $A\in \Delta_n$ is of the form
$\gamma \sigma_a \abcd{a}{b}{0}{n/a}$,
with $\gamma\in\Gamma_1(N)$ and $a,b$ suitably bounded.  This
proves the second claim.
\end{proof}

Let any element $\delta=\abcd{a}{b}{c}{d}\in\GL_2(\Q)$
act on the left on modular symbols $\M_k\tensor\Q$ by
$$
\delta(P\{\alpha,\beta\}) = P(dX-bY,-cX+aY) \{\delta(\alpha),\delta(\beta)\}.
$$
(Until now we had only defined an action of $\SL_2(\Z)$
on modular symbols.)
For $g=\abcd{a}{b}{c}{d}\in \GL_2(\Q)$, let
\begin{equation}\label{eqn:gtilde}
 \tilde{g} = \mtwo{\hfill{}d}{-b}{-c}{\hfill{}a}
           = \det(g)\cdot g^{-1}.
         \end{equation}
Note that $\tilde{\tilde{g}} = g$.  
Also, $\delta P(X,Y)=(P\circ \tilde{g})(X,Y)$,
where we set $$\tilde{g}(X,Y) = (dX-bY,-cX+aY).$$


Suppose $\Gamma$ and $\Delta$ are as above.    Fix
a finite set $R$ of representatives for $\Gamma\backslash \Delta$.
Let
$$
 T_{\Delta} : \M_k(\Gamma) \to \M_k(\Gamma)
$$
be the linear map
$$
  T_{\Delta}(x) = \sum_{\delta\in R} \delta x.
$$

This map is well defined because 
if $\gamma\in\Gamma$ and $x\in \M_k(\Gamma)$,
then 
$$\sum_{\delta\in R} \delta \gamma  x = 
\sum_{\text{certain $\delta'$}} \gamma \delta' x
 = \sum_{\text{certain $\delta'$}} \delta' x
 =\sum_{\delta\in R} \delta x,
$$
where we have used that $\Delta\Gamma = \Gamma\Delta$, and 
$\Gamma$ acts trivially on $\M_k(\Gamma)$.

Let $\Gamma=\Gamma_1(N)$ and $\Delta=\Delta_n$.  Then the $n$th
Hecke operator $T_n$ is $T_{\Delta_n}$, and by Lemma~\ref{lem:deltan},
$$
 T_n(x) = \sum_{a,b} \sigma_a \mtwo{a}{b}{0}{n/a} \cdot x,
$$
where $a,b$ are as in Lemma~\ref{lem:deltan}.

Given this definition, we can compute the Hecke operators on
$M_k(\Gamma_1(N))$ as follows.  Write $x$ as a modular symbol
$P\{\alpha,\beta\}$, compute $T_n(x)$ as a modular symbol,
and then convert  to Manin symbols using continued fractions
expansions.  This is extremely inefficient; fortunately Lo\"{i}c
Merel (following \cite{mazur:symboles}) found a much better way, which
we now describe (see \cite{merel:1585}).

\subsection{Hecke Operators on Manin Symbols}\label{sec:hecke_on_manin}
If $S$ is a subset of $\GL_2(\Q)$, 
let $$\tilde{S} = \{\tilde{g} : g \in  S\},$$
where $\tilde{g}$ is as in \eqref{eqn:gtilde}.
Also, for any ring $R$ and any subset $S\subset \Mat_2(\Z)$, 
let $R[S]$ denote the free $R$-module with basis the elements of $S$,
so the elements of $R[S]$ are the finite $R$-linear combinations
of the elements of $S$.

One of the main theorems of \cite{merel:1585} is that for any
$\Gamma,\Delta$ satisfying the condition at the beginning of
Section~\ref{sec:gendefn}, if we can find $\sum u_M M \in
\C[\Mat_2(\Z)]$ and a map
$$\phi: \tilde{\Delta} \SL_2(\Z) \to \SL_2(\Z)$$
that satisfies certain conditions, then for any Manin
symbol $[P,g] \in \M_k(\Gamma)$, we have 
$$ 
  T_\Delta([P,g]) =
\sum_{gM\in \tilde{\Delta}\SL_2(\Z)\text{ with }M\in\SL_2(\Z)} u_M
[\tilde{M}\cdot P,\,\, \phi(gM)].
$$ 
The paper
\cite{merel:1585} contains many examples of $\phi$ and $\sum u_M M \in
\C[\Mat_2(\Z)]$ that satisfy all of the conditions.
 
When $\Gamma=\Gamma_1(N)$, the complicated list of conditions becomes
simpler.    Let \sym{\Mat_2(\Z)_n} be the set of $2\times 2$
matrices with determinant $n$.  
An element
$$
  h = \sum u_M [M] \in \C[\Mat_2(\Z)_n]
$$
\defn{satisfies condition $C_n$}
if for every $K\in \Mat_2(\Z)_n / \SL_2(\Z)$, we have that 
\begin{equation}\label{eqn:cn}
  \sum_{M \in K} u_M([M\infty] - [M 0 ]) = [\infty] - [0]
\in \C[P^1(\Q)].
\end{equation}
If $h$ satisfies condition $C_n$, then for any Manin symbol
$[P,g] \in M_k(\Gamma_1(N))$, Merel proves that
\begin{equation}\label{eqn:tnmanin}
  T_n([P,(u,v)]) = \sum_{M} u_M [P(aX+bY, \, cX+dY), (u,v)M].
\end{equation}
Here $(u,v)\in (\Z/N\Z)^2$ corresponds via Proposition~\ref{prop:gamma1cosets}
to a coset of $\Gamma_1(N)$ in $\SL_2(\Z)$, and
if $(u',v')=(u,v)M \in (\Z/N\Z)^2$ and $\gcd(u',v',N)\neq 1$,
then we omit the corresponding summand.

For example, we will now check directly that the element
$$
 h_2 =  \left[\mtwo{2}{0}{0}{1}\right] + 
        \left[\mtwo{1}{0}{0}{2}\right] + 
        \left[\mtwo{2}{1}{0}{1}\right] + 
        \left[\mtwo{1}{0}{1}{2}\right]
$$
satisfies condition $C_2$. 
We have, as in the proof of Lemma~\ref{lem:deltan}
(but using elementary column operations), that
\begin{align*}
  \Mat_2(\Z)_2/\SL_2(\Z) &= \left\{
\mtwo{a}{0}{b}{2/a} \SL_2(\Z) : a=1,2\text{ and }0\leq b < 2/a
\right\} \\
&\hspace{-2ex}= \left\{ \mtwo{1}{0}{0}{2} \SL_2(\Z),\,\,\,\,
       \mtwo{1}{0}{1}{2} \SL_2(\Z),\,\,\,\,
      \mtwo{2}{0}{0}{1} \SL_2(\Z) \right\}.
  \end{align*}
To verify condition $C_2$, we consider in turn each of the three
elements of $\Mat_2(\Z)_2/\SL_2(\Z)$ and check that \eqref{eqn:cn}
holds.  We have that 
$$
\mtwo{1}{0}{0}{2} \in \mtwo{1}{0}{0}{2} \SL_2(\Z),
$$
$$
\mtwo{2}{1}{0}{1}, \mtwo{1}{0}{1}{2} \in \mtwo{1}{0}{1}{2} \SL_2(\Z),
$$
and
$$
\mtwo{2}{0}{0}{1}\in \mtwo{2}{0}{0}{1} \SL_2(\Z).
$$
Thus if $K = \abcd{1}{0}{0}{2}\SL_2(\Z)$, the left
sum of \eqref{eqn:cn} is 
$[\abcd{1}{0}{0}{2}(\infty)] - [\abcd{1}{0}{0}{2}(0)]  = [\infty]-[0]$,
as required.  If $K=\abcd{1}{0}{1}{2} \SL_2(\Z)$, then
the left side of \eqref{eqn:cn} is 
\begin{align*}
[\abcd{2}{1}{0}{1}(\infty)] &- [\abcd{2}{1}{0}{1}(0)] 
 + [\abcd{1}{0}{1}{2}(\infty)] - [\abcd{1}{0}{1}{2}(0)] \\
  &= [\infty] - [1] + [1] - [0] = [\infty] - [0].
\end{align*}
Finally, for $K = \abcd{2}{0}{0}{1}\SL_2(\Z)$
we also have
$[\abcd{2}{0}{0}{1}(\infty)] - [\abcd{2}{0}{0}{1}(0)]  = [\infty]-[0]$,
as required.
Thus by \eqref{eqn:tnmanin} we can compute $T_2$ on {\em any} Manin
symbol, by summing over the action 
 of the four matrices
$\abcd{2}{0}{0}{1}, \abcd{1}{0}{0}{2}, 
\abcd{2}{1}{0}{1}, \abcd{1}{0}{1}{2}$.

\begin{proposition}[Merel]\label{prop:heilbronn}
The element
$$\sum_{\substack{a>b\geq 0\\d>c\geq 0\\ad-bc=n}}
\left[\mtwo{a}{b}{c}{d}\right] \in \Z[\Mat_2(\Z)_n]
$$
satisfies condition $C_n$.
\end{proposition}
Merel's two-page proof of Proposition~\ref{prop:heilbronn}
is fairly elementary. 

\begin{remark}
  In \cite[\S2.4]{cremona:algs}, Cremona discusses the work of Merel
  and Mazur on Heilbronn matrices\index{Heilbronn matrices} in the special cases
  $\Gamma=\Gamma_0(N)$ and weight~$2$.  He gives a simple proof
  that the action of $T_p$ on Manin symbols can be computed by summing
  the action of some set $R_p$ of matrices of determinant~$p$.  He
  then describes the set $R_p$ and gives an efficient continued
  fractions algorithm for computing it (but he does not prove
  that his $R_p$ satisfy Merel's hypothesis). 
% (Note:
%  The author's experience is that Cremona's set $R_p$ is significantly smaller
%  than the sets appearing in Merel's paper, but when I've tried to use
%  $R_p$ to do certain more general higher-weight computations that are
%  correct using Merel's sets, they do not work.)
\end{remark}

\subsection{Remarks on Complexity}
Merel gives another family $\mathcal{S}_n$ of matrices
that satisfy condition $C_n$, and he proves that
as $n\to \infty$,
$$\#\mathcal{S}_n \sim \frac{12 \log(2)}{\pi^2} \cdot 
  \sigma_1(n)\log(n),$$
where $\sigma_1(n)$ is the sum of the divisors of $n$.
Thus for a fixed space $\sM_k(\Gamma)$ of modular symbols,
one can compute $T_n$ using
$O(\sigma_1(n)\log(n))$ arithmetic operations.  
Note that we have fixed $\sM_k(\Gamma)$, so we ignore 
the linear algebra involved in computation of a presentation; also, adding
elements takes a bounded number of field operations when the space
is fixed.  Thus, using Manin symbols 
the complexity of computing $T_p$, for~$p$ prime, is $O((p+1)\log(p))$
field operations, which is {\em exponential} in the number of digits
of~$p$.     

\subsection{Basmaji's Trick}\label{sec:basmaji}\index{Basmaji's trick}
There is a trick of Basmaji (see \cite{basmaji:thesis}) for computing
a matrix of $T_n$ on $\sM_k(\Gamma)$, when $n$ is very large, and it is more
efficient than one might naively expect.  Basmaji's trick does not
improve the big-oh complexity for a fixed space, but it
does improve the
complexity by a constant factor of the dimension of $\sM_k(\Gamma;\Q)$.
Suppose we are interested in computing the matrix for $T_n$ for some massive
integer~$n$ and that $\sM_k(\Gamma;\Q)$ has large dimension.
The trick is as follows.  Choose a list 
$$x_1 =[P_1,g_1],\ldots,x_r = [P_r,g_r] \in V = \sM_k(\Gamma;\Q)$$ 
of Manin symbols  such that the
map $\Psi:\T \to V^r$
 given by $$t \mapsto (t x_1, \ldots, t x_r)$$
is injective.  In practice, it is often possible to do this with $r$ ``very small''.
Also, we emphasize that $V^r$ is a $\Q$-vector space of dimension 
$r\cdot \dim(V)$.

Next find Hecke operators $T_i$, with $i$ small, 
whose images form a basis for the image of $\Psi$.
Now with the above data precomputed, which only required working with
Hecke operators $T_i$ for small $i$, we are ready to compute $T_n$
with $n$ huge.  Compute $y_i = T_n(x_i)$, for each $i=1,\ldots, r$,
which we can compute using Heilbronn matrices\index{Heilbronn matrices} since each
$x_i=[P_i,g_i]$ is a Manin symbol.  We thus obtain $\Psi(T_n) \in V^r.$
Since we have precomputed Hecke operators $T_j$ such that
$\Psi(T_j)$ generate $V^r$, we can find $a_j$ such that
$\sum a_j \Psi(T_j) = \Psi(T_n)$.  Then since $\Psi$
is injective, we have $T_n = \sum a_j T_j$, which gives
the full matrix of $T_n$ on $M_k(\Gamma;\Q)$.



\section{Cuspidal Modular Symbols}\label{sec:cuspmodsym}
Let $\sB$ be the free abelian group on symbols $\{\alpha\}$,
for $\alpha\in\P^1(\Q)$, and set
$$
  \sB_k = \Z[X,Y]_{k-2}\tensor\sB.
$$
Define a \defn{left action of $\SL_2(\Z)$} on $\sB_k$ by
$$
  g  (P \{\alpha\}) = (gP) \{g(\alpha)\},
$$
for $g\in \SL_2(\Z)$.
For any finite index subgroup $\Gamma\subset \SL_2(\Z)$,
let 
$\sB_k(\Gamma)$ be the quotient of $\sB_k$ by the relations
$x- gx$ for all $g\in \Gamma$ and by any torsion.  Thus
$\sB_k(\Gamma)$ is a torsion-free abelian group.

The \defn{boundary map} is the map
$$
  b:\sM_k(\Gamma) \to \sB_k(\Gamma)
$$
given by extending the map
$$
  b(P\{\alpha,\beta\}) = P\{\beta\} - P\{\alpha\}
$$
linearly.
The space \sym{\sS_k(\Gamma)} of \defn{cuspidal modular symbols}
is the kernel\index{cusp forms!for $\Gamma$}
$$
  \sS_k(\Gamma) = \ker(\sM_k(\Gamma) \to \sB_k(\Gamma)),
$$
so we have an exact sequence
$$
  0 \to \sS_k(\Gamma) \to \sM_k(\Gamma) \to \sB_k(\Gamma).
$$
One can prove that when $k>2$,
this sequence is exact on the right.  

Next we give a presentation of $\sB_k(\Gamma)$ in terms of
``boundary Manin symbols''. 


\subsection{Boundary Manin Symbols}%
\label{sec:computeboundary}%
\index{Manin symbols!and cuspidal subspace}%
\index{Manin symbols!and boundary space}%
\index{cuspidal modular symbols!and Manin symbols}%
\index{boundary modular symbols!and Manin symbols}%

We give an explicit description of the boundary
map in terms of Manin symbols for 
$\Gamma=\Gamma_1(N)$, then describe an
efficient way to compute the boundary map.

Let~$\cR$ be the equivalence relation 
on $\Gamma\backslash\Q^2$ given by
$$
 [\Gamma\smallvtwo{\lambda u}{\lambda v}] \sim   
  \sign(\lambda)^k[\Gamma\smallvtwo{u}{v}],
$$
for any $\lambda\in\Q^*$.  Denote by $B_k(\Gamma)$
the finite-dimensional $\Q$-vector space
with basis the equivalence classes
$(\Gamma\backslash\Q^2)/\cR$. 
The following two propositions are proved in \cite{merel:1585}.
\begin{proposition}
The map
$$\mu:\sB_k(\Gamma)\ra B_k(\Gamma),
\qquad P\left\{\frac{u}{v}\right\}\mapsto
    P(u,v)\left[\Gamma\vtwo{u}{v}\right]$$
is well defined and injective.  
Here $u$ and $v$ are assumed coprime.
\end{proposition}
Thus the kernel of $\delta:\sS_k(\Gamma)\ra \sB_k(\Gamma)$
is the same as the kernel of $\mu\circ \delta$. 
\begin{proposition}\label{prop:boundary}
Let $P\in V_{k-2}$ and $g=\abcd{a}{b}{c}{d}\in\sltwoz$.  We have
$$\mu\circ\delta([P,(c,d)])
      = P(1,0)[\Gamma\smallvtwo{a}{c}]
       -P(0,1)[\Gamma\smallvtwo{b}{d}].$$
\end{proposition}

\index{boundary map}We next describe how to explicitly compute 
$\mu\circ \delta:\sM_k(N,\eps)\ra B_k(N,\eps)$
by generalizing the algorithm of \cite[\S2.2]{cremona:algs}.
To compute the image of $[P,(c,d)]$, with 
$g=\abcd{a}{b}{c}{d}\in\sltwoz$,
we must compute the class of $[\smallvtwo{a}{c}]$ and of 
$[\smallvtwo{b}{d}]$.
Instead of finding a canonical form for cusps, we 
use a quick test for equivalence modulo scalars. 
In the following algorithm, by the $i$th 
symbol\index{cusps!and boundary map} we mean
the $i$th basis vector for a basis of $B_k(N,\eps)$.  This
basis is constructed as the algorithm is called successively.
We first give the algorithm, and then prove the facts
used by the algorithm in testing equivalence.

\begin{algorithm}{Cusp Representation}\label{alg:cusplist}
  Given a boundary Manin symbol
  $[\smallvtwo{u}{v}]$, this algorithm outputs an integer~$i$ and a
  scalar~$\alp$ such that $[\smallvtwo{u}{v}]$ is equivalent to~$\alp$
  times the $i$th symbol found so far.  (We call this algorithm
  repeatedly and maintain a running list of cusps seen so far.)
\begin{steps}
\item Use Proposition~\ref{prop:cusp1g0} to check
whether or not $[\smallvtwo{u}{v}]$ is equivalent, modulo scalars, to
any cusp found.  If so, return the representative, its
index, and the scalar.  If not, record $\smallvtwo{u}{v}$ in the
representative list.
\item Using Proposition~\ref{prop:cuspdies},
check whether or not $[\smallvtwo{u}{v}]$ is forced to equal $0$ by the
relations.  If it does not equal $0$, return its position in the list
and the scalar~$1$.  If it equals $0$, return the scalar~$0$ and the
position~$1$; keep $\smallvtwo{u}{v}$ in the list, and record that it
is equivalent to $0$.
\end{steps}
\end{algorithm}

The case considered in Cremona's book \cite{cremona:algs} only
involve the trivial character, so no cusp classes
are forced to vanish.  Cremona gives the following two
criteria for equivalence.
\begin{proposition}[Cremona]\label{prop:cusp1}
Consider $\smallvtwo{u_i}{v_i}$, $i=1,2$, with $u_i,v_i$ integers such
that $\gcd(u_i,v_i)=1$ for each $i$.
\begin{enumerate}
\item There exists $g\in\Gamma_0(N)$ such that 
    $g\smallvtwo{u_1}{v_1}=\smallvtwo{u_2}{v_2}$ if and only if
 $$s_1 v_2 \con s_2 v_1 \pmod{\gcd(v_1 v_2,N)},\,
\text{ where $s_j$ satisfies $u_j s_j\con 1\!\!\!\pmod{v_j}$}.$$
\item There exists $g\in\Gamma_1(N)$ such that 
    $g\smallvtwo{u_1}{v_1}=\smallvtwo{u_2}{v_2}$ if and only if
 $$v_2 \con v_1 \pmod{N}\text{ and } u_2 \con u_1 \pmod{\gcd(v_1,N)}.$$
\end{enumerate}
\end{proposition}
\begin{proof}
The first statement is \cite[Prop.~2.2.3]{cremona:algs}, and
the second is \cite[Lem.~3.2]{cremona:gammaone}.
\end{proof}

\begin{algorithm}{Explicit Cusp Equivalence}\label{alg:cusp1}%
Suppose $\smallvtwo{u_1}{v_1}$ and 
$\smallvtwo{u_2}{v_2}$ 
are equivalent modulo $\Gamma_0(N)$. 
This algorithm computes a matrix $g\in\Gamma_0(N)$ such
that $g\smallvtwo{u_1}{v_1}=\smallvtwo{u_2}{v_2}$.
\begin{steps}
\item Let $s_1, s_2, r_1, r_2$ be solutions to
$s_1 u_1 -r_1 v_1 =1$ and
$s_2 u_2 -r_2 v_2 =1$.
\item Find integers $x_0$ and $y_0$ such
that $x_0v_1v_2+y_0N=1$.
\item Let $x=-x_0(s_1v_2-s_2v_1)/(v_1v_2,N)$
and $s_1' = s_1 + xv_1$.
\item Output $g=\abcd{u_2}{r_2}{v_2}{s_2}
       \cdot \abcd{u_1}{r_1}{v_1}{s_1'}^{-1}$,
which sends $\smallvtwo{u_1}{v_1}$ to $\smallvtwo{u_2}{v_2}$.
\end{steps}
\end{algorithm}
\begin{proof}
  See the proof of
 \cite[Prop.~\ref{prop:cusp1}]{cremona:algs}.
\end{proof}

The~$\eps$ relations can make the situation more
complicated, since it is possible that $\eps(\alp)\neq
\eps(\beta)$ but
$$\eps(\alp)\left[\vtwo{u}{v}\right] =\left[\gamma_\alp \vtwo{u}{v}\right]= 
  \left[\gamma_\beta \vtwo{u}{v}\right]=\eps(\beta)\left[\vtwo{u}{v}\right].$$
One way out of this difficulty  is to construct 
the cusp classes for $\Gamma_1(N)$, and then quotient
out by the additional~$\eps$ relations using
Gaussian elimination. This is far too
inefficient to be useful in practice because the number of
$\Gamma_1(N)$ cusp classes can be unreasonably large.  
Instead, we give a quick test to determine whether or not
a cusp vanishes modulo the $\eps$-relations.

\begin{lemma}\label{lem:canlift}
Suppose $\alp$ and $\alp'$ are integers
such that $\gcd(\alp,\alp',N)=1$.
Then there exist integers $\beta$ and $\beta'$,
congruent to $\alp$ and $\alp'$ modulo $N$, respectively,
 such that $\gcd(\beta,\beta')=1$.
\end{lemma}
\begin{proof}
By \exref{ch:modsym}{ex:surjred} the map
$\SL_2(\Z)\ra\SL_2(\Z/N\Z)$ is surjective.
By the Euclidean algorithm, there exist
integers $x$, $y$ and $z$ such that
$x\alp + y\alp' + zN = 1$.  
Consider the matrix
$\abcd{y}{-x}{\alp}{\hfill\alp'}\in \SL_2(\Z/N\Z)$
and take $\beta$, $\beta'$ to be the bottom
row of a lift of this matrix to $\SL_2(\Z)$.
\end{proof}

\begin{proposition}\label{prop:cuspdies}\index{cusps!criterion for vanishing}
Let~$N$ be a positive integer and~$\eps$ a Dirichlet
character\index{Dirichlet character!and cusps} of modulus~$N$. 
Suppose $\smallvtwo{u}{v}$ is a cusp with $u$ and $v$ coprime.
Then $\smallvtwo{u}{v}$ vanishes modulo the relations
$$
\left[\gamma\smallvtwo{u}{v}\right]=
\eps(\gamma)\left[\smallvtwo{u}{v}\right],\qquad
\text{all $\gamma\in\Gamma_0(N)$},
$$
if and only if there exists $\alp\in(\Z/N\Z)^*$,
with $\eps(\alp)\neq 1$, such that
\begin{align*}
 v &\con \alp v \pmod{N},\\
 u &\con \alp u \pmod{\gcd(v,N)}.
\end{align*}
\end{proposition}
\begin{proof}
First suppose such an~$\alp$ exists.
By Lemma~\ref{lem:canlift}
there exists $\beta, \beta'\in\Z$ lifting
$\alp,\alp^{-1}$ such that $\gcd(\beta,\beta')=1$. 
The cusp $\smallvtwo{\beta u}{\beta' v}$ 
has coprime coordinates so,
by Proposition~\ref{prop:cusp1} and our
congruence conditions on~$\alp$, the cusps
$\smallvtwo{\beta{}u}{\beta'{}v}$ 
and $\smallvtwo{u}{v}$ are equivalent by
an element of $\Gamma_1(N)$.
This implies that $\left[\smallvtwo{\beta{}u}{\beta'{}v}\right]
   =\left[\smallvtwo{u}{v}\right]$.  
Since $\left[\smallvtwo{\beta{}u}{\beta'{}v}\right]
  = \eps(\alp)\left[\smallvtwo{u}{v}\right]$
and $\eps(\alp)\neq 1$, we have $\left[\smallvtwo{u}{v}\right]=0$. 

Conversely, suppose $\left[\smallvtwo{u}{v}\right]=0$. 
Because all relations are two-term relations and the
$\Gamma_1(N)$-relations identify $\Gamma_1(N)$-orbits,
there must exists $\alp$ and $\beta$ with
  $$\left[\gamma_\alp \vtwo{u}{v}\right]
      =\left[\gamma_\beta \vtwo{u}{v}\right]
   \qquad\text{ and }\eps(\alp)\ne \eps(\beta).$$
Indeed, if this did not occur,
then we could mod out by the $\eps$ relations by writing
each $\left[\gamma_\alp \smallvtwo{u}{v} \right]$
in terms of  $\left[\smallvtwo{u}{v}\right]$, and there would
be no further relations left to kill 
$\left[\smallvtwo{u}{v}\right]$.
Next observe that
\begin{align*}
\left[\gamma_{\beta^{-1}\alp}
      \vtwo{u}{v}\right]
  &= \left[\gamma_{\beta^{-1}}\gamma_\alp
      \vtwo{u}{v}\right]\\
 &= \eps(\beta^{-1})\left[\gamma_\alp
      \vtwo{u}{v}\right]
 = \eps(\beta^{-1})\left[\gamma_\beta
      \vtwo{u}{v}\right]
 = \left[\vtwo{u}{v}\right].
\end{align*}
Applying Proposition~\ref{prop:cusp1} and
noting that $\eps(\beta^{-1}\alp)\neq 1$ shows 
that $\beta^{-1}\alp$ satisfies the properties
of the ``$\alp$'' in the statement of the
proposition.
\end{proof}

We enumerate the possible~$\alp$ appearing 
in Proposition~\ref{prop:cuspdies} as follows. 
Let $g=(v,N)$ and list the
$\alp=v\cdot\frac{N}{g}\cdot{}a+1$, for $a=0,\ldots,g-1$,
such that $\eps(\alp)\neq 0$. 


\section{Pairing Modular Symbols and Modular Forms}%
\label{sec:pairing}%
In this section we define a pairing between modular symbols and
modular forms that the Hecke operators respect.  
We also define an involution on modular symbols and study
its relationship with the pairing.  This pairing is crucial in much
that follows, because it gives rise to period maps from modular
symbols to certain complex vector spaces.

Fix an integer weight $k\geq 2$ and a finite index subgroup
$\Gamma$ of $\SL_2(\Z)$.  Let $M_k(\Gamma)$ denote the space
of holomorphic modular forms of weight~$k$ for $\Gamma$,
and let $S_k(\Gamma)$ denote its cuspidal subspace.
Following \cite[\S1.5]{merel:1585}, let 
\begin{equation}\label{eqn:sbar}
 \Sbar_k(\Gamma) = \{\overline{f} : f \in S_k(\Gamma)\}
\end{equation}
denote the space of \defn{antiholomorphic} cusp forms.  Here
$\overline{f}$ is the function on $\h^*$ given by
$\overline{f}(z) = \overline{f(z)}$.

Define a pairing
\begin{equation}\label{eqn:intpairing}
 (S_k(\Gamma) \oplus \Sbar_k(\Gamma)) 
   \times \sM_k(\Gamma) \to \C
 \end{equation}
 by letting
$$
  \langle (f_1, f_2), P\{\alpha,\beta\} \rangle 
   = \int_{\alpha}^{\beta} f_1(z) P(z,1)\dz +
    \int_{\alpha}^{\beta} f_2(z) P(\zbar,1) \dzbar
$$
and extending linearly.   Here the integral is a complex
path integral along a simple path from~$\alpha$ 
to~$\beta$  in $\h$ (so, e.g., write $z(t)=x(t) + iy(t)$,
where $(x(t),y(t))$ traces out the path and consider
two real integrals).

\begin{proposition}\label{prop:pairwd}
The integration pairing is well defined, i.e., if we
replace $P\{\alpha,\beta\}$ by an equivalent modular symbol
(equivalent modulo the left action of~$\Gamma$), then the integral
is the same.  
\end{proposition}
\begin{proof}
This follows from the change of variables formulas
for integration and the fact that $f_1 \in S_k(\Gamma)$ and
$f_2 \in \Sbar_k(\Gamma)$.  For example, if $k=2$,
$g\in \Gamma$ and $f\in S_k(\Gamma)$, then
\begin{align*}
  \langle f, g\{\alpha,\beta\}\rangle
  &= \langle f, \{g(\alpha),g(\beta)\}\rangle \\
  &= \int_{g(\alpha)}^{g(\beta)} f(z)dz \\
  &= \int_{\alpha}^{\beta} f(g(z))d g(z) \\
  &= \int_{\alpha}^{\beta} f(z)dz  = 
\langle f, \{\alpha, \beta\}\rangle,
\end{align*}
where $f(g(z))d g(z) = f(z) dz$ because
$f$ is a weight~$2$ modular form.
For the case of arbitrary weight $k>2$, see 
\exref{ch:modsym}{ex:pairingwd}.
\end{proof}

The integration pairing is very relevant to the
study of special values of
$L$-functions.  The $L$-function
of a cusp form $f=\sum a_n q^n \in S_k(\Gamma_1(N))$ is
\begin{align}\label{eqn:lfs}
  L(f,s) &= (2\pi)^{s}\Gamma(s)^{-1} \int_{0}^{\infty} f(i t)t^s \frac{dt}{t}\\
&=\sum_{n=1}^{\infty} \frac{a_n}{n^s} \qquad \text{ for }\Re(s)\gg 0.
\end{align}
The equality of the integral and the
Dirichlet series follows 
by switching the order of summation and integration,
which is justified using standard estimates on $|a_n|$
(see, e.g., \cite[Section~VIII.5]{knapp:elliptic}).

For each integer $j$ with $1 \leq j \leq k-1$, we have,
setting $s=j$ and making the change of variables $t\mapsto -it$
in \eqref{eqn:lfs}, that
\begin{equation}\label{eqn:lfunc_ms}
L(f,j) = \frac{(-2\pi i)^{j}}{(j-1)!}
              \cdot \left\langle f,\,\,
       X^{j-1}Y^{k-2-(j-1)} \{0,\infty\} \right\rangle.
   \end{equation}
The integers $j$ as above are called \defn{critical integers}.
When $f$ is an eigenform, they have deep 
conjectural significance (see \cite{bloch-kato, scholl:motivesinvent}).
One can approximate $L(f,j)$ to any desired precision
by computing the above pairing explicitly using
the method described in Chapter~\ref{ch:periods}.
Alternatively, \cite{dokchitser:lfun} contains
methods for computing special values of very 
general $L$-functions, which can be used for
approximating $L(f,s)$ for arbitrary~$s$, in addition
to just the critical integers $1,\ldots,k-1$.

\begin{theorem}[Shokoruv]\label{thm:shokoruvpairing}
The pairing 
$$
\langle \cdot\, , \, \cdot \rangle: (S_k(\Gamma) \oplus \Sbar_k(\Gamma)) 
   \times \sS_k(\Gamma,\C) \to \C
$$
is a nondegenerate pairing of complex vector spaces.
\end{theorem}
\begin{proof}
This is \cite[Thm.~0.2]{shokurov:kuga}
and \cite[\S1.5]{merel:1585}. 
% This is the higher weight
%analogue of Theorem~\ref{thm:pairing2}.
\end{proof}

\begin{corollary}
We have 
$$
  \dim_\C\sS_k(\Gamma,\C) = 2\dim_\C S_2(\Gamma).
$$
\end{corollary}

The pairing is also compatible with Hecke operators.  Before proving
this, we define an \defn{action of Hecke operators} on $M_k(\Gamma_1(N))$ and
on $\Sbar_k(\Gamma_1(N))$.  The definition is similar to the one
we gave in Section~\ref{sec:hecke_one} for modular forms of level~$1$.
For a positive integer~$n$, let $R_n$ be a set of coset
representatives for $\Gamma_1(N)\backslash \Delta_n$ from
Lemma~\ref{lem:deltan}.  For any $\gamma=\abcd{a}{b}{c}{d} \in \GL_2(\Q)$
and $f\in M_k(\Gamma_1(N))$ set
$$
  f^{[\gamma]_k} = \det(\gamma)^{k-1} (cz+d)^{-k} f(\gamma(z)).
$$
Also, for $f\in \Sbar_k(\Gamma_1(N))$, set
$$
  f^{[\gamma]_k'} = \det(\gamma)^{k-1} (c\zbar+d)^{-k} f(\gamma(z)).
$$
Then for $f\in M_k(\Gamma_1(N))$, 
$$
  T_n(f) = \sum_{\gamma\in R_n} f^{[\gamma]_k}
$$
and for $f\in \Sbar_k(\Gamma_1(N))$, 
$$
  T_n(f) = \sum_{\gamma\in R_n} f^{[\gamma]_k'}.
$$
This agrees with the definition from Section~\ref{sec:hecke_one}
when $N=1$.

\begin{remark}
If $\Gamma$ is an arbitrary finite index subgroup of $\SL_2(\Z)$,
then we can define operators $T_\Delta$ on $M_k(\Gamma)$
for any $\Delta$ with $\Delta\Gamma=\Gamma\Delta = \Delta$
and $\Gamma\backslash \Delta$ finite.   For concreteness we
do not do the general case here or in the theorem below, but
the proof is exactly the same (see \cite[\S1.5]{merel:1585}).
\end{remark}

Finally we prove the promised Hecke compatibility of the pairing.
This proof should convince you that the definition of modular symbols
is sensible, in that they are natural objects to
integrate against modular forms.
\begin{theorem}\label{thm:tequivar}
If $$f=(f_1,f_2)\in S_k(\Gamma_1(N)) \oplus \Sbar_k(\Gamma_1(N))$$ and
$x \in \sM_k(\Gamma_1(N))$, then for any $n$, 
$$
  \langle T_n(f), x \rangle = \langle f, T_n(x) \rangle.
$$
\end{theorem}
\begin{proof}
  We follow \cite[\S2.1]{merel:1585} (but with more details) and will
  only prove the theorem when $f=f_1\in S_k(\Gamma_1(N))$, the proof
  in the general case being the same.

Let $\alpha, \beta \in \P^1(\Q)$, $P\in \Z[X,Y]_{k-2}$,
and for $g=\abcd{a}{b}{c}{d}\in\GL_2(\Q)$, set 
$$j(g,z) = cz+d.$$  Let $n$ be any positive integer,
and let $R_n$ be a set of coset
representatives for $\Gamma_1(N)\backslash \Delta_n$ from
Lemma~\ref{lem:deltan}.

We have
\begin{align*}
\langle T_n(f), P\{\alpha,\beta\}\rangle &= 
    \int_{\alpha}^{\beta} T_n(f) P(z,1) \dz\\
   &= \sum_{\delta\in R} \int_{\alpha}^{\beta}
            \det(\delta)^{k-1} f(\delta(z)) j(\delta,z)^{-k} P(z,1) \dz.
\end{align*}
Now for each summand corresponding to the $\delta\in R$,
make the change of variables $u=\delta z$.  Thus we 
make $\# R$ change of variables.  
Also, we will use the notation $$
\tilde{g} = \mtwo{\hfill{}d}{-b}{-c}{\hfill{}a}
           = \det(g)\cdot g^{-1}$$
for $g\in\GL_2(\Q)$.
We have
\begin{align*}
\langle T_n(f), P\{\alpha,\beta\}\rangle &=\\
   &\!\!\!\!\!\!\!\!\!\!\!\!\!\!\!\!\!\!\!\!
\sum_{\delta \in R}\int_{\delta(\alpha)}^{\delta(\beta)}
             \det(\delta)^{k-1} f(u) j(\delta, \delta^{-1}(u))^{-k}
             P(\delta^{-1}(u), 1) d(\delta^{-1}(u)).
\end{align*}
We have $\delta^{-1}(u) = \tilde{\delta}(u)$,
since a linear fractional transformation is unchanged
by a nonzero rescaling of a matrix that induces it.
Thus by the quotient rule, using that $\tilde{\delta}$
has determinant $\det(\delta)$, we see that
$$
   d(\delta^{-1}(u)) = d(\tilde{\delta}(u))
   = \frac{\det(\delta) du}{j(\tilde{\delta},u)^2}.
$$
We next show that
\begin{equation}\label{eqn:bigstuff}
j(\delta, \delta^{-1}(u))^{-k} P(\delta^{-1}(u),1)
=
j(\tilde{\delta},u)^k \det(\delta)^{-k} P(\tilde{\delta}(u),1).
\end{equation}
From the definitions, and again using that $\delta^{-1}(u)=\tilde{\delta}(u)$, we see that
$$
  j(\delta, \delta^{-1}(u)) = \frac{\det(\delta)}{j(\tilde{\delta}, u)},
$$
which proves that \eqref{eqn:bigstuff} holds.
Thus
\begin{align*}
\left\langle T_n(f), P\{\alpha,\beta\}\right\rangle &= \\
&\!\!\!\!\!\!\!\!\!\!\!\!\!\!\!\!\!\!\!\!\!\!\!
   \sum_{\delta \in R}\int_{\delta(\alpha)}^{\delta(\beta)}
             \det(\delta)^{k-1} f(u) j(\tilde{\delta},u)^{k}
             \det(\delta)^{-k}
             P(\tilde{\delta}(u), 1) \frac{\det(\delta) du}
                                        {j(\tilde{\delta}, u)^2}.
\end{align*}
Next we use that 
$$
  (\delta P)(u,1) = j(\tilde{\delta}, u)^{k-2}
                    P(\tilde{\delta}(u), 1).
$$
To see this, note that $P(X,Y) = P(X/Y, 1)\cdot Y^{k-2}$.
Using this we see that
\begin{align*}
(\delta P)(X,Y) &= (P\circ \tilde{\delta})(X,Y)\\
      &= P\left(\tilde{\delta}\left(\frac{X}{Y}\right),1\right) 
       \left(-c \cdot \frac{X}{Y} + a\right)^{k-2} \cdot Y^{k-2}.
\end{align*}
Now substituting $(u,1)$ for $(X,1)$, we see that
$$ (\delta P)(u,1) = 
  P(\tilde{\delta}(u), 1) \cdot (-c u + a)^{k-2},$$
as required.
Thus finally
\begin{align*}
\langle T_n(f), P\{\alpha,\beta\}\rangle 
  &= \sum_{\delta \in R}\int_{\delta(\alpha)}^{\delta(\beta)}
             f(u) j(\tilde{\delta},u)^{k-2}
             P(\tilde{\delta}(u), 1) du\\
  &= \sum_{\delta \in R} \int_{\delta(\alpha)}^{\delta(\beta)}
           f(u) \cdot ((\delta P)(u,1)) du\\
  &= \langle f, \, T_n(P\{\alpha,\beta\})\rangle.
\end{align*}
\end{proof}


Suppose  that $\Gamma$ is a finite index
subgroup of $\SL_2(\Z)$ such that 
if $\eta = \abcd{-1}{0}{\hfill 0}{1}$, then 
$$
 \eta \Gamma \eta = \Gamma.
$$
For example, $\Gamma=\Gamma_1(N)$ satisfies this condition.
There is an involution $\iota^*$ on $\sM_k(\Gamma)$ given by
\begin{equation}\label{eqn:star}
  \iota^*(P(X,Y)\{\alpha,\beta\})
    = -P(X,-Y)\{-\alpha,-\beta\},
\end{equation}
which we call the \defn{star involution}.  
On Manin symbols, $\iota^*$ is 
\begin{equation}\label{eqn:iota}
  \iota^*[P,(u,v)] = -[P(-X,Y), (-u,v)].
\end{equation}
Let $\sS_k(\Gamma)^+$ be the $+1$ eigenspace for
$\iota^*$ on $\sS_k(\Gamma)$, and let $\sS_k(\Gamma)^-$ be the $-1$ eigenspace.
There is also a map $\iota$ on modular forms, which
is adjoint to $\iota^*$.
\begin{remark}
  Notice the minus sign in front of $-P(X,-Y)\{-\alpha,-\beta\}$ in
  \eqref{eqn:star}.  This sign is missing in \cite{cremona:algs},
  which is a potential source of confusion (this is not a mistake,
  but  a different choice of convention).
\end{remark}

We now state the final result about the pairing, which
explains how modular symbols and modular forms
are related.
\begin{theorem}\label{thm:modsymformpairing}
  The integration pairing $\langle \cdot\, , \, \cdot \rangle$ induces
  nondegenerate Hecke compatible  bilinear pairings 
$$
 \sS_k(\Gamma)^+ \times S_k(\Gamma) \to \C
\qquad \text{and}\qquad
\sS_k(\Gamma)^- \times \Sbar_k(\Gamma) \to \C.
$$
\end{theorem}


\begin{remark}
We make some remarks about computing
the boundary map of Section~\ref{sec:computeboundary} when working
in the $\pm 1$ quotient.
Let~$s$ be a sign, either~$+1$
or~$-1$.  To compute $\sS_k(N,\eps)_s$, it is necessary to replace
$B_k(N,\eps)$ by its quotient modulo the additional relations $\left[
  \smallvtwo{-u}{\hfill v}\right] = s \left[\smallvtwo{u}{v}\right]$
for all cusps $\smallvtwo{u}{v}$.  Algorithm~\ref{alg:cusplist} can be
modified to deal with this situation as follows.  Given a cusp
$x=\smallvtwo{u}{v}$, proceed as in Algorithm~\ref{alg:cusplist} and
check if either $\smallvtwo{u}{v}$ or $\smallvtwo{-u}{\hfill v}$ is
equivalent (modulo scalars) to any cusp seen so far.  If not, use the
following trick to determine whether the $\eps$ and $s$-relations kill
the class of $\smallvtwo{u}{v}$: use the unmodified
Algorithm~\ref{alg:cusplist} to compute the scalars $\alp_1, \alp_2$
and indices $i_1$, $i_2$ associated to $\smallvtwo{u}{v}$ and
$\smallvtwo{-u}{\hfill v}$, respectively.  The $s$-relation kills the
class of $\smallvtwo{u}{v}$ if and only if $i_1=i_2$ but $\alp_1\neq
s\alp_2$.
\end{remark}


\section{Degeneracy Maps}\label{sec:degen}
In this section, we describe natural maps between spaces of
modular symbols with character of different levels. 
We consider spaces with character, since
they are so important in applications. 

Fix a positive integer~$N$ and a Dirichlet 
character\index{Dirichlet character} 
$\eps : (\Z/N\Z)^*\ra \C^*$.  Let~$M$ be a positive divisor
of~$N$ that is divisible by the conductor of~$\eps$, in the sense
that~$\eps$ factors through $(\Z/M\Z)^*$ via the natural map
$(\Z/N\Z)^*\ra (\Z/M\Z)^*$ composed with some uniquely defined
character $\eps':(\Z/M\Z)^*\ra\C^*$.  For any positive divisor~$t$ of
$N/M$, let $T=\abcd{1}{0}{0}{t}$ and fix a choice $D_t=\{T\gamma_i :
i=1,\ldots, n\}$ of coset representatives for $\Gamma_0(N)\backslash
T\Gamma_0(M)$.

\begin{remark}
Note that \cite[\S2.6]{merel:1585} contains a typo:
 The quotient ``$\Gamma_1(N)\backslash\Gamma_1(M)T$'' should be replaced 
by ``$\Gamma_1(N)\backslash T\Gamma_1(M)$''.
\end{remark}
\begin{proposition}
For each divisor~$t$ of $N/M$ there are well-defined linear maps
\begin{align*}
\alp_t : \sM_k(N,\eps) \ra \sM_k(M,\eps'),&\qquad
      \alp_t(x) = (tT^{-1})x = \mtwo{t}{0}{0}{1} x,\\
\beta_t : \sM_k(M,\eps') \ra \sM_k(N,\eps),&\qquad
      \beta_t(x) = \sum_{T\gamma_i\in D_t} \eps'(\gamma_i)^{-1}T\gamma_i{} x.
\end{align*}
Furthermore,
  $\alp_t\circ \beta_t$ is multiplication by 
  $t^{k-2}\cdot [\Gamma_0(M) : \Gamma_0(N)].$
\end{proposition}
\begin{proof}
To show that~$\alp_t$ is well defined, we must show that for
each $x\in\sM_k(N,\eps)$ and $\gamma=\abcd{a}{b}{c}{d}\in\Gamma_0(N)$, 
we have
  $$\alp_t(\gamma x  -\eps(\gamma) x)=0\in\sM_k(M,\eps').$$
We have
$$\alp_t(\gamma{}x) = \mtwo{t}{0}{0}{1}\gamma{}x
        = \mtwo{a}{tb}{c/t}{d}\mtwo{t}{0}{0}{1} x
        = \eps'(a)\mtwo{t}{0}{0}{1} x,$$
so 
$$\alp_t(\gamma x  -\eps(\gamma) x)
  = \eps'(a)\alp_t(x) - \eps(\gamma)\alp_t(x) = 0.$$

We next verify that~$\beta_t$ is well defined.
Suppose that $x\in\sM_k(M,\eps')$ and $\gamma\in\Gamma_0(M)$;
then $\eps'(\gamma)^{-1}\gamma{}x = x$, so
\begin{align*}
\beta_t(x) 
    &= \sum_{T\gamma_i\in D_t} 
        \eps'(\gamma_i)^{-1}T\gamma_i{}\eps'(\gamma)^{-1}\gamma{} x\\
    &= \sum_{T\gamma_i\gamma\in D_t} 
        \eps'(\gamma_i\gamma)^{-1}T\gamma_i{}\gamma{} x.
\end{align*}
This computation shows that the definition of~$\beta_t$ 
does not depend on the choice~$D_t$ of coset representatives.
To finish the proof that~$\beta_t$ is well defined,
we must show that, for $\gamma\in\Gamma_0(M)$, we have
$\beta_t(\gamma{}x) = \eps'(\gamma)\beta_t(x)$ so that $\beta_t$
respects the relations that define $\sM_k(M,\eps)$. 
Using that~$\beta_t$ does not depend on the choice of 
coset representative, we find that for $\gamma\in\Gamma_0(M)$,
\begin{align*}
  \beta_t(\gamma{}x) 
     &= \sum_{T\gamma_i\in D_t} \eps'(\gamma_i)^{-1}T\gamma_i{} \gamma{} x\\
     &= \sum_{T\gamma_i\gamma^{-1}\in D_t} 
         \eps'(\gamma_i\gamma^{-1})^{-1}T\gamma_i{}\gamma{}^{-1} \gamma{} x\\
     &= \eps'(\gamma)\beta_t(x).\\
\end{align*}
To compute $\alp_t\circ\beta_t$, we use 
that $\#D_t = [\Gamma_0(N) : \Gamma_0(M)]$:
\begin{align*}
 \alp_t(\beta_t(x)) &=
    \alp_t \left(\sum_{T\gamma_i}
        \eps'(\gamma_i)^{-1}T\gamma_i x\right)\\
  &= \sum_{T\gamma_i}
        \eps'(\gamma_i)^{-1}(tT^{-1})T\gamma_i x\\
  &= t^{k-2}\sum_{T\gamma_i}
        \eps'(\gamma_i)^{-1}\gamma_i x\\
  &= t^{k-2}\sum_{T\gamma_i} x \\
  &= t^{k-2} \cdot [\Gamma_0(N) : \Gamma_0(M)] \cdot x.
\end{align*}
The scalar factor of $t^{k-2}$ appears instead 
of~$t$, because~$t$ is acting on~$x$ as an element of $\GL_2(\Q)$
and {\em not} as an an element of~$\Q$.
\end{proof}

\begin{definition}[New and Old Modular Symbols]%
\label{def:newandoldsymbols}%
\index{new modular symbols|textit}%
\index{old modular symbols|textit}%
\index{modular symbols!new and old subspace of|textit}%
The space $\sM_k(N,\eps)_{\new}$ 
of \defn{new modular symbols} is the
intersection of the kernels of the $\alp_t$ as~$t$ 
runs through all positive divisors of $N/M$ and~$M$
runs through positive divisors of~$M$ strictly less than~$N$
and divisible by the conductor of~$\eps$.
The subspace $\sM_k(N,\eps)_{\old}$ 
of \defn{old modular symbols}
is the subspace generated by the images of the $\beta_t$
where~$t$ runs through all positive divisors of $N/M$ and~$M$
runs through positive divisors of~$M$ strictly less than~$N$
and divisible by the conductor of~$\eps$.
The new and old subspaces of cuspidal modular
symbols are the intersections of the above spaces
with $\sS_k(N,\eps)$.
\end{definition}

\begin{example}
  The new and old subspaces need not be disjoint, as the following
  example illustrates!  (This contradicts ~\cite[pg.~80]{merel:1585}.)
  Consider, for example, the case $N=6$, $k=2$, and trivial character.
  The spaces $\sM_2(\Gamma_0(2))$ and $\sM_2(\Gamma_0(3))$ are each of
  dimension~$1$, and each is generated by the modular symbol
  $\{\infty,0\}$.  The space $\sM_2(\Gamma_0(6))$ is of dimension~$3$
  and is generated by the three modular symbols $\{\infty, 0\}$,
  $\{-1/4, 0\}$, and $\{-1/2, -1/3\}$.  The space generated by the two
  images of $\sM_2(\Gamma_0(2))$ under the two degeneracy maps has
  dimension~$2$, and likewise for $\sM_2(\Gamma_0(3))$.  Together
  these images generate $\sM_2(\Gamma_0(6))$, so $\sM_2(\Gamma_0(6))$
  is equal to its old subspace.  However, the new subspace is
  nontrivial because the two degeneracy maps $\sM_2(\Gamma_0(6)) \ra
  \sM_2(\Gamma_0(2))$ are equal, as are the two degeneracy maps
  $$\sM_2(\Gamma_0(6)) \ra \sM_2(\Gamma_0(3)).$$  In particular, the
  intersection of the kernels of the degeneracy maps has dimension at
  least~$1$ (in fact, it equals~$1$).
%begin_sage
We verify some of the above claims using \sage.
\begin{verbatim}
sage: M = ModularSymbols(Gamma0(6)); M
Modular Symbols space of dimension 3 for Gamma_0(6) 
of weight 2 with sign 0 over Rational Field
sage: M.new_subspace()
Modular Symbols subspace of dimension 1 of Modular 
Symbols space of dimension 3 for Gamma_0(6) of weight 
2 with sign 0 over Rational Field
sage: M.old_subspace()
Modular Symbols subspace of dimension 3 of Modular 
Symbols space of dimension 3 for Gamma_0(6) of weight 
2 with sign 0 over Rational Field
\end{verbatim}
%end_sage
\end{example}


\section{Explicitly Computing $\sM_k(\Gamma_0(N))$}\label{sec:compg0n}
In this section we explicitly compute $\sM_k(\Gamma_0(N))$ for various
$k$ and $N$.  We represent Manin symbols for $\Gamma_0(N)$ as triples
of integers $(i, u,v)$, where $(u,v)\in\P^1(\Z/N\Z)$, and $(i, u,v)$ corresponds
to $[X^i Y^{k-2-i}, (u,v)]$ in the usual notation.  Also, recall from
Proposition~\ref{prop:p1bij} that $(u,v)$ corresponds to the right
coset of $\Gamma_0(N)$ that contains a matrix
$\abcd{a}{b}{c}{d}$ with $(u,v)\con (c,d)$ as elements of
$\P^1(\Z/N\Z)$, i.e., up to rescaling by an element of $(\Z/N\Z)^*$.

\subsection{Computing $\P^1(\Z/N\Z)$}\label{sec:p1rep}
In this section we give an algorithm to compute a canonical
representative for each element of $\P^1(\Z/N\Z)$.  This algorithm is
extremely important because modular symbols implementations use it a
huge number of times.  A more naive approach would be to store all
pairs $(u,v)\in (\Z/N\Z)^2$ and a fixed reduced representative, but
this wastes a huge amount of memory.  For example, if $N=10^5$,
we would store an array of
$$
  2 \cdot 10^5 \cdot 10^5 = 20 \text{ billion integers}.
$$

Another approach to enumerating $\P^1(\Z/N\Z)$ is described at the end
of \cite[\S2.2]{cremona:algs}. It uses the fact that is easy to
test whether two pairs $(u_0,v_0), (u_1,v_1)$ define the same element
of $\P^1(\Z/N\Z)$; they do if and only if we have equality of cross
terms $u_0 v_1 = v_0 u_1\pmod{N}$ (see
\cite[Prop.~2.2.1]{cremona:algs}).  So we consider the
$0$-based list of elements 
\begin{equation}\label{eqn:zerobased}
 (1,0), (1,1), \ldots, (1,N-1), (0,1)
\end{equation}
concatenated with the list of nonequivalent elements
$(d,a)$ for $d\mid N$ and
$a=1,\ldots, N-1$, checking each time we add a new element to our
list (of $(d,a)$) whether we have already seen it.  

Given a random pair $(u,v)$, the problem
is then to find the index of the element of our list of the equivalent 
representative in $\P^1(\Z/N\Z)$. 
We use the following
algorithm, which finds a canonical representative for each element of
$\P^1(\Z/N\Z)$.  Given an arbitrary
$(u,v)$, we first find the canonical equivalent elements $(u',v')$.
If $u'=1$, then the index is $v'$.  If
$u'\neq 1$, we find the corresponding element
in an explicit sorted list, e.g., using binary search.


In the following algorithm, $a \pmod{N}$ denotes the residue of~$a$
modulo~$N$ that satisfies $0 \leq a < N$.
Note that we {\em never} create and store the
list \eqref{eqn:zerobased} itself in memory.
\begin{algorithm}{Reduction in $\P^1(\Z/N\Z)$ 
to Canonical Form}\label{alg:p1rep}
Given $u$ and $v$ and a positive integer~$N$,
this algorithm outputs 
a pair $u_0, v_0$ such that $(u,v) \con (u_0,v_0)$ as
  elements of $\P^1(\Z/N\Z)$ and $s\in\Z$ such that $(u,v) = (su_0,
  sv_0) \pmod{\Z/n\Z}$.  
Moreover, the element $(u_0,v_0)$ does not
  depend on the class of $(u,v)$, i.e., for any~$s$ with $\gcd(N,s)=1$
  the input $(su,sv)$ also outputs $(u_0,v_0)$.  
  If $(u,v)$ is not in $\P^1(\Z/N\Z)$, this algorithm
  outputs $(0,0), 0$.
\begin{steps}
\item{}[Reduce] Reduce both $u$ and $v$ modulo~$N$.
\item{}[Easy $(0,1)$ Case] If $u=0$, check that $\gcd(v,N)=1$.
If so, return $s=1$ and $(0,1)$;  otherwise return $0$.
\item{}[GCD] Compute $g=\gcd(u,N)$ and $s,t\in\Z$ such that $g=su+tN$.
\item{}[Not in $P^1$?] We have $\gcd(u,v,N)=\gcd(g,v)$, so if 
$\gcd(g,v)>1$, then $(u,v)\not\in\P^1(\Z/N\Z)$,
and we return $0$.
\item{}[Pseudo-Inverse]
Now $g = su + tN$, so we may think of~$s$ as ``pseudo-inverse''
  of $u \pmod{N}$, in the sense that $su$ is as close as possible to
  being $1$ modulo~$N$. Note that since $g\mid u$, changing~$s$ modulo
  $N/g$ does not change $su\pmod{N}$.  We can adjust~$s$ modulo $N/g$
  so it is coprime to $N$ (by adding multiples of
$N/g$ to $s$).  (This is because $1 = s u/g + tN/g$, so
  $s$ is a unit mod $N/g$, and the map $(\Z/N\Z)^* \to (\Z/(N/g)\Z)^*$
 is surjective, e.g., as we saw in the proof of 
  Algorithm~\ref{alg:restrict}.)

\item{}[Multiply by $s$] Multiply $(u,v)$ by $s$, and replace  
$(u,v)$  by the
   equivalent element $(g,sv)$ of $\P^1(\Z/N\Z)$.

\item{}[Normalize] 
Compute the unique pair $(g,v')$ equivalent to $(g,v)$ that
minimizes $v$, as follows:
\begin{steps}
\item{}[Easy Case] If $g=1$, this pair is $(1,v)$.

\item{}[Enumerate and Find Best]
 Otherwise, note that if $1\neq t\in(\Z/N\Z)^*$ and
$tg \con g\pmod{N}$, then $(t-1)g \con 0\pmod{N}$, so $t-1 = kN/g$
for some $k$ with $1\leq k \leq g-1$.   Then for $t=1+kN/g$ coprime to~$N$, 
we have $(gt, vt) = (g, v+kvN/g)$.  So we compute
all pairs $(g, v+kvN/g)$ and pick out the one that minimizes the 
least nonnegative residue of $vt$ modulo~$N$.

\item{}[Invert $s$ and Output]
The $s$ that we computed in the above steps multiplies
the input $(u,v)$ to give the output $(u_0,v_0)$.  Thus we
invert it, since the scalar  we output
multiplies $(u_0,v_0)$ to give $(u,v)$.
\end{steps}

\end{steps}

\end{algorithm}

\begin{remark}
  In the above algorithm, there are many gcd's with $N$ so one should
  create a table of the gcd's of $0,1,2,\ldots, N-1$ with~$N$.
\end{remark}
\begin{remark}
Another approach is to instead use that
 $$\P^1(\Z/N\Z)\isom \prod_{p\mid N} \P^1(\Z/p^{\nu_p}\Z),$$
where $\nu_p = \ord_p(N)$,
 and that it is relatively easy to enumerate the
 elements of $\P^1(\Z/p^n\Z)$ for a prime power $p^n$.
\end{remark}

\begin{algorithm}{List $\P^1(\Z/N\Z)$}
  Given an integer $N>1$, this algorithm makes a sorted list of the
  distinct representatives $(c,d)$ of $\P^1(\Z/N\Z)$ with 
  $c\neq 0,1$, as output by Algorithm~\ref{alg:p1rep}.
\begin{steps}
\item  For each $c=1,\ldots, N-1$ with $g = \gcd(c,N) > 1$ do the
following:
\begin{steps}
  \item Use Algorithm~\ref{alg:p1rep} to compute
the canonical representative $(u',v')$ equivalent to $(c,1)$,
and include it in the list.
  \item If $g=c$, for each $d=2,\ldots, N-1$ with $\gcd(d,N)>1$
and $\gcd(c,d)=1$, append the normalized representative
of $(c,d)$ to the list.
\end{steps}
  
\item Sort the list.

\item Pass through the sorted list and delete any duplicates.

\end{steps}

\end{algorithm}

\section{Explicit Examples}
Explicit detailed examples are crucial when implementing modular
symbols algorithms from scratch.  This section contains a number of
such examples.
\subsection{Examples of Computation of $\sM_k(\Gamma_0(N))$}
In this section, we compute $\sM_k(\Gamma_0(N))$ explicitly in a few cases.
\begin{example}\label{example:m4_1}
We compute $V=\sM_4(\Gamma_0(1))$.   Because $S_k(\Gamma_0(1))=0$
and $M_k(\Gamma_0(1))=\C E_4$, we expect $V$ to have dimension~$1$
and for each integer $n$ the Hecke operator $T_n$ to have eigenvalue the
sum $\sigma_3(n)$ of the cubes of positive divisors of $n$.

The Manin symbols are 
$$
  x_0 = (0, 0, 0), \quad x_1 = (1, 0, 0), \quad x_2 = (2, 0, 0).
$$
The relation matrix is
$$
\left(\begin{matrix}
{}1&{}0&{}1\\
{}0&{}0&{}0\\
\hline
{}2&{}-2&{}2\\
{}1&{}-1&{}1\\
{}2&{}-2&{}2
\end{matrix}\right),
$$
where the first two rows correspond to $S$-relations and the second
three to $T$-relations. 
Note that we do not include all $S$-relations, since it is obvious
that some are redundant, e.g., $x+xS=0$ and $(xS) + (xS)S=xS + x=0$
are the same since~$S$ has order~$2$.  

The echelon form of the relation matrix
is
$$
\left(\begin{matrix}
\hfill{}1&\hfill{}0&\hfill{}1\\
\hfill{}0&\hfill{}1&\hfill{}0\\
\end{matrix}\right),
$$
where we have deleted the zero rows from the bottom.
Thus we may replace the above complicated list
of relations with the following simpler list of relations:
\begin{align*}
   x_0 + x_2  &= 0,\\
    x_1 &= 0
\end{align*}
from which we immediately read off that the second generator $x_1$ is
$0$ and $x_0=-x_2$.  Thus $\sM_4(\Gamma_0(1))$ has dimension~$1$, with
basis the equivalence class of $x_2$ (or of $x_0$).

Next we compute the Hecke operator $T_2$ on $\sM_4(\Gamma_0(1))$.
The Heilbronn matrices\index{Heilbronn matrices}
of determinant $2$ from Proposition~\ref{prop:heilbronn}
are
\begin{align*}
h_0&=\left(\begin{matrix}
\hfill{}1&\hfill{}0\\
\hfill{}0&\hfill{}2
\end{matrix}\right),\quad{}\\
h_1&=\left(\begin{matrix}
\hfill{}1&\hfill{}0\\
\hfill{}1&\hfill{}2
\end{matrix}\right),\quad{}\\
h_2&=\left(\begin{matrix}
\hfill{}2&\hfill{}0\\
\hfill{}0&\hfill{}1
\end{matrix}\right),\quad{}\\
h_3&=\left(\begin{matrix}
\hfill{}2&\hfill{}1\\
\hfill{}0&\hfill{}1
\end{matrix}\right).
\end{align*}
To compute $T_2$, we apply each of these matrices to $x_0$,
then reduce modulo the relations.
We have
\begin{align*}
  x_2\mtwo{1}{0}{0}{2} &= [X^2,(0,0)] \mtwo{1}{0}{0}{2}
         x_2,\\
  x_2 \mtwo{1}{0}{1}{2} &= [X^2,(0,0)] = x_2,\\
  x_2 \mtwo{2}{0}{0}{1} &= [(2X)^2,(0,0)] = 4x_2,\\
  x_2 \mtwo{2}{1}{0}{1} &= [(2X+1)^2,(0,0)] = x_0 + 4x_1 + 4x_2
                                            \sim 3x_2.
\end{align*}
Summing we see that $T_2(x_2) \sim 9x_2$ in $\sM_4(\Gamma_0(1))$.
Notice that 
$$9=1^3+2^3 = \sigma_3(2).$$

The Heilbronn matrices\index{Heilbronn matrices} of determinant $3$ from 
Proposition~\ref{prop:heilbronn} are
\begin{align*}
h_0&=\left(\begin{matrix}
\hfill{}1&\hfill{}0\\
\hfill{}0&\hfill{}3
\end{matrix}\right),\quad{}
h_1=\left(\begin{matrix}
\hfill{}1&\hfill{}0\\
\hfill{}1&\hfill{}3
\end{matrix}\right),\quad{}\\
h_2&=\left(\begin{matrix}
\hfill{}1&\hfill{}0\\
\hfill{}2&\hfill{}3
\end{matrix}\right),\quad{}
h_3=\left(\begin{matrix}
\hfill{}2&\hfill{}1\\
\hfill{}1&\hfill{}2
\end{matrix}\right),\quad{}\\
h_4&=\left(\begin{matrix}
\hfill{}3&\hfill{}0\\
\hfill{}0&\hfill{}1
\end{matrix}\right),\quad{}
h_5=\left(\begin{matrix}
\hfill{}3&\hfill{}1\\
\hfill{}0&\hfill{}1
\end{matrix}\right),\quad{}\\
h_6&=\left(\begin{matrix}
\hfill{}3&\hfill{}2\\
\hfill{}0&\hfill{}1
\end{matrix}\right).
\end{align*}

We have
\begin{align*}
  x_2 \mtwo{1}{0}{0}{3} &= [X^2,(0,0)] \mtwo{1}{0}{0}{3} =
         x_2,\\
  x_2 \mtwo{1}{0}{1}{3} &= [X^2,(0,0)] = x_2,\\
  x_2 \mtwo{1}{0}{2}{3} &= [X^2,(0,0)] = x_2,\\
  x_2 \mtwo{2}{1}{2}{2} &= [(2X+1)^2,(0,0)] = x_0 + 4x_1 + 4x_2
                                            \sim 3x_2,\\
  x_2 \mtwo{3}{0}{0}{1} &= [(3X)^2,(0,0)] = 9x_2,\\
  x_2 \mtwo{3}{1}{0}{1} &= [(3X+1)^2,(0,0)] = x_0 + 6x_1 + 9x_2 \sim 8x_2,\\
  x_2 \mtwo{3}{2}{0}{1} &= [(3X+2)^2,(0,0)] = 4x_0 + 12x_1 + 9x_2 \sim 5x_2.\\
\end{align*}
Summing we see that 
$$T_3(x_2) \sim x_2 + x_2 + x_2 + 3x_2 + 9x_2 + 8x_2 + 5x_2
           = 28 x_2.
$$
Notice that 
$$28=1^3+3^3 = \sigma_3(3).$$

\end{example}

\begin{example}
Next we compute $\sM_2(\Gamma_0(11))$ explicitly.
The Manin symbol generators are 
$$
x_0= (0, 1),\quad\!\!\!\!
x_1= (1, 0),\quad\!\!\!\!
x_2= (1, 1),\quad\!\!\!\!
x_3= (1, 2),\quad\!\!\!\!
x_4= (1, 3),\quad\!\!\!\!
x_5= (1, 4),\quad\!\!
$$
$$
x_6= (1, 5),\quad\!\!\!\!
x_7= (1, 6),\quad\!\!\!\!
x_8= (1, 7),\quad\!\!\!\!
x_9= (1, 8),\quad\!\!\!\!
x_{10}= (1, 9),\quad\!\!\!\!
x_{11}= (1, 10).
$$
The relation matrix is as follows, where the $S$-relations 
are above the line and the $T$-relations
are below it:
$$
\left(\begin{array}{ccccccccccccc}
1&1&0&0&0&0&0&0&0&0&0&0\\
0&0&1&0&0&0&0&0&0&0&0&1\\
0&0&0&1&0&0&1&0&0&0&0&0\\
0&0&0&0&1&0&0&0&1&0&0&0\\
0&0&0&0&0&1&0&0&0&1&0&0\\
0&0&0&0&0&0&0&1&0&0&1&0\\
\hline
1&1&0&0&0&0&0&0&0&0&0&1\\
0&0&1&0&0&0&1&0&0&0&1&0\\
0&0&0&1&0&1&0&0&1&0&0&0\\
0&0&0&0&1&0&0&1&0&1&0&0
\end{array}\right).
$$
In weight $2$, two out of three $T$-relations are redundant, so we
do not include them.
The reduced row echelon form of the relation matrix is
$$
\left(\begin{array}{cccccccccccc}
1&1&0&0&0&0&0&0&0&0&0&0\\
0&0&1&0&0&0&0&0&0&0&0&0\\
0&0&0&1&0&0&0&0&0&0&-1&0\\
0&0&0&0&1&0&0&0&0&1&-1&0\\
0&0&0&0&0&1&0&0&0&1&0&0\\
0&0&0&0&0&0&1&0&0&0&1&0\\
0&0&0&0&0&0&0&1&0&0&1&0\\
0&0&0&0&0&0&0&0&1&-1&1&0\\
0&0&0&0&0&0&0&0&0&0&0&1\\
0&0&0&0&0&0&0&0&0&0&0&0
\end{array}\right).
$$
From the echelon form we see that every symbol
is equivalent to a combination of
$x_1=(1,0)$, $x_9=(1,8)$, and $x_{10}=(1,9)$.
(Notice that columns $1$, $9$, and $10$ are the pivot
columns, where we index columns starting at $0$.)
%Explicitly, if $(a,b,c)$ is the $i$th row of the following
%matrix, then  $x_i = a x_1 + b x_9 + c x_{10}$:
%$$\left(\begin{array}{ccc}
%-1&0&0\\
%1&0&0\\
%0&0&0\\
%0&0&1\\
%0&-1&1\\
%0&-1&0\\
%0&0&-1\\
%0&0&-1\\
%0&1&-1\\
%0&1&0\\
%0&0&1\\
%0&0&0
%\end{array}\right).$$

To compute $T_2$, we apply each of the Heilbronn
matrices\index{Heilbronn matrices} of determinant~$2$ from Proposition~\ref{prop:heilbronn}
to $x_1$, then to $x_9$, and finally to $x_{10}$.
The matrices are as in Example~\ref{example:m4_1} above.
We have
$$
  T_2(x_1) = 3(1,0) + (1,6) \sim 3x_1 - x_{10}.
$$
Applying $T_2$ to $x_9=(1,8)$, we get
$$
 T_2(x_9) =  (1,3)  + (1,4) + (1,5) + (1,10) \sim -2x_9.
$$
Applying $T_2$ to $x_{10}=(1,9)$, we get
$$
T_2(x_{10})= (1,4)+(1,5) + (1,7) + (1,10) \sim -x_1 - 2x_{10}.
$$
Thus the matrix of $T_2$ with respect to this basis is
$$
T_2 = 
\left(\begin{array}{ccc}
3&0&0\\
0&-2&0\\
-1&0&-2
\end{array}\right),
$$
where we write the matrix as an operator on the left on vectors
written in terms of $x_1$, $x_9$, and $x_{10}$.
The matrix $T_2$ has characteristic polynomial
$$
  (x-3) (x+2)^2.
$$
The $(x-3)$ factor corresponds to the weight $2$ Eisenstein
series, and the $x+2$ factor corresponds to the elliptic
curve $E=X_0(11)$, which has $$a_2 = -2 = 2+1 - \#E(\F_2).$$

%We have
%\begin{align*}
% T_3(x_1) &= 4 (1,0) + (1,4) + (1,6) + (1,8) \sim 4 x_1 - x_{10},\\
% T_3(x_9) &= (1,2)  + (1,3) + (1,4)+ (1,5) + (1,7) + 2(1,10) \sim -x_9,\\
% T_3(x_{10}) &= (0,1) + (1,0) + (1,2) + (1,3) + (1,5) + (1,6) +(1,7)
%    \sim -x_{10},
%\end{align*}
%so
%$$
%T_3 = 
%\left(\begin{array}{ccc}
%4&0&0\\
%0&-1&0\\
%-1&0&-1
%\end{array}\right).
%$$
%The characteristic polynomial of $T_3$ is 
%$(x-4)(x+1)^2$.
\end{example}

\begin{example}
In this example, we compute $\sM_6(\Gamma_0(3))$, 
which illustrates both weight greater than $2$
and level greater than $1$.
We have the following generating Manin symbols:
\begin{align*}
x_{0} &= [XY^4,(0,1)],\qquad x_{1} = [XY^4,(1,0)],\\
x_{2} &= [XY^4,(1,1)], \qquad x_{3} = [XY^4,(1,2)],\\
x_{4} &= [XY^3,(0,1)], \qquad x_{5} = [XY^3,(1,0)],\\
x_{6} &= [XY^3,(1,1)], \qquad 
x_{7} = [XY^3,(1,2)],\\
x_{8} &= [X^2Y^2,(0,1)], \qquad 
\!\!x_{9} = [X^2Y^2,(1,0)],\\
x_{10} &= [X^2Y^2,(1,1)], \qquad 
\!\!\!\!x_{11} = [X^2Y^2,(1,2)],\\
x_{12} &= [X^3Y,(0,1)], \qquad 
x_{13} = [X^3Y,(1,0)],\\
x_{14} &= [X^3Y,(1,1)], \qquad 
x_{15} = [X^3Y,(1,2)],\\
x_{16} &= [X^4Y,(0,1)], \qquad 
x_{17} = [X^4Y,(1,0)],\\
x_{18} &= [X^4Y,(1,1)], \qquad 
x_{19} = [X^4Y,(1,2)].\\
\end{align*}
The relation matrix is already very large
for $\M_6(\Gamma_0(3))$.  It is as follows, where the $S$-relations
are before the line and the $T$-relations after it:
{\mbox{}\tiny
$$
\left(\begin{array}{cccccccccccccccccccc}
1&0&0&0&0&0&0&0&0&0&0&0&0&0&0&0&0&1&0&0\\
0&1&0&0&0&0&0&0&0&0&0&0&0&0&0&0&1&0&0&0\\
0&0&1&0&0&0&0&0&0&0&0&0&0&0&0&0&0&0&0&1\\
0&0&0&1&0&0&0&0&0&0&0&0&0&0&0&0&0&0&1&0\\
0&0&0&0&1&0&0&0&0&0&0&0&0&-1&0&0&0&0&0&0\\
0&0&0&0&0&1&0&0&0&0&0&0&-1&0&0&0&0&0&0&0\\
0&0&0&0&0&0&1&0&0&0&0&0&0&0&0&-1&0&0&0&0\\
0&0&0&0&0&0&0&1&0&0&0&0&0&0&-1&0&0&0&0&0\\
0&0&0&0&0&0&0&0&1&1&0&0&0&0&0&0&0&0&0&0\\
0&0&0&0&0&0&0&0&0&0&1&1&0&0&0&0&0&0&0&0\\\hline
1&0&0&1&0&0&0&-4&0&0&0&6&0&0&0&-4&0&1&0&1\\
1&1&0&0&-4&0&0&0&6&0&0&0&-4&0&0&0&1&0&0&1\\
0&0&2&0&0&0&-4&0&0&0&6&0&0&0&-4&0&0&0&2&0\\
0&1&0&1&0&-4&0&0&0&6&0&0&0&-4&0&0&1&1&0&0\\
0&0&0&1&1&0&0&-3&0&0&0&3&0&-1&0&-1&0&1&0&0\\
1&0&0&0&-3&1&0&0&3&0&0&0&-1&0&0&-1&0&0&0&1\\
0&0&1&0&0&0&-2&0&0&0&3&0&0&0&-2&0&0&0&1&0\\
0&1&0&0&0&-3&0&1&0&3&0&0&-1&-1&0&0&1&0&0&0\\
0&0&0&1&0&0&0&-2&1&1&0&1&0&-2&0&0&0&1&0&0\\
1&0&0&0&-2&0&0&0&1&1&0&1&0&0&0&-2&0&0&0&1\\
0&0&1&0&0&0&-2&0&0&0&3&0&0&0&-2&0&0&0&1&0\\
0&1&0&0&0&-2&0&0&1&1&0&1&-2&0&0&0&1&0&0&0\\
0&0&0&1&0&-1&0&-1&0&3&0&0&1&-3&0&0&0&1&0&0\\
1&0&0&0&-1&0&0&-1&0&0&0&3&0&1&0&-3&0&0&0&1\\
0&0&1&0&0&0&-2&0&0&0&3&0&0&0&-2&0&0&0&1&0\\
0&1&0&0&-1&-1&0&0&3&0&0&0&-3&0&0&1&1&0&0&0\\
0&1&0&1&0&-4&0&0&0&6&0&0&0&-4&0&0&1&1&0&0\\
1&0&0&1&0&0&0&-4&0&0&0&6&0&0&0&-4&0&1&0&1\\
0&0&2&0&0&0&-4&0&0&0&6&0&0&0&-4&0&0&0&2&0\\
1&1&0&0&-4&0&0&0&6&0&0&0&-4&0&0&0&1&0&0&1
\end{array}\right).
$$}The reduced row echelon form of the relations matrix,
with zero rows removed is
{\tiny
$$
\hspace{-1em}\left(\begin{array}{cccccccccccccccccccc}
1&0&0&0&0&0&0&0&0&0&0&0&0&0&0&0&0&1&0&0\\
0&1&0&0&0&0&0&0&0&0&0&0&0&0&0&0&1&0&0&0\\
0&0&1&0&0&0&0&0&0&0&0&0&0&0&0&0&0&0&0&1\\
0&0&0&1&0&0&0&0&0&0&0&0&0&0&0&0&0&0&1&0\\
0&0&0&0&1&0&0&0&0&0&0&0&0&0&0&0&0&0&3/16&-3/16\\
0&0&0&0&0&1&0&0&0&0&0&0&0&0&0&0&0&0&-1/16&1/16\\
0&0&0&0&0&0&1&0&0&0&0&0&0&0&0&0&0&1/2&-5/16&-3/16\\
0&0&0&0&0&0&0&1&0&0&0&0&0&0&0&0&0&-1/2&3/16&5/16\\
0&0&0&0&0&0&0&0&1&0&0&0&0&0&0&0&0&-1/6&1/12&1/12\\
0&0&0&0&0&0&0&0&0&1&0&0&0&0&0&0&0&1/6&-1/12&-1/12\\
0&0&0&0&0&0&0&0&0&0&1&0&0&0&0&0&0&0&1/4&-1/4\\
0&0&0&0&0&0&0&0&0&0&0&1&0&0&0&0&0&0&-1/4&1/4\\
0&0&0&0&0&0&0&0&0&0&0&0&1&0&0&0&0&0&-1/16&1/16\\
0&0&0&0&0&0&0&0&0&0&0&0&0&1&0&0&0&0&3/16&-3/16\\
0&0&0&0&0&0&0&0&0&0&0&0&0&0&1&0&0&-1/2&3/16&5/16\\
0&0&0&0&0&0&0&0&0&0&0&0&0&0&0&1&0&1/2&-5/16&-3/16\\
\end{array}\right).
$$}Since these relations are equivalent to the original
relations, we see how
$x_0, \ldots, x_{15}$ can be expressed in 
terms of $x_{16}$, $x_{17}$, $x_{18}$, and $x_{19}$.
Thus $\sM_6(\Gamma_0(3))$ has dimension $4$.
For example, 
$$x_{15} \sim \frac{1}{2} x_{17} -\frac{5}{16}x_{18}
        - \frac{3}{16} x_{19}.
$$


Notice that the number of relations is already quite large. It is
perhaps surprising how complicated the presentation is already for
$\sM_6(\Gamma_0(3))$.  Because there are denominators in the
relations, the above calculation is only a computation of
$\sM_6(\Gamma_0(3);\Q)$.  Computing $\sM_6(\Gamma_0(3);\Z)$ involves
finding a $\Z$-basis for the kernel of the relation matrix
(see \exref{ch:linalg}{ex:hnf}).


As before, we find that with respect to the basis 
$x_{16}$, $x_{17}$, $x_{18}$, and $x_{19}$
$$
T_2 = 
\left(\begin{array}{cccc}
33&0&0&0\\
3&6&12&12\\
-3/2&27/2&15/2&27/2\\
-3/2&27/2&27/2&15/2
\end{array}\right).
$$
Notice that there are denominators in the matrix for
$T_2$ with respect to this basis.  It is clear from
the definition of $T_2$ acting on Manin symbols that
$T_2$ preserves the $\Z$-module $\sM_6(\Gamma_0(3))$,
so there is some basis for $\sM_6(\Gamma_0(3))$
such that $T_2$ is given by an integer matrix.  Thus
the characteristic polynomial $f_2$ of $T_2$ will have integer
coefficients; indeed, 
$$
  f_2 = (x-33)^2 \cdot (x+6)^2.
$$
Note the factor $(x-33)^2$, which comes from the two images
of the Eisenstein series $E_4$ of level~$1$.  The
factor $x+6$ comes from the cusp form 
$$
  g = q - 6q^2 + \cdots \in S_6(\Gamma_0(3)).
$$
By computing more Hecke operators $T_n$, we
can find more coefficients of~$g$.
For example, the charpoly of $T_3$ is 
$(x-1)(x-243)(x-9)^2$, and the matrix of $T_5$
is 
$$
T_5 = \left(\begin{array}{cccc}
3126&0&0&0\\
240&966&960&960\\
-120&1080&1086&1080\\
-120&1080&1080&1086
\end{array}\right),
$$
which has characteristic polynomial
$$
  f_5 = (x-3126)^2 (x-6)^2.
$$
The matrix of $T_7$ is 
$$T_7 = \left(\begin{array}{cccc}
16808&0&0&0\\
1296&5144&5184&5184\\
-648&5832&5792&5832\\
-648&5832&5832&5792
\end{array}\right),
$$
with characteristic polynomial
$$
 f_7 = (x -16808)^2 (x + 40)^2.
$$
One can put this information together to deduce that
$$
  g = q - 6q^2 + 9q^3 + 4q^4 + 6q^5 - 54q^6 - 40q^7 +\cdots.
$$
\end{example}

\begin{example}
Consider $\M_2(\Gamma_0(43))$, which has dimension $7$.  
With respect to the symbols 
\begin{align*}
x_1&=(1,0),\quad x_{32}=(1,31),\quad
 x_{33} = (1,32),\\ 
x_{39} &= (1,38),\quad
 x_{40} = (1,39), \quad
x_{41} = (1,40),\quad
  x_{42} = (1,41),
\end{align*}
the matrix of $T_2$ is
$$
T_2=\left(\begin{array}{ccccccc}
3&0&0&0&0&0&0\\
0&-2&-1&-1&-1&0&0\\
0&1&1&0&0&-2&-1\\
0&0&1&-1&1&0&0\\
0&0&0&2&1&2&1\\
0&0&-1&-1&-1&-2&0\\
-1&0&0&1&1&1&-1
\end{array}\right),
$$
which has characteristic polynomial
$$
  (x-3)(x+2)^2 (x^2-2)^2.
$$
There is one Eisenstein series and there are three cusp forms
with $a_2 = -2$ and $a_2 = \pm \sqrt{2}$.
\end{example}

\begin{example}
To compute
$\sM_2(\Gamma_0(2004);\Q)$, we first make
a list of the 
$$
4032 = (2^2 + 2) \cdot (3+1) \cdot (167 + 1)
$$
elements $(a,b)\in\P^1(\Z/2004\Z)$ using Algorithm~\ref{alg:p1rep}.
The list looks like this:
$$
  (0, 1), (1, 0), (1, 1), (1, 2),
  \ldots,   (668, 1), (668, 3), (668, 5), (1002, 1).
$$
For each of the symbols $x_i$, we consider the $S$-relations
and $T$-relations.
Ignoring the redundant relations, we find $2016$ $S$-relations and
$1344$ $T$-relations.  It is simple to quotient out by the
$S$-relations, e.g., by identifying $x_i$ with $-x_i S = -x_j$ for
some $j$ (or setting $x_i=0$ if $x_i S = x_i$).  Once we have taken
the quotient by the $S$-relations, we take the {\em image} of all
$1344$ of the $T$-relations modulo the $S$-relations and quotient out by
those relations.  Because $S$ and $T$ do not commutate, we cannot
only quotient out by $T$-relations $x_i + x_i T + x_i T^2 = 0$ where
the $x_i$ are the basis after quotienting out by the $S$-relations.
The relation matrix has rank $3359$, so $\sM_2(\Gamma_0(2004);\Q)$ has
dimension $673$.

If we instead compute the quotient $\sM_2(\Gamma_0(2004);\Q)^+$
of $\sM_2(\Gamma_0(2004);\Q)$ by the subspace of elements $x-\eta^*(x)$,
we include relations $x_i+x_iI=0$, where $I=\abcd{-1}{0}{0}{1}$.
There are now 2016 $S$-relations, 2024 $I$-relations,
and 1344 $T$-relations.  Again, it is relatively easy to quotient
out by the $S$-relations by identifying $x_i$ and $-x_i S$.
We then take the image of all $2024$ $I$-relations modulo the
$S$-relations, and again directly quotient out by the $I$-relations
by identifying $[x_i]$ with $-[x_iI]=-[x_j]$ for some $j$, where
by $[x_i]$ we mean the class of $x_i$ modulo the $S$-relations.
Finally, we quotient out by the $1344$ $T$-relations, which involves
sparse Gauss elimination on a matrix with $1344$ rows
and at most three nonzero entries per row.    
The dimension of $\sM_2(\Gamma_0(2004);\Q)^+$ is $331$.
\end{example}

\section{Refined Algorithm for the Presentation}
\begin{algorithm}{Modular Symbols Presentation}\label{alg:modsympresent}
This is an algorithm to compute $\sM_k(\Gamma_0(N);\Q)$ or 
$\sM_k(\Gamma_0(N);\Q)^{\pm}$, which
only requires doing generic sparse linear algebra to deal
with the three term $T$-relations. 
\begin{steps}
\item Let $x_0,\ldots, x_n$ by a list of all Manin symbols.
\item Quotient out the two-term $S$-relations and if the $\pm $
quotient is desired, by the two-term $\eta$-relations.
(Note that this is more subtle than 
just ``identifying symbols in pairs'',
since complicated relations can cause generators to surprisingly
equal $0$.)
       Let $[x_i]$ denote the
          class of $x_i$ after this quotienting process.
\item Create a sparse matrix $A$ with $m$ columns, whose rows 
      encode the relations
$$
                 [x_i] + [x_iT] + [x_iT^2] = 0.
$$
          For example, there are about $n/3$ such rows when 
          $k=2$.
          The number of nonzero entries per row is at most $3(k-1)$.
          Note that we must include rows for {\em all} $i$, since
          even if $[x_i] = [x_j]$, it need not be the case that
          $[x_i T] = [x_jT]$, since the matrices $S$ and $T$ do not commute.
          However, we have an a priori formula for the dimension
          of the quotient by all these relations, so we could omit
          many relations and just check that there are enough at the
          end---if there are not, we add in more.

        \item Compute the reduced row echelon form of $A$ using
          Algorithm~\ref{alg:modech}.  For $k=2$, this is the echelon
          form of a matrix with size about $n/3 \times n/4$.

 \item Use what we have done above to read off a sparse matrix $R$
          that expresses each of the $n$ Manin symbols in terms of
          a basis of Manin symbols, modulo the relations. 
        \end{steps}
      \end{algorithm}












    



\section{Applications}
\subsection{Later in This Book}
We sketch some of the ways in which we will apply the
modular symbols algorithms of this chapter later in this book.

Cuspidal modular symbols are in Hecke-equivariant duality with
cuspidal modular forms, and as such we can compute modular forms by
computing systems of eigenvalues for the Hecke operators acting on
modular symbols.  By the Atkin-Lehner-Li theory of newforms\index{newform} (see,
e.g., Theorem~\ref{thm:ALL}), we can construct $S_k(N,\eps)$ for
any~$N$, any $\eps$, and $k\geq 2$ using this method.  See
Chapter~\ref{ch:modform} for more details.

Once we can compute spaces of modular symbols, we move to computing
the corresponding modular forms.  We
define inclusion and trace maps from modular symbols of one level $N$
to modular symbols of level a multiple or divisor of $N$.  Using these,
we compute the quotient $V$ of the new subspace of cuspidal modular
symbols on which a ``star involution'' acts as $+1$.  The Hecke
operators act by diagonalizable commuting matrices on this space, and
computing the systems of Hecke eigenvalues is equivalent
to computing\index{newform} newforms $\sum a_n q^n$.  In this way, we
obtain a list of {\em all} newforms (normalized eigenforms) in
$S_k(N,\eps)$ for any $N$, $\eps$, and $k\geq 2$.

In Chapter~\ref{ch:periods}, we compute with the period mapping from
modular symbols to $\C$ attached to a newform $f \in S_k(N,\eps)$.
When $k=2, \eps=1$ and~$f$ has rational Fourier coefficients, this
gives a method to compute the period lattice associated to a modular
elliptic curve attached to a newform (see
Section~\ref{sec:findall}).  In general, computation of this map
is important when finding equations for modular $\Q$-curves, CM
curves, and curves with a given modular Jacobian.  It is also
important for computing special values of the $L$-function $L(f,s)$ at
integer points in the critical strip.

\subsection{Discussion of the Literature and Research}
Modular symbols were introduced by Birch \cite{birch:bsd} for
computations in support of the Birch and Swinnerton-Dyer conjecture.
Manin \cite{manin:parabolic} used modular symbols to
prove rationality results about special values of $L$-functions.

Merel's paper \cite{merel:1585} builds on work of Shokurov (mainly
\cite{sokurov:shimura}), which develops a higher-weight
generalization of Manin's work partly to understand rationality
properties of special values of $L$-functions. Cremona's book
\cite{cremona:algs} discusses how to compute the space of
weight~$2$ modular symbols for $\Gamma_0(N)$, in connection with the
problem of enumerating all elliptic curves of given conductor, and his
article \cite{cremona:gammaone} discusses the $\Gamma_1(N)$ case and
computation of modular symbols with character.

There have been several Ph.D. theses about modular symbols.
Basmaji's thesis \cite{basmaji:thesis} contains tricks to efficiently
compute Hecke operators $T_p$, with $p$ very large (see
Section~\ref{sec:basmaji}), and also discusses how to compute spaces
of half integral weight modular forms building on what one can get
from modular symbols of integral weight. The author's Ph.D. thesis
\cite{stein:phd} discusses higher-weight modular
symbols and applies modular symbols to study Shafarevich-Tate groups
(see also \cite{agashe:phd}).  Martin's thesis \cite{martin:thesis} is
about an attempt to study an analogue of analytic modular symbols for
weight~$1$.  Gabor Wiese's thesis \cite{gabor:phd} uses modular
symbols methods to study weight 1 modular forms modulo~$p$.  Lemelin's
thesis \cite{lemelin:dominic} discusses modular symbols for quadratic
imaginary fields in the context of $p$-adic analogues of the Birch and
Swinnerton-Dyer conjecture.  See also the survey paper
\cite{frey-muller}, which discusses computation with weight~$2$
modular symbols in the context of modular abelian
varieties.

The appendix of this book is about
analogues of modular symbols for groups besides
finite index subgroups of $\SL_2(\Z)$, e.g., for subgroup
of higher rank groups such as $\SL_3(\Z)$.  
There has also been work on computing Hilbert
modular forms, e.g., by Lassina Dembel\'e \cite{dembele:hilbert5} Hilbert modular forms are functions on a product of copies of~$\h$, and
$\SL_2(\Z)$ is replaced by a group of matrices with entries in a
totally real field.  

Glenn Stevens, Robert Pollack and Henri Darmon (see
\cite{darmon-pollack}) have worked for many years to develop an
analogue of modular symbols in a rigid analytic context, which is 
helpful for questions about computing with overconvergent
$p$-adic modular forms or proving results about $p$-adic
$L$-functions.

Finally we mention that Barry Mazur and some other authors 
use the term ``modular symbol''
in a different way than we do.
They use the term in a way that is dual
to our usage; for example, they attach a ``modular symbol'' to
a modular form or elliptic curve.   
See \cite{mtt} for an extensive discussion
of modular symbols from this point of view, where they are used to
construct $p$-adic $L$-functions.


\begin{exercises}
\item \label{ex:surjredmod}
Suppose $M$ is an integer multiple of $N$.  Prove
that the natural map $(\Z/M\Z)^* \to (\Z/N\Z)^*$ is surjective.

\item \label{ex:surjred}
Prove that $\SL_2(\Z)\to \SL_2(\Z/N\Z)$ is surjective
(see Lemma~\ref{lem:canlift}).

\item \label{ex:m3g13}
Compute $\sM_{3}(\Gamma_1(3))$.  List
each Manin symbol the relations they satisfy, compute
the quotient, etc.  Find the matrix of $T_2$.
(Check: The dimension of $\sM_{3}(\Gamma_1(3))$ is $2$, and
the characteristic polynomial of $T_2$ is $(x-3)(x+3)$.)

\item \label{ex:pairingwd} 
Finish the proof of Proposition~\ref{prop:pairwd}.

\item \label{ex:etaconj}
\begin{enumerate}
\item Show that if $\eta = \abcd{-1}{0}{0}{1}$, then
$\eta \Gamma \eta = \Gamma$ for $\Gamma=\Gamma_0(N)$
and $\Gamma=\Gamma_1(N)$.  
\item (*) Give an example of a finite index
subgroup $\Gamma$ such that $\eta \Gamma \eta \neq \Gamma$.
\end{enumerate}


\end{exercises}




















\chapter{Computing with Newforms}\label{ch:modformcomp}\index{newform!computing}

In this chapter we pull together results and algorithms from
Chapter~\ref{ch:modform2}, \ref{ch:dirichlet}, \ref{ch:linalg}, and
\ref{ch:modsym} and explain how to use linear algebra techniques to
compute cusp forms and eigenforms using modular symbols.

We first discuss in Section~\ref{sec:decompdirchar} how to decompose
$M_k(\Gamma_1(N))$ as a direct sum of subspaces corresponding to
Dirichlet characters.  Next in Section~\ref{sec:ALL} we state the main
theorems of Atkin-Lehner-Li theory, which 
decomposes $S_k(\Gamma_1(N))$ into subspaces on which the
Hecke operators act diagonalizably with ``multiplicity one''.
In Section~\ref{sec:compcuspforms} we describe two
algorithms for computing modular forms. One algorithm finds
a basis of $q$-expansions, and the other computes eigenvalues
of newforms. 

\section{Dirichlet Character Decomposition}
\label{sec:decompdirchar}
The group $(\Z/N\Z)^*$ acts on $M_k(\Gamma_1(N))$ through 
\defn{diamond-bracket operators} $\langle d \rangle$, as follows.
For $d\in(\Z/N\Z)^*$, define
$$
   f|\dbd{d} = f^{\left[\abcd{a}{b}{c}{d'}\right]_k},
$$
where $\abcd{a}{b}{c}{d'} \in \SL_2(\Z)$ is congruent to
$\abcd{d^{-1}}{0}{0\hfill }{d}\pmod{N}$.  Note that the map
$\SL_2(\Z)\to \SL_2(\Z/N\Z)$ is surjective (see \exref{ch:modsym}{ex:surjred}),
so the matrix $\abcd{a}{b}{c}{d'}$ exists.
To prove that $\dbd{d}$ preserves $M_k(\Gamma_1(N))$,
we prove the more general fact that $\Gamma_1(N)$ is a 
normal subgroup of
$$  \Gamma_0(N) = \left\{ \mtwo{a}{b}{c}{d} \in \SL_2(\Z) : 
                    \mtwo{a}{b}{c}{d}\con \mtwo{*}{*}{0}{*}\pmod{N}
                \right\}.
$$
This will imply that $\dbd{d}$ preserves $M_k(\Gamma_1(N))$
since $\abcd{a}{b}{c}{d'} \in \Gamma_0(N)$.

\begin{lemma}
The group $\Gamma_1(N)$ is a normal subgroup of $\Gamma_0(N)$,
and the quotient $\Gamma_0(N)/\Gamma_1(N)$ is
isomorphic to $(\Z/N\Z)^*$.
\end{lemma}
\begin{proof}
See \exref{ch:modformcomp}{ex:g0g1quo}.
\end{proof}

Alternatively, one can prove that $\dbd{d}$ preserves
$M_k(\Gamma_1(N))$ by showing that $\dbd{d}\in \T$ and noting that
$M_k(\Gamma_1(N))$ is preserved by~$\T$ (see
Remark~\ref{rem:dbdinhecke}).

The \defn{diamond-bracket action} is the action of
$\Gamma_0(N)/\Gamma_1(N)\isom(\Z/N\Z)^*$ on $M_k(\Gamma_1(N))$.  Since
$M_k(\Gamma_1(N))$ is a finite-dimensional vector space over $\C$, the
$\dbd{d}$ action breaks $M_k(\Gamma_1(N))$ up as a direct sum of
factors corresponding to the Dirichlet characters $D(N,\C)$ of
modulus~$N$.
\begin{proposition}\label{prop:dcdecomp}
We have
$$
  M_k(\Gamma_1(N)) = \bigoplus_{\eps \in D(N,\C)} M_k(N,\eps),
$$
where
$$
  M_k(N,\eps) = \bigl\{ f \in \M_k(\Gamma_1(N)) : f|\dbd{d} = \eps(d) f,
                    \text{ all }d \in (\Z/N\Z)^*\bigr\}.
$$
\end{proposition}
\begin{proof}
  The linear transformations $\dbd{d}$, for the $d\in(\Z/N\Z)^*$, all
  commute, since $\dbd{d}$ acts through the abelian group
  $\Gamma_0(N)/\Gamma_1(N)$.  Also, if $e$ is the exponent of
  $(\Z/N\Z)^*$, then $\dbd{d}^e = \dbd{d^e} = \dbd{1} = 1$, so the
  matrix of $\dbd{d}$ is diagonalizable.  It is a standard fact from
  linear algebra that any commuting family of diagonalizable linear
  transformations is simultaneously diagonalizable (see
  \exref{ch:eisen}{ex:diag}), so there is a basis $f_1,\ldots, f_n$
  for $M_k(\Gamma_1(N))$ such that all $\dbd{d}$ act by diagonal
  matrices.  The system of eigenvalues of the action of $(\Z/N\Z)^*$ on a fixed
  $f_i$ defines a Dirichlet character, i.e., each $f_i$ has
  the property that $f_i |\dbd{d} = \eps_i(d)$, for all
  $d\in(\Z/N\Z)^*$ and some Dirichlet character $\eps_i$.  The $f_i$
  for a given $\eps$ then span $M_k(N,\eps)$, and taken together the
  $M_k(N,\eps)$ must span $M_k(\Gamma_1(N))$.
\end{proof}

\begin{definition}[Character of Modular Form]
If $f\in M_k(N,\eps)$, we say that~$\eps$ is 
the \defn{character of the modular form}~$f$.
\end{definition}

The spaces $M_k(N,\eps)$ are a direct sum of subspaces
$S_k(N,\eps)$ and $E_k(N,\eps)$, where $S_k(N,\eps)$ is the
subspace of cusp forms, i.e., forms that vanish at {\em all}
cusps (elements of $\Q \union \{\infty\}$), and $E_k(N,\eps)$
is the subspace of Eisenstein series, which is the unique
subspace of $M_k(N,\eps)$ that is invariant under all Hecke
operators and is such that 
$M_k(N,\eps) = S_k(N,\eps) \oplus E_k(N,\eps)$.
The space $E_k(N,\eps)$ can also be defined as the space
spanned by all Eisenstein series of weight $k$ and level~$N$,
as defined in Chapter~\ref{ch:eisen}. The space 
$E_k(N,\eps)$ can be defined in a third way using the Petersson
inner product (see \cite[\S{}VII.5]{lang:modular}).\index{Petersson
inner product}

%begin_sage
The diamond-bracket operators preserve cusp forms, so
the isomorphism of Proposition~\ref{prop:dcdecomp} restricts to an
isomorphism of the corresponding cuspidal subspaces.  
We illustrate how to use \sage{}
to make a table of dimension of $M_k(\Gamma_1(N))$ 
and $M_k(N,\eps)$ for $N=13$.
\sageindex{dimension with character}
\begin{verbatim}
sage: G = DirichletGroup(13)
sage: G
Group of Dirichlet characters of modulus 13 over 
Cyclotomic Field of order 12 and degree 4
sage: dimension_modular_forms(Gamma1(13),2)
13
sage: [dimension_modular_forms(e,2) for e in G]
[1, 0, 3, 0, 2, 0, 2, 0, 2, 0, 3, 0]
\end{verbatim}
Next we do the same for $N=100$.
\begin{verbatim}
sage: G = DirichletGroup(100)
sage: G
Group of Dirichlet characters of modulus 100 over 
Cyclotomic Field of order 20 and degree 8
sage: dimension_modular_forms(Gamma1(100),2)
370
sage: v = [dimension_modular_forms(e,2) for e in G]; v
[24, 0, 0, 17, 18, 0, 0, 17, 18, 0, 0, 21, 18, 0, 0, 17, 
  18, 0, 0, 17, 24, 0, 0, 17, 18, 0, 0, 17, 18, 0, 0, 21, 
  18, 0, 0, 17, 18, 0, 0, 17]
sage: sum(v)
370
\end{verbatim}
%end_sage


\section{Atkin-Lehner-Li Theory}\label{sec:ALL}
In Section~\ref{sec:degen} we defined maps between modular symbols
spaces of different level.  There are similar maps between spaces of
cusp forms.  Suppose $N$ and $M$ are positive integers with $M\mid N$
and that $t$ is a divisor of $N/M$.  Let
\begin{equation}\label{eqn:degenmf}
  \alpha_t : S_k(\Gamma_1(M)) \to S_k(\Gamma_1(N))
\end{equation}
be the \defn{degeneracy map}, which is given by $f(q) \mapsto f(q^t)$.  
There are also maps $\beta_t$ in the other direction; see
\cite[Ch.~VIII]{lang:modular}.

The \defn{old subspace} of $S_k(\Gamma_1(N))$, denoted $S_k(\Gamma_1(N))_{\old}$,
is the sum of the images of all maps $\alpha_t$ with~$M$ a proper divisor
of $N$ and $t$ any divisor of $N/M$ (note that $\alpha_t$ depends on 
$t$, $N$, and $M$, so there is a slight abuse of notation).
The \defn{new subspace} of $S_k(\Gamma_1(N))$, which we denote by
$S_k(\Gamma_1(N))_{\new}$, is the intersection
of the kernel of all maps $\beta_t$ with $M$ a proper divisor
of $N$.  One can use the Petersson inner product to show
that 
$$
  S_k(\Gamma_1(N)) = S_k(\Gamma_1(N))_{\new} \oplus S_k(\Gamma_1(N))_{\old}.
$$
Moreover, the new and old subspaces are preserved by all
Hecke operators. 

Let $\T = \Z[T_1, T_2, \ldots ]$ be the commutative polynomial ring
in infinitely many indeterminates $T_n$. This
ring acts (via $T_n$ acting as the $n$th Hecke operator)
on $S_k(\Gamma_1(N))$ for every integer $N$.  Let
$\T^{(N)}$ be the subring of $\T$ generated by the $T_n$ with
$\gcd(n,N)=1$.

\begin{theorem}[Atkin, Lehner, Li]\label{thm:ALL}
We have a decomposition
\begin{equation}\label{eqn:decomp}
  S_k(\Gamma_1(N)) = 
   \bigoplus_{M\mid N} \bigoplus_{d\mid N/M} \alpha_d(S_k(\Gamma_1(M))_{\new}).
\end{equation}
Each space $S_k(\Gamma_1(M))_{\new}$ is a direct
sum of distinct (nonisomorphic) simple $\T^{(N)}_{\C}$-modules.
\end{theorem}
\begin{proof}
The complete proof is in \cite{winnie:newforms}.
See also \cite[Ch.~5]{diamond-shurman} for a beautiful
modern treatment of this and related results.
\end{proof}
The analogue of Theorem~\ref{thm:ALL} with $\Gamma_1$ replaced by
$\Gamma_0$ is also true (this is what was proved in
\cite{atkin-lehner}).  The analogue for $S_k(N,\eps)$ is also valid,
as long as we omit the spaces $S_k(\Gamma_1(M),\eps)$ for which
$\cond(\eps) \nmid M$.

\begin{example}
If $N$ is prime and $k\leq 11$, then 
$S_k(\Gamma_1(N))_{\new} = S_k(\Gamma_1(N))$,
since $S_k(\Gamma_1(1))=0$.   
\end{example}

One can prove using the Petersson inner product that the 
operators $T_n$ on $S_k(\Gamma_1(N))$, with $\gcd(n,N)=1$, are
diagonalizable.  Another result of Atkin-Lehner-Li theory is that the
ring of endomorphisms of $S_k(\Gamma_1(N))_{\new}$ generated by all
Hecke operators equals the ring generated by the Hecke operators $T_n$
with $(n,N)=1$.  This statement need not be true if we do not
restrict to the new subspace, as the following example shows. 
\begin{example}\label{ex:onlyold}
We have
$$
  S_2(\Gamma_0(22)) = S_2(\Gamma_0(11)) \oplus \alpha_2(S_2(\Gamma_0(11))),
$$
where each of the spaces $S_2(\Gamma_0(11))$ has dimension $1$.  Thus
$S_2(\Gamma_0(22))_{\new} = 0$.  The Hecke operator $T_2$ on 
$S_2(\Gamma_0(22))$ has characteristic polynomial $x^2+2x+2$,
which is irreducible.  Since $\alpha_2$ commutes with all Hecke 
operators $T_n$, with $\gcd(n,2)=1$, the subring $\T^{(22)}$ of the Hecke
algebra generated by operators $T_n$ with~$n$ odd is isomorphic
to~$\Z$ (the $2\times 2$ scalar matrices).  Thus on the full space
$S_2(\Gamma_0(22))$, we do not have $\T^{(22)}=\T$.  However, on the new subspace 
we do have this equality, since the new subspace has dimension~$0$.
\end{example}

\begin{example}\label{ex:new_old}
The space $S_2(\Gamma_0(45))$ has dimension $3$ and basis
\begin{align*}
f_{0} &= q - q^{4} - q^{10} - 2q^{13} - q^{16} + 4q^{19} + \cdots, \\
f_{1} &= q^{2} - q^{5} - 3q^{8} + 4q^{11} - 2q^{17} + \cdots, \\
f_{2} &= q^{3} - q^{6} - q^{9} - q^{12} + q^{15} + q^{18} + \cdots. \\
\end{align*}
The new subspace $S_2(\Gamma_0(45))_{\new}$ is spanned by the single cusp form
$$
q + q^{2} - q^{4} - q^{5} - 3q^{8} - q^{10} + 4q^{11} - 2q^{13} + \cdots.
$$
We have $S_2(\Gamma_0(45/5)) = 0$ and $S_2(\Gamma_0(15))$ 
has dimension $1$ with basis
$$
q - q^{2} - q^{3} - q^{4} + q^{5} + q^{6} + 3q^{8} + q^{9} - q^{10} - 4q^{11} + q^{12} - 2q^{13} + \cdots.
$$

We use \sage{} to verify the above assertions.
\begin{verbatim}
sage: S = CuspForms(Gamma0(45), 2, prec=14); S
Cuspidal subspace of dimension 3 of Modular Forms space 
of dimension 10 for Congruence Subgroup Gamma0(45) of 
weight 2 over Rational Field
sage: S.basis()
[
q - q^4 - q^10 - 2*q^13 + O(q^14),
q^2 - q^5 - 3*q^8 + 4*q^11 + O(q^14),
q^3 - q^6 - q^9 - q^12 + O(q^14)
]
sage: S.new_subspace().basis()
(q - q^4 - q^10 - 2*q^13 + O(q^14),)
sage: CuspForms(Gamma0(9),2)
Cuspidal subspace of dimension 0 of Modular Forms space 
of dimension 3 for Congruence Subgroup Gamma0(9) of 
weight 2 over Rational Field
sage: CuspForms(Gamma0(15),2, prec=10).basis()
[
q - q^2 - q^3 - q^4 + q^5 + q^6 + 3*q^8 + q^9 + O(q^10)
]
\end{verbatim}
\end{example}

\begin{example}
This example is similar to Example~\ref{ex:onlyold}, except that there
are newforms.
We have
$$
  S_2(\Gamma_0(55)) = S_2(\Gamma_0(11)) \oplus \alpha_5(S_2(\Gamma_0(11)))
                     \oplus S_2(\Gamma_0(55))_{\new},
$$
where $S_2(\Gamma_0(11))$ has dimension $1$ and
$S_2(\Gamma_0(55))_{\new}$ has dimension~$3$.
The Hecke operator $T_5$ on 
$S_2(\Gamma_0(55))_{\new}$ acts  via the matrix
$$
\mthree
{-2}{\hfill 2}{\hfill -1}
{-1}{\hfill 1}{-1}
{\hfill 1}{-2}{\hfill 0}
$$
with respect to some basis.
This matrix has eigenvalues $1$ and $-1$.
Atkin-Lehner theory asserts that $T_5$ must be a
linear combination of $T_n$, with 
$\gcd(n,55)=1$.
Upon computing the matrix for $T_2$,
we find by simple linear algebra that $T_5 = 2T_2 - T_4$.
\end{example}




\begin{definition}[Newform]\label{def:newform}
A \defn{newform} is a $\T$-eigenform $f\in S_k(\Gamma_1(N))_{\new}$
that is normalized so that the coefficient of $q$ is $1$.
\end{definition}
We now motivate this definition by explaining why any $\T$-eigenform
can be normalized so that the coefficient of~$q$ is~$1$ and how such
an eigenform has the property that its Fourier
coefficients are exactly the Hecke eigenvalues.
\begin{proposition}\label{prop:eigncoeffs}
If $f=\sum_{n=0}^{\infty} a_n q^n \in M_k(N,\eps)$ 
is an eigenvector for all Hecke operators
$T_n$ normalized so that $a_1=1$, 
then $T_n(f) = a_n f$.
\end{proposition}
\begin{proof}
If $\eps=1$, then $f\in M_k(\Gamma_0(N))$ and this
is Lemma~\ref{lem:hecketn}.   However, we have not
yet considered Hecke operators on $q$-expansions for
more general spaces of modular forms. 

The Hecke operators $T_p$, for $p$ prime, act on $S_k(N,\eps)$
by 
$$
T_p\left(\sum_{n=0}^{\infty} a_n q^n\right) =
\sum_{n=0}^{\infty} \left( a_{np}q^n + \eps(p) p^{k-1}
  a_{n}q^{np}\right),
$$
and there is a similar formula for $T_m$ with $m$ composite.
If $f=\sum_{n=0}^{\infty} a_n q^n$ is an eigenform for all
$T_p$, with eigenvalues $\lambda_p$, then by the above formula 
\begin{equation}\label{eqn:eigenexp}
\lambda_p f = \lambda_p a_1 q + \lambda_p a_2 q^2 + \cdots = T_p(f) = a_{p} q +
\text{higher terms}.  
\end{equation}
Equating coefficients of $q$, we see that if
$a_1=0$, then $a_p=0$ for all~$p$; hence
$a_n=0$ for all~$n$, because of the multiplicativity of Fourier
coefficients and the recurrence $$a_{p^r} = a_{p^{r-1}} a_p - \eps(p)
p^{k-1} a_{p^{r-2}}.$$ 
This would mean that $f=0$, a contradiction.  Thus $a_1\neq 0$,
and it makes sense to normalize $f$ so that $a_1=1$.
With this normalization, \eqref{eqn:eigenexp} implies that $\lambda_p = a_p$,
as desired.  
\end{proof}

\begin{remark}\label{rem:dbdinhecke}
  The Hecke algebra $\T_\Q$ on $M_k(\Gamma_1(N))$ contains the
  operators $\dbd{d}$, since they satisfy the relation $T_{p^2} =
  T_p^2 - \dbd{p}p^{k-1}$.  Thus any $\T$-eigenform in
  $M_k(\Gamma_1(N))$ lies in a subspace $M_k(N,\eps)$ for some
  Dirichlet character $\eps$. Also, one can even prove that
  $\dbd{d} \in \Z[\ldots,T_n,\ldots]$ (see
  \exref{ch:modformcomp}{ex:dbdinhecke}).
\end{remark}

\section{Computing Cusp Forms}\label{sec:compcuspforms}
Let $\sS_k(N,\eps;\C)$ be the space of cuspidal modular symbols as in
Chapter~\ref{ch:modsym}.  Let $\iota^*$ be the map of 
\eqref{eqn:iota}, and let $\sS_k(N,\eps;\C)^+$
 be the \defn{plus one quotient} of cuspidal modular symbols,
i.e., the quotient of $\sS_k(N,\eps;\C)$ by the image of $\iota^*-1$. 
It follows from
Theorem~\ref{thm:modsymformpairing} and compatibility of the
degeneracy maps (for modular symbols they are defined
in Section~\ref{sec:degen}) that the $\T$-modules $S_k(N,\eps)_{\new}$
and $\sS_k(N,\eps;\C)^+_{\new}$ are dual as $\T$-modules.
Thus finding the systems of $\T$-eigenvalues on cusp forms is the same
as finding the systems of $\T$-eigenvalues on cuspidal modular
symbols.

Our strategy to compute $S_k(N,\eps)$ is to first compute spaces
$S_k(N,\eps)_{\new}$ using the Atkin-Lehner-Li decomposition
\eqref{eqn:decomp}.  To compute $S_k(N,\eps)_{\new}$ to a given
precision, we compute the systems of eigenvalues of the Hecke
operators $T_p$ on the space $V=\sS_k(N,\eps;\C)^+_{\new}$, which we
will define below.  Using Proposition~\ref{prop:eigncoeffs}, we then
recover a basis of $q$-expansions for newforms.  Note that we only
need to compute Hecke eigenvalues $T_p$, for $p$ prime, not the $T_n$
for $n$ composite, since the $a_n$ can be quickly recovered in terms
of the $a_p$ using multiplicativity and the recurrence.

For some problems, e.g., construction of models for modular curves,
having a basis of $q$-expansions is enough.  For many other problems,
e.g., enumeration of modular abelian varieties, one is really
interested in the newforms.  We next discuss algorithms aimed at each
of these problems.

\subsection{A Basis of $q$-Expansions}
The following algorithm generalizes Algorithm~\ref{alg:s2basis}.
It computes $S_k(N,\eps)$ without finding any eigenspaces.
\begin{algorithm}{Merel's Algorithm for Computing a Basis}\label{alg:merelqexp}
Given integers $m$, $N$ and $k$ and a Dirichlet character~$\eps$ with
modulus~$N$, this algorithm computes a basis of $q$-expansions
for $S_k(N,\eps)$ to precision $O(q^{m+1})$.
\begin{steps}
\item{}[Compute Modular Symbols] 
Use  Algorithm~\ref{alg:modsympresent} to
compute $$V = \sS_k(N,\eps)^+\tensor \Q(\eps),$$
viewed as a $K=\Q(\eps)$ vector space,
with an action of the $T_n$.

\item{}[Basis for Linear Dual]
Write down a basis for $V^* = \Hom(V,\Q(\eps))$.  E.g., if we identify
$V$ with $K^n$ viewed as column vectors, then $V^*$ is the space
of row vectors of length $n$, and the pairing is the row $\times$ column product.

\item{}[Find Generator]
Find $x\in V$ such that
$\T{} x = V$ by choosing random~$x$ until we find
one that generates.  The set of~$x$ that fail
to generate lie in a union of a finite number of
proper subspaces. 
\item{}[Compute Basis]  
The set of power series
$$
  f_i = \sum_{n=1}^{m} \psi_i(T_n(x)) q^n + O(q^{m+1})
$$
forms a basis for $S_k(N,\eps)$ to precision $m$.
\end{steps}
\end{algorithm}

In practice Algorithm~\ref{alg:merelqexp} seems slower than the
eigenspace algorithm that we will describe in the rest of this
chapter.  The theoretical complexity of Algorithm~\ref{alg:merelqexp}
{\em may} be better, because it is not necessary to factor any
polynomials. Polynomial factorization is difficult from the worst-case
complexity point of view, though it is usually fast in practice.  The
eigenvalue algorithm only requires computing a few images $T_p(x)$ for
$p$ prime and $x$ a Manin symbol on which $T_p$ can easily be
computed.  The Merel algorithm involves computing $T_n(x)$ for
all~$n$ and for a fairly easy~$x$, which is potentially more work.

\begin{remark}
By ``easy $x$'', I mean that computing $T_n(x)$ is easier on~$x$
than on a completely random element of $\sS_k(N,\eps)^+$,
e.g.,~$x$ could be a Manin symbol.
\end{remark}

\subsection{Newforms: Systems of Eigenvalues}\label{sec:syseigen}\index{newform!system of eigenvalues}
In this section we describe an algorithm for computing
the system of Hecke eigenvalues associated to a simple
subspace of a space of modular symbols.   This algorithm
is better than  doing linear algebra directly
over the number field generated by the eigenvalues.  It
only involves linear algebra over the base field and also
yields a compact representation for the answer, which
is better than writing the eigenvalues in terms of
a power basis for a number field.
In order to use this algorithm, it is necessary to
decompose the space of cuspidal modular symbols as a
direct sum of simples, e.g., using Algorithm~\ref{alg:decomp}.



Fix~$N$ and a Dirichlet character~$\eps$ of modulus~$N$,
and let 
$$
  V=\sM_k(N,\eps)^+
$$
be the $+1$ quotient of
modular symbols (see equation \eqref{eqn:iota}).

\begin{algorithm}{System of Eigenvalues}\label{alg:eigsys}%
Given a $\T$-simple subspace $W\subset V$ of modular symbols,
this algorithm outputs maps~$\psi$ and~$e$, where 
$
   \psi : \T_K \to W
$ 
is a $K$-linear map  and $e:W \isom L$ is an isomorphism of $W$ with a
number field~$L$, such that $a_n = e (\psi(T_n))$ is
the eigenvalue of the $n$th Hecke operator acting on a fixed
$\T$-eigenvector in $W\tensor\Qbar$.  (Thus 
$f = \sum_{n=1}^{\infty} e(\psi(T_n)) q^n$ is a newform.)
\begin{steps}

\item{}[Compute Projection] 
Let $\vphi:V \to W'$ be any surjective linear map
such that $\ker(\vphi)$ equals the kernel
of the $\T$-invariant projection onto $W$.
For example, compute~$\vphi$ by finding a simple submodule
of $V^*=\Hom(V,K)$ that is isomorphic to~$W$, e.g., by applying
Algorithm~\ref{alg:decomp} to $V^*$ with $T$ replaced by
the transpose of $T$.   

\item{}[Choose $v$]\label{step:eig:v} Choose a nonzero element
  $v\in{}V$ such that $\pi(v)\neq 0$ and computation of $T_n(v)$ is
  ``easy'', e.g., choose~$v$ to be a Manin symbol.

\item{}[Map from Hecke Ring] Let $\psi$ be the map $\T \to W'$, given 
by $\psi(t) = \pi(t v)$.  Note that computation of $\psi$
is relatively easy, because $v$ was chosen so that
$tv$ is relatively easy to compute.  In particular, 
if $t=T_p$, we do not need to compute the full matrix
of $T_p$ on~$V$; instead we just compute $T_p(v)$.


\item{}[Find Generator]\label{alg:ap:gen} 
Find a random $T \in \T$ such that the iterates 
$$\psi(T^0), \quad \psi(T), \quad \psi(T^2), \quad \ldots, \quad \psi(T^{d-1})$$
are a basis for~$W'$, where $W$ has dimension~$d$. 

\item{}[Characteristic Polynomial] Compute the characteristic polynomial 
$f$ of $T|_W$, and let $L=K[x]/(f)$.
Because
  of how we chose $T$ in step~(\ref{alg:ap:gen}), the minimal
  and characteristic polynomials of $T|_W$ are equal, and
  both are irreducible, so~$L$ is an extension of~$K$ 
  of degree~$d=\dim(W)$. 

\item{}[Field Structure] 
In this step we endow $W'$ with a field structure.
Let $e:W'\to L$ be the unique $K$-linear isomorphism such that 
$$
  e(\psi(T^i)) \con x^i \pmod{f}
$$
for $i = 0,1,2,...,\deg(f)-1$. 
The map $e$ is uniquely determined since the $\psi(T^i)$ are a 
basis for~$W'$.  To compute $e$, we compute
the change of basis matrix from the standard basis for~$W'$ 
to the basis $\{\psi(T^i)\}$.  
This change of basis matrix is the inverse
of the matrix whose rows are the $\psi(T^i)$ 
for $i=0,...,\deg(f)-1$.

\item{}[Hecke Eigenvalues] Finally for each integer
$n\geq 1$, we have
$$       
          a_n = e(\psi(T_n)) = e(\pi(T_n(v))),
$$        
where $a_n$ is the eigenvalue of $T_n$.
Output the maps $\psi$ and $e$ and terminate. 
\end{steps}
\end{algorithm}

One reason we separate $\psi$ and~$e$ is that when $\dim(W)$ is large,
the values $\psi(T_n)$ take less space to store and are easier to
compute, whereas each one of the values $e(\psi(n))$ is
huge.\footnote{John Cremona initially suggested to me the idea of
  separating these two maps.}  The function $e$ typically involves
large numbers if $\dim(W)$ is large, since $e$ is obtained from the
iterates of a single vector.  For many applications, e.g., databases,
it is better to store a matrix that defines~$e$ and the images under
$\psi$ of many $T_n$.





\begin{example}\label{ex:23}
The space $S_2(\Gamma_0(23))$ of cusp forms has dimension~$2$
and is spanned by two $\Gal(\Qbar/\Q)$-conjugate newforms, 
one of which is
$$
f = q + aq^2 + (-2a - 1)q^3 + (-a - 1)q^4 + 2aq^5 + \cdots,
$$
where $a=(-1+\sqrt{5})/2$.
We will use Algorithm~\ref{alg:eigsys} to compute a few
of these coefficients.

The space $\sM_2(\Gamma_0(23))^+$ of modular symbols has dimension~$3$.
It has the following basis of Manin symbols:
$$
[(0,0)],\quad [(1,0)], \quad [(0,1)],
$$
where we use square brackets to differentiate Manin symbols
from vectors. 
The Hecke operator 
$$
T_2 = \left(\begin{array}{ccc}
3&0&0\\
0&0&2\\
-1&1/2&-1
\end{array}\right)
$$ has characteristic polynomial 
$(x-3)(x^2 + x - 1)$.  The kernel of $T_2-3$ corresponds to the span of
the Eisenstein series of
level~$23$ and weight~$2$, and the kernel $V$ of $T_2^2+T_2-1$
corresponds to $S_2(\Gamma_0(23))$.  (We could also 
have computed~$V$ as the kernel of the boundary map 
$\sM_2(\Gamma_0(23))^+\to \sB_2(\Gamma_0(23))^+$.)
Each of the following steps corresponds to the 
step of Algorithm~\ref{alg:eigsys} with the same number.
\begin{enumerate}
\item{}[Compute Projection] We compute projection onto $V$ (this
will suffice to give us a map $\phi$ as in the algorithm).  The
  matrix whose first two columns are the echelon basis for $V$ and
  whose last column is the echelon basis for the Eisenstein subspace
  is
$$
\left(\begin{array}{ccc}
0&0&1\\
1&0&-2/11\\
0&1&-3/11
\end{array}\right)
$$
and 
$$
 B^{-1} = \left(\begin{array}{ccc}
2/11&1&0\\
3/11&0&1\\
1&0&0
\end{array}\right),
$$
so projection onto $V$ is given by the first two rows:
$$
  \pi = \left(\begin{array}{ccc}
2/11&1&0\\
3/11&0&1
\end{array}\right).
$$


\item{}[Choose~$v$] Let $v=(0,1,0)^t$. Notice that $\pi(v)=(1,0)^t\neq 0$, 
and $v=[(1,0)]$ is a sum of
only one Manin symbol.

\item{}[Map from Hecke Ring] 
This step is purely conceptual, since no actual work needs to be
done.  We illustrate it by computing $\psi(T_1)$ and $\psi(T_2)$. 
 We have
$$\psi(T_1) = \pi(v) = (1,0)^t$$
and 
$$\psi(T_2) = \pi(T_2(v)) = \pi((0,0,1/2)^t) = (0,1/2)^t.$$

\item{}[Find Generator] We have 
$$\psi(T_2^0) = \psi(T_1) = (1,0)^t,$$
which is clearly independent from $\psi(T_2)=(0,1/2)^t$.  Thus we
find that the image of the powers of $T=T_2$ generate $V$.

\item{}[Characteristic Polynomial]  The matrix
of $T_2|_V$  is $\abcd{0}{2}{1/2}{-1}$, which has characteristic
polynomial $f=x^2+x-1$.   Of course, we already knew this because
we computed $V$ as the kernel of $T_2^2+T_2-1$.

\item{}[Field Structure] We have 
$$\psi(T_2^0) = \pi(v) = (1,0)^t
\text{ and } 
\psi(T_2) = (0,1/2).$$
The matrix with rows the $\psi(T_2^i)$
is $\abcd{1}{0}{0}{1/2}$, which has inverse $e=\abcd{1}{0}{0}{2}$.
The matrix $e$ defines an isomorphism between $V$ and
the field $$L=\Q[x]/(f) = \Q((-1+\sqrt{5})/2).$$
I.e., $e((1,0))=1$ and $e((0,1))=2x$, where $x=(-1+\sqrt{5})/2$.

\item{}[Hecke Eigenvalues] We have $a_n = e(\Psi(T_n))$.  For example,
\begin{align*}
\qquad\qquad  a_1 &= e(\Psi(T_1)) = e((1,0)) = 1,\\
  a_2 &= e(\Psi(T_2)) = e((0,1/2)) = x,\\
  a_3 &= e(\Psi(T_3)) = e(\pi(T_3(v))) \!=\! e(\pi((0,-1,-1)^t))\\
      &= e((-1,-1)^t)\!=\! -1-2x,\\
  a_4 &= e(\Psi(T_4)) = e(\pi((0,-1,-1/2)^t)) = e((-1,-1/2)^t) = -1-x,\\
  a_5 &= e(\Psi(T_5)) = e(\pi((0,0,1)^t)) = e((0,1)^t) = 2x,\\
  a_{23} &=   e(\Psi(T_{23})) = e(\pi((0,1,0)^t)) = e((1,0)^t) = 1,\\
  a_{97} &= e(\Psi(T_{23})) = e(\pi((0,14,3)^t)) = e((14,3)^t) = 14+6x.
 \end{align*}
\end{enumerate}
\end{example}

\begin{example}
  It is easier to appreciate Algorithm~\ref{alg:eigsys} after seeing
  how big the coefficients of the power series expansion of a newform
  typically are, when the newform is defined over a large field. For
  example, there is a newform
$$
  f = \sum_{n=1}^{\infty} a_n q^n\in S_2(\Gamma_0(389))
$$ 
such that if $\alpha = a_2$, then
\begin{align*}
1097385680 \cdot &a_3(f) = 
-20146763x^{19} + 102331615x^{18} + 479539092x^{17} \\
 &- 3014444212x^{16} - 3813583550x^{15} + 36114755350x^{14} \\
 &+ 6349339639x^{13} - 227515736964x^{12} + 71555185319x^{11} \\
 &+ 816654992625x^{10} - 446376673498x^{9} - 1698789732650x^{8} \\
 &+ 1063778499268x^{7} + 1996558922610x^{6} - 1167579836501x^{5}\\
 & - 1238356001958x^{4} + 523532113822x^{3} + 352838824320x^{2} \\
 &- 58584308844x - 25674258672.
\end{align*}

In contrast, if we take $v=\{0,\infty\} = (0,1) \in
\sM_2(\Gamma_0(389))^+$, then 
$$
  T_3(v) = -4(1,0) + 2(1,291) - 2(1,294) - 2(1,310) + 2(1,313) + 2(1,383).
$$
Storing $T_3(v), T_5(v), \ldots$ as vectors 
is more compact than storing
$a_3(f)$, $a_5(f)$, $\ldots$ directly as polynomials in $a_2$!
\end{example}


\section{Congruences between Newforms}
This section is about congruences between modular forms.
Understanding congruences is crucial for studying Serre's
conjectures, Galois representations, and explicit construction of
Hecke algebras.  We assume more background in algebraic number theory
here than elsewhere in this book.

\subsection{Congruences between Modular Forms}
Let $\Gamma$ be an arbitrary congruence subgroup of $\SL_2(\Z)$,
and suppose $f\in M_k(\Gamma)$ is a modular form of integer weight~$k$
for $\Gamma$. Since $\abcd{1}{N}{0}{1}\in\Gamma$ for some integer~$N$,
the form $f$ has a Fourier expansion in
nonnegative powers of $q^{1/N}$.  For a rational
number $n$, let $a_n(f)$ be the coefficient of $q^n$ in
the Fourier expansion of~$f$.
Put
$$
  \ord_q(f) = \min\{n\in \Q : a_n \neq 0\},
$$
where by convention we take $\min \emptyset = +\infty$,
so $\ord_q(0)=+\infty$.

\subsubsection{The $j$-invariant}
Let
$$
  j = \frac{1}{q} + 744 + 196884 q + \cdots
$$
be the $j$-function,\index{$j$-function} 
which is a weight~$0$ modular function
that is holomorphic except for a simple pole at $\infty$ and
has integer Fourier coefficients (see, e.g., 
\cite[Section~VIII.3.3]{serre:arithmetic}).
\begin{lemma}\label{lem:jform}
Suppose $g$ is a weight~$0$ level~$1$
modular function that is holomorphic except possibly 
with a pole of order~$n$ at $\infty$.
Then $g$ is a polynomial in $j$ of degree at most~$n$.
Moreover, the coefficients of this polynomial lie in 
the ideal $I$ generated by the coefficients $a_m(g)$
with $m\leq  0$.
\end{lemma}
\begin{proof}
If $n=0$, then $g\in M_0(\SL_2(\Z))=\C$, so $g$ is constant with
constant term in~$I$, so the statement is true.  Next suppose $n>0$ and the lemma
has been proved for all functions with smaller order poles.
Let $\alpha = a_n(g)$, and note that
$$
  \ord_q(g - \alpha j^n) = 
  \ord_q\left(g - \alpha \cdot \left(\frac{1}{q} + 744 + 196884 q + \cdots \right)^n\right) > -n.
$$
Thus by induction $h=g - \alpha j^{n}$ is a polynomial in $j$
of degree $<n$ with coefficients in the ideal generated 
by the coefficients $a_m(g)$ with $m< 0$.  It follows
that $g=\alpha \cdot j^n - h$ satisfies the conclusion
of the lemma.
\end{proof}

\subsubsection{Sturm's Theorem}
If $\O$ is the ring of integers of a number field, $\m$
is a maximal ideal of~$\O$, and $f=\sum a_n q^n \in \O[[q^{1/N}]]$
for some integer~$N$,
let
$$
 \ord_\m(f) = \ord_q(f \!\!\mod{\m}) = \min\{n\in\Q : a_n \not\in \m\}.
$$
Note that $\ord_\m(fg) = \ord_\m(f) + \ord_\m(g)$.
The following theorem was first proved in \cite{sturm:cong}.
%and our proof is an expanded version of the one in \cite{sturm:cong}.

\begin{theorem}[Sturm]\label{thm:sturm}
  Let $\m$ be a prime ideal in the ring of integers $\O$ of a
  number field~$K$, and let $\Gamma$ be a congruence subgroup
of $\SL_2(\Z)$ of index~$m$ and level~$N$.
Suppose $f\in M_k(\Gamma,\O)$ is a modular form
and 
$$
  \ord_\m(f) > \frac{km}{12}
$$
or $f \in S_k(\Gamma,\O)$ is a cusp form
and 
$$
\ord_\m(f) > \frac{km}{12} - \frac{m-1}{N}.
$$
Then $f\con 0\pmod{\m}$.
\end{theorem}
\begin{proof}
\noindent{\bf Case 1: First we assume $\Gamma=\SL_2(\Z)$.}\\
Let 
$$
 \Delta  = q + 24q^2 + \cdots \in S_{12}(\SL_2(\Z),\Z)
$$
be the $\Delta$ function. 
Since $\ord_\m(f) > k/12$, we have $\ord_\m(f^{12}) > k$.
We have 
\begin{equation}\label{eqn:sturm1}
  \ord_q(f^{12}\cdot\Delta^{-k}) 
 = 12 \cdot \ord_q(f) - k\cdot \ord_q(\Delta) \geq -k,
\end{equation}
since $f$ is holomorphic at infinity and $\Delta$ has a zero of order~$1$.
Also
\begin{equation}\label{eqn:sturm2}
  \ord_\m(f^{12}\cdot \Delta^{-k}) =  \ord_\m(f^{12}) - k\cdot \ord_\m(\Delta)
          > k - k = 0.
\end{equation}
Combining \eqref{eqn:sturm1} and \eqref{eqn:sturm2}, we see that
$$
  f^{12} \cdot \Delta^{-k} = \sum_{n\geq -k} b_n q^n,
$$
with $b_n \in \O$ and $b_n\in \m$ if $n\leq 0$. 

By Lemma~\ref{lem:jform}, 
$$
   f^{12} \cdot \Delta^{-k} \in \m[j]
$$
is a polynomial in~$j$ of degree at most~$k$ with coefficients
in $\m$.
Thus 
$$
  f^{12} \in \m[j]\cdot \Delta^k,
$$
so since the coefficients of $\Delta$ are integers, 
every coefficient of $f^{12}$ is in $\m$.
Thus
$
\ord_\m(f^{12}) = +\infty,
$
hence 
$\ord_\m(f) = +\infty$, so $f=0$, as claimed.

\vspace{2ex}
\noindent{\bf Case 2: $\Gamma$ Arbitrary}
\newcommand{\wM}{\mathfrak{M}}


Let $N$ be such that $\Gamma(N)\subset \Gamma$, so also
$f \in M_k(\Gamma(N))$.  
If $g\in M_k(\Gamma(N))$ is arbitrary, then
because $\Gamma(N)$ is a normal subgroup of $\SL_2(\Z)$,
we have that for any $\gamma \in\Gamma(N)$ and $\delta \in \SL_2(\Z)$,
$$
  (g^{[\delta]_k})^{[\gamma]_k} = g^{[\delta\gamma]_k}
          = g^{[\gamma'\delta]_k} = (g^{[\gamma']_k})^{[\delta]_k}
           = g^{[\delta]_k},
$$
where $\gamma'\in\SL_2(\Z)$.  Thus for any $\delta\in \SL_2(\Z)$,
we have that $g^{[\delta]_k} \in M_k(\Gamma(N))$, so $\SL_2(\Z)$
acts on $M_k(\Gamma(N))$.

It is a standard (but nontrivial) fact about modular forms, which
comes from the geometry of the modular curve $X(N)$ over $\Q(\zeta_N)$
and $\Z[\zeta_N]$, that $M_k(\Gamma(N))$ has a basis with Fourier
expansions in $\Z[\zeta_N][[q^{1/N}]]$ and that the action of
$\SL_2(\Z)$ on $M_k(\Gamma(N))$ preserves $$
M_k(\Gamma(N),\Q(\zeta_N)) = M_k(\Gamma(N)) \cap{}
(\Q(\zeta_N)[[q^{1/N}]])$$
and the cuspidal subspace
$S_k(\Gamma(N),\Q(\zeta_N))$.  In particular, for any
$\gamma\in\SL_2(\Z)$, $$
f^{[\gamma]_k} \in M_k(\Gamma(N), K(\zeta_N))
$$
Moreover, the denominators of $f^{[\gamma]_k}$ are bounded, since~$f$
is an $\O[\zeta_N]$-linear combination of a basis for
$M_k(\Gamma(N),\Z[\zeta_N])$, and the denominators of $f^{[\gamma]_k}$
divide the product of the denominators of the images of each of these
basis vectors under $[\gamma]_k$.
  
Let $L=K(\zeta_N)$.
Let $\wM$ be a prime of $\O_L$ that divides
$\m \O_L$.   We will now show that for each $\gamma\in\SL_2(\Z)$, the
Chinese Remainder Theorem implies that there is an element $A_\gamma\in L^*$
such that
\begin{equation}\label{eqn:sturm_crt}
A_\gamma \cdot f^{[\gamma]_k} \in M_k(\Gamma(N), \O_L)
\qquad \text{and}
\qquad \ord_{\wM} (A_{\gamma} \cdot f^{[\gamma]_k}) < \infty.
\end{equation}
First find $A \in L^*$ such that 
$A\cdot f^{[\gamma]_k}$ has coefficients in $\O_L$.
Choose $\alpha\in \wM$ with $\alpha \not\in \wM^2$, and find a negative power
$\alpha^t$ such that $\alpha^t\cdot A\cdot f^{[\gamma]_k}$
has $\wM$-integral coefficients and finite valuation.  
This is possible because we assumed that $f$ is nonzero.
Use the Chinese Remainder Theorem to find $\beta\in \O_L$
such that $\beta\con 1\pmod{\wM}$ and $\beta\con 0\pmod{\wp}$
for each prime $\wp\neq \wM$ that divides $(\alpha)$.
Then for some $s$ we have 
$$
 \beta^s \cdot \alpha^t \cdot A\cdot f^{[\gamma]_k}  
   = A_\gamma \cdot f^{[\gamma]_k}
   \in M_k(\Gamma(N),\O_L)
$$
and $\ord_{\wM}(A_\gamma \cdot f^{[\gamma]_k}) < \infty$.



Write 
$$
 \SL_2(\Z) = \bigcup_{i=1}^m \Gamma{} \gamma_i
$$
with $\gamma_1 = \abcd{1}{0}{0}{1}$, and let
$$
  F = f \cdot \prod_{i=2}^m A_{\gamma_i} \cdot f^{[\gamma_i]_k}.
$$
Then 
$
  F \in M_{km}(\SL_2(\Z))
$
and since $\wM\cap \O_K = \m$, we have $\ord_{\wM}(f) = \ord_{\m}(f)$,
so
$$
 \ord_{\wM}(F) \geq \ord_{\wM}(f) = \ord_{\m}(f) > \frac{km}{12}.
$$
Thus we can apply Case 1 to conclude that
$$
 \ord_{\wM}(F) = +\infty.
$$
Thus
\begin{equation}\label{eqn:sturm:ords}
  \infty = \ord_{\wM}(F) = \ord_\m(f) + \sum_{i=2}^m \ord_{\wM}(A_{\gamma_i} 
  f^{[\gamma]_k}),
\end{equation}
so $\ord_\m(f)=+\infty$, because of \eqref{eqn:sturm_crt}.
    
We next obtain a better bound when~$f$ is a cusp form.
Since $[\gamma]_k$ preserves
cusp forms, $\ord_{\wM}(A_{\gamma_i} f^{[\gamma]_k}) \geq \frac{1}{N}$
for each $i$.
Thus
$$
 \ord_{\wM}(F) \geq \ord_{\wM}(f)  + \frac{m-1}{N} = \ord_{\m}(f) + \frac{m-1}{N} > \frac{km}{12},
$$
since now we are merely assuming that 
$$
  \ord_{\m}(f) > \frac{km}{12} 
     - \frac{m - 1}{N}.
$$
Thus we again apply Case 1 to conclude that
$
 \ord_{\wM}(F) = +\infty,
$
and using \eqref{eqn:sturm:ords}, conclude that
$\ord_\m(f)=+\infty$.
\end{proof}

\begin{corollary}
  Let $\m$ be a prime ideal in the ring of integers $\O$ of a
  number field.  Suppose $f,g\in M_k(\Gamma,\O)$ are modular
  forms and 
$$
  a_n(f) \con  a_n(g) \pmod{\m}
$$
for all 
$$
  n\leq 
\begin{cases}
\ds\frac{km}{12} - \frac{m-1}{N} & \text{ if }f-g \in S_k(\Gamma,\O),\vspace{2pt}\\
\ds\frac{km}{12} & \text{otherwise},
\end{cases}
$$
where $m=[\SL_2(\Z):\Gamma]$.
Then $f\con g\pmod{\m}$.
\end{corollary}

Buzzard proved the following corollary, which is extremely useful
in practical computations.
It asserts that the Sturm bound for modular forms with character
is the same as the Sturm bound for $\Gamma_0(N)$.
\begin{corollary}[Buzzard]
  Let~$\m$ be a prime ideal in the ring of integers~$\O$ of a
  number field.  
  Suppose $f,g\in M_k(N,\eps,\O)$ are modular
  forms with Dirichlet character~$\eps:(\Z/N\Z)^* \to \C^*$ and assume
that
  $$
  a_n(f) \con  a_n(g) \pmod{\m}
\qquad\text{for all}\qquad
  n \leq  \frac{km}{12},
$$
where 
$$m=[\SL_2(\Z):\Gamma_0(N)] = 
   \#\P^1(\Z/N\Z) = N \cdot \prod_{p|N}\left(1+\frac{1}{p}\right).$$
Then $f\con g\pmod{\m}$.
\end{corollary}
\begin{proof}
Let $h=f-g$ and 
let $r=km/12$,
so $\ord_\m(h) > r$.  Let $s$ be the order of the
Dirichlet character $\eps$.  Then $h^s \in M_{ks}(\Gamma_0(N))$ 
and 
$$
\ord_\m(h^s) > sr = \frac{ksm}{12}.
$$
By Theorem~\ref{thm:sturm}, we have $\ord_\m(h^s) = \infty$,
so $\ord_\m(h) = \infty$.  It follows that $f\con g\pmod{\m}$.
\end{proof}

\subsubsection{Congruence for Newforms}
Sturm's paper \cite{sturm:cong} also applies some results of Asai on
$q$-expansions at various cusps to obtain a more refined result for
newforms.  

\begin{theorem}[Sturm]
Let $N$ be a positive integer that is square-free, and 
suppose~$f$ and~$g$ are two newforms in 
$S_k(N,\eps,\O)$,
where $\O$ is the ring of integers of a number field, 
and suppose that $\m$ is a maximal ideal of $\O$.
Let~$I$ be an arbitrary subset of the prime divisors of~$N$.
If 
$
 a_p(f) = a_p(g)
$
for all $p\in I$ and if
$$
 a_p(f) \con a_p(g)\pmod{\m}
$$
for all primes 
$$
  p \leq \frac{k\cdot [\SL_2(\Z):\Gamma_0(N)]}{12 \cdot 2^{\#I}},
$$ 
then $f\con g\pmod{\m}$.
\end{theorem}


The paper \cite{buzzard-stein:artin} contains a similar  result about
congruences between newforms, which does not require that the level be
square-free.  Recall from Definition~\ref{defn:conductordir} that the
conductor of a Dirichlet character $\eps$ is the largest divisor~$c$
of $N$ such that~$\eps$ factors through $(\Z/c\Z)^\times$.

\begin{theorem}
Let $N>4$ be any integer, and
suppose $f$ and $g$ are two normalized eigenforms in $S_k(N,\eps;\O)$,
where $\O$ is the ring of integers of a number field, 
and suppose that $\m$ is a maximal ideal of $\O$.
Let~$I$ be the set of prime divisors of $N$ that
do not divide~$\frac{N}{\cond(\eps)}$.
If 
$$
a_p(f) \con a_p(g)\pmod{\m}
$$
for all primes $p\in I$ and for all primes
$$
  p \leq \frac{k\cdot [\SL_2(\Z):\Gamma_0(N)]}{12 \cdot 2^{\#I}},
$$ 
then $f\con g\pmod{\m}$.
\end{theorem}
For the proof, see Lemma~1.4 and Corollary 1.7 in \cite[\S1.3]{buzzard-stein:artin}.




\subsection{Generating the Hecke Algebra}

The following theorem appeared in \cite[Appendix]{lario-schoof}, except that
we give a better bound here.  It is a nice application of the
congruence result above, which makes possible explicit computations
with Hecke rings $\T$.

\begin{theorem}\label{thm:heckegen}
Suppose $\Gamma$ is a congruence subgroup that contains
$\Gamma_1(N)$ and let 
\begin{equation}\label{eqn:perfectR}
  r=\frac{km}{12} - \frac{m - 1}{N},
\end{equation}
where $m=[\SL_2(\Z):\Gamma]$.
Then the Hecke algebra\index{Hecke algebra!generators over $\Z$} 
$$\T=\Z[\ldots, T_n,\ldots] \subset \End(S_k(\Gamma))$$
is generated as a $\Z$-module by the Hecke operators
$T_n$ for $n\leq r$.
\end{theorem}
\begin{proof}
  For any ring~$R$, let $S_k(N,R) = S_k(N;\Z)\otimes{}R$, where
  $S_k(N;\Z)\subset \Z[[q]]$ is the submodule of cusp forms with
  integer Fourier expansion at the cusp~$\infty$, and let $\T{}_R =
  \T{}\otimes_\Z R$.  For any ring~$R$, there is a perfect pairing 
$$S_k(N,R) \otimes_R \T{}_R \rightarrow R$$ 
given by $\langle f,T\rangle\mapsto a_1(T(f))$ (this is true
for~$R=\Z$, hence for any $R$).

  Let~$M$ be the submodule of $\T{}$ generated by
  $T_1,T_2,\ldots,T_r$, where~$r$ is the largest integer $\leq
  {\frac{kN}{12}}\cdot [\SL_2(\Z):\Gamma]$.
  Consider the exact sequence of additive abelian groups
  $$0\rightarrow M \mathop{\rightarrow}\limits^{i}\ \T{}
  \rightarrow \T{}/M \rightarrow 0.$$
  Let~$p$ be a prime and use the fact that
  tensor product is right exact to obtain an exact sequence
  $$M\otimes \F_p \mathop{\rightarrow}\limits^{\overline{i}}\ 
    \T \otimes\F_p \rightarrow (\T{}/M)\otimes\F_p\rightarrow 0.$$
  Suppose that $f\in S_k(N,\F_p)$ pairs to~$0$ with each of
  $T_1,\ldots, T_r$.  Then $$a_m(f)=a_1(T_m f)=\langle f, T_m\rangle =
  0$$ in $\F_p$ for each $m\leq r$.  By Theorem~\ref{thm:sturm}, it
  follows that $f = 0$. Thus the pairing restricted to the image of
  $M\otimes\F_p$ in $\T_{\F_p}$ is nondegenerate, so
because \eqref{eqn:perfectR} is perfect, it follows that
  $$
  \dim_{\F_p} \overline{i}(M\otimes\F_p) =  \dim_{\F_p} S_k(N,\F_p).
  $$
  Thus $(\T{}/M) \tensor \F_p = 0$. 
  Repeating the argument for all primes~$p$ shows that $\T{}/M=0$, 
  as claimed.
\end{proof}


\begin{remark}
In general, the conclusion of 
Theorem~\ref{thm:heckegen} is not true if one considers only
 $T_n$ where~$n$ runs over the primes less than the bound.
Consider, for example, $S_2(11)$, where the bound is~$1$ and
there are no primes $\leq 1$. 
However, the Hecke algebra is generated as an algebra
by operators $T_p$ with $p\leq r$.
\end{remark}




\begin{exercises}
\item \label{ex:g0g1quo}
  Prove that the group $\Gamma_1(N)$ is a normal subgroup of
  $\Gamma_0(N)$ and that the quotient $\Gamma_0(N)/\Gamma_1(N)$ is
  isomorphic to $(\Z/N\Z)^*$.

\item \label{ex:dbdinhecke} Prove that the operators $\dbd{d}$ are elements
of $\Z[\ldots,T_n,\ldots]$.  [Hint: Use Dirichlet's theorem
on primes in arithmetic progression.]

\item \label{exer:onlyold} Find an example like
  Example~\ref{ex:onlyold} but in which the new subspace is nonzero.
  More precisely, find an integer $N$ such that the Hecke ring on
  $S_2(\Gamma_0(N))$ is not equal to the ring generated by Hecke
  operators $T_n$ with $\gcd(n,N)=1$ and $S_2(\Gamma_0(N))_{\new}\neq
  0$.

\item \label{ex:31} 
\begin{enumerate}
\item  Following Example~\ref{ex:23}, compute a
basis for $S_2(\Gamma_0(31))$. 
\item Use Algorithm~\ref{alg:merelqexp} to compute a basis
for $S_2(\Gamma_0(31))$. 
\end{enumerate}

\end{exercises}







\chapter{Computing Periods}\label{ch:periods}


This chapter is about computing period maps associated to 
newforms.\index{newform!associated period map}
We assume you have read Chapters~\ref{ch:modsym} and
\ref{ch:modformcomp} and that you are familiar with abelian
varieties at the level of \cite{rosen:abvars}.

In Section~\ref{sec:periodintro} we introduce the period map and give
some examples of situations in which computing it is relevant.
Section~\ref{sec:newformabvar} is about how to use the period mapping
to attach an abelian variety to any newform.  In
Section~\ref{defn:extendedmodsyms}, we introduce extended modular
symbols, which are the key computational tool for quickly computing
periods of modular symbols.  We turn to numerical computation of
period integrals in Section~\ref{sec:numper}, and in
Section~\ref{sec:wntrick} we explain how to use Atkin-Lehner operators
to speed convergence. In Section~\ref{sec:computephi} we explain how
to compute the full period map with a minimum amount of work.

Section~\ref{sec:findall} briefly sketches three approaches
to computing all elliptic curves of a given conductor.


This chapter was inspired by \cite{cremona:algs}, which contains
similar algorithms in the special case of a
newform $f=\sum a_n q^n \in S_2(\Gamma_0(N))$ with $a_n\in\Z$. 

See also \cite{dokchitser:lfun} for algorithmic methods to compute
special values of very general $L$-functions, which can be used for
approximating $L(f,s)$ for arbitrary~$s$.

\section{The Period Map}\label{sec:periodintro}
Let $\Gamma$ be a subgroup of $\SL_2(\Z)$ that contains
$\Gamma_1(N)$ for some $N$, and 
suppose
$$
  f =\sum_{n\geq 1} a_n q^n\in S_k(\Gamma)
$$
is a newform (see Definition~\ref{def:newform}).
In this chapter we describe how to approximately 
compute the complex period mapping
$$
   \Phi_f : \sM_k(\Gamma) \to \C,
$$
given by 
$$
  \Phi_f(P\{\alpha,\beta\}) = \langle f,\, P\{\alpha,\beta\} \rangle
  = \int_{\alpha}^{\beta} f(z) P(z,1) d z,
$$
as in Section~\ref{sec:pairing}.  As an application, we can
approximate the special values $L(f,j)$, for $j=1,2,\ldots, k-1$ using
\eqref{eqn:lfunc_ms}.  We can also
compute the period lattice attached to a modular abelian variety,
which is an important step, e.g., in enumeration of 
$\Q$-curves (see, e.g., \cite{gonz_lario_quer}) or computation of 
a curve whose Jacobian is a modular abelian variety $A_f$
(see, e.g., \cite{wang:factors}).



\section{Abelian Varieties Attached to Newforms}
\label{sec:newformabvar}
Fix a newform $f\in S_k(\Gamma)$, where $\Gamma_1(N)\subset \Gamma$
for some~$N$.  Let $f_1,\ldots, f_d$ be the
$\Gal(\Qbar/\Q)$-conjugates of $f$, where $\Gal(\Qbar/\Q)$ acts via
its action on the Fourier coefficients, which are algebraic integers
(since they are the eigenvalues of matrices with integer entries).
Let
\begin{equation}\label{eqn:vf}
  V_f = \C f_1 \oplus \cdots \oplus \C f_d \subset S_k(\Gamma)
\end{equation}
be the subspace of cusp forms spanned by the 
$\Gal(\Qbar/\Q)$-conjugates of~$f$.
One can show using the results discussed in Section~\ref{sec:ALL}
that the above sum is direct, i.e., that $V_f$ has dimension~$d$.

The integration pairing induces a $\T$-equivariant homomorphism 
$$ 
  \Phi_f : \sM_k(\Gamma) \ra V_f^* = \Hom_\C(V_f,\C)
$$
from modular symbols to the $\C$-linear dual $V_f^*$ of $V_f$.
Here $\T$ acts on $V_f^*$ via $(\vphi t)(x) = \vphi(tx)$,
and this homomorphism is $\T$-stable by Theorem~\ref{thm:tequivar}.
The \defn{abelian variety attached to $f$} is the
quotient 
$$
  A_f(\C) = V_f^*/\Phi_f(\sS_k(\Gamma;\Z)).
$$
Here $\sS_k(\Gamma;\Z)=\sS_k(\Gamma)$, and we include
the $\Z$ in the notation to emphasize that
these are integral modular symbols.
See \cite{shimura:surles} for a proof that $A_f(\C)$
is an abelian variety (in particular, 
$\Phi_f(\sS_k(\Gamma;\Z))$ is a lattice, and 
$V_f^*$ is equipped with a nondegenerate Riemann form).

When $k=2$, we can also construct $A_f$ as a quotient
of the modular Jacobian $\Jac(X_\Gamma)$, so $A_f$ is an 
abelian variety canonically defined over~$\Q$.

In general, we have an exact sequence
$$
  0 \to \Ker(\Phi_f) \to  \sS_k(\Gamma) \to V_f^* \to A_f(\C) \to 0.
$$


\begin{remark}
  When $k=2$, the abelian variety $A_f$ has a canonical structure of
  abelian variety over $\Q$.  Moreover, there is a conjecture of Ribet
  and Serre in \cite{ribet:abvars} that describes the simple abelian
  varieties~$A$ over $\Q$ that should arise via this construction.  In
  particular, the conjecture is that $A$ is isogenous to some abelian
  variety $A_f$ if and only if $\End(A/\Q)\tensor\Q$ is a number field
  of degree $\dim(A)$.  The abelian varieties $A_f$ have this property
  since $\Q(\ldots, a_n(f),\ldots)$ embeds in $\End(A/\Q)\tensor\Q$
  and the endomorphism ring over~$\Q$ has degree at most $\dim(A)$
  (see \cite{ribet:abvars} for details).  Ribet proves that his
  conjecture is a consequence of Serre's conjecture
  \cite{serre:conjectures} on modularity of mod~$p$ odd irreducible
  Galois representations (see Section~\ref{sec:apps}).  Much of Serre's
  conjecture has been proved by Khare and Wintenberger (not
  published).  In particular, it is a theorem that if~$A$ is a simple
  abelian variety over~$\Q$ with $\End(A/\Q)\tensor\Q$ a number field
  of degree $\dim(A)$ and if $A$ has good reduction at $2$, then $A$ is
  isogenous to some abelian variety $A_f$.
\end{remark}

\begin{remark}
  When $k>2$, there is an object called a
  \defn{Grothendieck motive} that is attached to $f$ and has a canonical
  ``structure over $\Q$''.  See \cite{scholl:motivesinvent}.
\end{remark}





\section{Extended Modular Symbols}\label{defn:extendedmodsyms}%
In this section, we extend the notion of modular
symbols to allows symbols of the form
$P\{w,z\}$ where~$w$ and~$z$ are
arbitrary elements of $\h^*=\h\union\P^1(\Q)$.  
\begin{definition}[Extended Modular Symbols]
  The abelian group \sym{\esM_k} 
  of \defn{extended modular symbols} of weight~$k$
is the $\Z$-span of symbols $P\{w,z\}$, with $P\in
  V_{k-2}$ a homogeneous polynomial of degree $k-2$ with integer
  coefficients, modulo the relations $$P\cdot (\{w,y\} + \{y,z\} +
  \{z,w\}) = 0$$
  and modulo any torsion.
\end{definition}

Fix a finite index subgroup $\Gamma\subset \SL_2(\Z)$.
Just as for usual modular symbols, $\esM_k$ is equipped with
an action of $\Gamma$, and we
define the space of \defn{extended modular symbols} of weight~$k$
for $\Gamma$ to be the quotient
$$
  \esM_k(\Gamma) = (\esM_k/\langle \gamma x - x : \gamma\in\Gamma, x \in \esM_k\rangle)/\tor.
$$
The quotient $\esM_k(\Gamma)$ is torsion-free and fixed by $\Gamma$.

The integration  pairing extends naturally to a pairing
\begin{equation}\label{eqn:extpair}
      \left(S_k(\Gamma) \oplus \overline{S}_k(\Gamma)\right) 
\times 
\esM_k(\Gamma) 
\to \C,
    \end{equation}
where we recall from \eqref{eqn:sbar} that  $\overline{S}_k(\Gamma)$ denotes 
the space of antiholomorphic cusp forms.
Moreover, if
$$
  \iota: \sM_k(\Gamma)\rightarrow \esM_k(\Gamma)
$$
is the natural map, then $\iota$
respects \eqref{eqn:extpair} in the sense that
for all $f \in S_k(\Gamma) \oplus \overline{S}_k(\Gamma)$
and $x\in \sM_k(\Gamma)$, we have
$$
   \langle f, x\rangle = \langle f , \iota(x)\rangle.
$$
As we will see soon, it
is often useful to replace
$x\in\sM_k(\Gamma)$ first by $\iota(x)$ and then 
by an equivalent sum $\sum y_i$ of symbols 
$y_i\in \esM_k(N,\eps)$ such that
$\langle f, \sum y_i \rangle$ is easier to 
compute numerically than $\langle f, x\rangle$.


Let $\eps$ be a Dirichlet character of modulus~$N$.
If $\gamma=\abcd{a}{b}{c}{d} \in \SL_2(\Z)$, let $\eps(\gamma)=\eps(d)$.\index{$\eps(\gamma)$}
Let \sym{\esM_k(N,\eps)} be the quotient of
$\esM_k(N,\Z[\eps])$ by the relations $\gamma(x) -
\eps(\gamma) x$, for all $x \in \sM_k(N,\Z[\eps])$,
$\gamma \in \Gamma_0(N)$, and modulo any torsion.  


\section{Approximating Period Integrals}
\label{sec:numper}
In this section we assume $\Gamma$ is a congruence subgroup
of $\SL_2(\Z)$ that contains $\Gamma_1(N)$ for some~$N$.
Suppose~$\alp\in \h$, so $\Im(\alpha)>0$ and $m$ is
an integer such that $0\leq m\leq k-2$,
and consider the extended modular symbol
$X^mY^{k-2-m}\{\alp,\infty\}$.
Let $\langle \cdot, \cdot \rangle$ denote 
the integration pairing from Section~\ref{sec:pairing}.
Given an arbitrary cusp form $f =\sum_{n=1}^{\infty} a_n q^n\in S_k(\Gamma)$,
we have
\begin{align}
\Phi_f(X^mY^{k-2-m}\{\alpha,\infty\}) &=
\left\langle f, \, X^mY^{k-2-m}\{\alpha,\infty\}\right\rangle \\
&= \int_{\alpha}^{\infty} f(z)z^m dz \\
&=  \sum_{n=1}^{\oo} a_n \int_{\alpha}^{\infty} e^{2\pi i n z} z^m dz.
\label{eqn:intsum}
\end{align}
The reversal of summation and integration is justified because
the imaginary part of~$\alp$ is positive so that the sum
converges absolutely.  The following
lemma is useful for computing the above infinite sum.
\begin{lemma}\label{lem:intexp}
\begin{equation}\label{intexp}
\int_{\alpha}^{\infty} e^{2\pi i n z} z^m dz
   \,\,=\,\, e^{2\pi i n \alpha} 
      \sum_{s=0}^m \left(
          \frac{(-1)^s \alpha^{m-s}}
              {(2\pi i n)^{s+1}}
         \prod_{j=(m+1)-s}^m j\right).
\end{equation}
\end{lemma}
\begin{proof}
See \exref{ch:periods}{ex:intexp}
\end{proof}


In practice we will usually be interested in computing the period map
$\Phi_f$ when $f \in S_k(\Gamma)$ is a newform.  Since $f$ is a
newform, there is a Dirichlet character $\eps$ such that $f \in
S_k(N,\eps)$.  The period map $\Phi_f : \sM_k(\Gamma) \to \C$ then
factors through the quotient $\sM_k(N,\eps)$, so it suffices to
compute the period map on modular symbols in $\sM_k(N,\eps)$.

The following proposition is an
analogue of \cite[Prop. 2.1.1(5)]{cremona:algs}.
\begin{proposition}\label{modsym-errorterm}
For any $\gamma\in \Gamma_0(N)$, $P\in V_{k-2}$ and $\alp\in\h^*$,
we have the following relation in $\esM_k(N,\eps)$:
\begin{eqnarray}
P\{\oo, \gamma(\oo)\}
  &=& P\{\alp,\gamma(\alp)\} + (P - \eps(\gamma)\gamma^{-1}P)\{\oo,\alp\}\\
  &=& \eps(\gamma)(\gamma^{-1}P)\{\alp, \oo\} - P\{\gamma(\alp),\oo\}.
\end{eqnarray}
\end{proposition}
\begin{proof}
By definition, if $x\in\sM_k(N,\eps)$ is a modular symbol 
and $\gamma\in\Gamma_0(N)$, then $\gamma{}x=\eps(\gamma)x$.
Thus $\eps(\gamma)\gamma^{-1}x=x$, so 
\begin{eqnarray*}
P\{\oo, \gamma(\oo)\}
  &=& P\{\oo,\alp\} + P\{\alp,\gamma(\alp)\} + P\{\gamma(\alp),\gamma(\oo)\}\\
  &=& P\{\oo,\alp\} + P\{\alp,\gamma(\alp)\} + \eps(\gamma)\gamma^{-1}(P\{\gamma(\alp),\gamma(\oo)\})\\
  &=& P\{\oo,\alp\} + P\{\alp,\gamma(\alp)\} + \eps(\gamma)(\gamma^{-1}P)\{\alp, \oo\}\\
  &=& P\{\alp,\gamma(\alp)\} + P\{\oo,\alp\}  - \eps(\gamma)(\gamma^{-1}P)\{\oo, \alp\}\\
  &=& P\{\alp,\gamma(\alp)\} + (P - \eps(\gamma)\gamma^{-1}P)\{\oo,\alp\}.
\end{eqnarray*}
The second equality in the statement of the proposition now follows easily.
\end{proof}
In the case of weight $2$ and trivial character,
the ``error term'' 
\begin{equation}\label{eqn:erroreqn}
(P - \eps(\gamma)\gamma^{-1}P)\{\oo,\alp\}
\end{equation}
vanishes since $P$ is constant and $\eps(\gamma)=1$.  In general this term does 
not vanish.  However, we can suitably modify the formulas 
found in \cite[2.10]{cremona:algs} and still obtain an algorithm
for computing period integrals.

\begin{algorithm}{Period Integrals}
Given $\gamma\in\Gamma_0(N)$,  $P\in V_{k-2}$ and 
 $f\in S_k(N,\eps)$ 
presented as a $q$-expansion to some precision,
this algorithm outputs an approximation to
 the period integral
$\langle f, \,P\{\oo, \gamma(\oo)\}\rangle$.
\begin{steps}
\item Write $\gamma=\abcd{\hfill a}{b}{cN}{d}\in\gzero$, with $a,b,c,d\in\Z$,
and set $\alpha = \frac{-d+i}{cN}$ in Proposition~\ref{modsym-errorterm}.
\item 
Replacing~$\gamma$ by $-\gamma$ if necessary,
we find that the imaginary parts of~$\alpha$ and 
$\gamma(\alpha)=\frac{a+i}{cN}$
are both equal to the positive number $\frac{1}{cN}$.
\item 
Use \eqref{eqn:intsum} and Lemma~\ref{lem:intexp} to 
compute the integrals that appear in
Proposition~\ref{modsym-errorterm}.
\end{steps}
\end{algorithm}

It would be nice if the modular symbols
of the form $P\{\oo,\gamma(\oo)\}$ for $P\in V_{k-2}$ and
$\gamma\in \Gamma_0(N)$ were to generate a large subspace
of $\sM_k(N,\eps)\tensor\Q$.
When $k=2$ and $\eps=1$, 
Manin proved in~\cite{manin:parabolic}
that the map $\Gamma_0(N)\ra H_1(X_0(N),\Z)$ 
sending~$\gamma$ to $\{0,\gamma(0)\}$ is a surjective
group homomorphism.  When $k>2$, the author 
does not know a similar
group-theoretic statement.  However, we have the following
theorem.

\begin{theorem}\label{onlyoo}
Any element of $\sS_k(N,\eps)$ can be written in the form
$$\sum_{i=1}^n P_{i}\{\infty,\gamma_i(\infty)\}$$
for some $P_i\in V_{k-2}$ and $\gamma_i\in\gzero.$
Moreover, $P_i$ and $\gamma_i$ can be chosen so that 
$\sum  P_i = \sum \eps(\gamma_i)\gamma_i^{-1}(P_i)$,
so the error term \eqref{eqn:erroreqn} vanishes.
\end{theorem}
\begin{figure}
\begin{center}
\begin{picture}(230,150)(0,0)
\put(-10,10){\line(1,0){200}}
\put(10,10){\circle*{3}}
\put(10,0){$\infty$}
\qbezier(10,10)(50,40)(75,10)
\put(49,25){\vector(1,0){1}}
\put(45,28){$P$}
\put(75,10){\circle*{3}}
\put(75,0){$\gamma \infty$}
\qbezier(75,10)(130,50)(140,10)
\put(120,30){\vector(1,0){1}}
\put(115,34){$Q$}
\put(140,10){\circle*{3}}
\put(140,0){$\beta \infty$}
\put(20,90){\circle*{3}}
\put(17,96){$\alpha$}
\qbezier(20,90)(50,120)(55,60)
\put(40,100){\vector(1,0){1}}
\put(40,104){$P$}
\put(55,60){\circle*{3}}
\put(60,54){$\gamma \alpha$}
\qbezier(55,60)(70,180)(100,140)
\put(87,150){\vector(1,0){1}}
\put(85,154){$Q$}
\put(100,140){\circle*{3}}
\put(100,146){$\beta \alpha$}
\comment{
\qbezier(10,10)(-10,30)(10,50)
\qbezier(10,50)(40,70)(20,90)
\put(5,45){\vector(2,3){1}}
\qbezier(75,10)(90,30)(75,40)
\qbezier(75,40)(60,50)(55,60)
\put(69,45){\vector(1,-1){1}}
\put(80,35){\vector(-1,2){1}}
\qbezier(140,10)(160,30)(140,50)
\qbezier(140,50)(90,80)(110,90)
\qbezier(110,90)(130,100)(100,120)
\qbezier(100,120)(90,130)(100,140)
\put(110,90){\vector(-1,-1){1}}
}
\put(100,50){\vector(-1,4){10}}
\put(50,40){\vector(-1,4){10}}
\put(140,155){\parbox[t][4in][l]{2in}{
The modular symbol 
$$P\{\infty ,\gamma \infty \}+Q\{\gamma \infty ,\beta \infty \}$$
$$\hspace{2em}=P\{\infty ,\gamma \infty \}+Q\{\infty ,\beta \infty \}-Q\{\infty ,\gamma \infty \}$$
can be ``transported'' to 
$$P\{\alpha,\gamma \alpha\}+Q\{\gamma \alpha,\beta \alpha\},$$
provided that\\
$$P\,+\,Q\,-\,Q=\gamma ^{-1}P\,+\,\beta ^{-1}Q\,-\,\gamma ^{-1}Q.$$\\
}}
\end{picture}
\end{center}
\caption{``Transporting'' a transportable modular symbol.}
\label{fig:trans}
\end{figure}

The author and Helena Verrill prove this theorem 
in \cite{stein-verrill:transportable}. 
The condition
that the error term vanishes means that one can replace
$\infty$ by any $\alpha$ in the expression for the modular symbol
and obtain an equivalent modular symbol.  For this reason,
we call such modular symbols \defn{transportable}, as illustrated
in Figure~\ref{fig:trans}.

Note that in general not every element of the form
$P\{\oo,\gamma(\oo)\}$ must lie in $\sS_k(N,\eps)$.  However, if
$\gamma P = P$, then $P\{\oo,\gamma(\oo)\}$ does lie in
$\sS_k(N,\eps)$.  It would be interesting to know under what
circumstances $\sS_k(N,\eps)$ is generated by symbols of the form
$P\{\oo,\gamma(\oo)\}$ with $\gamma P = P$.  This sometimes fails for
$k$ odd; for example, when $k=3$, the condition $\gamma P = P$ implies
that $\gamma\in\gzero$ has an eigenvector with eigenvalue~$1$, and hence
is of finite order.  When~$k$ is even, the author can see no
obstruction to generating $\sS_k(N,\eps)$ using such symbols.


\section{Speeding Convergence Using Atkin-Lehner}\label{sec:wntrick}
Let $w_N = \abcd{0}{-1}{N}{\hfill 0} \in \Mat_2(\Z)$.
Consider the Atkin-Lehner involution $W_N$ on $M_k(\Gamma_1(N))$,
which is defined by
\begin{align*}
   W_N(f) &= N^{(2-k)/2} \cdot f|_{[w_N]_k}\\
          &= N^{(2-k)/2} \cdot f\left(-\frac{1}{Nz}\right) \cdot N^{k-1} \cdot (Nz)^{-k}\\
          &= N^{-k/2} \cdot z^{-k} \cdot f\left(-\frac{1}{Nz}\right).
\end{align*}
Here we take the positive square root if $k$ is odd.
Then $W_N^2 = (-1)^k$ is an involution when $k$ is even.

There is an operator on modular symbols, which we also denote
$W_N$, which is given by
\begin{align*}
 W_N(P \{\alpha, \beta\}) &= N^{(2-k)/2} \cdot w_N(P) \{w_N(\alpha), w_N(\beta)\}\\
  &=N^{(2-k)/2} \cdot P(-Y,NX) \left\{-\frac{1}{\alpha N}, -\frac{1}{\beta N}\right\},
 \end{align*}
and one has that if $f \in S_k(\Gamma_1(N))$ and $x\in\sM_k(\Gamma_1(N))$, then
$$
  \langle W_N(f), x\rangle = \langle f, W_N(x) \rangle.
$$

If $\eps$ is a Dirichlet character of modulus~$N$, then
the operator $W_N$ sends $S_k(N,\eps)$ to
$S_k(\Gamma_1(N),\overline{\eps})$.  Thus if $\eps^2=1$,
then $W_N$ preserves $S_k(N,\eps)$.  In particular,
$W_N$ acts on $S_k(\Gamma_0(N))$. 

 
The next proposition shows how to compute the pairing
$\langle f, P\{\oo,\gamma(\oo)\} \rangle$ under
certain restrictive assumptions. 
It generalizes a result of~\cite{cremona:periods} to 
higher weight.  %(Compare also~\cite{goldfeld:complexity}.)

\begin{proposition}\label{wntrick}
Let $f \in S_k(N,\eps)$ be a cusp form which is
an eigenform for the Atkin-Lehner operator~$W_N$ 
having eigenvalue $w\in \{\pm 1\}$ (thus $\eps^2=1$ and~$k$ is even). 
Then for any $\gamma\in\Gamma_0(N)$ and any 
$P\in V_{k-2}$, with the property that $\gamma P = \eps(\gamma)P$, we have 
the following formula, valid for any $\alpha\in\h$:
\begin{align*}
\langle f, P\{\oo,\gamma(\oo)\}\rangle &= 
 \Bigl\langle f, \quad w \frac{P(Y,-NX)}{N^{k/2-1}}\{w_N(\alp),\oo\}\\
      &\quad +
   \left(P - w \frac{P(Y,-NX)}{N^{k/2-1}} \right)\left\{i/\sqrt{N},\oo\right\} 
        -P\left\{\gamma(\alp),\oo\right\} \Bigr\rangle.
\end{align*}
Here $\ds w_N(\alpha) = -\frac{1}{N\alpha}$. 
\end{proposition}
\begin{proof}
By Proposition~\ref{modsym-errorterm} our condition on~$P$
implies that $P\{\oo,\gamma(\oo)\}= P\{\alp,\gamma(\alp)\}$. 
We describe the steps of the following computation below.\vspace{1ex}\\
\begin{align*}
&\Bigl\langle f,\quad P\{\alp,\gamma(\alp)\}\Bigr\rangle \\
  &=\left\langle 
  f,\quad\!\! P\{\alp,i/\sqrt{N}\} + P\{i/\sqrt{N},W(\alp)\}+P\{W(\alp),\gamma(\alp)\}
           \right\rangle \\
  &=\left\langle f,\quad\!\! w \frac{W(P)}{N^{k/2-1}}
        \{W(\alp),i/\sqrt{N}\} + P\{i/\sqrt{N},W(\alp)\}+P\{W(\alp),\gamma(\alp)\}
           \right\rangle.
      \end{align*}
For the first equality, we break the path into three paths,
and in the second, we apply the $W$-involution to the first
term and use that the action of~$W$ is compatible with
the pairing $\langle \,,\, \rangle$ and that~$f$ is an eigenvector
with eigenvalue $w$. 
In the following sequence of equalities we combine the first two terms
and break up the third; then we replace $\{ W(\alp), i/\sqrt{N}\}$ by
$\{W(\alp),\infty\}+\{\infty,i/\sqrt{N}\}$ and regroup:
\begin{align*}
&w \frac{W(P)}{N^{k/2-1}}
        \{W(\alp),i/\sqrt{N}\} + P\{i/\sqrt{N},W(\alp)\}+P\{W(\alp),\gamma(\alp)\}\\
 &= \left(w \frac{W(P)}{N^{k/2-1}}-P\right)
        \{W(\alp),i/\sqrt{N}\} +P\{W(\alp),\oo\} - P\{\gamma(\alp),\oo\}\\
  &=  w \frac{W(P)}{N^{k/2-1}}\{W(\alp),\oo\}
        +\left(P - w \frac{W(P)}{N^{k/2-1}}\right)\{i/\sqrt{N},\oo\} 
        -P\{\gamma(\alp),\oo\}.
\end{align*}
\end{proof}

A good choice for~$\alp$ is 
$\alp=\gamma^{-1}\left(\frac{b}{d}+\frac{i}{d\sqrt{N}}\right)$,
so that $W(\alp) = \frac{c}{d}+\frac{i}{d\sqrt{N}}$.
This maximizes the minimum of the imaginary parts 
of~$\alp$ and~$W(\alp)$, which results in series that converge
more quickly.

Let $\gamma=\abcd{a}{b}{c}{d}\in \Gamma_0(N)$.  
The polynomial
$$P(X,Y) = (cX^2 + (d-a)XY - bY^2)^{\frac{k-2}{2}}$$
satisfies $\gamma(P)=P$. 
We obtained this formula by viewing $V_{k-2}$ as
the $(k-2)$th symmetric product of the $2$-dimensional
space on which $\gzero$ acts naturally.  For example,
observe that since $\det(\gamma)=1$,
the symmetric product of two eigenvectors for~$\gamma$ is an eigenvector
in $V_{2}$ having eigenvalue~$1$. 
For the same reason, if $\eps(\gamma)\neq 1$, there need
not be a polynomial $P(X,Y)$ such that $\gamma(P)=\eps(\gamma) P$.
One remedy is to choose another~$\gamma$ so that $\eps(\gamma)=1$.

Since the imaginary parts of the terms $i/\sqrt{N}$, $\alp$ and
$W(\alp)$ in the proposition are all relatively large, the sums
appearing at the beginning of Section~\ref{sec:numper} converge
quickly if~$d$ is small.  It is {\em important} to
choose~$\gamma$ in Proposition~\ref{wntrick} with~$d$ small; otherwise
the series will converge very slowly.

\begin{remark}
  Is there a generalization of Proposition~\ref{wntrick}
  without the restrictions that $\eps^2=1$ and $k$ is even?
\end{remark}

\subsection{Another Atkin-Lehner Trick}
  Suppose $E$ is an elliptic curve and let $L(E,s)$ be the
  corresponding $L$-function.  Let $\eps\in\{\pm 1\}$ be the root
  number of $E$, i.e., the sign of the functional equation for
  $L(E,s)$, so $\Lambda(E,s) = \eps\Lambda(E,2-s)$, where
  $\Lambda(E,s) = N^{s/2} (2\pi)^{-s} \Gamma(s) L(E,s)$.
Let $f=f_E$ be the modular form associated to~$E$ (which
exists by \cite{wiles:fermat, breuil-conrad-diamond-taylor}).
  If $W_N(f) = w f$, then $\eps = -w$ (see \exref{ch:periods}{ex:funceqn}).
  We have
\begin{align*}
   L(E,1) &= -2\pi \int_{0}^{\infty} f(z)\, dz \\
          &= -2\pi i \, \left\langle f,\, \{0,\infty\}\right\rangle\\
    &=-2\pi i\, \left\langle f, \{0,i/\sqrt{N}\} + \{i/\sqrt{N}, \infty\}\right\rangle\\
    &=-2\pi i\, \left\langle wf, \{w_N(0), w_N(i/\sqrt{N})\} 
                    + \{i/\sqrt{N}, \infty\}\right\rangle\\
    &=-2\pi i\, \left\langle wf, \{\infty , i/\sqrt{N}\} 
                    + \{i/\sqrt{N}, \infty\}\right\rangle\\
    &=-2\pi i\, (w-1) \left\langle f, \{\infty , i/\sqrt{N}\} \right\rangle.\\
\end{align*}
If $w=1$, then $L(E,1)=0$.  If $w= -1$, then 
\begin{equation}\label{eqn:lone}
 L(E,1) = 4\pi i \left\langle f, \{\infty , i/\sqrt{N}\} \right\rangle
    =  2 \sum_{n=1}^{\infty} \frac{a_n}{n} e^{-2\pi n/\sqrt{N}}.
  \end{equation}
  
  For more about computing with $L$-functions of elliptic curves,
  including a trick for computing $\eps$ quickly without directly
  computing $W_N$, see \cite[\S7.5]{cohen:course_ant} and
  \cite[\S2.11]{cremona:algs}.  One can also find higher derivatives
  $L^{(r)}(E,1)$ by a formula similar to \eqref{eqn:lone} (see
  \cite[\S2.13]{cremona:algs}).  The methods in this chapter for
  obtaining rapidly converging series are not just of computational
  interest; see, e.g., \cite{greenberg:bsd} for a nontrivial
  theoretical application to the Birch and Swinnerton-Dyer conjecture.


\section{Computing the Period Mapping}\label{sec:computephi}%
\index{period mapping!computation of}%
Fix a newform $f=\sum a_nq^n\in S_k(\Gamma)$, 
where $\Gamma_1(N)\subset \Gamma$
for some~$N$.
Let $V_f$ be as in \eqref{eqn:vf}.


Let $\Theta_f:M_k(\Gamma;\Q) \to V$ be {\em any} $\Q$-linear map with the same kernel
as $\Phi_f$; we call any such map a \defn{rational period mapping}
associated to~$f$.  Let $\Phi_f$ be the period mapping associated
to the $\Gal(\Qbar/\Q)$-conjugates of~$f$.  We
have a commutative diagram
$$\xymatrix{
 {\sM_k(\Gamma;\Q)}\ar[dr]_{\Theta_f}\ar[rr]^{\Phi_f} 
   &          & \Hom_\C(V_f,\C) \\
    &     V\ar@{^(->}[ur]^{i_f}
}
$$
Recall from Section~\ref{sec:newformabvar} that the cokernel of
$\Phi_f$ is the abelian variety $A_f(\C)$.

The Hecke algebra $\T$ acts on the linear dual
$$\sM_k(\Gamma;\Q)^* = \Hom (\sM_k(\Gamma), \Q)$$
by $(t \vphi)(x) = \vphi(tx)$.
Let $I=I_f\subset \T$ be the kernel of the 
ring homomorphism $\T\ra \Z[a_2, a_3, \ldots]$ that sends~$T_n$ to~$a_n$. 
Let $$ 
  \sM_k(\Gamma;\Q)^*[I] = \{\vphi \in \sM_k(\Gamma;\Q)^* : t \vphi = 
   0 \text{ all }t \in I\}.
$$
Since $f$ is a newform, one can show that
$\sM_k(\Gamma;\Q)^*[I]$ has dimension $d$.
Let $\theta_1, \ldots, \theta_d$ be a basis
for $\sM_k(\Gamma;\Q)^*[I]$, so
$$
  \Ker(\Phi_f) = \Ker(\theta_1) \oplus \cdots \oplus \Ker(\theta_d).
$$
We can thus compute $\Ker(\Phi_f)$, hence a choice of $\Theta_f$.
To compute $\Phi_f$, it remains to compute~$i_f$. 

Let $S_k(\Gamma;\Q)$ denote the space of cusp forms with
$q$-expansion in $\Q[[q]]$.  
By \exref{ch:periods}{ex:qvs}
$$
  S_k(\Gamma; \Q)[I] = S_k(\Gamma)[I] \cap \Q[[q]]
$$
is a $\Q$-vector space of dimension~$d$.
Let $g_1,\ldots,g_d$ be a basis for this $\Q$-vector space.
We will compute $\Phi_f$ with respect to the basis of
$\Hom_\Q(S_k(\Gamma;\Q)[I];\C)$ dual to this basis. 
Choose elements $x_1,\ldots,x_d\in \sM_k(\Gamma)$
with the following properties:
\begin{enumerate}
\item Using Proposition~\ref{modsym-errorterm} or Proposition~\ref{wntrick},
it is possible to compute the period integrals 
$\langle g_i, x_j \rangle$, $i,j\in\{1,\ldots, d\}$,
efficiently.
\item The $2d$ elements $v+\eta(v)$ and $v-\eta(v)$ for $v=\Theta_f(x_1),\ldots,\Theta_f(x_d)$ 
span a space of dimension $2d$ (i.e., they span $\sM_k(\Gamma)/\Ker(\Phi_f)$).
\end{enumerate}
Given this data, we can compute
$$i_f (v+\eta(v)) = 
   2\Re(\langle g_1, x_i\rangle, \ldots, \langle g_d, x_i\rangle)$$
and
$$i_f(v-\eta(v)) = 
   2i\Im(\langle g_1, x_i\rangle, \ldots, \langle g_d, x_i\rangle).$$
We break the integrals into real and imaginary parts because this
increases the precision of our answers.
Since the vectors $v_n+\eta(v_n)$ and $v_n-\eta(v_n)$, $n=1,\ldots,d$, span 
$\sM_k(N,\eps;\Q)/\Ker(\Phi_f)$,  we have computed~$i_f$. 

\begin{remark}
We want to find symbols~$x_i$ satisfying the
conditions of Proposition~\ref{wntrick}.  This is usually possible
when~$d$ is very small, but in practice it is difficult  when~$d$ is large. 
\end{remark}

\begin{remark}
The above strategy was motivated by
\cite[\S2.10]{cremona:algs}.
\end{remark}






\section{All Elliptic Curves of Given Conductor}\label{sec:findall}
\label{sec:ecmodsym}
Using modular symbols and the period map, we can compute all elliptic
curves over~$\Q$ of conductor~$N$, up to isogeny.  The algorithm in
this section gives all \defn{modular elliptic curves} (up to isogeny),
i.e., elliptic curves attached to modular forms, of conductor~$N$.
Fortunately, it is now known by \cite{wiles:fermat,
  breuil-conrad-diamond-taylor, taylor-wiles:fermat} that every
elliptic curve over~$\Q$ is modular, so the procedure of this section
gives all elliptic curves (up to isogeny) of given conductor.  See
\cite{cremona:history} for a nice historical discussion of this
problem.

\begin{algorithm}{Elliptic Curves of Conductor~$N$}
Given~$N> 0$, this
algorithm outputs equations for all
elliptic curves of conductor~$N$, up to isogeny.
\begin{steps}
\item{}[Modular Symbols] 
  Compute $\sM_2(\Gamma_0(N))$ using Section~\ref{sec:compg0n}.

\item{}[Find Rational Eigenspaces] \label{step:rateigen}
Find the $2$-dimensional eigenspaces $V$ in 
$\sM_2(\Gamma_0(N))_{\new}$ that correspond to elliptic curves.
Do {\em not} use the algorithm for decomposition from 
Section~\ref{sec:decompmodsym}, which is too complicated
and gives more information than we need.  Instead, for the first
few primes $p\nmid N$, compute all eigenspaces $\Ker(T_p - a)$, where
$a$ runs through integers with $-2\sqrt{p} < a  <  2\sqrt{p}$.
Intersect these eigenspaces to find the eigenspaces that 
correspond to elliptic curves.  To find just the new ones,
either compute the degeneracy maps to lower level or find
all the rational eigenspaces of all levels that strictly
divide~$N$ and exclude them.

\item{}[Find Newforms]\label{step:ratnew}
Use Algorithm~\ref{alg:eigsys}
to compute to some precision each newform $f=\sum_{n=1}^{\infty} a_n q^n \in \Z[[q]]$ associated
to each eigenspace $V$ found in step~(\ref{step:rateigen}).

\item{}[Find Each Curve] 
For each newform~$f$ 
found in step~(\ref{step:ratnew}), do the following:

\begin{steps}
\item{}[Period Lattice] Compute the corresponding
period lattice $\Lambda=\Z\omega_1 + \Z\omega_2$ 
by computing the image of $\Phi_f$, 
as described in Section~\ref{sec:computephi}.

\item{}[Compute $\tau$]\label{step:tau}
Let $\tau=\omega_1/\omega_2$.  If $\Im(\tau)<0$, swap
$\omega_1$ and $\omega_2$, so $\Im(\tau)>0$.  
By successively applying generators of $\SL_2(\Z)$,
we find an $\SL_2(\Z)$ equivalent element $\tau'$
in $\cF$, i.e.,
$|\Re(\tau')|\leq 1/2$ and $|\tau|\geq 1$.

\item{}[$c$-invariants] Compute the invariants $c_4$
and $c_6$ of the lattice~$\Lambda$ using the following 
rapidly convergent series:
\begin{align*}
c_4 &= \left(\frac{2\pi}{\omega_2}\right)^4 \cdot
      \left(1+240\sum_{n=1}^{\infty}\frac{n^3q^n}{1-q^n}\right),\\
c_6 &= \left(\frac{2\pi}{\omega_2}\right)^6 \cdot
      \left(1-504\sum_{n=1}^{\infty}\frac{n^5q^n}{1-q^n}\right),
    \end{align*}
where $q=e^{2\pi i \tau'}$, where $\tau'$ is as in step~(\ref{step:tau}).
A theorem of Edixhoven
(that the Manin constant is an integer) implies that
the invariants $c_4$ and $c_6$ of $\Lambda$ are integers,
so it is only necessary to compute $\Lambda$ to large
precision to completely determine them.

\item{}[Elliptic Curve] An elliptic curve with
invariants $c_4$ and $c_6$ is 
$$
E: \quad y^2 = x^3 -\frac{c_4}{48}x  -\frac{c_6}{864}.
$$

\item{}[Prove Correctness]\label{step:correct}  
Using Tate's algorithm, find
the conductor of $E$. If the conductor is not~$N$, then recompute $c_4$ and
$c_6$ using more terms of $f$ and real numbers
 to larger precision, etc.  If the conductor is~$N$, 
compute the coefficients $b_p$ of the modular form $g=g_E$
attached to the elliptic curve $E$, for $p\leq \#\P^1(\Z/N\Z)/6$.
Verify that $a_p = b_p$, where $a_p$ are the coefficients
of~$f$.   If this equality holds, then $E$ must be isogenous
to the elliptic curve attached to~$f$, by the Sturm
bound (Theorem~\ref{thm:sturm}) and Faltings's isogeny theorem.
If the equality fails for some~$p$, recompute $c_4$ and $c_6$
to larger precision.
\end{steps}
\end{steps}
\end{algorithm}
There are numerous tricks to optimize the above algorithm.
For example, often one can work separately with
 $\sM_k(\Gamma_0(N))_{\new}^+$
and 
$\sM_k(\Gamma_0(N))_{\new}^-$ and
get enough information to find $E$, up to isogeny
(see \cite{cremona:periods}).


Once we have one curve from each isogeny class of curves of
conductor~$N$, we find each curve in each isogeny class (which is
another interesting problem discussed in \cite{cremona:algs}), hence
all curves of conductor~$N$.  If $E/\Q$ is an elliptic curve, then any
curve isogenous to $E$ is isogenous via a chain of isogenies of prime
degree.  There is an {\em a priori} bound on the degrees of these
isogenies due to Mazur.  Also, there are various methods for finding
all isogenies of a given degree with domain $E$.   See
\cite[\S3.8]{cremona:algs} for more details.



\subsection{Finding Curves: $S$-Integral Points}
In this section we briefly survey an alternative approach to
finding curves of a given conductor by finding integral points on
other elliptic curves.

Cremona and others have developed a complementary
approach to the problem of computing all elliptic curves of given
conductor (see \cite{cremona-lingham}).  Instead of computing all
curves of given conductor, we instead consider the seemingly more
difficult problem of finding all curves with good reduction outside a
finite set $S$ of primes.  Since one can compute the conductor of a
curve using Tate's algorithm \cite[\S3.2]{tate:antwerpiv,
  cremona:algs}, if we know all curves with good reduction
outside~$S$, we can find all curves of conductor~$N$ by letting $S$ be
the set of prime divisors of~$N$.

There is a strategy for finding all curves with good reduction
outside~$S$.  It is not an algorithm, in the sense
that it is always guaranteed to terminate (the modular symbols method
above {\em is} an algorithm), but in practice it often works.  Also,
this strategy makes sense over any number field, whereas the modular
symbols method does not (there are generalizations of modular
symbols to other number fields).

Fix a finite set~$S$ of primes of a number field~$K$.  It is a theorem
of Shafarevich that there are only finitely many elliptic curves with
good reduction outside~$S$ (see \cite[Section~IX.6]{silverman:aec}).  His
proof uses that the group of $S$-units in $K$ is finite and Siegel's
theorem that there are only finitely many $S$-integral points on an
elliptic curve.  One can make all this explicit, and sometimes in
practice one can compute all these $S$-integral points.  

The problem of finding all elliptic curves with good reduction outside
of~$S$ can be broken into several subproblems, the main ones being
\begin{enumerate}
\item determine the following finite subgroup of $K^*/(K^*)^m$:
$$
  K(S,m) = \{x \in K^*/(K^*)^m \, : \, m \mid \ord_\p(x) \text{ all } \p\not\in S\};
$$
\item find all $S$-integral points on certain elliptic curves
$y^2=x^3 + k$.
\end{enumerate}

In \cite{cremona-lingham}, there is one example, where they find all
curves of conductor $N=2^8\cdot 17^2 = 73984$ by finding all curves
with good reduction outside $\{2, 17\}$.  They finds $32$ curves of
conductor $73984$ that divide into $16$ isogeny classes. 
(Note that $\dim S_2(\Gamma_0(N)) = 9577$.)

\subsection{Finding Curves: Enumeration}
One can also find curves by simply enumerating Weierstrass equations.
For example, the paper \cite{stein-watkins:ants5} discusses a database
that the author and Watkins created that contains hundreds of millions
of elliptic curves.  It was constructed by enumerating Weierstrass
equations of a certain form.  This database does not contain {\em
  every} curve of each conductor included in the database.  It is,
however, fairly complete in some cases.  For example, using the Mestre
method of graphs \cite{mestre:graphs}, we verified in \cite{rank4} that
the database contains all elliptic curve of prime conductor $< 234446$, which implies that the smallest conductor
rank $4$ curve is composite.
















\begin{exercises}

\item \label{ex:intexp} Prove Lemma~\ref{lem:intexp}. 

\item \label{ex:funceqn} Suppose $f \in S_2(\Gamma_0(N))$ is a newform
and that $W_N(f) = w f$.  Let $\Lambda(E,s) = N^{s/2} (2\pi)^{-s} \Gamma(s) L(E,s)$.
Prove that $$\Lambda(E,s) = -w \Lambda(E,2-s).$$
[Hint: Show that $\Lambda(f,s) = \int_{0,\infty} f(iy/\sqrt{N})y^{s-1}\, dy$.
Then substitute $1/y$ for $y$.]

\item \label{ex:qvs} Let $f=\sum a_n q^n \in \C[[q]]$ be a power
  series whose coefficients $a_n$ together generate a number field $K$ of
  degree $d$ over $\Q$.  Let $V_f$ be the complex vector space spanned
  by the $\Gal(\Qbar/\Q)$-conjugates of~$f$. 
\begin{enumerate}
\item Give an example to show that $V_f$ need not have dimension~$d$.
\item Suppose $V_f$ has dimension~$d$.  Prove that $V_f \cap
  \Q[[q]]$ is a $\Q$-vector space of dimension $d$. 
\end{enumerate}

\item \label{ex:11} Find an elliptic curve of conductor $11$
using Section~\ref{sec:ecmodsym}.

\end{exercises}




























\chapter{Solutions to Selected Exercises}\label{ch:solutions}
\section{Chapter~\ref{ch:modform}}%ack: r. bradshaw wrote all these!
\begin{enumerate}
\item\exref{ch:modform}{ex:upperhalfpres}.
Suppose $\gamma = \abcd{a}{nb}{c}{d} \in \GL_2(\R)$ is a matrix
with positive determinant.
Then $\gamma$ is a linear fractional transformation that fixes the
real line, so it must either fix or swap the upper and
lower half planes. Now
$$\gamma(i) = \frac{ai+b}{ci+d} = \frac{ac+bd + (ad-bc)i}{d^2+c^2},$$
so since $\det \gamma = ad-bc > 0$, the imaginary part of $\gamma(i)$
is positive; hence $\gamma$ sends the upper half plane to itself.

\item\exref{ch:modform}{ex:meromorphic}.
Avoiding poles, the quotient rule for differentiation goes
  through exactly as in the real case, so any rational
  function $f(z) = p(z)/q(z)$ ($p,q \in \C[z]$) is holomorphic on
  $\C-\{\alpha : q(\alpha) = 0 \}$. By the fundamental theorem of
  algebra, this set of poles is finite, and hence it is discrete.
  Write $q (z) = a_n(z-\alpha_1)^{r_1}\cdots(z-\alpha_k)^{r_k}$ for
  each $\alpha_i$ and let $q_i(z)=q(z)/(z-\alpha_i)^{r_i}$ which is a
  polynomial nonzero at $\alpha_i$. Thus for each $i$ we have
  $(z-\alpha_i)^{r_i}f(z) = p(z)/q'(z)$ is holomorphic at $\alpha_i$
  and hence $f(z)$ is meromorphic on $\C$.

\item \exref{ch:modform}{ex:wmfprod}.
\begin{enumerate}
\item The product $fg$ of two meromorphic functions on the upper half
  plane is itself meromorphic. Also, for all $\gamma \in \SL_2(\Z)$
  we have
\begin{align*}
(fg)^{[\gamma]_{k+j}} 
&= \frac{1}{(cz+d)^{k+j}}((fg) \circ \gamma)\\
&= \frac{1}{(cz+d)^k}(f \circ \gamma) \frac{1}{(cz+d)^j}(g \circ \gamma)
= fg,
\end{align*} so $fg$ is weakly modular. 

\item If $f$ is meromorphic on the upper half plane, then so is $1/f$.
 Now
$$
  \qquad\qquad\frac{1}{f} = \frac{1}{(cz+d)^{-k} f \circ \gamma} = (cz+d)^k ((1/f) \circ \gamma) = \frac{1}{f}^{[\gamma]_{-k}},
$$
so $1/f$ is a weakly modular form of weight $-k$. 

\item Let $f$ and $g$ be modular functions. Then, as above, $fg$ is a
  weakly modular function.  Let $\sum_{n=m}^\infty a_nq^n$ and
  $\sum_{n=m'}^\infty b_nq^n$ be their $q$-expansions around any
  $\alpha \in \P^1(\Q)$; then their formal product is the
  $q$-expansion of $fg$. But the formal product of two Laurent series
  about the same point is itself a Laurent series with convergence in
  the intersection of the convergent domains of the original series,
  so $fg$ has a meromorphic $q$-expansion at each $\alpha \in
  \P^1(\Q)$ and hence at each cusp.

\item We are in exactly the same case as in part (c), but because $f$
  and $g$ are modular functions, $m, m'\ge 0$ and hence the function
  is holomorphic at each of its cusps.

\end{enumerate}

\item \exref{ch:modform}{ex:nomodformodd}.
Let $f$ be a weakly modular function of odd weight
$k$. Since
$\gamma 
   = \abcd{-1}{0}{0}{-1} \in \SL_2(\Z)$, we have $f(z) =
(-1)^{-k}f(\gamma(z)) = -f(z)$ so $f = 0$.

\item  \exref{ch:modform}{ex:allsame}.
  Because $\SL_2(\Z/1\Z)$ is the trivial group, $\Gamma(1) =
  \ker(\SL_2(\Z) \rightarrow \SL_2(\Z/1\Z))$ must be all of
  $\SL_2(\Z)$.  As $\SL_2(\Z) = \Gamma(1) \subset \Gamma_1(1)
  \subset \Gamma_0(1) \subset \SL_2(\Z)$, we must have $\Gamma(1) =
  \Gamma_1(1) = \Gamma_0(1) = \SL_2(\Z)$.

\item \exref{ch:modform}{ex:conggamma1}.
\begin{enumerate}
\item The group $\Gamma_1(N)$ is the inverse image of
the subgroup of $\SL_2(\Z/N\Z)$ generated by $\abcd{1}{1}{0}{1}$, and
the inverse image of a group (under a group homomorphism)
is a group.
\item The group contains the kernel of the homomorphism
$\SL_2(\Z) \to \SL_2(\Z/N\Z)$, and 
that kernel has finite index since the quotient is
contained in $\SL_2(\Z/N\Z)$, which is finite. 
\item Same argument as previous part. 
\item The level is at most $N$ since both groups contain $\Gamma(N)$.
  It can be no greater than $N$ since $\abcd{1}{N}{0}{1}$ is in both
  groups.
\end{enumerate}

\item \exref{ch:modform}{ex:grpact}.
See \cite[Lemma~1.2.2]{diamond-shurman}.

\item \exref{ch:modform}{ex:sl2ztrans}.
   Let $\alpha = p/q \in \Q$,
  where $p$ and $q$ are relatively prime. By the Euclidean
  algorithm, we can find $x, y \in \Z$ such that $px + qy = 1$. Let
  $\gamma_\alpha = \abcd{p}{-y}{q}{x}$. Note that $\gamma_\alpha \in
  \SL_2(\Z)$ and $\gamma_{\alpha}(\infty) = \alpha$. Also let
  $\gamma_\infty$ be the identity map on $\P^1(\Q)$. Now
  $\gamma_\beta^{-1}$ sends $\beta$ to $\infty$ so we have
  $\gamma_\alpha \circ \gamma_\beta^{-1}$ which sends $\alpha$ to
  $\beta$.

\end{enumerate}

\section{Chapter~\ref{ch:level1}}

\begin{enumerate}
\item \exref{ch:level1}{ex:zetaval}.
We have
$$\zeta(26) = \frac{1315862 \cdot \pi^{26}}{11094481976030578125}.$$
Variation: Compute $\zeta(28)$.

\item \exref{ch:level1}{ex:odd_bernoulli}. Omitted.

\item \exref{ch:level1}{ex:expeis}.
\begin{align*}
 E_8 &=  -\frac{B_8}{16} + q + \sum_{n=2}^{\infty} \sigma_7(n) q^n\\
     &= \frac{1}{480} + q + 129q^{2} + 2188q^{3} + \cdots.
   \end{align*}
Variation: Compute $E_{10}$.


\item \exref{ch:level1}{ex:g4g6quo}. Omitted.

\item  \exref{ch:level1}{ex:vm}.
We have $d=\dim S_{28} = 2$.  A choice of
$a,b$ with $4a+6b\leq 14$ and $4a+6b \con 4\pmod{12}$
is $a=1, b=0$.
A basis for $S_{28}$ is then 
\begin{align*}
\qquad \qquad g_1 &=  \Delta F_6^{2(2-1) + 0} F_4 =q - 792q^{2} - 324q^{3} + 67590208q^{4} + \cdots,\\
g_2 &=  \Delta^2 F_6^{2(2-2)+0} F_4 = q^{2} + 192q^{3} - 8280q^{4} + \cdots.
\end{align*}
The Victor Miller basis is then
\begin{align*}
  f_1 &= g_1 +729g_2 = q + 151740q^{3} + 61032448q^{4} + \cdots, \\
  f_2 &= g_2 = q^{2} + 192q^{3} - 8280q^{4} + \cdots.
\end{align*}
Variation: Compute the Victor Miller basis for $S_{30}$.

\item \exref{ch:level1}{ex:elkies28}.
From the previous exercise we have $f = \Delta^2 F_4$.
Then
\begin{align*}
\qquad \qquad f = \Delta^2 F_4 &= \left(\frac{F_4^3 - F_6^2}{-1728}\right)^2 \cdot F_4 \\
      &= \left(\frac{\left(-\frac{8}{B_4} E_4\right)^3 - 
              \left(-\frac{12}{B_6} E_6\right)^2}{-1728}\right)^2
           \cdot \left(-\frac{8}{B_4} E_4\right) \\
      &= 5186160 E_4E_6^{4} - 564480000 E_4^{4}E_6^{2} + 15360000000 E_4^{7}.
\end{align*}


\item \exref{ch:level1}{ex:fkint}.  
No, it is not always integral.
For example, for $k=12$, the coefficient of $q$
is $-2\cdot 12/B_{12} = 65520/691 \not\in\Z$.
Variation: Find, with proof, the set of all $k$ such that the normalized
series $F_k$ {\em is} integral (use that $B_k$ is eventually very
large compared to $2k$).

\item \exref{ch:level1}{ex:vm32t2}. 
We compute the Victor Miller basis to precision 
great enough to determine $T_2$. This means
we need up to $O(q^5)$. 
\begin{align*}
 f_0 &= 1 + 2611200q^{3} + 19524758400q^{4} + \cdots,\\
 f_1 &= q + 50220q^{3} + 87866368q^{4} +\cdots,\\
 f_2 &= q^{2} + 432q^{3} + 39960q^{4} + \cdots.
\end{align*}
Then the matrix of $T_2$ on this basis is 
$$
\left(\begin{array}{rrr}
2147483649&0&19524758400\\
0&0&2235350016\\
0&1&39960
\end{array}\right).
$$
(The rows of this matrix are the linear combinations that
give the images of the $f_i$ under $T_2$.)
This matrix has characteristic polynomial
$$
   (x - 2147483649) \cdot (x^{2} - 39960x - 2235350016).
$$


    
\end{enumerate}

\section{Chapter~\ref{ch:modform2}}

\begin{enumerate}

\item \exref{ch:modform2}{ex:x0n}. 
Write $g=\abcd{a}{b}{c}{d}$, so 
$\lambda' = \frac{a\lambda + b}{c\lambda + d}$.
Let $f$ be the isomorphism $\C/\Lambda \to \C/\Lambda'$
given by $f(z) = z/(c\lambda + d)$.
We have
$$
 \qquad\quad f\left(\frac{1}{N}\right) = 
    \frac{1}{N(c\lambda + d)} = \frac{a}{N} - \frac{c}{N} \cdot \frac{a \lambda + b}{c\lambda + d}
   \cong \frac{a}{N} \pmod{\Z + \Z \lambda'},
$$
where the second equality can be verified easily by expanding out each side,
and for the congruence we use that $N\mid c$.
Thus the subgroup of $\C/\Lambda$ generated by $\frac{1}{N}$ is
taken isomorphically to the subgroup of $\C/\Lambda'$ generated
by $\frac{1}{N}$.

\item \exref{ch:modform2}{ex:modsymzero}. 
For any integer $r$, 
we have $\abcd{1}{r}{0}{1}\in \Gamma_0(N)$,
so $\{0,\infty\} = \{r, \infty\}$.
Thus
$$
  \qquad  0 = \{0,\infty\} - \{0,\infty\} = 
           \{n,\infty\} - \{m,\infty\} = 
               \{n,\infty\} + \{\infty, m\} = \{n,m\}.
$$

\item \exref{ch:modform2}{ex:bij}.
\begin{enumerate}
\item $(0:1), (1:0), (1:1), \ldots, (1,p-1)$.
\item $p+1$.
\item See \cite[Prop.~2.2.1]{cremona:algs}.
\end{enumerate}

\item \exref{ch:modform2}{ex:intermsofmanin}.
We start with $b=4$, $a=7$.  Then $4\cdot 2 \con 1\pmod{7}$.
Let $\delta_1 = \abcd{4}{1}{7}{2}\in\SL_2(\Z)$.  Since 
$\delta_1\in \Gamma_0(7)$,
we use the right coset representative $\abcd{1}{0}{0}{1}$ and see that
$$
   \{0,4/7\} = \{0,1/2\} + \abcd{1}{0}{0}{1}\{0,\infty\}.
$$
Repeating the process, we have $\delta_2 = \abcd{1}{1}{2}{0}$, which is
in the same coset at $\abcd{0}{-1}{1}{\hfill0}$.  Thus
$$
\{0,1/2\} = \abcd{0}{6}{1}{0} \{0,\infty\} + \{0,0\}.
$$
Putting it together gives 
$$
  \{0,4/7\} = \abcd{1}{0}{0}{1} \{0,\infty\} + \abcd{0}{6}{1}{0}\{0,\infty\}
             = [(0,1)] + [(1,0)].
$$

\item \exref{ch:modform2}{ex:heckes2g03}.
%soln from R. Bradshaw.
\begin{enumerate}
\item Coset representatives for $\Gamma_0(3)$ in $\SL_2(\Z)$ are 
$$
\mtwo{ 1 }{ 0  }{ 0  }{ 1 },\quad
\mtwo{ 1 }{ 0 }{ 1  }{ 1 },\quad
\mtwo{ 1 }{ 0  }{ 2  }{ 1 },\quad
\mtwo{ 0 }{ -1  }{ 1  }{ 0 },
$$
which we refer to below as $[r_0], [r_1], [r_2], $ and $[r_3]$, respectively. 

\item 
In terms of representatives we have 
%
% T := Matrix([[1,-1],[1,0]]);
% S := Matrix([[0,-1],[1,0]]);
% r[0] := Matrix([[1,0],[0,1]]);
% r[1] := Matrix([[1,0],[1,1]]);
% r[2] := Matrix([[1,0],[2,1]]);
% r[3] := Matrix([[0,-1],[1,0]]);

%for i from 0 to 3 do  
%A := multiply(r[i],T^2);
%printf("r[%d] * X  = ", i);
%print(evalm(A));
%for j from 0 to 3 do
%if ((multiply(A, r[j]^(-1))[2,1] mod 3)=0) then printf("in class r[%d]\n", j);
%printf("r[%d] * X * r[%d]^(-1)\n", i, j);
%print(evalm(multiply(A, r[j]^(-1))));
%fi;
%od;
%od:
%
$$
\begin{array}{ccc}
& [r_0]+[r_3]=0, & [r_0]+[r_3]+[r_2]=0, \\
& [r_1]+[r_2]=0, & [r_1]+[r_1]+[r_1]=0, \\
& [r_2]+[r_1]=0, & [r_2]+[r_0]+[r_3]=0, \\
& [r_3]+[r_0]=0, & [r_3]+[r_2]+[r_0]=0.
\end{array}
$$

\item By the first three relations we have $[r_2] = [r_1] = 0 = 0[r_0]$ and $[r_3] = -1[r_0]$. 

\item \begin{align*}
\qquad\quad T_2([r_0]) & =  
[r_0] \abcd{1}{0}{0}{2}
+ [r_0]\abcd{2 }{ 0 }{ 0 }{ 1}
+ [r_0]\abcd{2 }{ 1 }{ 0 }{ 1}
+ [r_0]\abcd{1 }{ 0 }{ 1 }{ 2} \\
& = 
\left[ \abcd{1 }{ 0 }{ 0 }{ 2} \right]
+ \left[ \abcd{2 }{ 0 }{ 0 }{ 1} \right]
+ \left[ \abcd{2 }{ 1 }{ 0 }{ 1} \right]
+ \left[ \abcd{1 }{ 0 }{ 1 }{ 2} \right] \\
& =  [r_0] + [r_0] + [r_0] + [r_2] \\
& =  3[r_0].
\end{align*}

\end{enumerate}


\end{enumerate}

\section{Chapter~\ref{ch:dirichlet}}

\begin{enumerate}
\item \exref{ch:dirichlet}{ex:notdir}.
Suppose~$f$ is a Dirichlet character with modulus~$N$.
Then $-1=f(-1) = f(-1+N)=1$, a contradiction.

\item \exref{ch:dirichlet}{ex:cyclic}.
\begin{enumerate}
\item Any finite subgroup of the multiplicative group of a field
is cyclic (since the number of roots of a polynomial over a field
is at most its degree), so $(\Z/p\Z)^*$ is cyclic.  Let $g$ be an
integer that reduces to a generator of $(\Z/p\Z)^*$.
Let $x=1+p \in (\Z/p^n\Z)^*$; by the binomial theorem 
$$\qquad\qquad x^{p^{n-2}} = 1 + p^{n-2} \cdot p + \cdots
         \con 1 + p^{n-1} \not\con 0 \pmod{p^{n}},
$$
so $x$ has order $p^{n-1}$. Since $p$ is odd, $\gcd(p^{n-1},p-1) = 1$,
so $xg$ has order $p^{n-1}\cdot (p-1) = \vphi(p^n)$; hence
$(\Z/p^n\Z)^*$ is cyclic.  
\item By the binomial theorem $(1+2^2)^{2^{n-3}} \not\con 1 \pmod{2^n}$,
so $5$ has order $2^{n-2}$ in $(\Z/2^n\Z)^*$, and clearly $-1$ has
order $2$. Since $5\con 1\pmod{4}$, $-1$ is not a power of $5$ in 
$(\Z/2^n\Z)^*$.  Thus the subgroups $\langle -1 \rangle$
and $\langle 5 \rangle$ have trivial intersection.  
The product of their orders is $2^{n-1} = \vphi(2^n) = \#(\Z/2^n\Z)^*$,
so the claim follows. 
\end{enumerate}

\item\exref{ch:dirichlet}{ex:orderalg}.
Write $n=\prod p_i^{e_i}$.
The order of $g$ divides $n$, so 
the condition implies that $p_i^{e_i}$ divides the order of $g$
for each $i$.   Thus the order of $g$ is divisible by the
least common multiple of the $p_i^{e_i}$, i.e., by $n$.


\item\exref{ch:dirichlet}{ex:dlogadd}.
\begin{enumerate}
\item 
The bijection given by
$1+p^{n-1}a\pmod{p^n} \mapsto a\pmod{p}$
 is a homomorphism
since $$\hspace{6em}(1+p^{n-1}a)(1+p^{n-1}b) \con 1 + p^{n-1}(a+b) \pmod{p^n}.$$
\item We have an exact sequence
$$
 \qquad\qquad  1 \to 1 + p \Z/p^{n}\Z \to (\Z/p^n\Z)^* \to (\Z/p\Z)^*\to 1,
$$
so it suffices to solve the discrete log problem in the kernel
and cokernel. We prove by induction on $n$ that we can solve
the discrete log problem
in the kernel easily (compared to known methods for solving
the discrete log problem in $(\Z/p\Z)^*$).
We have an exact sequence
$$
\qquad\qquad  1 \to 1 + p^{n-1} \Z/p^{n}\Z \to (\Z/p^n\Z)^*  \to 
             (\Z/p^{n-1}\Z)^* \to 1.$$
The
first part of this problem shows that we can solve the
discrete log problem in the kernel, and by induction
we can solve it in the cokernel.  This completes
the proof.
\end{enumerate}

\item \exref{ch:dirichlet}{ex:cond2}.
If $\eps(5)=1$, then since $\eps$ is 
nontrivial, \exref{ch:dirichlet}{ex:cyclic} implies
that $\eps$ factors through $(\Z/4\Z)^*$, hence has
conductor $4 = 2^{1 + 1}$, as claimed. 
If $\eps(5)\neq 1$, then again from \exref{ch:dirichlet}{ex:cyclic}
we see that if $\eps$ has order~$r$, then~$\eps$ 
factors through $(\Z/2^{r+2}\Z)^*$ but nothing smaller. 

\item \exref{ch:dirichlet}{ex:charmod5}.
\begin{enumerate}
\item Take $f=x^2+2$.
\item The element $2$ has order $4$.
\item A minimal generator for $(\Z/25\Z)^*$ is $2$,
and the characters are $[1]$, $[2]$, $[3]$, $[4]$.
\item Each of the four Galois orbits has size $1$.
\end{enumerate}


\end{enumerate}

\section{Chapter~\ref{ch:eisen}}

\begin{enumerate}
\item \exref{ch:eisen}{ex:diag}.
The eigenspace $E_{\lambda}$ of $A$ with eigenvalue $\lambda$ is
preserved by~$B$, since if $v\in E_{\lambda}$, then 
$$ABv = BAv = B(\lambda v) = \lambda B v.$$
Because $B$ is diagonalizable, its minimal polynomial
equals its characteristic polynomial; hence the
same is true for the restriction of $B$ to $E_{\lambda}$,
i.e., the restriction of $B$ is diagonalizable.  
Choose basis for all $E_{\lambda}$ so that the
restrictions of $B$ to these eigenspaces
is diagonal with respect to these bases. Then
the concatenation of these bases is a basis that
simultaneously diagonalizes $A$ and $B$.

\item \exref{ch:eisen}{ex:bern_triv}.
When $\eps$ is the trivial character, the
$B_{k,\eps}$ are defined by
$$
  \qquad \sum_{a=1}^{1} \frac{\eps(a) x e^{ax}}{e^{x}-1}
   = \frac{x e^x}{e^x - 1} = x + \frac{x}{e^x - 1}
   = \sum_{k=0}^{\infty} B_{k,\eps} \frac{x^k}{k!}.
$$
Thus $B_{1,\eps} = 1 + B_1 = \frac{1}{2}$,
and for $k>1$, we have $B_{k,\eps} = B_k$.

\item \exref{ch:eisen}{ex:gbparity}. Omitted.

\item \exref{ch:eisen}{ex:dime3}.
The  Eisenstein series in our
basis for $E_3(\Gamma_1(13))$ are of 
the form $E_{3,1,\eps}$ or $E_{3,\eps,1}$
with $\eps(-1)=(-1)^3=-1$.
There are six characters~$\eps$
with modulus $13$ such that  $\eps(-1)=-1$,
and we have the two series $E_{3,1,\eps}$ and
$E_{3,\eps,1}$ associated to each of these.
This gives a dimension of $12$.


\end{enumerate}

\section{Chapter~\ref{ch:dim}}
\begin{enumerate}
\item \exref{ch:dim}{ex:mu0n}.
\begin{enumerate}
\item By Proposition~\ref{prop:p1bij}, we have 
$[\SL_2(\Z):\Gamma_0(N)] = \#\P^1(\Z/N\Z)$.  
By the Chinese Remainder Theorem, 
$$\#\P^1(\Z/N\Z) = \prod_{p\mid N} \#\P^1(\Z/p^{\ord_p(N)}\Z).$$
So we are reduced to computing $\#\P^1(\Z/p^{\ord_p(N)}\Z)$.
We have $(a,b)\in (\Z/p^n\Z)^2$ with $\gcd(a,b,p)=1$
if and only if $(a,b)\not\in (p\Z/p^n\Z)^2$, so
there are $p^{2n} - p^{2(n-1)}$ such pairs. 
The unit group $(\Z/p^n\Z)^*$ has order $\vphi(p^n) = p^n - p^{n-1}$.
It follows that 
$$
  \#\P^1(\Z/N\Z) = \frac{p^{2n} - p^{2(n-1)}}{p^n - p^{n-1}}
   = p^{n} + p^{n-1}.
$$
\item Omitted.
\end{enumerate}

\item \exref{ch:dim}{ex:formsl2z}.  Omitted.

\item \exref{ch:dim}{ex:newmatch}. Omitted.

\item \exref{ch:dim}{ex:suma4}. Omitted.

\item \exref{ch:dim}{ex:sagedim}. See the source
code to \sage.  

\end{enumerate}


\section{Chapter~\ref{ch:linalg}}
\begin{enumerate}
\item \exref{ch:linalg}{ex:matrixwithkernel}.
Take a basis of $W$ and let $G$ be the matrix
 whose rows are these basis elements. Let
 $B$ be the row echelon form of $G$. 
After a permutation $p$ of columns, we may
 write $B = p_i(I | C)$, where $I$ is the identity matrix.
 The matrix $A = p^{-1}(-C^t | I)$, where $I$ is a different
 sized identity matrix, has the property that
 $W = \Ker(A)$.

\item \exref{ch:linalg}{ex:rrefmod}.
The answer is no.  For example if $A=nI$ is $n$ times the
identity matrix and if $p \mid n$, then $\rref(A\pmod{p})=0$
but $\rref(A)\pmod{p}=I$.

\item \exref{ch:linalg}{ex:rrefsmodp}.
Let $T=\prod E_i$ be an invertible matrix such that $TA = E$ is
in (reduced) echelon form and the $E_i$ are elementary matrices,
i.e., the result of applying an elementary row operation to the
identity matrix.  If $p$ is a prime
that does not divide any of the nonzero numerators
or denominators of the entries of $A$ and any $E_i$, then 
$\rref(A \pmod{p}) = \rref(A) \pmod{p}$.  This is
because $E\pmod{p}$ is in echelon form and $A\pmod{p}$
can be transformed to $E\pmod{p}$ via a series of
elementary row operations modulo~$p$.

\item \exref{ch:linalg}{ex:linalg1to9}.
\begin{enumerate}
\item The echelon form (over $\Q$) is 
$$\left(\begin{array}{rrr}
1&0&-1\\
0&1&2\\
0&0&0
\end{array}\right).$$
\item The kernel is the $1$-dimensional span of 
$\left(1,-2,1\right)$.
\item The characteristic polynomial is $x \cdot (x^{2} - 15x - 18)$.
\end{enumerate}

%\item \exref{ch:linalg}{ex:transmat}.

\item \exref{ch:linalg}{ex:hnf}.
\begin{enumerate}
\item The answer is given in the problem.
\item See \cite[\S2.4]{cohen:course_ant}. 
\end{enumerate}

\end{enumerate}

\section{Chapter~\ref{ch:modsym}}
\begin{enumerate}
\item \exref{ch:modsym}{ex:surjredmod}.  
Using the Chinese Remainder Theorem we immediately reduce
to proving the statement when both $M=p^r$ and $N=p^s$ are powers
of a prime $p$. Then $[a]\in(\Z/p^s\Z)^*$ is represented by an
integer $a$ with $\gcd(a,p)=1$.  That same integer $a$
defines an element of $(\Z/p^r\Z)^*$ that reduces modulo
$p^s$ to $[a]$.



\item \exref{ch:modsym}{ex:surjred}. See \cite[Lemma 1.38]{shimura:intro}.


\item \exref{ch:modsym}{ex:m3g13}.
Coset representatives for $\Gamma_1(3)$ are in bijection with $(c,d)$ where 
$c,d \in \Z/3\Z$ and gcd$(c,d,N)=1$, so the following are representatives:
$$ 
\abcd{1 }{0}{0}{1},\,
\abcd{2 }{0}{0}{2}, \,
\abcd{0 }{2}{1}{0},\,
\abcd{1 }{0}{1}{1},\,
\abcd{2 }{0}{1}{2},\,
\abcd{1 }{1}{2}{0},\,
\abcd{1 }{0}{2}{1},\,
\abcd{2 }{0}{2}{2},
$$
which we call $r_1, \ldots, r_8$, respectively.
Now our Manin symbols are of the form $[X,r_i]$ 
and $[Y, r_i]$ for $1 \le i \le 8$ 
modulo the relations 
$$x+x\sigma = 0, \quad x+x\tau+x\tau^2 = 0, \quad\text{ and } x-xJ = 0.$$
First, note that $J$ acts trivially on Manin symbols of odd weight because 
it sends $X$ to $-X$, $Y$ to $-Y$ and $r_i$ to $-r_i$, so $$[z,g]J = [-z,-g] = [z,g].$$ Thus the last relation is trivially true. 

Now $\sigma^{-1}X = -Y$ and $\sigma^{-1}Y = X$. 
Also $\tau^{-1}X = -Y, \tau^{-1}Y = X-Y, \tau^{-2}X = -X+Y$ 
and $\tau^{-2}Y = -X$. 

The first relation on the first symbol says that
$$[X, r_1] = -[-Y, r_3] = [Y, r_3]$$
and the second relation tells us that
$$[X, r_1] + [-Y, r_5] + [-X+Y, r_6] = 0.$$



\item \exref{ch:modsym}{ex:pairingwd}.
  Let $f \in S_k(\Gamma)$ and $g \in \Gamma$. All that remains to be
  shown is that this pairing respects the relation $x=xg$ for all
  modular symbols $x$.  By linearity it suffices to show the invariance of
  $\langle f, X^{k-i-2}Y^i\{\alpha, \beta\} \rangle$.  We have
\begin{align*}
\qquad&\left\langle f, (X^{k-2-i}Y^i \{\alpha, \beta\})g^{-1} \right\rangle\\
& =  \left \langle f, (aX+bY)^{k-i-2}(cX+dY)^i \{g^{-1}(\alpha), g^{-1}(\beta)\} \right \rangle \\
& =  \int_{g^{-1}(\alpha)}^{g^{-1}(\beta)} f(z) (az+b)^{k-i-2}(cz+d)^i \; dz \\
& =  \int_{g^{-1}(\alpha)}^{g^{-1}(\beta)} f(z) \frac{(az+b)^{k-i-2}}{(cz+d)^{k-i-2}}(cz+d)^{k-2} \; dz \\
& =  \int_{g^{-1}(\alpha)}^{g^{-1}(\beta)} f(z) \, g(z)^{k-i-2} (cz+d)^{k-2} \; dz \\
& =  \int_{\alpha}^{\beta} f(g^{-1}(z)) \, g(g^{-1}(z))^{k-i-2} (cg^{-1}(z)+d)^{k-2} \; d(g^{-1}(z)) \\
& =  \int_{\alpha}^{\beta} f(g^{-1}(z)) \, z^{k-i-2} (cg^{-1}(z)+d)^{k-2} \; (cg^{-1}(z)+d)^2 \; dz \\
& =  \int_{\alpha}^{\beta} f(z) \, z^{k-i-2} \; dz \\
& =  \left \langle f, X^{k-i-2}Y^i \{ \alpha , \beta \} \right \rangle,
\end{align*}
where the second to last simplification is due to invariance under
$[g]_k$, i.e.,
$$\qquad f(g^{-1}(z)) = f^{[g]_k}(g^{-1}(z)) = (cg^{-1}(z)+d)^{-k} f(g(g^{-1}(z))).$$
(The proof for $f \in \overline{S}_k(\Gamma)$ works in
 exactly the same way.) 

\item \exref{ch:modsym}{ex:etaconj}.
\begin{enumerate}
\item Let $\eta = {\scriptsize \left(\begin{array}{rr}-1 & 0 \\ 0 & 1
      \end{array}\right)} $. For any $\gamma = {\scriptsize
    \left(\begin{array}{cc}a & b \\c & d \end{array}\right)}$ we have
$$\qquad\qquad \gamma \eta =  {\scriptsize  \left(\begin{array}{rr}-a & b \\ -c & d \end{array}\right)},~~~~
\eta \gamma = {\scriptsize \left(\begin{array}{rr}-a & -b \\ c & d
    \end{array}\right)},~~~~ \mbox{and}~~~~~ \eta \gamma \eta =
{\scriptsize \left(\begin{array}{rr} a & -b \\ -c & d
    \end{array}\right)}.
 $$
 First, if $\gamma \in SL_2(\Z)$, then $\eta \gamma \eta \in GL_2(\Z)$
 and $$\det (\eta \gamma \eta) = \det \eta \det \gamma \det \eta =
 (-1)(1)(-1) = 1$$ so $\eta \gamma \eta \in SL_2(\Z)$. As $\eta^2 = 1$,
 conjugation by $\eta$ is self-inverse, so it must be a bijection.

 Now if $\gamma \in \Gamma_0(N)$, then $c \equiv 0$ (mod $N$), so $-c
 \equiv 0$ (mod $N$), and so $\eta \gamma \eta \in \Gamma_0(N)$. Thus
 $\eta\Gamma_0(N)\eta = \Gamma_0(N)$.

 If $\gamma \in \Gamma_1(N)$, then $-c \equiv 0$ (mod $N$) as before
 and also $a \equiv d \equiv 1$ (mod $N$), so $\eta \gamma \eta \in
 \Gamma_1(N)$. Thus $\eta\Gamma_1(N)\eta = \Gamma_1(N)$.

\item Omitted.

\end{enumerate}


\end{enumerate}


\section{Chapter~\ref{ch:modformcomp}}
\begin{enumerate}
\item \exref{ch:modformcomp}{ex:g0g1quo}.
Consider the surjective homomorphism $$r:\SL_2(\Z)\to \SL_2(\Z/N\Z).$$
Notice that $\Gamma_1(N)$ is the exact inverse image of the subgroup $H$ of
matrices of the form $\abcd{1}{*}{0}{1}$ and $\Gamma_0(N)$
is the inverse image of the subgroup~$T$ of upper triangular matrices.
It thus suffices to observe that $H$ is normal in $T$, which is clear.
Finally, the quotient $T/H$ is isomorphic to the group of diagonal
matrices in $\SL_2(\Z/N\Z)^*$, which is isomorphic to $(\Z/N\Z)^*$.

\item \exref{ch:modformcomp}{ex:dbdinhecke}.
It is enough to show $\dbd{p}\in\Z[\ldots,T_n,\ldots]$ for
primes~$p$, since each $\dbd{d}$ can be written in terms of the
$\dbd{p}$. Since $p\nmid N$, we have that
\[
  T_{p^2} = T_p^2 - \dbd{p}p^{k-1},
\]
so $$\dbd{p}p^{k-1} = T_p^2 - T_{p^2}.$$
By Dirichlet's theorem on primes in arithmetic progression, 
there is a prime~$q\neq p$ congruent
to~$p$ mod $N$. Since $p^{k-1}$ and $q^{k-1}$ are relatively
prime, there exist integers~$a$ and~$b$ such that $a p^{k-1} + b
q^{k-1} = 1$. Then
$$
\qquad\dbd{p}=\dbd{p}(a p^{k-1} + b q^{k-1})
       = a({T_p}^2-T_{p^2}) + b({T_q}^2-T_{q^2})
       \in \Z[\ldots, T_n,\ldots].
$$

\item \exref{ch:modformcomp}{exer:onlyold}.
Take $N=33$. The space $S_2(\Gamma_0(33))$ is a direct
sum of the two old subspaces coming from $S_2(\Gamma_0(11))$
and the new subspace, which has dimension $1$.
If $f$ is a basis for $S_2(\Gamma_0(11))$ and
$g$ is a basis for $S_2(\Gamma_0(33))_{\new}$,
then $\alpha_1(f), \alpha_3(f), g$ is a basis
for $S_2(\Gamma_0(33))$ on which all Hecke
operators $T_n$, with $\gcd(n,33)=1$, have
diagonal matrix.
However, the operator
$T_3$ on $S_2(\Gamma_0(33))$ does
not act as a scalar on $\alpha_1(f)$,
so it cannot be in the ring generated
by all operators $T_n$ with $\gcd(n,33)=1$.

\item \exref{ch:modformcomp}{ex:31}. Omitted.

\end{enumerate}

\section{Chapter~\ref{ch:periods}}
\begin{enumerate}
\item \exref{ch:periods}{ex:intexp}. Hint: Use either repeated
integration by parts or a change of variables that relates the integral
to the $\Gamma$ function.

\item \exref{ch:periods}{ex:funceqn}. See \cite[\S2.8]{cremona:algs}.


\item \exref{ch:periods}{ex:qvs}.
\begin{enumerate}
\item Let $f = \sqrt{-1} \sum q^n$.  Then $d=2$, but the nontrivial 
conjugate of $f$ is $-f$, so $V_f$ has dimension $1$.
\item Choose $\alpha \in K$ such that $K=\Q(\alpha)$.
 Write 
\begin{equation}\label{eqn:fidag}
\qquad\qquad f=\sum_{i=0}^{d-1} \alpha^i g_i
\end{equation}
with $g_i \in \Q[[q]]$.
Let $W_g$ be the $\Q$-span of the $g_i$, and let
$W_f = V_f \cap \Q[[q]]$. 
 By considering the $\Gal(\Qbar/\Q)$ conjugates of (\ref{eqn:fidag}),
we see that the Galois conjugates of $f$ are in the $\C$-span 
of the $g_i$, so
\begin{equation}\label{eqn:dcvfqw}
\qquad\qquad d = \dim_\C V_f \leq \dim_\Q W_g.
\end{equation}
Likewise, taking the above modulo $O(q^n)$ for any~$n$, 
we obtain a matrix equation 
$$
\qquad\qquad  F = A G,
$$
where the columns of $F$ are the $\Gal(\Qbar/\Q)$-conjugates of
$f$, the matrix $A$ is the Vandermonde matrix corresponding to
$\alpha$ (and its $\Gal(\Qbar/\Q)$ conjugates),  and~$G$ has
columns~$g_i$.  Since~$A$ is a Vandermonde matrix, it is 
invertible, so $A^{-1}F = G$.  Taking the limit as $n$ goes
to infinity, we see that each $g_i$ is a linear combination
of the $f_i$, hence an element of $V_f$.   
Thus $W_g \subset W_f$, so (\ref{eqn:dcvfqw})
implies that $\dim_\Q W_f \geq d$.  But $W_f \tensor_\Q \C \subset V_f$
so finally
$$
\qquad\qquad  d\leq \dim_\Q W_f = \dim_\C (W_f \tensor_\Q \C) \leq \dim_\C V_f = d.
$$

\end{enumerate}

\item \exref{ch:periods}{ex:11}. See the appendix to Chapter II in
  \cite{cremona:algs}, where this example is worked out in complete
  detail.  

\end{enumerate}

%%% Local Variables: 
%%% mode: latex
%%% TeX-master: t
%%% End: 

\appendix

\newcommand{\G}{{\mathscr G}}
\newcommand{\V}{{\mathscr V}}
\newcommand{\bG}{\mathbf{G}}
\newcommand{\bH}{\mathbf{H}}
\newcommand{\bR}{\mathbf{R}}
\newcommand{\bC}{\mathbf{C}}
\newcommand{\bP}{\mathbf{P}}
\newcommand{\Proj}{{\mathbb P}}
\renewcommand{\P}{{\mathbb P}}
\newcommand{\PPP}{\mathscr P}
\newcommand{\MMM}{\mathscr M}
\newcommand{\WWW}{\mathscr W}
\newcommand{\sA}{\mathscr A}
\newcommand{\sC}{\mathscr C}
\newcommand{\sO}{\mathscr O}
\newcommand{\sV}{\mathscr V}
\renewcommand{\sS}{\mathscr S}
\newcommand{\fg}{\mathfrak g}
\newcommand{\aaa}{\mathbf{a}}
\renewcommand{\H}{\mathscr H}
\newcommand{\fh}{\mathfrak h}
\renewcommand{\emptyset}{\varnothing}
\newcommand{\cusp}{\text{cusp}}
\newcommand{\bv}{\mathbf{v}}
\newcommand{\bu}{\mathbf{u}}
\newcommand{\borser}[1]{#1^{\text{BS}}}
\newcommand{\closure}[1]{(#1)^{-}}
\renewcommand{\norm}[1]{\|#1\|}
\newcommand{\Vor}{Vorono\v{\i}\xspace}
\newcommand{\vv}{{\mathbf {v}}}
\renewcommand{\aa}{{\mathbf {a}}}
\newcommand{\bb}{{\mathbf {b}}}
\newcommand{\ww}{{\mathbf {w}}}
\newcommand{\xx}{{\mathbf {x}}}
\newcommand{\uu}{{\mathbf {u}}}
\newcommand{\yy}{{\mathbf {y}}}
\newcommand{\qq}{{\mathbf {q}}}
\newcommand{\sets}[1]{[\![#1]\!]}
\newcommand{\coinv}[1]{(S_{#1})_{\Gamma }}
\renewcommand{\riso}{\xrightarrow{\sim}}


\chapter{Computing in Higher Rank}\label{ch:appendix}
\noindent{\Large\bf by Paul E. Gunnells}

%%%%%%%%%%%%%%%%%%
%                %
%  introduction  %
%                %
%%%%%%%%%%%%%%%%%%
\section{Introduction}\label{introduction}
%
%\subsection*{}
This book has addressed the theoretical and practical problems of
performing computations with modular forms.  Modular forms are the
simplest examples of the general theory of automorphic forms attached
to a reductive algebraic group $G$ with an arithmetic
subgroup~$\Gamma$; they are the case $G=\SL_{2} (\R)$ with $\Gamma$ a
congruence subgroup of $\SL_{2} (\Z)$.  For such pairs $(G,\Gamma)$
the Langlands philosophy asserts that there should be deep connections
between automorphic forms and arithmetic, connections that are
revealed through the action of the Hecke operators on spaces of
automorphic forms.  There have been many profound advances in recent
years in our understanding of these phenomena, for example:
\begin{itemize}
\item the establishment of the modularity of
elliptic curves defined over $\Q$ \cite{wiles:fermat, taylor-wiles:fermat, diamond,
cdt, breuil-conrad-diamond-taylor}, 
\item the proof by Harris--Taylor of the local Langlands
correspondence \cite{harris.taylor}, and  
\item Lafforgue's proof of the
global Langlands correspondence for function fields \cite{lafforgue}.
\end{itemize}
Nevertheless, we are still far from seeing that the links between
automorphic forms and arithmetic hold in the broad scope in which 
they are generally believed.  Hence one has the natural problem of
studying spaces of automorphic forms computationally.

%\subsection*{}
%
The goal of this appendix is to describe some computational techniques
for automorphic forms.  We focus on the case $G= \SL_{n} (\R)$ and
$\Gamma \subset \SL_{n} (\Z)$, since the automorphic forms that arise
are one natural generalization of modular forms, and since this is the
setting for which we have the most tools available.  In fact, we do not
work directly with automorphic forms, but rather with the cohomology
of the arithmetic group $\Gamma$ with certain coefficient modules.
This is the most natural generalization of the tools developed in
previous chapters.

Here is a brief overview of the contents.  Section \ref{aut.forms}
gives background on automorphic forms and the cohomology of arithmetic
groups and explains why the two are related.  In Section
\ref{comb.models} we describe the basic topological tools used to
compute the cohomology of $\Gamma$ explicitly.  Section
\ref{mod.symb.section} defines the Hecke operators, describes the
generalization of the modular symbols from Chapter \ref{ch:modsym} to
higher rank, and explains how to compute the action of the Hecke
operators on the top degree cohomology group.  Section
\ref{other.coho} discusses computation of the Hecke action on
cohomology groups below the top degree.  Finally, Section \ref{apps}
briefly discusses some related material and presents some open
problems.

\subsection{}
%
The theory of automorphic forms is notorious for the difficulty of its
prerequisites.  Even if one is only interested in the cohomology of
arithmetic groups---a small part of the full theory---one needs
considerable background in algebraic groups, algebraic topology, and
representation theory.  This is somewhat reflected in our
presentation, which falls far short of being self-contained.  Indeed,
a complete account would require a long book of its own.  We have
chosen to sketch the foundational material and to provide many
pointers to the literature; good general references are
\cite{borel.wallach, harder.notes, schwermer.article, vogan}.  We hope
that the energetic reader will follow the references and fill many
gaps on his/her own.

The choice of topics presented here is heavily influenced (as usual)
by the author's interests and expertise.  There are many computational
topics in the cohomology of arithmetic groups we have completely
omitted, including the trace formula in its many incarnations
\cite{gross.pollack}, the explicit Jacquet--Langlands correspondence
\cite{dembele,whitehouse}, and moduli space techniques
\cite{faber-van.der.geer, van.der.geer}.  We encourage the reader to investigate
these extremely interesting and useful techniques.

\subsection{Acknowledgements}
I thank Avner Ash, John Cremona, Mark McConnell, and Dan Yasaki for helpful
comments.  I also thank the NSF for support.  

\section{Automorphic Forms and Arithmetic Groups}\label{aut.forms}
%
%\subsection{Modular forms and cohomology}
\subsection{}
%
Let $\Gamma = \Gamma_{0} (N)\subset \SL_{2} (\Z)$ be the usual Hecke
congruence subgroup of matrices upper-triangular mod $N$. Let
$Y_{0} (N)$ be the modular curve $\Gamma \backslash \fh$, and let
$X_{0} (N)$ be its canonical compactification obtained by adjoining
cusps.  For any integer $k\geq 2$, let $S_{k} (N)$ be the space of
weight $k$ holomorphic cuspidal modular forms on $\Gamma$.  According
to Eichler--Shimura \cite[Chapter 8]{shimura:intro}, we have the isomorphism
\begin{equation}\label{eq:es}
H^{1} (X_{0} (N); \C)\riso S_{2} (N)\oplus \overline{S_{2} (N)},
\end{equation} 
where the bar denotes complex conjugation and where the isomorphism
is one of Hecke modules. 

More generally, for any integer $n\geq 0$, let $M_{n}\subset \C [x,y]$
be the subspace of degree $n$ homogeneous polynomials.  The space
$M_{n}$ admits a representation of $\Gamma $ by the ``change of
variables'' map
\begin{equation}\label{eq:rep}
\left(\begin{array}{cc}
a&b\\
c&d
\end{array} \right)\cdot p (x,y) = p (dx-by, -cx+ay).
\end{equation}
This induces a local system $\widetilde{M}_{n}$ on the curve $X_{0}
(N)$.\footnote{The classic references for cohomology with local
systems are \cite[Section~31]{steenrod} and \cite[Ch. V]{eilenberg}.  A more
recent exposition (in the language of \v Cech cohomology and locally
constant sheaves) can be found in \cite[II.13]{bott.tu}.  For an
exposition tailored to our needs, 
see \cite[Section~2.9]{harder.notes}.}  Then the
analogue of \eqref{eq:es} for higher-weight modular forms is the
isomorphism
\begin{equation}\label{eq:eshw}
H^{1} (X_{0} (N); \widetilde{M}_{k-2})\riso S_{k} (N)\oplus
\overline{S_{k} (N)}.  
\end{equation}
Note that \eqref{eq:eshw} reduces to \eqref{eq:es} when $k=2$.

Similar considerations apply if we work with the open curve $Y_{0}
(N)$ instead, except that Eisenstein series also contribute to the
cohomology.  More precisely, let $\Eis_{k} (N)$ be the space of weight
$k$ Eisenstein series on $\Gamma_{0} (N)$.  Then \eqref{eq:eshw}
becomes
\begin{equation}\label{eq:eseis}
H^{1} (Y_{0} (N); \widetilde{M}_{k-2})\riso S_{k} (N)\oplus
\overline{S_{k} (N)}\oplus \Eis_{k} (N).  
\end{equation}
These isomorphisms lie at the heart of the modular symbols method.

%\subsection{Modular forms and functions on $\SL_{2}
%(\R)$}\label{ss:pre.aut.forms}
\subsection{}\label{ss:pre.aut.forms}
%
The first step on the path to general automorphic forms is a
reinterpretation of modular forms in terms of functions on $\SL_{2}
(\R)$.  Let $\Gamma \subset \SL_{2}
(\Z)$ be a congruence subgroup.  A weight $k$ modular form on $\Gamma$
is a holomorphic function $f\colon \fh \rightarrow \C$ satisfying
the transformation property
\[
f ((az+b)/ (cz+d)) = j (\gamma, z)^{k}f (z), \quad \gamma = \left(\begin{array}{cc}
a&b\\
c&d
\end{array} \right)\in \Gamma, \quad z \in \fh.
\]  
Here $j (\gamma , z)$ is the \defn{automorphy factor} $cz+d$.  There
are some additional conditions $f$ must satisfy at the cusps of $\fh
$, but these are not so important for our discussion.

The group $G=\SL_{2} (\R)$ acts transitively on $\fh$, with the
subgroup $K=\SO (2)$ fixing $i$.  Thus $\fh$ can be written as the
quotient $G/K$.  From this, we see that $f$ can be viewed as a function
$G\rightarrow \C$ that is \emph{$K$-invariant on the right} and that
satisfies a certain symmetry condition with respect to the
\emph{$\Gamma$-action on the left}.  Of course not every $f$ with these
properties is a modular form: some extra data is needed to take the
role of holomorphicity and to handle the behavior at the cusps.
Again, this can be ignored right now.

We can turn this interpretation around as follows.  Suppose
$\varphi$ is a function $G\rightarrow \C$ that is \defn{$\Gamma
$-invariant on the left}, that is, $\varphi (\gamma g) = \varphi (g)$
for all $\gamma \in \Gamma$.  Hence $\varphi$ can be thought of as a
function $\varphi \colon \Gamma \backslash G\rightarrow \C$.  We
further suppose that $\varphi$ satisfies a certain symmetry condition
with respect to the \emph{$K$-action on the right.}  In particular,
any matrix $m\in K$ can be written 
\begin{equation}\label{eq:m-theta}
m = \left(\begin{array}{cc}
\cos \theta & -\sin \theta \\
\sin \theta & \cos \theta \\
\end{array} \right), \quad \theta \in \R ,
\end{equation}
with $\theta$ uniquely determined modulo $2\pi$.  Let $\zeta_{m}$ be
the complex number $e^{i\theta}$.  Then the $K$-symmetry we require is 
\[
\varphi (gm) = \zeta_{m}^{-k}\varphi(z), \quad m\in K,
\]
where $k$ is some fixed nonnegative integer.

It turns out that such functions $\varphi$ are very closely related to
modular forms: \emph{any} $f\in S_{k} (\Gamma)$ uniquely determines such a
function $\varphi_{f} \colon \Gamma \backslash G \rightarrow \C$.  The
correspondence is very simple.  Given a weight $k$ modular form $f$, define
\begin{equation}\label{eq:corresp}
\varphi_{f} (g) := f (g\cdot i)j (g,i)^{-k}.
\end{equation}
We claim $\varphi_{f}$ is left $\Gamma$-invariant and satisfies the 
desired $K$-symmetry on the right.
Indeed, since $j$ satisfies the
cocycle property $$j (gh, z) = j (g,h\cdot z)j (h,z),$$ we have 
\[
\varphi_{f} (\gamma g) = f
((\gamma g)\cdot i)j (\gamma g,i)^{-k} = j (\gamma ,g\cdot i)^{k}f (g\cdot
i)j (\gamma ,g\cdot i)^{-k}j (g,i)^{-k} = \varphi_{f} (g).
\]
Moreover, any
$m\in K$ stabilizes $i$.  Hence 
\[
\varphi_{f} (gm) = f ((gm)\cdot i) j
(gm,i)^{-k} = f (g\cdot i) j (m,i)^{-k}j (g,m\cdot i)^{-k}.
\]
From
\eqref{eq:m-theta} we have $j (m,i)^{-k} = (\cos \theta + i\sin
\theta)^{-k} = \zeta_{m}^{-k}$, and thus $\varphi_{f} (gm) =
 \zeta^{-k}_{m}\varphi_{f} (g)$.

Hence in \eqref{eq:corresp} the weight and the automorphy factor
``untwist'' the $\Gamma$-action to make $\varphi_{f}$ left
$\Gamma$-invariant.  The upshot is that we can study modular forms by
studying the spaces of functions that arise through the construction
\eqref{eq:corresp}.

Of course, not every $\varphi \colon \Gamma \backslash G \rightarrow
\C$ will arise as $\varphi_{f}$ for some $f\in S_{K} (\Gamma)$: after
all, $f$ is holomorphic and satisfies rather stringent growth
conditions.  Pinning down all the requirements is somewhat technical
and is (mostly) done in the sequel.

%\subsection{Arithmetic groups}\label{ss:a.gps}
\subsection{}\label{ss:a.gps}
%
Before we define automorphic forms, we need to find the correct
generalizations of our groups $\SL_{2} (\R)$ and $\Gamma_{0} (N)$.
The correct setup is rather technical, but this really reflects the
power of the general theory, which handles so many different
situations (e.g., Maass forms, Hilbert modular forms, Siegel
modular forms, etc.).

Let $G$ be a connected Lie group, and let $K\subset G$ be a maximal
compact subgroup.  We assume that $G$ is the set of real points of a
connected semisimple algebraic group $\bG $ defined over $\Q$.  These
conditions mean the following \cite[\S2.1.1]{plat.rapin}:
\begin{enumerate}
\item The group $\bG$ has the structure of an affine algebraic variety given by an ideal
$I$ in the ring $R = \C [x_{ij}, D^{-1} ]$, where the variables $\{
x_{ij} \mid 1\leq i,j\leq n\}$ should be interpreted as the entries of
an ``indeterminate matrix,'' and $D$ is the polynomial $\det
(x_{ij})$.  Both the group multiplication $\bG \times \bG \rightarrow \bG$
 and inversion $\bG \rightarrow \bG$ are required to be morphisms of
algebraic varieties.

The ring $R$ is the coordinate ring of the algebraic group $\GL_{n}$.
Hence this condition means that $\bG$ can be essentially viewed as a
subgroup of $\GL_{n} (\C)$ defined by polynomial equations in the
matrix entries of the latter.
\item \defn{Defined over $\Q$} means that $I$ is generated by
polynomials with rational coefficients.
\item \defn{Connected} means that $\bG $ is connected as an algebraic variety.
\item \defn{Set of real points} means that $G$ is the set of real
solutions to the equations determined by $I$. We write $G = \bG (\R)$.
\item \defn{Semisimple} means that the maximal connected 
solvable normal subgroup of $\bG$ is trivial.
\end{enumerate}

\begin{example}
The most important example for our purposes is 
the \defn{split
form of $\SL_{n}$}.  For this choice we have 
\[
\text{$G=\SL_{n} (\R)$ and $K=\SO (n)$.}
\]
\end{example}

\begin{example}\label{ex:resofscalars}
Let $F/\Q$ be a number field.  Then there is a $\Q$-group $\bG$ such
that $\bG (\Q) = \SL_{n} (F)$.  The group $\bG$ is constructed as
$\bR_{F/\Q } (\SL_{n})$, where $\bR_{F/\Q }$ denotes the
\defn{restriction of scalars} from $F$ to $\Q$
\cite[\S2.1.2]{plat.rapin}.  For example, if $F$ is totally real, the
group $\bR_{F/\Q} (\SL_{2})$ appears when one studies Hilbert modular
forms.

Let $(r,s)$ be the signature of the field $F$, so that $F\otimes \R
\simeq \R^{r}\times \C^s$.  Then $G = \SL_n (\R)^{r} \times \SL_n
(\C)^{s}$ and $K = \SO (n)^{r}\times \SU(n)^{s}$.
\end{example}


\begin{example}
Another important example is the \defn{split symplectic group}
$\SP_{2n}$.  This is the group that arises when one studies Siegel
modular forms.  The group of real points $\SP_{2n} (\R)$ is the
subgroup of $\SL_{2n} (\R)$ preserving a fixed nondegenerate
alternating bilinear form on $\R^{2n}$.  We have $K = \UN (n)$.
\end{example}

\subsection{}
%
To generalize $\Gamma_{0} (N)$, we need the notion of an
\defn{arithmetic group}.  This is a discrete subgroup $\Gamma$ of the
group of rational points $\bG (\Q)$ that is commensurable with the set
of integral points $\bG (\Z)$.  Here commensurable simply means that
$\Gamma \cap \bG (\Z)$ is a finite index subgroup of both $\Gamma$ and
$\bG (\Z)$; in particular $\bG (\Z)$ itself is an arithmetic group.

\begin{example}\label{ex:cong.subgp}
For the split form of $\SL_{n}$ we have $\bG (\Z) = \SL_{n} (\Z)
\subset \bG (\Q) = \SL_{n} (\Q)$.  A trivial way to obtain other
arithmetic groups is by conjugation: if $g\in \SL_{n} (\Q)$, then
$g\cdot \SL_{n} (\Z) \cdot g^{-1}$ is also arithmetic.

A more interesting collection of examples is given by the congruence
subgroups.  The \defn{principal congruence subgroup} $\Gamma (N)$ is
the group of matrices congruent to the identity modulo $N$ for some
fixed integer $N\geq 1$.  A \defn{congruence subgroup} is a group
containing $\Gamma (N)$ for some $N$.

In higher dimensions there are many candidates to generalize the Hecke
subgroup $\Gamma_{0} (N)$.  For example, one can take the
subgroup of $\SL_{n} (\Z)$ that is upper-triangular mod $N$.  From
a computational perspective, this choice is not so good since its
index in $\SL_{n} (\Z)$ is large.  A better choice, and the one that
usually appears in the literature, is to define $\Gamma_{0} (N)$ to be
the subgroup of $\SL_{n} (\Z)$ with bottom row congruent to $(0,\dotsc
,0,*) \mod N$.
\end{example}

%\subsection{Automorphic forms}\label{ss:aut.forms} We are finally
\subsection{}\label{ss:aut.forms} 
%
We are almost ready to define automorphic forms.  Let $\fg$ be the Lie
algebra of $G$, and let $U (\fg)$ be its universal enveloping algebra
over $\C$.  Geometrically, $\fg$ is just the tangent space at the
identity of the smooth manifold $G$.  The algebra $U (\fg)$ is a
certain complex associative algebra canonically built from $\fg$.  The
usual definition would lead us a bit far afield, so we will settle for an
equivalent characterization: $U (\fg)$ can be realized as a certain
subalgebra of the ring of differential operators on $C^{\infty} (G)$,
the space of smooth functions on $G$.

In particular, $G$ acts on $C^{\infty} (G)$ by \defn{left
translations}: given $g\in G$ and $f\in C^{\infty} (G)$, we define
\[
L_{g} (f) (x) := f (g^{-1}x).
\]
Then $U (\fg)$ can be identified with the ring of all differential
operators on $C^{\infty} (G)$ that are invariant under left
translation.  For our purposes the most important part of $U (\fg)$ is
its center $Z (\fg)$.  In
terms of differential operators, $Z (\fg)$ consists of those operators
that are also invariant under \defn{right translation}:
\[
R_{g} (f) (x) := f (xg).
\]

\begin{definition}\label{def:autform}
An \defn{automorphic form} on $G$ with
respect to $\Gamma$ is a function $\varphi \colon G\rightarrow \C$ satisfying
\begin{enumerate}
\item $\varphi (\gamma g) = \varphi (g)$ for all $\gamma \in \Gamma $,
\item the right translates $\{\varphi  (gk)\mid k\in K \}$ span a finite-dimensional
space $\xi$ of functions,\label{ktype}
\item there exists an ideal $J\subset Z (\fg)$ of finite codimension
such that $J$ annihilates $\varphi $, and 
\item $\varphi $ satisfies a certain growth condition that we do not
wish to make precise.  (In the literature, $\varphi $ is said to be
\defn{slowly increasing}.)
\end{enumerate}
\end{definition}

For fixed $\xi$ and $J$, we denote by $\sA(\Gamma ,\xi ,J,K)$ the
space of all functions
satisfying the above four conditions.  It is a basic theorem, due to
Harish-Chandra \cite{harish-chandra}, that $\sA(\Gamma ,\xi ,J,K)$ is
finite-dimensional.

\begin{example}\label{ex:holo.aut}
We can identify the cuspidal modular forms $S_{k} (N)$ in the language
of Definition \ref{def:autform}.  Given a modular form $f$, let
$\varphi_{f}\in C^{\infty} (\SL_{2} (\R))$ be the function from  
\eqref{eq:corresp}.  Then the map $f\mapsto \varphi_{f}$ identifies
$S_{k} (N)$ with the subspace $\sA_{k} (N)$ of functions $\varphi$ satisfying
\begin{enumerate}
\item $\varphi (\gamma g) = \varphi (g)$ for all $\gamma \in \Gamma_{0} (N)$,
\item $\varphi (gm) = \zeta_{m}^{-k}\varphi (g)$ for all $m\in \SO (2)$,
\item $(\Delta +\lambda_{k})\varphi = 0$, where $\Delta \in Z (\fg)$
is the \defn{Laplace--Beltrami--Casimir operator} and 
\[
\lambda_{k} = \frac{k}{2}\left(\frac{k}{2}-1 \right),
\]
\item $\varphi$ is slowly increasing, and 
\item $\varphi$ is \defn{cuspidal}.
\end{enumerate}

% the casimir operator is
% h^{2}/2 - h + 2ef
% it should be the same .. .. check it

The first four conditions parallel Definition \ref{def:autform}.  Item (1) is the
$\Gamma$-invariance.  Item (2) implies that the right translates of
$\varphi$ by $\SO (2)$ lie in a fixed finite-dimensional
representation of $\SO (2)$.  Item (3) is how holomorphicity appears,
namely that $\varphi$ is killed by a certain differential operator.
Finally, item (4) is the usual growth condition.

The only condition missing from the general
definition is (5), which is an extra constraint placed on $\varphi$ to
ensure that it comes from a cusp form.  This condition can be
expressed by the vanishing of certain integrals (``constant terms'');
for details we refer to \cite{bump.autforms, gelbart}.
\end{example}

\begin{example}
Another important example appears when we set $k=0$ in (2) in Example
\ref{ex:holo.aut} and relax (3) by requiring only that $(\Delta
-\lambda)\varphi =0$ for \emph{some} nonzero $\lambda \in \R$.  Such
automorphic forms cannot possibly arise from modular forms, since
there are no nontrivial cusp forms of weight 0.  However, there are
plenty of solutions to these conditions: they correspond to
\defn{real-analytic} cuspidal modular forms of weight 0 and
are known as \defn{Maass forms}.  Traditionally one writes $\lambda =
(1-s^{2})/4$.  The positivity of $\Delta$ implies that $s\in
(-1,1)$ or is purely imaginary.

Maass forms are highly elusive objects.  Selberg proved that there are
infinitely many linearly independent Maass forms of full level
(i.e., on $\SL_{2} (\Z)$), but to this date no explicit construction of
a single one is known.  (Selberg's argument is indirect and relies on
the trace formula; for an exposition see \cite{sarnak}.)  For higher
levels some explicit examples can be constructed using theta series
attached to indefinite quadratic forms \cite{vigneras}.  Numerically
Maass forms have been well studied; see for example \cite{farmer}.

In general the arithmetic nature of the eigenvalues
$\lambda$ that correspond to Maass forms is unknown, although a famous
conjecture of Selberg states that for congruence subgroups they
satisfy the inequality $\lambda \geq 1/4$ (in other words, only purely
imaginary $s$ appear above).  The truth of this
conjecture would have far-reaching consequences, from analytic number
theory to graph theory \cite{lubotzky}.
\end{example}

\subsection{}\label{ss:cusp.eis}
As Example \ref{ex:holo.aut} indicates, there is a notion of \defn{cuspidal
automorphic form}.  The exact definition is too technical to state
here, but it involves an appropriate generalization of the notion of
constant term familiar from modular forms.

There are also \defn{Eisenstein series} \cite{langlands.e.s,
arthur.e.s}.  Again the complete definition is technical; we only
mention that there are different types of Eisenstein series
corresponding to certain subgroups of $G$.  The Eisenstein series that
are easiest to understand are those built from cusp forms on lower
rank groups.  Very explicit formulas for Eisenstein series on
$\GL_{3}$ can be seen in \cite{bump.gl3}.  For a down-to-earth
exposition of some of the Eisenstein series on $\GL_{n}$, we refer to
\cite{goldfeld}.

The decomposition of $M_{k} (\Gamma_{0} (N))$ into cusp forms and
Eisenstein series also generalizes to a general group $G$, although
the statement is much more complicated.  The result is a theorem of
Langlands \cite{langlands} known as the {\em spectral decomposition
of $L^{2} (\Gamma \backslash G)$}.  A thorough recent presentation of
this can be found in \cite{mw}.

%\subsection{Cohomology of arithmetic groups}\label{ss:coho}
\subsection{}\label{ss:coho}
%
Let $\sA = \sA (\Gamma ,K)$ be the space of all automorphic forms,
where $\xi$ and $J$ range over all possibilities.  The space $\sA$ is
huge, and the arithmetic significance of much of it is unknown.  This
is already apparent for $G=\SL_{2} (\R)$.  The automorphic forms
directly connected with arithmetic are the holomorphic modular forms,
not the Maass forms\footnote{However, Maass forms play a very
important \emph{indirect} role in arithmetic.}. Thus the question
arises: which automorphic forms in $\sA$ are the most natural
generalization of the modular forms?

One answer is provided by the isomorphisms \eqref{eq:es},
\eqref{eq:eshw}, \eqref{eq:eseis}.  These show that modular forms
appear naturally in the cohomology of modular curves.  Hence a
reasonable approach is to generalize the \emph{left} of \eqref{eq:es},
\eqref{eq:eshw}, \eqref{eq:eseis}, and to study the resulting
cohomology groups.  This is the approach we will take.  One drawback
is that it is not obvious that our generalization has anything to do
with automorphic forms, but we will see eventually that it certainly
does.  So we begin by looking for an appropriate generalization of the
modular curve $Y_{0} (N)$.

Let $G$ and $K$ be as in Section~\ref{ss:a.gps}, and let $X$ be the quotient
$G/K$.  This is a global Riemannian symmetric space \cite{helgason}.  One can prove
that $X$ is contractible.  Any arithmetic group $\Gamma\subset G$ acts
on $X$ properly discontinuously.  In particular, if $\Gamma$ is
torsion-free, then the quotient $\Gamma \backslash X$ is a smooth
manifold.

Unlike the modular curves, $\Gamma \backslash X$ will not have a
complex structure in general\footnote{The symmetric spaces that have
a complex structure are known as \defn{bounded domains}, or
\defn{Hermitian symmetric spaces} \cite{helgason}.}; nevertheless,
$\Gamma \backslash X$ is a very nice space.  In particular, if
$\Gamma$ is torsion-free, it is an \defn{Eilenberg--Mac~Lane} space for
$\Gamma$, otherwise known as a $K (\Gamma , 1)$.  This means that the
only nontrivial homotopy group of $\Gamma \backslash X$ is its
fundamental group, which is isomorphic to $\Gamma$, and that the
universal cover of $\Gamma \backslash X$ is contractible.  Hence
$\Gamma \backslash X$ is in some sense a ``topological 
incarnation''\footnote{This apt phrase is due to Vogan \cite{vogan}.}
of~$\Gamma$.

This leads us to the notion of the \defn{group cohomology} $H^{*}
(\Gamma ; \C)$ of $\Gamma$ with trivial complex coefficients.  In the
early days of algebraic topology, this was defined to be the complex
cohomology of an Eilenberg--Mac~Lane space for $\Gamma$
\cite[Introduction, I.4]{brown}:
\begin{equation}\label{eq:coho.iso}
H^{*} (\Gamma ;\C) = H^{*} (\Gamma \backslash X; \C).
\end{equation}
Today there are purely algebraic approaches to $H^{*} (\Gamma ;\C)$
\cite[III.1]{brown}, but for our purposes \eqref{eq:coho.iso} is
exactly what we need.  In fact, the group cohomology $H^{*} (\Gamma ;
\C)$ can be identified with the cohomology of the quotient $\Gamma
\backslash X$ even if $\Gamma$ has torsion, since we are working with
complex coefficients.  The cohomology groups $H^{*} (\Gamma ; \C)$,
where $\Gamma$ is an arithmetic group, are our proposed generalization
for the weight 2 modular forms.

What about higher weights?  For this we must replace the trivial
coefficient module $\C$ with local systems, just as we did in
\eqref{eq:eshw}.  For our purposes it is enough to let $\MMM$ be a
rational finite-dimensional representation of $G$ over the complex
numbers.  Any such $\MMM$ gives a representation of $\Gamma \subset
G$ and thus induces a local system $\widetilde{\MMM}$ on $\Gamma
\backslash X$.  As before, the group cohomology $H^{*} (\Gamma ;
\MMM)$ is the cohomology $H^{*} (\Gamma\backslash X;
\widetilde{\MMM})$.  In \eqref{eq:eshw} we took $\MMM =M_{n}$, the
$n$th symmetric power of the standard representation.  For a general
group $G$ there are many kinds of representations to consider.  In any
case, we contend that the cohomology spaces
\[
H^{*} (\Gamma ; \MMM) =  H^{*}(\Gamma \backslash X; \widetilde{\MMM}) 
\]
are a good generalization of the spaces of modular forms.

\subsection{}\label{ss:franke}
%
It is certainly not obvious that the cohomology groups $H^{*} (\Gamma
; \MMM)$ have \emph{anything} to do with automorphic forms, although
the isomorphisms \eqref{eq:es}, \eqref{eq:eshw}, \eqref{eq:eseis} look
promising.  

The connection is provided by a deep theorem of Franke \cite{franke},
which asserts
that
\begin{enumerate}
\item the cohomology groups $H^{*} (\Gamma ; \MMM)$ can be directly
computed in terms of certain automorphic forms (the automorphic forms
of ``cohomological type,'' also known as those with ``nonvanishing
$(\fg , K)$ cohomology'' \cite{vog.zuck}); and \label{bc}
\item there is a direct sum decomposition 
\begin{equation}\label{franke.decomp}
H^{*} (\Gamma ; \MMM) = H^{*}_{\text{cusp}} (\Gamma ;
\MMM)\oplus\bigoplus_{\{P \}}H^{*}_{\{P \}} (\Gamma ; \MMM),
\end{equation}
where the sum is taken over the set of classes of \defn{associate proper
$\Q$-parabolic subgroups of $G$}. \label{ld}
\end{enumerate}
The precise version of statement \eqref{bc} is known in the literature
as the \defn{Borel conjecture}.  Statement \eqref{ld} parallels
Langlands's spectral decomposition of $L^{2} (\Gamma \backslash G)$.

\begin{example}
For $\Gamma = \Gamma_{0} (N)\subset \SL_{2} (\Z)$, the decomposition
\eqref{franke.decomp} is exactly \eqref{eq:eseis}.  The cusp forms
$S_{k} (N)\oplus \overline{S_{k} (N)}$ correspond to the summand
$H^{1}_{\text{cusp}} (\Gamma ; \MMM)$.  There is one class of proper
$\Q$-parabolic subgroups in $\SL_{2} (\R)$, represented by the Borel subgroup
of upper-triangular matrices.  Hence only one term appears in big
direct sum on the right of \eqref{franke.decomp}, which is the
Eisenstein term
$\Eis_{k}$.
%For general $G$,  \eqref{franke.decomp}
%parallels the Hence for these $\MMM$ the cohomology $H^{*} (\Gamma ;
%\MMM)$ can be completely described in terms of automorphic forms.
\end{example}


The summand $H^{*}_{\text{cusp}} (\Gamma ; \MMM)$ of
\eqref{franke.decomp} is called the \defn{cuspidal cohomology}; this
is the subspace of classes represented by cuspidal automorphic forms.
The remaining summands constitute the \defn{Eisenstein cohomology} of
$\Gamma$ \cite{harder.icm}.  In particular the summand indexed by $\{P \}$ is
constructed using Eisenstein series attached to certain cuspidal
automorphic forms on lower rank groups.  Hence $H^{*}_{\text{cusp}}
(\Gamma ; \MMM)$ is in some sense the most important part of the
cohomology: all the rest can be built systematically from cuspidal
cohomology on lower rank groups\footnote{This is a bit of an
oversimplification, since it is a highly nontrivial problem to decide
when cusp cohomology from lower rank groups appears in $\Gamma$.
However, many results are known; as a selection we mention
\cite{harder.icm, harder.gl2, li.schwermer}}. This leads us to our
basic computational problem:

\begin{problem}
Develop tools to compute explicitly the cohomology spaces $H^{*}
(\Gamma ; \MMM)$ and to identify the cuspidal
subspace $H^{*}_{\text{cusp}} (\Gamma ; {\MMM})$. 
\end{problem}

%%%%%%%%%%%%%
%           %
%  retract  %
%           %
%%%%%%%%%%%%%

\section{Combinatorial Models for Group Cohomolo\-gy}\label{comb.models}
%
\subsection{}
%
In this section, we restrict attention to $G=\SL_{n} (\R)$ and
$\Gamma$, a congruence subgroup of $\SL_{n} (\Z)$.  By the previous
section, we can study the group cohomology $H^{*} (\Gamma ; \MMM)$ by
studying the cohomology $H^{*} (\Gamma \backslash X;
\widetilde{\MMM})$.  The latter spaces can be studied using standard
topological techniques, such as taking the cohomology of complexes
associated to cellular decompositions of $\Gamma \backslash X$.  For
$\SL_{n} (\R ) $, one can construct such decompositions using a
version of explicit reduction theory of real positive-definite
quadratic forms due to \Vor\ \cite{vor1}.  The goal of this section is
to explain how this is done.  We also discuss how the cohomology can
be explicitly studied for congruence subgroups of $\SL_{3} (\Z)$.

%\subsection{A linear model for the symmetric
%space}\label{vcd.section}
\subsection{}\label{vcd.section}
%
Let $V$ be the $\R $-vector space of all symmetric $n\times n$
matrices, and let $C\subset V$ be the subset of positive-definite
matrices.  The space $C$ can be identified with the space of all real
positive-definite quadratic forms in $n$ variables: in coordinates, if $x=
(x_{1},\dotsc ,x_{n})^{t}\in \R^{n}$ (column vector), then the matrix
$A\in C$ induces the 
quadratic form  
\[
x \longmapsto x^{t}Ax,
\]
and it is well known that any positive-definite quadratic form arises
in this way.  The space $C$ is a cone, in that it is preserved by
homotheties: if $x\in C$, then $\lambda x\in C$ for all $\lambda \in
\R_{>0}$.  It is also convex: if $x_{1}, x_{2}\in C$, then $tx_{1}+
(1-t)x_{2}\in C$ for $t\in [0,1]$.  Let $D$ be the quotient of $C$ by
homotheties.

\begin{example}
The case $n=2$ is illustrative.  We can take
coordinates on $V \simeq \R^{3}$ by representing any matrix in $V$ as 
\[
\left(\begin{array}{cc}
x&y\\
y&z
\end{array} \right), \quad x,y,z\in \R.  
\] 
The subset of singular matrices $Q = \{xz-y^{2}=0 \}$ is a quadric
cone in $V$ dividing the complement $V\smallsetminus Q$ into three
connected components.  The component containing the identity matrix is
the cone $C$ of positive-definite matrices. The quotient $D$ can be
identified with an open $2$-disk.
\end{example}

The group $G$ acts on $C$ on the left by
\[
(g,c)\longmapsto gcg^{t}.
\]
This action commutes with that of the homotheties and thus
descends to a $G$-action on $D$.  One can show that $G$ acts transitively
on $D$ and that the stabilizer of the image of the identity matrix is
$K=\SO (n)$.  Hence we may identify $D$ with our symmetric space $X =
\SL_{n} (\R )/\SO (n)$.  We will do this in the sequel, using the
notation $D$ when we want to emphasize the coordinates coming from the
linear structure of $C\subset V$ and using the notation $X$ for the
quotient $G/K$.

We can make the identification $D\simeq X$ more explicit.  If $g\in
\SL_{n} (\R)$, then the map 
\begin{equation}\label{chmap}
g \longmapsto gg^{t}
\end{equation}
takes $g$ to a symmetric positive-definite matrix.  Any coset $gK$ is
taken to the same matrix since $KK^{t} = \Id$.  Thus \eqref{chmap}
identifies $G/K$ with a subset $C_{1}$ of $C$, namely those
positive-definite symmetric matrices with determinant $1$.  It is
easy to see that $C_{1}$ maps diffeomorphically onto $D$.  

The inverse map $C_{1}\rightarrow G/K$ is more complicated. Given a
determinant $1$ positive-definite symmetric matrix $A$, one must find
$g\in \SL_{n} (\R)$ such that $gg^{t}=A$.  Such a representation
always exists, with $g$ determined uniquely up to right multiplication
by an element of $K$.  In computational linear algebra, such a $g$ can
be constructed through \defn{Cholesky decomposition} of $A$.

The group $\SL_{n} (\Z )$ acts on $C$ via the $G$-action and does so
properly discontinuously.  This is the ``unimodular change of
variables'' action on quadratic forms \cite[V.1.1]{serre:arithmetic}.
Under our identification of $D$ with $X$, this is the usual action of
$\SL_{n} (\Z)$ by left translation from Section~\ref{ss:coho}.

\subsection{}\label{ss:vcd}
%
Now consider the group cohomology $H^{*} (\Gamma ; \MMM) = H^{*}
(\Gamma \backslash X; \widetilde{\MMM})$.  The identification $D\simeq
X$ shows that the dimension of $X$ is $n (n+1)/2 - 1$.  Hence $H^{i}
(\Gamma ; \MMM)$ vanishes if $i > n (n+1)/2 - 1$.  Since $\dim X$
grows quadratically in $n$, there are many potentially interesting
cohomology groups to study.

However, it turns out that there is some additional vanishing of the
cohomology for deeper (topological) reasons.  For $n=2$, this is easy
to see.  The quotient $\Gamma \backslash \fh$ is homeomorphic to a
topological surface with punctures, corresponding to the cusps of
$\Gamma$.  Any such surface $S$ can be retracted onto a finite graph
simply by ``stretching'' $S$ along its punctures.  Thus $H^{2} (\Gamma
; \MMM) = 0$, even though $\dim \Gamma \backslash \fh = 2$.

For $\Gamma \subset \SL_{n} (\Z )$, a theorem of Borel--Serre implies
that $H^{i} (\Gamma ;{\MMM})$ vanishes if $i> \dim X -n+1 = n(n-1)/2$
\cite[Theorem 11.4.4]{borel.serre}.  The number $\nu = n (n-1)/2$ is
called the \defn{virtual cohomological dimension} of $\Gamma $ and is
denoted $\vcd \Gamma$.  Thus we only need to consider cohomology in
degrees $i\leq \nu$.

Moreover we know from Section~\ref{ss:franke} that the most interesting part
of the cohomology is the cuspidal cohomology.  In what degrees can it
live?  For $n=2$, there is only one interesting cohomology group
$H^{1} (\Gamma ; \MMM)$, and it contains the cuspidal cohomology.  For
higher dimensions, the situation is quite different: for most $i$, the
subspace $H^{i}_{\text{cusp}} (\Gamma ; \MMM)$ vanishes!  In fact in
the late 1970's Borel, Wallach, and Zuckerman observed that the cuspidal
cohomology can only live in the cohomological degrees lying in an
interval around $(\dim X)/2$ of size linear in $n$.  An explicit
description of this interval is given in \cite[Proposition
3.5]{schwermer}; one can also look at Table \ref{dimtable}, from which
the precise statement is easy to determine.

Another feature of Table \ref{dimtable} deserves to be mentioned.
There are exactly two values of $n$, namely $n=2, 3$, such that
virtual cohomological dimension equals the upper limit of the cuspidal
range.  This will have implications later, when we study the action of
the Hecke operators on the cohomology.

\begin{table}[htb]
\begin{center}
\begin{tabular}{c||c|c|c|c|c|c|c|c}
$n$&2&3&4&5&6&7&8&9\cr
\hline
\hline
$\dim X$&2 & 5 & 9 & 14 & 20 & 27 & 35 & 44 \cr
$\vcd \Gamma $ & 1 & 3 & 6 & 10 & 15 & 21 & 28 & 36\cr
top degree of $H^{*}_{\text{cusp}} $& 1 & 3 & 5 & 8 & 11 & 15 & 19& 24 \cr
bottom degree of $H^{*}_{\text{cusp}} $& 1 & 2 & 4 & 6 & 9 & 12 & 16 & 20 \cr
\end{tabular}
\end{center}
\medskip
\caption{The virtual cohomological dimension and the cuspidal range
for subgroups of $\SL_{n}(\Z)$.\label{dimtable}}
\end{table}

%\subsection{The \Vor\  polyhedron}\label{vor.poly.section}
\subsection{}\label{vor.poly.section}
%
Recall that a point in $\Z ^{n}$ is said to be \defn{primitive} if the
greatest common divisor of its coordinates is $1$.  In particular, a
primitive point is nonzero.  Let $\PPP\subset \Z ^{n}$ be the set of
primitive points.  Any $v\in \PPP$, written as a column vector,
determines a rank-$1$ symmetric matrix $q (v)$ in the closure $\bar C$
via $q ( v) = vv^{t}$.  The \defn{\Vor \ polyhedron} $\Pi $ is defined
to be the closed convex hull in $\bar C$ of the points $q (v)$, as $v$
ranges over $\PPP$.  Note that by construction, $\SL_{n} (\Z )$ acts
on $\Pi $, since $\SL_{n} (\Z)$ preserves the set $\{q (v) \}$ and
acts linearly on $V$.

\begin{example}\label{ex:vor.poly}
Figure \ref{tessfig} represents a crude attempt to show what $\Pi$
looks like for $n=2$.  These images were constructed by computing a
large subset of the points $q (v)$ and taking the convex hull (we took
all points $v\in \PPP$ such that $\Trace q (v) < N$ for some large
integer $N$).  From a distance, the polyhedron $\Pi$ looks almost
indistinguishable from the cone $C$; this is somewhat conveyed by the
right of Figure \ref{tessfig}.  Unfortunately $\Pi$ is not locally
finite, so we really cannot produce an accurate picture.  To get a more
accurate image, the reader should imagine that each vertex meets
infinitely many edges.  On the other hand, $\Pi$ is not hopelessly
complex: each maximal face is a triangle, as the pictures suggest.


\begin{figure}[htb] \centering
\subfigure[\label{fig.origin}]{\includegraphics[scale=0.2]{diagrams/tess1}}
\quad\quad
\subfigure[\label{fig.side}]{\includegraphics[scale=0.2]{diagrams/tess2}}
\caption{\label{tessfig}The polyhedron $\Pi$ for $\SL_{2} (\Z)$.  In~(a) we see $\Pi$ from the origin, in~(b) from
the side.  The small triangle at the right center of~(a) is the
facet with vertices $\{q (e_{1}), q (e_{2}), q (e_{1}+e_{2}) \}$,
where $\{e_{1},e_{2}\}$ is the standard basis of $\Z^{2}$.  In~(b) the $x$-axis runs along the top from left to right, and
the $z$-axis runs down the left side.  The facet from~(a) is the little triangle at the top left corner of~(b).}
\end{figure}
\end{example}


\subsection{}
%
The polyhedron $\Pi$ is quite complicated: it has infinitely many
faces and is not locally finite.  However, one of \Vor 's great
insights is that $\Pi$ is actually not as complicated as it seems.

For any $A\in C$, let $\mu (A)$ be the minimum value attained by $A$
on $\PPP$ and let $M (A)\subset \PPP$ be the set on which $A$ attains
$\mu (A)$.  Note that $\mu (A)> 0 $ and $M (A)$ is finite since $A$ is
positive-definite.  Then $A$ is called \defn{perfect} if it is
recoverable from the knowledge of the pair $(\mu (A), M (A) )$.  In
other words, given $(\mu (A), M (A) )$, we can write a system of
linear equations
\begin{equation}\label{system}
mZm^{t} = \mu (A), \quad m\in M (A),
\end{equation}
where $Z = (z_{ij})$ is a symmetric matrix of variables.  Then $A$ is
perfect if and only if $A$ is the unique solution to the system
\eqref{system}. 

\begin{example}\label{ex:A2}
The quadratic form $Q (x,y) =  x^{2}-xy+y^{2}$ is perfect.  The smallest
nontrivial value
it attains on $\Z^{2}$ is $\mu (Q) = 1$, and it does so on the columns of 
\[
M (Q) =\left(\begin{array}{ccc}
1&0&1\\
0&1&1
\end{array} \right)
\]   
and their negatives.   Letting $\alpha x^{2}+\beta xy + \gamma
y^{2}$ be an undetermined quadratic form and applying the data $(\mu
(Q), M (Q))$, we are led to the system of linear equations 
\[
\alpha =1, \quad \gamma =1, \quad \alpha +\beta +\gamma =1.
\]
From this we recover $Q (x,y)$.
\end{example}

\begin{example}\label{ex:cubicform}
The quadratic form $Q' (x,y) = x^{2}+y^{2}$ is not perfect.  Again the
smallest nontrivial value of $Q'$ on $\Z^{2}$ is $m (Q')= 1$, attained on the
columns of 
\[
M (Q')=\left (\begin{array}{cc}
1&0\\
0&1
\end{array}\right)
\]
and their negatives.  But every member of the one-parameter family of quadratic forms 
\begin{equation}\label{}
x^{2}+\alpha xy+y^{2}, \quad \alpha \in (-1,1)
\end{equation}
has the same
set of minimal vectors, and so $Q'$ cannot be recovered from the
knowledge of $m (Q'), M (Q')$.
\end{example}

\begin{example}\label{ex:An}
Example \ref{ex:A2} generalizes to all $n$.  Define
\begin{equation}\label{anform}
A_{n} (x) := \sum_{i=1}^{n}x_{i}^{2} - \sum_{1\leq i<j\leq n} x_{i}x_{j}.
\end{equation}
Then $A_{n}$ is perfect for all $n$.  We have $\mu
(A_{n})=1$, and $M (A_{n})$ consists of all points of the form
\[
\pm (e_{i}+e_{i+1}+\dotsb + e_{i+k}), \quad 1\leq i\leq n, \quad i\leq i+k\leq n,
\]
where $\{e_{i} \}$ is the standard basis of $\Z^{n}$.  This quadratic
form is closely related to the $A_{n}$ root lattice
\cite{fulton.harris}, which explains its name.  It is one of two
infinite families of perfect forms studied by \Vor (the other is
related to the $D_{n}$ root lattice).
\end{example}

We can now summarize \Vor 's main results:
\begin{enumerate}
\item There are finitely many equivalence classes of perfect forms
modulo the action of $\SL_{n} (\Z)$.  \Vor \ even gave an explicit
algorithm to determine all the perfect forms of a given dimension.
\item \label{item:minvec} The facets of $\Pi$, in other words the
codimension $1$ faces, are in bijection with the rank $n$ perfect
quadratic forms.  Under this correspondence the minimal vectors $M
(A)$ determine a facet $F_{A}$ by taking the convex hull in $\bar C$
of the finite point set $\{q (m)\mid m\in M (A) \}$.  Hence there are
finitely many faces of $\Pi$ modulo $\SL_{n} (\Z)$ and thus finitely
many modulo any finite index subgroup $\Gamma$.
\item Let $\V$ be the set of cones over the faces of $\Pi $.  Then $\V$
is a \defn{fan}, which means (i) if $\sigma \in \V$, then any face of
$\sigma$ is also in $\V$; and (ii) if $\sigma ,\sigma '\in \V$, then
$\sigma \cap \sigma '$ is a common face of each\footnote{Strictly
speaking, \Vor \ actually showed that every codimension 1 cone is
contained in two top-dimensional cones.}.  The fan $\V$ provides a
reduction theory for $C$ in the following sense: any point $x\in C$ is
contained in a unique $\sigma (x) \in \V$, and the set $\{\gamma \in
\SL _{n} (\Z ) \mid \gamma \cdot \sigma (x) = \sigma (x) \}$ is
finite.  \Vor\ also gave an explicit algorithm to determine $\sigma
(x)$ given $x$, the \defn{\Vor \ reduction algorithm}.
\end{enumerate}

The number $N_{\text{perf}}$ of equivalence classes of perfect forms modulo the
action of $\GL_{n}(\Z)$ grows rapidly with
$n$ (Table \ref{vortable}); the complete classification is known
only for $n\leq 8$.  For a list of perfect forms up to $n = 7$, see
\cite{cs}.  For a recent comprehensive treatment of perfect forms,
with many historical remarks, see
\cite{martinet}.

%\begin{table}[htb]
%\begin{center}
%\begin{tabular}{c|ccccccc}
%$n$&2&3&4&5&6&7&8\\
%\hline
%$N$&1&1&2&3&?&33&10916\\
%\end{tabular}
%\end{center}
%\medskip
%\caption{\label{vortable}}
%\end{table}

\begin{table}[htb]
\begin{center}
\begin{tabular}{c|c|c}
Dimension&$N_{\text{perf}}$&Authors\\
\hline
\hline
$2$&$1$&\Vor \cite{vor1}\\
$3$&$1$& \emph{ibid.}\\
$4$&$2$& \emph{ibid.}\\
$5$&$3$& \emph{ibid.}\\
$6$&$7$&Barnes \cite{barnes}\\
$7$&$33$&Jaquet-Chiffelle \cite{jaquet, jaquet-chif}\\
$8$&$10916$&Dutour--Sch\"urmann--Vallentin \cite{dutour}\\
\end{tabular}
\end{center}
\medskip
\caption{The number $N_{\text{perf}}$ of equivalence 
classes of perfect forms.\label{vortable}}
\end{table}


%\subsection{Cellular decompositions and the well-rounded
%retract}\label{wrr.section}
\subsection{}\label{wrr.section}
%
Our goal now is to describe how the \Vor \
fan $\V$ can be used to compute the cohomology $H^{*} (\Gamma ;
\MMM)$.  The idea is to use the cones in $\V$ to chop the quotient $D
$ into pieces.  

For any $\sigma \in \V$, let $\sigma^{\circ}$ be the open cone
obtained by taking the complement in $\sigma$ of its proper faces.
Then after taking the quotient by homotheties, the cones $\{\sigma^{\circ }
\cap C \mid \sigma \in \V \}$ pass to locally closed subsets of $D$.
Let $\sC$ be the set of these images.

Any $c\in \sC$ is a \defn{topological cell}, i.e., it is homeomorphic to
an open ball, since $c$ is homeomorphic to a face of $\Pi$.  Because
$\sC$ comes from the fan $\V$, the cells in $\sC$ have good incidence
properties: the closure in $D$ of any $c\in \sC$ can be written as a
finite disjoint union of elements of $\sC$.  Moreover, $\sC$ is
locally finite: by taking quotients of all the $\sigma^{\circ }$
meeting $C$, we have eliminated the open cones lying in $\bar C$, and
it is these cones that are responsible for the failure of local
finiteness of $\V$.  We summarize these properties by saying that
$\sC$ gives a \defn{cellular decomposition} of $D$.  Clearly $\SL_{n}
(\Z)$ acts on $\sC$, since $\sC$ is constructed using the fan $\V$.
Thus we obtain a cellular decomposition\footnote{If $\Gamma$ has torsion, then
cells in $\sC$ can have nontrivial stabilizers in $\Gamma$, and thus
$\Gamma \backslash \sC$ should be considered as an ``orbifold''
cellular decomposition.} of $\Gamma \backslash D $ for
any torsion-free $\Gamma$.  We call $\sC$ the \defn{\Vor \
decomposition} of $D$.

Some care must be taken in using these cells to perform topological
computations.  The problem is that even though the individual pieces
are homeomorphic to balls and are glued together nicely, the
boundaries of the closures of the pieces are not homeomorphic to
spheres in general.  (If they were, then the \Vor \ decomposition
would give rise to a \defn{regular} cell complex \cite{cooke.finney},
which can be used as a substitute for a simplicial or CW complex in
homology computations.)  Nevertheless, there is a way to remedy this.

Recall that a subspace $A$ of a topological space $B$ is a
\defn{strong deformation retract} if there is a continuous map
$f\colon B\times [0,1]\rightarrow B$ such that $f (b,0)=b$, $f
(b,1)\in A$, and $f (a,t) = a$ for all $a\in A$.  For such pairs
$A\subset B$ we have $H^{*} (A) = H^{*} (B)$.  One can show that there
is a strong deformation retraction from $C$ to itself equivariant
under the actions of both $\SL_n(\Z)$ and the homotheties and that
the image of the retraction modulo homotheties, denoted $W$, is
naturally a locally finite regular cell complex of dimension $\nu $.
Moreover, the cells in $W$ are in bijective, inclusion-reversing
correspondence with the cells in $\sC $.  In particular, if a cell in
$\sC$ has \defn{codimension} $d$, the corresponding cell in $W$ has
\defn{dimension} $d$.  Thus, for example, the vertices of $W$ modulo
$\SL_{n} (\Z)$ are in bijection with the top-dimensional cells in
$\sC$, which are in bijection with equivalence classes of perfect
forms.

In the literature $W$ is called the \defn{well-rounded retract}.  The
subspace $W\subset D \simeq X$ has a beautiful geometric
interpretation.  The quotient
\[
\SL_{n} (\Z)\backslash X = \SL_{n} (\Z)
\backslash \SL_{n} (\R)/ \SO (n)
\]
can be interpreted as the moduli space of lattices in $\R^{n}$ modulo
the equivalence relation of rotation and positive scaling
(cf. \cite{ash.gross}; for $n=2$ one can also see \cite[VII,
Proposition 3]{serre:arithmetic}).  Then $W$ corresponds to those lattices whose
shortest nonzero vectors span $\R^{n}$.  This is the origin of the
name: the shortest vectors of such a lattice are ``more round'' than
those of a generic lattice.

The space $W$ was known classically for $n=2$ and was constructed for
$n\geq 3$ by Lannes and Soul\'e, although Soul\'e only published the
case $n=3$ \cite{soule}.  The construction for all $n$ appears in work
of Ash \cite{Ash80, ash.wrr}, who also generalized $W$ to a much
larger class of groups.  Explicit computations of the cell structure
of $W$ have only been performed up to $n=6$ \cite{gangl}.  Certainly
computing $W$ explicitly for $n=8$ seems very difficult, as Table
\ref{vortable} indicates.

\begin{example}\label{sl2}
Figure \ref{tess} illustrates $\sC$ and $W$ for $\SL_{2} (\Z)$.  As in
Example \ref{ex:vor.poly}, the polyhedron $\Pi$ is 3-dimensional, and
so the \Vor \ fan $\V$ has cones of dimensions $0,1,2,3$.  The
$1$-cones of $\V$, which correspond to the vertices of $\Pi$, pass to
infinitely many points on the boundary $\partial \bar D = \bar D
\smallsetminus D$.  The $3$-cones become triangles in $\bar D$ with
vertices on $\partial \bar D$.  In fact, the identifications $D\simeq
\SL_{2} (\R)/\SO (2)\simeq \fh$ realize $D$ as the Klein model for the
hyperbolic plane, in which geodesics are represented by Euclidean line
segments.  Hence, the images of the $1$-cones of $\V$ are none other
than the usual cusps of $\fh$, and the triangles are the $\SL_{2}
(\Z)$-translates of the ideal triangle with vertices $\{0,1,\infty
\}$.  These triangles form a tessellation of $\fh$ sometimes known as
the \defn{Farey tessellation}. The edges of the \Vor are the $\SL_{2}
(\Z)$-translates of the ideal geodesic between $0$ and $\infty$.
After adjoining cusps and passing to the quotient $X_{0} (N)$, these
edges become the supports of the Manin symbols from Section~\ref{sec:manin}
(cf.~Figure \ref{figure:modsym}).  This example also shows how the
\Vor \ decomposition fails to be a regular cell complex: the
boundaries of the closures of the triangles in $D$ do not contain the
vertices and thus are not homeomorphic to circles.

The virtual cohomological dimension of $\SL_{2} (\Z)$ is 1.  Hence the
well-rounded retract $W$ is a graph (Figures \ref{tess} and
\ref{retract2.fig}).  Note that $W$ is not a manifold.  The vertices
of $W$ are in bijection with the Farey triangles---each vertex lies at
the center of the corresponding triangle---and the edges are in
bijection with the Manin symbols.  Under the map $D\rightarrow \fh$,
the graph $W$ becomes the familiar ``$\PSL_{2}$-tree'' embedded in
$\fh$, with vertices at the order 3 elliptic points (Figure
\ref{retract2.fig}).

\begin{figure}[htb]
\begin{center}
\includegraphics[scale=0.4]{diagrams/retract3}
\end{center}
\caption{\label{tess}The \Vor \ decomposition and the retract in $D$.}
\end{figure}

\begin{figure}[htb]
\begin{center}
\includegraphics[scale=0.4, trim=0 330 0 0]{diagrams/retract2}
\end{center}
\caption{\label{retract2.fig}The \Vor \ decomposition and the retract in $\fh$.}
\end{figure}

\end{example}

%\subsection{Example: $\SL_{3} (\Z)$}\label{sl3}
\subsection{}\label{sl3}

We now discuss the example $\SL_{3} (\Z)$ in some detail.  This
example gives a good feeling for how the general situation compares to
the case $n=2$.

We begin with the \Vor \ fan $\V$.  The cone $C$ is 6-dimensional, and
the quotient $D$ is 5-dimensional.  There is one equivalence class of
perfect forms modulo the action of $\SL_{3} (\Z )$, represented by the
form \eqref{anform}.  Hence there are 12 minimal vectors; six are the
columns of the matrix
\begin{equation}\label{sl3minvecs}
\left(
\begin{array}{cccccc}
1&0&0&1&0&1\\
0&1&0&1&1&1\\
0&0&1&0&1&1
\end{array}
 \right),
\end{equation}
and the remaining six are the negatives of these.  This implies that
the cone $\sigma$ corresponding to this form is 6-dimensional and
simplicial.  The latter implies that the faces of $\sigma$ are the
cones generated by $\{q (v)\mid v\in S \}$, where $S$ ranges over all
subsets of \eqref{sl3minvecs}.  To get the full structure of the fan,
one must determine the $\SL_{3} (\Z)$ orbits of faces, as well as
which faces lie in the boundary $\partial \bar C = \bar C
\smallsetminus C$.  After some pleasant computation, one finds:
\begin{enumerate}
\item There is one equivalence class modulo $\SL_{3} (\Z)$ for each of the
6-, 5-, 2-, and
1-dimensional cones.
\item There are two equivalence classes of the 4-dimensional cones,
represented by the sets of minimal vectors
\[
\left(\begin{array}{cccc}
1&0&0&1\\
0&1&0&1\\
0&0&1&1
\end{array} \right)\quad\text{and}\quad 
\left(\begin{array}{cccc}
1&0&0&1\\
0&1&0&1\\
0&0&1&0
\end{array} \right).\label{3faces}
\]
\item There are two equivalence classes of the 3-dimensional cones,
represented by the sets of minimal vectors
\[
\left(\begin{array}{cccc}
1&0&0\\
0&1&0\\
0&0&1
\end{array} \right)\quad\text{and}\quad 
\left(\begin{array}{cccc}
1&0&1\\
0&1&1\\
0&0&0
\end{array} \right).
\]
The second type of 3-cone lies in $\partial \bar C$ and thus does 
not determine a cell in $\sC$.

\item The 2- and 1-dimensional cones lie entirely in
$\partial \bar C$ and do not determine cells in $\sC$.
\end{enumerate}

After passing from $C$ to $D$, the cones of dimension $k$ determine
cells of dimension $k-1$.
Therefore, modulo the action of $\SL_{3} (\Z)$ there are five types of
cells in the \Vor \ decomposition $\sC$, with dimensions from $5$ to $2$.
We denote these cell types by $c_{5}$, $c_{4}$, $c_{3a}$, $c_{3b}$,
and $c_{2}$.  Here $c_{3a}$ corresponds to the first type of 4-cone in item
\eqref{3faces} above, and $c_{3b}$ to the second.  For a beautiful way
to index the cells of $\sC$ using configurations in projective spaces,
see \cite{M91}.

The virtual cohomological dimension of $\SL_{3} (\Z)$ is 3, which
means that the retract $W$ is a 3-dimensional cell complex.  The
closures of the top-dimensional cells in $W$, which are in bijection
with the \Vor \ cells of type $c_{2}$, are homeomorphic to solid cubes
truncated along two pairs of opposite corners (Figure
\ref{soulecube}).  To compute this, one must see how many \Vor cells
of a given type contain a fixed cell of type $c_{2}$ (since the
inclusions of cells in $W$ are the \defn{opposite} of those in $\sC$).

A table of the incidence relations between the cells of $\sC$ and $W$
is given in Table \ref{incidence}.  To interpret the table, let $m =
m(X,Y)$ be the integer in row $X$ and column $Y$.
\begin{itemize}
\item If $m$ is below the diagonal, then the boundary of a cell
of type $Y$ contains $m$ cells of type $X$.
\item If $m$ is above the diagonal, then a cell of type $Y$
appears in the boundary of $m$ cells of type $X$.
\end{itemize}
For instance, the entry $16$ in row $c_{5}$ and column $c_{2}$ means
that a \Vor cell of type $c_{2}$ meets the boundaries of $16$ cells
of type $c_{5}$.  This is the same as the number of vertices in the
Soul\'e cube (Figure
\ref{soulecube}).  Investigation of the table shows that the triangular
(respectively, hexagonal) faces of the Soul\'e cube correspond to the
\Vor \ cells of type $c_{3a}$ (resp., $c_{3b}$).

Figure \ref{souleschlegel} shows a Schlegel diagram for the Soul\'e
cube.  One vertex is at infinity; this is indicated by the arrows on
three of the edges.  This Soul\'e cube is dual to the \Vor cell $C$ of
type $c_{2}$ with minimal vectors given by the columns of the identity
matrix.  The labels on the $2$-faces are additional minimal vectors
that show which \Vor cells contain $C$.  For example, the central
triangle labelled with $(1,1,1)^{t}$ is dual to the \Vor cell of type
$c_{3a}$ with minimal vectors given by those of $C$ together with
$(1,1,1)^{t}$.  Cells of type $c_{4}$ containing $C$ in their closure
correspond to the edges of the figure; the minimal vectors for a given
edge are those of $C$ together with the two vectors on the $2$-faces
containing the edge.  Similarly, one can read off the minimal vectors
of the top-dimensional \Vor cells containing $C$, which correspond to
the vertices of Figure \ref{souleschlegel}.


\begin{table}[htb]
\begin{center}
\begin{tabular*}{0.75\textwidth}{@{\extracolsep{\fill}}c||c|c|c|c|c}
&$c_{5}$&	$c_{4}$&	$c_{3a}$&	$c_{3b}$&	$c_{2}$\\
\hline
\hline
$c_{5}$&	$\bullet$&	$2$&	$3$&	$6$&	$16$\\
$c_{4}$&	$6$&	$\bullet$&	$3$&    $6$&	$24$\\
$c_{3a}$&	$3$&	$1$&	$\bullet$&	$\bullet$&	$4$\\
$c_{3b}$&	$12$&	$4$&	$\bullet$&	$\bullet$&	$6$\\
$c_{2}$&	$12$&	$8$&	$4$&	$3$&	$\bullet$\\
\end{tabular*}
\end{center}
\medskip
\caption{Incidence relations in the \Vor decomposition and the retract
for $\SL_{3} (\Z)$.\label{incidence}}
\end{table}

\begin{figure}[htb]
\begin{center}
\includegraphics[scale=0.35]{diagrams/soulecube}
\end{center}
\caption{\label{soulecube}The Soul\'e cube.}
\end{figure}

\begin{figure}[htb]
\psfrag{111}{$(1,1,1)^t$}
\psfrag{-111}{$(-1,1,1)^t$}
\psfrag{1-11}{$(1,-1,1)^t$}
\psfrag{11-1}{$(1,1,-1)^t$}
\psfrag{110}{$(1,1,0)^t$}
\psfrag{1-10}{$(1,-1,0)^t$}
\psfrag{011}{$(0,1,1)^t$}
\psfrag{01-1}{$(0,1,-1)^t$}
\psfrag{101}{$(1,0,1)^t$}
\psfrag{10-1}{$(1,0,-1)^t$}
\begin{center}
\includegraphics[scale=0.5]{diagrams/souleschlegel}
\end{center}
\caption{A Schlegel diagram of a Soul\'e cube, showing the minimal
vectors that correspond to the $2$-faces.\label{souleschlegel}}
\end{figure}

\subsection{}\label{ss:cohocompnohecke}
%
Now let $p$ be a prime, and let $\Gamma = \Gamma_{0} (p)\subset
\SL_{3} (\Z)$ be the Hecke subgroup of matrices with bottom row
congruent to $(0,0,*) \mod p$ (Example \ref{ex:cong.subgp}).  The
virtual cohomological dimension of $\Gamma$ is $3$, and the cusp
cohomology with constant coefficients can appear in degrees $2$ and
$3$.  One can show that the cusp cohomology in degree $2$ is dual to
that in degree $3$, so for computational purposes it suffices to focus
on degree $3$.  

In terms of $W$, these will be cochains supported on the 3-cells.
Unfortunately we cannot work directly with the quotient $\Gamma
\backslash W$ since $\Gamma$ has torsion: there will be cells taken to
themselves by the $\Gamma$-action, and thus the cells of $W$ need to
be subdivided to induce the structure of a cell complex on $\Gamma
\backslash W$.  Thus when $\Gamma$ has torsion, the ``set of $3$-cells
modulo $\Gamma$'' unfortunately makes no sense.

To circumvent this problem, one can mimic the idea of Manin symbols.
The quotient $\Gamma \backslash \SL_{3} (\Z)$ is in bijection with the
finite projective plane $\Proj^{2} (\F_{p})$, where $\F_p$ is the
field with $p$ elements (cf. Proposition \ref{prop:p1bij}).
The group $\SL_{3} (\Z)$ acts transitively on the set of all $3$-cells
of $W$; if we fix one such cell $w$, its stabilizer $\Stab (w) =
\{\gamma \in \SL_{3} (\Z) \mid \gamma w = w \}$ is a finite subgroup
of $\SL_{3} (\Z)$.  Hence the set of $3$-cells modulo $\Gamma$ should
be interpreted as the set of orbits in $\Proj^{2} (\F_{p})$ of the
finite group $\Stab (w)$.  This suggests describing $H^{3}(\Gamma ;
\C)$ in terms of the space $\sS $ of complex-valued functions $f\colon
\Proj^{2} (\F_{p})\rightarrow \C$.  To carry this out, there are two problems:
\begin{enumerate}
\item How do we explicitly describe $H^{3}(\Gamma ;
\C)$ in terms of $\sS$?
\item How can we isolate the cuspidal subspace $H^{3}_{\text{cusp}}
(\Gamma ;\C) \subset H^{3} (\Gamma ; \C)$ in terms of our description?
\end{enumerate}
Fully describing the solutions to these problems is rather
complicated.  We content ourselves with presenting the following
theorem, which collects together several statements in \cite{agg}.
This result should be compared to Theorems \ref{thm:manin_present} and
\ref{thm:mansym}.

\begin{theorem}[Theorem 3.19 and Summary 3.23 of \cite{agg}]
\label{thm:agg}
  We have 
\[
\dim H^{3} (\Gamma_{0} (p); \C) = \dim H^{3}_{\text{cusp}} (\Gamma_{0}
(p);\C) + 2S_{p},
\]
where $S_{p}$ is the dimension of the space of weight $2$ holomorphic
cusp forms on $\Gamma_{0} (p)\subset \SL_{2} (\Z)$.  Moreover, 
the cuspidal cohomology $H^{3}_{\text{cusp}} (\Gamma_{0}
(p); \C)$ is isomorphic to the vector space of functions $f\colon \Proj^{2}
(\F_p)\rightarrow \C$ satisfying
\begin{enumerate}
\item $f (x,y,z) = f (z,x,y) = f (-x,y,z) = -f (y,x,z)$,
\item $f (x,y,z) + f (-y,x-y,z) + f (y-x,-x,z) = 0$,
\item $f (x,y,0)=0$, and 
\item $\sum_{z=1}^{p-1} f (x,y,z) = 0$.
\end{enumerate}
\end{theorem}

Unlike subgroups of $\SL_{2} (\Z)$, cuspidal cohomology is
apparently much rarer for $\Gamma_{0} (p)\subset \SL_{3} (\Z)$.  The
computations of \cite{agg, fourdutch} show that the only prime levels
$p\leq 337$ with nonvanishing cusp cohomology are $53$, $61$, $79$, $89$,
and $223$.  In all these examples, the cuspidal subspace is
$2$-dimensional. 

For more details of how to implement such computations, we refer to
\cite{agg, fourdutch}.  For further details about the additional
complications arising for higher rank groups, in particular subgroups
of $\SL_{4} (\Z)$, see \cite[Section~3]{computation}.


%%%%%%%%%%%%%%%%%%%%%%%%%%%%%%%
%                             %
%  hecke and modular symbols  %
%                             %
%%%%%%%%%%%%%%%%%%%%%%%%%%%%%%%

\section{Hecke Operators and Modular Symbols}\label{mod.symb.section}
%
\subsection{}
%
There is one ingredient missing so far in our discussion of the
cohomology of arithmetic groups, namely the Hecke operators.\index{Hecke operator}
These are an essential tool in the study of modular forms.  Indeed,
the forms with the most arithmetic significance are the Hecke
eigenforms, and the connection with arithmetic is revealed by the
Hecke eigenvalues. 

In higher rank the situation is similar.  There is an algebra of
Hecke operators acting on the cohomology spaces $H^{*} (\Gamma ;
\MMM)$. The eigenvalues of these operators are conjecturally related
to certain representations of the Galois group.  Just as in the case
$G=\SL_{2} (\R)$, we need tools to compute the Hecke action.

In this section we discuss this problem.  We begin with a general
description of the Hecke operators and how they act on cohomology.
Then we focus on one particular cohomology group, namely the top
degree $H^{\nu} (\Gamma ;\C)$, where $\nu =\vcd (\Gamma)$ and $\Gamma$
has finite index in $\SL_{n} (\Z)$.  This is the setting that
generalizes the modular symbols method from Chapter \ref{ch:modsym}.
We conclude by giving examples of Hecke eigenclasses in the cuspidal
cohomology of $\Gamma_{0} (p)\subset \SL_{3} (\Z)$.


%\subsection{Hecke operators}\label{hecke.ops}
\subsection{}\label{hecke.ops}
%
Let $g\in \SL_{n} (\Q)$.   The group $\Gamma ' = \Gamma \cap
g^{-1}\Gamma g$ has finite index in both $\Gamma$ and $g^{-1}\Gamma
g$.  The element $g$ determines a diagram $C (g)$
\[
\xymatrix{&{\Gamma'\backslash X}\ar[dl]_{s}\ar[dr]^{t}&\\
 {\Gamma\backslash X}&&{\Gamma\backslash X}}
\]
called a \defn{Hecke correspondence}.  The map $s$ is induced by the
inclusion $\Gamma '\subset \Gamma$, while $t$ is induced by the
inclusion $\Gamma '\subset g^{-1}\Gamma g$ followed by the
diffeomorphism $g^{-1}\Gamma g\backslash X\rightarrow \Gamma
\backslash X$ given by left multiplication by $g$.  Specifically, 
\[
s (\Gamma 'x) = \Gamma x, \quad t (\Gamma 'x) = \Gamma gx, \quad x\in X.
\]

The maps $s$ and $t$ are finite-to-one, since the indices $[\Gamma
':\Gamma]$ and $[\Gamma ':g^{-1}\Gamma g]$ are finite.   This implies
that we obtain maps on cohomology 
\[
s^{*}\colon H^{*} (\Gamma \backslash X) \rightarrow
H^{*} (\Gamma' \backslash X), \quad t_{*}\colon H^{*} (\Gamma' \backslash X) \rightarrow
H^{*} (\Gamma \backslash X).
\]
Here the map $s^{*}$ is the usual induced map on cohomology, while the
``wrong-way'' map\footnote{Under the identification $H^{*} (\Gamma
\backslash X; \widetilde{\MMM})\simeq H^{*} (\Gamma ;\MMM)$, the map
$t_{*}$ becomes the transfer map in group cohomology
\cite[III.9]{brown}.} $t_{*}$ is given by summing a class over the finite
fibers of $t$.  These maps can be composed to give a map
\[
T_{g} := t_{*}s^{*}\colon H^{*} (\Gamma \backslash X;
\widetilde{\MMM})\longrightarrow  H^{*} (\Gamma \backslash X;
\widetilde{\MMM}).
\]
This is called the \defn{Hecke operator} associated to $g$.  There is
an obvious notion of isomorphism of Hecke correspondences.  One can
show that up to isomorphism, the correspondence $C (g)$ and thus the
Hecke operator $T_{g}$ depend only on the double coset $\Gamma
g\Gamma$.  One can compose Hecke correspondences, and thus we obtain
an algebra of operators acting on the cohomology, just as in the
classical case.

\begin{example}
Let $n=2$, and let $\Gamma = \SL_{2} (\Z)$.  If we take $g=\diag
(1,p)$, where $p$ is a prime, then the action of $T_{g}$ on $H^{1}
(\Gamma ; M_{k-2})$ is the same as the action of the classical Hecke
operator $T_{p}$ on the weight $k$ holomorphic modular forms.  If we
take $\Gamma = \Gamma_{0} (N)$, we obtain an operator $T (p)$ for all
$p$ prime to $N$, and the algebra of Hecke operators coincides with
the (semisimple) Hecke algebra generated by the $T_{p}$, $(p,N)=1$.
For $p|N$, one can also describe the $U_{p}$ operators in this language.
\end{example}

\begin{example}\label{ex:heckeops}
Now let $n>2$ and let $\Gamma = \SL_{n} (\Z)$. The picture is very
similar, except that now there are several Hecke operators attached to
any prime $p$.  In fact there are $n-1$ operators $T (p,k)$,
$k=1,\dotsc ,n-1$.  The operator $T (p,k)$ is associated to the
correspondence $C (g)$, where $g = \diag (1,\dotsc ,1,p,\dotsc ,p)$
and where $p$ occurs $k$ times.  If we consider the congruence
subgroups $\Gamma_{0} (N)$, we have operators $T (p,k)$ for
$(p,N)=1$ and analogues of the $U_{p}$ operators for $p|N$.

Just as in the classical case, any double coset $\Gamma g \Gamma$ can
be written as a disjoint union of left cosets 
\[
\Gamma g\Gamma  = \coprod_{h\in \Omega} \Gamma h
\]
for a certain finite set of $n\times n$ integral matrices $\Omega$.
For the operator $T (p,k)$, the set $\Omega$ can be taken to be all
upper-triangular matrices of the form \cite[Proposition 7.2]{krieg}
\[
\left(\begin{array}{cccc}
p^{e_{1}}&&a_{ij}\\
&\ddots&&\\
&&p^{e_{n}}&\\
\end{array} \right),
\]
where
\begin{itemize}
\item $e_{i}\in \{0,1 \}$ and exactly $k$ of the $e_{i}$ are equal to
$1$ and
\item $a_{ij}=0$ unless $e_{i}=0$ and $e_{j}=1$, in which case
$a_{ij}$ satisfies $0\leq a_{ij}<p$.
\end{itemize}


\end{example}

\begin{remark}
The number of coset representatives for the operator $T (p,k)$ is the
same as the number of points in the finite Grassmannian $G (k,n)
(\F_{p})$.  A similar phenomenon is true for the Hecke operators for
any group $G$, although there are some subtleties \cite{gross}.  
\end{remark}



%\subsection{Modular symbols}\label{mod.symbs}
\subsection{}\label{mod.symbs}
%
Recall that in Section~\ref{wrr.section} we constructed the \Vor
decomposition $\sC$ and the well-rounded retract $W$ and that we can
use them to compute the cohomology $H^{*} (\Gamma ;\MMM)$.
Unfortunately, we cannot directly use them to compute the action of the
Hecke operators on cohomology, since the Hecke operators do not act
cellularly on $\sC$ or $W$.  The problem is that the Hecke image of a
cell in $\sC$ (or $W$) is usually not a union of cells in $\sC$ (or
$W$).  This is already apparent for $n=2$.  The edges of $\sC$ are the
$\SL_{2} (\Z)$-translates of the ideal geodesic $\tau$ from $0$ to
$\infty$ (Example \ref{sl2}).  Applying a Hecke operator takes such an
edge to a union of ideal geodesics, each with vertices at a pair of
cusps.  In general such geodesics are not an $\SL_{2} (\Z)$-translate
of~$\tau$.  

For $n=2$, one solution is to work with all possible ideal geodesics
with vertices at the cusps, in other words the space of modular
symbols $\M_{2}$ from Section~\ref{sec:modsym}.  Manin's trick (Proposition
\ref{prop:manintrick}) shows how to write any modular symbol as a
linear combination of unimodular symbols, by which we mean modular
symbols supported on the edges of $\sC$.  These are the ideas we now
generalize to all $n$.

\begin{definition}\label{def:modular.symbols}
Let $S_{0}$ be the $\Q$-vector space
spanned by the symbols $\vv = [v_{1},\dotsc ,v_{n}]$, where $v_{i}\in
\Q^{n}\smallsetminus \{0 \}$, modulo the following relations:
\begin{enumerate}
\item If $\tau $ is a permutation on $n$ letters, then
\[
[v_{1},\ldots,v_{n}] = \sign (\tau ) [\tau (v_{1}),\ldots,\tau
(v_{n})], 
\]
where $\sign (\tau )$ is the sign of $\tau $.
\item  \label{projpts} If $q\in \Q^{\times }$, then 
\[
[q v_{1},v_{2}\ldots,v_{n}] = [v_{1},\ldots,v_{n}].
\]
\item \label{lindep} If the points $v_{1},\dotsc ,v_{n}$ are linearly dependent, then $\vv = 0$.
\end{enumerate}
Let $B\subset S_{0}$ be the subspace generated by linear combinations
of the form 
\begin{equation}\label{fundrel}
\sum _{i=0}^{n}  (-1)^{i}
[v_{0},\ldots,\hat{v_{i}},\ldots,v_{n}],  
\end{equation}
where $v_{0},\dotsc ,v_{n}\in
\Q^{n}\smallsetminus \{0 \}$ and where $\hat{v_{i}}$ means to omit $v_{i}$.
\end{definition}

We call $S_{0}$ the space of \defn{modular symbols}.  We caution the
reader that there are some differences in what we call modular symbols
and those found in Section~\ref{sec:modsym} and Definition
\ref{defn:modsym}; we compare them in 
Section~\ref{sec:comparison}.  The group $\SL_{n} (\Q)$ acts on $S_{0}$ by
left multiplication: $g\cdot \vv = [gv_{1},\dotsc ,gv_{n}]$.  This
action preserves the subspace $B$ and thus induces an action on the
quotient $M = S_{0}/B$.  For $\Gamma \subset \SL_{n} (\Z)$ a finite
index subgroup, let $M_{\Gamma}$ be the space of $\Gamma$-coinvariants
in $M$.  In other words, $M_{\Gamma}$ is the quotient of $M$ by the
subspace generated by $\{ m - \gamma \cdot m \mid \gamma \in
\Gamma\}$.

The relationship between modular symbols and the cohomology of
$\Gamma$ is given by the following theorem, first proved for $\SL_{n}$
by Ash and Rudolph \cite{ash.rudolph} and by Ash for general $G$
\cite{ash.minmod}:

\begin{theorem}[\cite{ash.minmod, ash.rudolph}]\label{thm:modsymbiso}
 Let $\Gamma \subset \SL_{n} (\Z)$ be a
finite index subgroup.  There is an isomorphism
\begin{equation}\label{modsymbiso}
M_{\Gamma}\xrightarrow{\sim} H^{\nu} (\Gamma ; \Q),
%\xymatrix@1{S_{0,\Gamma}\ar[r]^{\sim}& H^{\nu} (\Gamma ; \Q)}
\end{equation}
where $\Gamma$ acts trivially on $\Q$ and where $\nu = \vcd (\Gamma )$.
\end{theorem}

We remark that Theorem \ref{thm:modsymbiso} remains true if $\Q$ is
replaced with nontrivial coefficients as in Section \ref{ss:coho}.
Moreover, if $\Gamma$ is assumed to be torsion-free then we can
replace $\Q$ with $\Z$.


The great virtue of $M_{\Gamma}$ is that it admits an action of the
Hecke operators.  Given a Hecke operator $T_{g}$, write 
the double coset $\Gamma g\Gamma$ as a disjoint union of left cosets
\begin{equation}\label{eq:coset.decomp}
\Gamma g\Gamma  = \coprod_{h\in \Omega} \Gamma h
\end{equation}
as in Example \ref{ex:heckeops}.  
Any class in $M_{\Gamma}$ can be lifted to a representative $\eta =
\sum q (\vv)\vv \in S_0$, where $q (\vv)\in \Q$
and almost all $q (\vv)$ vanish.  
Then we define
\begin{equation}\label{eq:heckeaction}
T_{g} (\vv) = \sum_{h\in \Omega} h\cdot \vv
\end{equation}
and extend to $\eta$ by linearity.  The right side of
\eqref{eq:heckeaction} depends on the choices of $\eta$ and $\Omega$,
but after taking quotients and coinvariants, we obtain a well-defined
action on cohomology via \eqref{modsymbiso}.

\subsection{}\label{sec:comparison}
%
The space $S_{0}$ is closely related to the space $\M_{2}$ from
Section~\ref{sec:modsym} and Section~\ref{sec:modsymgen}.  Indeed, $\M_{2}$ was
defined to be the quotient $(F/R)/ (F/R)_{\text{tor}}$, where $F$ is
the free abelian group generated by ordered pairs
\begin{equation}\label{symbol.def}
\{\alpha ,\beta \}, \quad \alpha ,\beta \in \Proj^{1} (\Q),
\end{equation}
and $R$ is the subgroup generated by elements of the form
\begin{equation}\label{fund.rel}
  \{\alpha,\beta\} + \{\beta,\gamma\} + \{\gamma,\alpha\}, \quad
\alpha ,\beta, \gamma  \in \Proj^{1} (\Q).
\end{equation}
The only new feature in Definition \ref{def:modular.symbols} is item
\eqref{lindep}.  For $n=2$ this corresponds to the condition $\{\alpha,\alpha
\}=0$, which follows from \eqref{fund.rel}.  
We have
\[
S_{0}/B \simeq \M_{2} \otimes \Q.
\]
Hence there are two differences between $S_{0}$ and $\M_{2}$: our
notion of modular symbols uses rational coefficients instead of
integral coefficients and is the space of symbols \emph{before}
dividing out by the subspace of relations $B$; we further caution the
reader that this is somewhat at
odds with the literature. 

We also remark that the general arbitrary weight definition of modular
symbols for a subgroup $\Gamma \subset \SL_{2} (\Z)$ given in
Section~\ref{sec:modsymgen} also includes taking $\Gamma$-coinvariants, as
well as extra data for a coefficient system.  We have not included the
latter data since our emphasis is trivial coefficients, although it
would be easy to do so in the spirit of Section~\ref{sec:modsymgen}.

Elements of $\M_{2}$ also have a geometric interpretation: the symbol
$\{\alpha ,\beta \}$ corresponds to the ideal geodesic in $\fh$ with
endpoints at the cusps $\alpha$ and $\beta$.  We have a similar
picture for the symbols $\vv = [v_{1},\dotsc ,v_{n}]$.  We can assume
that each $v_{i}$ is primitive, which means that each $v_{i}$
determines a vertex of the \Vor polyhedron $\Pi$.  The rational cone
generated by these vertices determines a subset $\Delta (\vv)\subset
D$, where $D$ is the linear model of the symmetric space $X = \SL_{n}
(\R)/\SO (n)$ from Section~\ref{vcd.section}.  This subset $\Delta (\vv)$ is
then an ``ideal simplex'' in $X$.  There is also a connection between
$\Delta (\vv)$ and torus orbits in $X$; we refer to \cite{ash.minmod}
for a related discussion.

%\subsection{The modular symbol  algorithm}
\subsection{}
%
Now we need a generalization of the Manin trick
(Section~\ref{sec:manintrick}).  This is known in the literature as the
\defn{modular symbols algorithm}.

We can define a kind of norm function on $S_{0}$ as follows.  Let $\vv
=[v_{1},\dotsc ,v_{n}]$ be a modular symbol.  For each $v_{i}$, 
choose $\lambda_{i}\in \Q^{\times}$ such that $\lambda_{i}v_{i}$ is
primitive.  Then we define 
\[
\norm{\vv} := |\det (\lambda_{1}v_{1},\dotsc ,\lambda_{n}v_{n})| \in \Z .
\]
Note that $\norm{\vv}$ is well defined, since the $\lambda_{i}$ are
unique up to sign, and permuting the $v_{i}$ only changes the
determinant by a sign.  We extend $\norm{\phantom{a}}$ to all of $S_{0}$
by taking the maximum of $\norm{\phantom{a}}$ over the support of any
$\eta \in S_{0}$: if $\eta =\sum q (\vv)\vv$, where $q (\vv)\in \Q$
and almost all $q (\vv)$ vanish, then we put
\[
\norm{\eta} = \Max_{q (\vv) \not =0} \norm{\vv}.
\]
We say a modular symbol $\eta $ is \defn{unimodular} if $\norm{\eta }
= 1$.  It is clear that the images of the unimodular symbols generate
a finite-dimensional subspace of $M_{\Gamma}$.  The next theorem shows
that this subspace is actually \emph{all} of $M_{\Gamma }$.

\begin{theorem}[\cite{ash.rudolph, barvinok}]
\label{thm:modsymbalg}
The space $M_{\Gamma }$ is spanned by
the images of the unimodular symbols.  More precisely, given any symbol $\vv\in S_{0}$ with
$\norm{\vv}>1$, 
\begin{enumerate}
\item in $S_{0}/B$ we may write 
\begin{equation}\label{relation}
\vv = \sum q (\ww)\ww , \quad q (\ww)\in \Z ,
\end{equation}
where if $q (\ww)\not =0$, then $\norm{\ww} = 1$, and 
\item the number of terms on the right side of \eqref{relation} is bounded
by a polynomial in $\log \|\vv\|$ that depends only on the dimension $n$.
\end{enumerate}
\end{theorem}

\begin{proof}
%
(Sketch)
Given a modular symbol $\vv =[v_{1},\dotsc ,v_{n}]$, we may assume
that the points $v_{i}$ are primitive.  We will show that if
$\norm{\vv}>1$, we can find a point $u$ such that when we apply the
relation \eqref{fundrel} using the points $u,v_{1},\dotsc ,v_{n}$, all
terms other than $\vv$ have norm less than $\norm{\vv}$.  We call such
a point a \defn{reducing point} for $\vv$.

Let $P\subset \R^{n}$ be the open parallelotope
\[
P:=\Bigl\{\sum \lambda _{i}v _{i}\Bigm | |\lambda_{i} |< \norm{\vv }^{-1/n}
\Bigr\}. 
\]
Then $P$ is an $n$-dimensional centrally symmetric convex body with
volume $2^{n}$.  By Minkowski's theorem from the geometry of numbers
(cf.~\cite[IV.2.6]{fr.tay}), $P\cap \Z ^{n}$ contains a nonzero point
$u$.  
Using \eqref{fundrel}, we find 
\begin{equation}\label{move}
\vv = \sum_{i=1}^{n} (-1)^{i-1}\vv_{i} (u),
\end{equation}
where $\vv_{i} (u)$ is the symbol
\[
\vv_{i} (u) = [v_{1},\dotsc ,v_{i-1},u, v_{i+1},\dotsc ,v_{n}].
\]
Moreover, it is easy to see that the new symbols satisfy
\begin{equation}\label{barvest}
0\leq \norm{\vv_{i} (u)}<\norm{\vv}^{(n-1)/n}, \quad i=1,\dotsc ,n.
\end{equation}
This completes the proof of the first statement.

To prove the second statement, we must estimate how many times
relations of the form \eqref{move} need to be applied to obtain
\eqref{relation}.  A nonunimodular symbol produces at most $n$ new
modular symbols after \eqref{move} is performed; we potentially have
to apply \eqref{move} again to each of the symbols that result, which
in turn could produce as many as $n$ new symbols for each.  Hence we
can visualize the process of constructing \eqref{relation} as building
a rooted tree, where the root is $\vv$, the leaves are the symbols
$\ww$, and where each node has at most $n$ children.  It is not hard
to see that the bound \eqref{barvest} implies that the depth of this
tree (i.e., the longest length of a path from the root to a leaf) is
$O(\log \log \norm{\vv})$.  From this the second statement follows
easily.
%
\end{proof}

Statement (1) of Theorem \ref{thm:modsymbalg} is due to Ash and Rudolph
\cite{ash.rudolph}.  Instead of $P$, they used the larger
parallelotope $P'$ defined by 
\[
P':=\Bigl\{\sum \lambda _{i}v _{i}\Bigm | |\lambda_{i} |< 1
\Bigr\},
\]
which has volume $2^{n}\norm{\vv}$.  The observation that $P'$ can be
replaced by $P$ and the proof of (2) are both due to Barvinok
\cite{barvinok}.

\subsection{}\label{sec:retractandmanin}
%
The relationship between Theorem \ref{thm:modsymbalg} and Manin's
trick should be clear.  For $\Gamma \subset \SL_{2} (\Z)$, the Manin
symbols correspond exactly to the unimodular symbols mod $\Gamma$.  So
Theorem \ref{thm:modsymbalg} implies that every modular symbol (in the
language of Section~\ref{sec:modsymgen}) is a linear combination of Manin
symbols.  This is exactly the conclusion of Proposition
\ref{prop:mangen}.

In higher rank the relationship between Manin symbols and unimodular
symbols is more subtle.  In fact there are two possible notions of
``Manin symbol,'' which agree for $\SL_{2} (\Z)$ but not in general.
One possibility is the obvious one: a Manin symbol is a unimodular
symbol.  

The other possibility is to define a Manin symbol to be a modular
symbol corresponding to a top-dimensional cell of the retract $W$.
But for $n\geq 5$, such modular symbols need not be unimodular.  In
particular, for $n=5$ there are two equivalence classes of
top-dimensional cells.  One class corresponds to the unimodular
symbols, the other to a set of modular symbols of norm $2$.  However,
Theorems~\ref{thm:modsymbiso} and~\ref{thm:modsymbalg} show that
$H^{\nu} (\Gamma ; \Q)$ is spanned by unimodular symbols.  Thus as far
as this cohomology group is concerned, the second class of symbols is
in some sense unnecessary.


\subsection{}\label{ss:cohocomphecke}
%
We return to the setting of Section~\ref{ss:cohocompnohecke} and give
examples of Hecke eigenclasses in the cusp cohomology of $\Gamma =
\Gamma_{0} (p)\subset \SL_{3} (\Z)$.  We closely follow \cite{agg,
fourdutch}.  Note that since the top of the cuspidal range for
$\SL_{3}$ is the same
as the virtual cohomological dimension $\nu$, we can use modular
symbols to compute the Hecke action on cuspidal classes.

Given a prime $l$ coprime to $p$, there are two Hecke operators of
interest $T (l,1)$ and $T (l,2)$.  We can compute the action of these
operators on $H^{3}_{\text{cusp}} (\Gamma ; \C)$ as follows.  Recall
that $H^{3}_{\text{cusp}} (\Gamma ; \C)$ can be identified with a
certain space of functions $f\colon \Proj^{2} (\F_{p})\rightarrow \C$
(Theorem \ref{thm:agg}).  Given $x\in \Proj^{2} (\F_p)$, let $Q_{x}\in
\SL_{3} (\Z)$ be a matrix such that $Q_{x}\mapsto x$ under the
identification $\Proj^{2} (\F_p) \xrightarrow{\sim} \Gamma \backslash
\SL_{3} (\Z)$.  Then $Q_{x}$ determines a unimodular symbol $[Q_{x}]$
by taking the $v_{i}$ to be the columns of $Q_{x}$.  Given any Hecke
operator $T_{g}$, we can find coset representatives $h_{i}$ such that
$\Gamma g \Gamma  = \coprod \Gamma h_{i}$ (explicit representatives
for $\Gamma = \Gamma_{0} (p)$ and $T_{g} = T (l,k)$ are given in
\cite{agg, fourdutch}).  The modular symbols $[h_{i}Q_{x}]$ are no
longer unimodular in general, but we can apply Theorem
\ref{thm:modsymbalg} to write 
\[
[h_{i}Q_{x}] = \sum_{j} [R_{ij}], \quad R_{ij}\in \SL_{3} (\Z).
\]
Then for $f\colon \Proj^{2} (\F_p)\rightarrow \C$ as in Theorem
\ref{thm:agg}, we have 
\[
(T_{g}f) (x) = \sum_{i,j} f (\overline{R_{ij}}), 
\]
where $\overline{R_{ij}}$ is the class of $R_{ij}$ in $\Proj^{2}
(\F_p)$. 

Now let $\xi\in H^{3}_{\text{cusp}} (\Gamma ; \C)$ be a simultaneous
eigenclass for all the Hecke operators $T (l,1)$, $T (l,2)$, as $l$
ranges over all primes coprime with $p$.  General considerations from
the theory of automorphic forms imply that the eigenvalues $a (l,1)$,
$a (l,2)$ are complex conjugates of one other.  Hence it suffices to
compute $a (l,1)$.  We give two examples
of cuspidal eigenclasses for two different prime levels.

\begin{example}
Let $p=53$.  Then $H^{3}_{\text{cusp}} (\Gamma_{0} (53); \C)$ is
$2$-dimensional.  Let $\eta = (1+\sqrt{-11})/2$.  One eigenclass is
given by the data 
\begin{center}
\begin{tabular*}{0.75\textwidth}{@{\extracolsep{\fill}}c||c|c|c|c|c|c}
$l$&2&3&5&7&11&13\\
\hline
$a (l,1)$&$-1-2\eta$&$-2+2\eta$&$1$&$-3$&$1$&$-2-12\eta$\\
\end{tabular*}
\end{center}
and the other is obtained by complex conjugation.
\end{example}

\begin{example}
Let $p=61$.  Then $H^{3}_{\text{cusp}} (\Gamma_{0} (61); \C)$ is
$2$-dimensional.  Let $\omega  = (1+\sqrt{-3})/2$.  One eigenclass is
given by the data 
\begin{center}
\begin{tabular*}{0.9\textwidth}{@{\extracolsep{\fill}}c||c|c|c|c|c|c}
$l$&2&3&5&7&11&13\\
\hline
$a (l,1)$&$1-2\omega$&$-5+4\omega $&$-2+4\omega $&$-6\omega $&$-2+2\omega $&$-2-4\omega$\\
\end{tabular*}
\end{center}
and the other is obtained by complex conjugation.
\end{example}

%%%%%%%%%%%%%%%%%%%%%%%%%%%%%
%                           %
%  other cohomology groups  %
%                           %
%%%%%%%%%%%%%%%%%%%%%%%%%%%%%

\section{Other Cohomology Groups}\label{other.coho}
%
\subsection{}
%
In Section~\ref{mod.symb.section} we saw how to compute the Hecke action on
the top cohomology group $H^{\nu} (\Gamma ; \C)$.  Unfortunately for
$n\geq 4$, this cohomology group does not contain any cuspidal
cohomology.  The first case is $\Gamma \subset \SL_{4} (\Z )$; we have
$\vcd (\Gamma) = 6$, and the cusp cohomology lives in degrees $4$ and
$5$.  One can show that the cusp cohomology in degree $4$ is dual to
that in degree $5$, so for computational purposes it suffices to be
able to compute the Hecke action on $H^{5} (\Gamma ; \C)$.  But
modular symbols do not help us here.

In this section we describe a technique to compute the Hecke action on
$H^{\nu -1} (\Gamma ; \C)$, following \cite{experimental}.  The
technique is an extension of the modular symbol algorithm to these
cohomology groups.  In principle the ideas in this section can be
modified to compute the Hecke action on other cohomology groups
$H^{\nu -k} (\Gamma ; \C)$, $k>1$, although this has not been
investigated\footnote{The first interesting case is $n=5$, for which
the cuspidal cohomology lives in $H^{\nu -2}$.}.  For $n=4$, we have
applied the algorithm in joint work with Ash and McConnell to
investigate computationally the cohomology $H^{5} (\Gamma ; \C)$,
where $\Gamma_{0} (N)\subset \SL_{4} (\Z)$ \cite{computation}.

%\subsection{The Sharbly complex}\label{sharbly.sect}
\subsection{}\label{sharbly.sect}
%
To begin, we need an analogue of Theorem \ref{thm:modsymbiso} for
lower degree cohomology groups.  In other words, we need a
generalization of the modular symbols for other cohomology groups.
This is achieved by the \defn{sharbly complex} $S_{*}$:

\begin{definition}[\cite{ash.sharb}] \label{sharbly.complex}
 Let 
$\left\{S_{*},\partial \right\}$ be the chain complex given by the following data:
\begin{enumerate}
\item For $k\geq 0$, $S_{k}$ is the $\Q$-vector space generated
by the symbols $\uu  = [v_{1},\ldots,v_{n+k}]$, where $v_{i}\in
\Q^{n}\smallsetminus \{0 \}$, modulo the relations:
\begin{enumerate}
\item If $\tau $ is a permutation on $(n+k)$ letters, then
\[
[v_{1},\ldots,v_{n+k}] = \sign (\tau ) [\tau (v_{1}),\ldots,\tau
(v_{n+k})], 
\]
where $\sign (\tau )$ is the sign of $\tau $.
\item  If $q \in \Q^{\times }$, then 
\[
[q v_{1},v_{2}\ldots,v_{n+k}] = [v_{1},\ldots,v_{n+k}].
\]
\item If the rank of the matrix
$(v_{1},\ldots,v_{n+k})$ is less than $n$, then $\uu = 0$.
\end{enumerate}
\item For $k>0$, the boundary map $\partial \colon S_{k}\rightarrow S_{k-1}$ is 
$$[v_{1},\ldots,v_{n+k}] \longmapsto \sum _{i=1}^{n+k}  (-1)^{i}
[v_{1},\ldots,\hat{v_{i}},\ldots,v_{n+k}].  $$
We define $\partial$ to be identically zero on $S_{0}$.
\end{enumerate}
\end{definition}

The elements $$\uu = [v_{1},\dots ,v_{n+k}]$$ are called
\defn{$k$-sharblies}\footnote{The terminology for $S_{*}$ is due to
Lee Rudolph, in honor of Lee and Szczarba.  They introduced a very
similar complex in \cite{l.s} for $\SL_{3} (\Z)$.}.  The $0$-sharblies
are exactly the modular symbols from Definition
\ref{def:modular.symbols}, and the subspace $B\subset S_{0}$ is the
image of the boundary map $\partial \colon S_{1}\rightarrow S_{0}$.

There is an obvious left action of $\Gamma $ on $S_{* }$ commuting
with $\partial $.  For any $k\geq 0$, let $S_{k,\Gamma }$ be the space
of $\Gamma $-coinvariants.  Since the boundary map $\partial$ commutes
with the $\Gamma $-action, we obtain a complex $(S_{*,\Gamma},
\partial_{\Gamma})$.  The following theorem shows that this complex
computes the cohomology of $\Gamma $:

\begin{theorem}[\cite{ash.sharb}]
\label{thm:sharbiso}
There is a natural isomorphism 
\[
H^{\nu - k}(\Gamma ; \C) \xrightarrow{\sim} H_{k}
(S_{*,\Gamma }\otimes \C ).
\]
\end{theorem}

\subsection{}
%
We can extend our norm function $\|\phantom{a}\|$ from modular
symbols to all of $S_{k}$ as follows.  Let $\uu = [v_{1},\dots
,v_{n+k}]$ be a $k$-sharbly, and let $Z(\uu)$ be the set of all
submodular symbols determined by $\uu$.  In other words, $Z (\uu)$
consists of the modular symbols of the form $[v_{i_{1}},\dotsc
,v_{i_{n}}]$, where $\{i_{1},\dotsc ,i_{n} \}$ ranges over all
$n$-fold subsets of $\{1,\dotsc ,n+k \}$.  Define $\|\uu \|$ by
\[
\|\uu\| = \Max_{\vv \in Z (\uu )} \| \vv \|.
\]
Note that $\| \uu \|$ is well defined modulo the relations in
Definition \ref{sharbly.complex}.  As for modular symbols, we extend
the norm to sharbly chains $\xi = \sum q (\uu)\uu $ taking the maximum
norm over the support.  Formally, we let $\support (\xi) = \{\uu \mid
q (\uu)\not =0 \}$ and $Z (\xi) = \bigcup_{\uu \in \support (\xi )} Z
(\uu)$, and then we define $\|\xi \|$ by
\[
\|\xi \| = \Max_{\vv \in Z (\xi )} \| \vv \|.
\]

We say that $\xi $ is \defn{reduced} if $\|\xi \| = 1$.  Hence $\xi$
is reduced if and only if all its submodular symbols are unimodular or
have determinant $0$.  Clearly there are only finitely many reduced
$k$-sharblies modulo $\Gamma$ for any $k$.

In general the cohomology groups $H^{*} (\Gamma ; \C)$ are \emph{not}
spanned by reduced sharblies.  However, it is known (cf. \cite{M91})
that for $\Gamma \subset \SL _{4} (\Z )$, the group $H^{5} (\Gamma ;
\C)$ is spanned by reduced $1$-sharbly cycles.  The best one can say
in general is that for each pair $n,k$, there is an integer $N=N
(n,k)$ such that for $\Gamma \subset \SL_{n} (\Z)$, $H^{\nu -k}
(\Gamma; \C)$ is spanned by $k$-sharblies of norm $\leq N$.  This set
of sharblies is also finite modulo $\Gamma$, although it is not known
how large $N$ must be for any given pair $n,k$.

%\subsection{Sharblies and Voronoi}\label{that.section}
\subsection{}\label{that.section}
%
Recall that the cells of the well-rounded retract $W$ are indexed by
sets of primitive vectors in $\Z^{n}$.  Since any primitive vector
determines a point in $\Q^{n}\smallsetminus \{0 \}$ and since sets of
such points index sharblies, it is clear that there is a close
relationship between $S_{*}$ and the chain complex associated to $W$,
although of course $S_{*}$ is much bigger.  In any case, both
complexes compute $H^{*} (\Gamma ; \C)$.

The main benefit of using the sharbly complex to compute cohomology is
that it admits a Hecke action.  Suppose $\xi = \sum q (\uu )\uu $ is a
sharbly cycle mod $\Gamma $, and consider a Hecke operator $T_{g}$.
Then we have
\begin{equation}\label{how.hecke.acts}
T_{g}(\xi ) = \sum _{h\in \Omega , \uu } n (\uu )h\cdot \uu, 
\end{equation}
where $\Omega$ is a set of coset representatives as in
\eqref{eq:coset.decomp}.  
Since $\Omega \not \subset \SL _{n} (\Z )$ in general, the 
Hecke image of a
reduced sharbly is not usually reduced.

\subsection{}
%
We are now ready to describe our algorithm for the computation of the
Hecke operators on $H^{\nu -1} (\Gamma ;\C )$.   It suffices to describe an
algorithm that takes as input a $1$-sharbly cycle $\xi $ and produces
as output a cycle $\xi '$ with
\begin{enumerate}
\item [(a)] the classes of $\xi $ and $\xi '$ in $H^{\nu -1} (\Gamma
;\C )$ the same, and 
\item [(b)] $\|\xi' \| < \|\xi \|$ if $\|\xi \|>1$.
\end{enumerate}

Below, we will present an algorithm satisfying (a).  In
\cite{experimental}, we conjectured (and presented evidence) that the
algorithm satisfies (b) for $n\leq 4$.  Further evidence is provided
by the computations in \cite{computation}, which relied on the
algorithm to compute the Hecke action on $H^{5} (\Gamma ; \C)$, where
$\Gamma = \Gamma_{0} (N)\subset \SL_{4} (\Z)$.

The idea behind the algorithm is simple: given a $1$-sharbly cycle
$\xi$ that is not reduced, (i) simultaneously apply the modular symbol
algorithm (Theorem \ref{thm:modsymbalg}) to each of its submodular
symbols, and then (ii) package the resulting data into a new $1$-sharbly
cycle.  Our experience in presenting this algorithm is that most
people find the geometry involved in (ii) daunting.  Hence we will
give details only for $n=2$ and will provide a sketch for $n>2$.
Full details are contained in \cite{experimental}.  Note that $n=2$ is
topologically and arithmetically uninteresting, since we are computing
the Hecke action on $H^{0} (\Gamma ; \C)$; nevertheless, the geometry
faithfully represents the situation for all $n$.

\subsection{}\label{sl2.start}
%
Fix $n=2$, let $\xi \in S_{1}$ be a $1$-sharbly cycle mod $\Gamma$ for
some $\Gamma \subset \SL_{2} (\Z )$, and suppose $\xi $ is not
reduced.  Assume $\Gamma$ is torsion-free to simplify the
presentation.  

Suppose first that all submodular symbols $\vv \in Z (\xi )$ are
nonunimodular.  Select reducing points for each $\vv \in Z (\xi )$
and make these choices $\Gamma $-equivariantly.  This means the
following.  Suppose $\uu,\uu '\in \support \xi $ and $\vv \in \support
(\partial \uu )$ and $\vv' \in \support (\partial \uu' )$ are modular
symbols such that $\vv = \gamma\cdot \vv ' $ for some $\gamma \in
\Gamma $.  Then we select reducing points $w$ for $\vv $ and $w'$ for
$\vv '$ such that $w = \gamma\cdot w'$.  (Note that since $\Gamma$ is
torsion-free, no modular symbol can be identified to itself by an
element of $\Gamma$; hence $\vv \not = \vv '$.)  This is possible
since if $\vv $ is a modular symbol and $w$ is a reducing point for
$\vv $, then $\gamma \cdot w $ is a reducing point for $\gamma \cdot
\vv$ for any $\gamma \in \Gamma $.  Because there are only finitely
many $\Gamma $-orbits in $Z (\xi )$, we can choose reducing points
$\Gamma $-equivariantly by selecting them for some set of orbit
representatives.

It is important to note that $\Gamma $-equivariance is the only
global criterion we use when selecting reducing.  In
particular, there is a priori no relationship among the three reducing
points chosen for any $\uu \in \support \xi $.

\subsection{}\label{ss:buildingxiprime}
%
Now we want to use the reducing points and the $1$-sharblies in $\xi$ to
build $\xi '$.  Choose $\uu = [v_{1},v_{2},v_{3}]\in
\support \xi $, and denote the reducing point for $[v_{i},v_{j}]$ by $w_{k}$,
where $\{i,j,k \} = \{1,2,3 \}$.  We use the $v_{i}$ and the
$w_{i}$ to build a $2$-sharbly chain $\eta (\uu )$ as follows.

Let $P$ be an octahedron in $\R ^{3}$.  Label the vertices of $P$ with
the $v_{i}$ and $w_{i}$ such that the vertex labeled $v_{i}$ is
opposite the vertex labeled $w_{i}$ (Figure \ref{octa.fig}).
Subdivide $P$ into four tetrahedra by connecting two opposite
vertices, say $v_{1}$ and $w_{1}$, with an edge (Figure
\ref{octa-hom.fig}).
For each tetrahedron $T$, take the labels of four vertices and arrange
them into a quadruple.  If we orient $P$, then we can use the induced
orientation on $T$ to order the four primitive points.  In this way,
each $T$ determines a $2$-sharbly, and $\eta (\uu)$ is defined to be
the sum.  For example, if we use the decomposition in Figure
\ref{octa-hom.fig}, we have
\begin{equation}\label{eta}
\eta (\uu ) = [v_{1},v_{3},
v_{2},w_{1}]+[v_{1},w_{2},v_{3},w_{1}]+[v_{1},w_{3},w_{2},w_{1}]+[v_{1},v_{2},w_{3},w_{1}].
\end{equation}
Repeat this construction for all $\uu \in \support \xi $, and let 
$\eta =\sum q (\uu ) \eta (\uu )$.  Finally, let $\xi ' = \xi +\partial \eta $.

\begin{figure}[htb]
\begin{center}
\psfrag{v1}{$v_{1}$}
\psfrag{v2}{$v_{2}$}
\psfrag{v3}{$v_{3}$}
\psfrag{w1}{$w_{1}$}
\psfrag{w2}{$w_{2}$}
\psfrag{w3}{$w_{3}$}
\includegraphics[scale=.5]{diagrams/octa}
\end{center}
\caption{\label{octa.fig}}
\end{figure}

\begin{figure}[htb]
\begin{center}
\psfrag{v1}{$v_{1}$}
\psfrag{v2}{$v_{2}$}
\psfrag{v3}{$v_{3}$}
\psfrag{w1}{$w_{1}$}
\psfrag{w2}{$w_{2}$}
\psfrag{w3}{$w_{3}$}
\includegraphics[scale=.5]{diagrams/octa-hom}
\end{center}
\caption{\label{octa-hom.fig}}
\end{figure}

\subsection{}\label{endofdiscussion}
%
By construction, $\xi' $ is a cycle mod $\Gamma$ in the same
class as $\xi $.  We claim in addition that no submodular symbol from
$\xi $ appears in $\xi '$.  To see this, consider $\partial \eta (\uu
)$.  From \eqref{eta}, we have
\begin{multline}\label{bound}
\partial \eta (\uu ) = - [v_{1},v_{2},v_{3}] + [v_{1},v_{2},w_{3}] +
[v_{1},w_{2},v_{3}] + [w_{1},v_{2},v_{3}] \\
- [v_{1},w_{2},w_{3}] -
[w_{1},v_{2},w_{3}] - [w_{1},w_{2},v_{3}] + [w_{1},w_{2},w_{3}].
\end{multline}
Note that this is the boundary in $S_{*}$, not in $S_{*,\Gamma }$.
Furthermore, $\partial \eta (\uu )$ is independent of which pair of
opposite vertices of $P$ we connected to build $\eta (\uu )$.

From \eqref{bound}, we see that in $\xi +\partial \eta $ the
$1$-sharbly $-[v_{1},v_{2},v_{3}]$ is canceled by $\uu \in \support
\xi $.  We also claim that $1$-sharblies in \eqref{bound} of the form
$[v_{i},v_{j},w_{k}]$ vanish in 
$\partial_{\Gamma } \eta $.

To see this, let $\uu $,$\uu '\in \support \xi $, and suppose $\vv =
[v_{1},v_{2}]\in \support \partial \uu $ equals $\gamma \cdot \vv '$
for some $\vv ' = [v_{1}',v_{2}']\in \support \partial \uu '$.  Since
the reducing points were chosen $\Gamma $-equivariantly, we have
$w=\gamma \cdot w'$.  This means that the $1$-sharbly $[v_{1},v_{2},
w]\in \partial \eta (\uu )$ will be canceled mod $\Gamma $ by
$[v_{1}',v_{2}', w']\in \partial \eta (\uu' )$.  Hence, in passing
from $\xi $ to $\xi '$, the effect in $\coinv{*}$ is to replace $\uu $
with \emph{four} $1$-sharblies in $\support \xi '$:
\begin{equation}\label{tfm}
[v_{1},v_{2},v_{3}]\longmapsto - [v_{1},w_{2},w_{3}] -
[w_{1},v_{2},w_{3}] - [w_{1},w_{2},v_{3}] + [w_{1},w_{2},w_{3}].
\end{equation}
Note that in \eqref{tfm}, there are no $1$-sharblies of the form
$[v_{i},v_{j},w_{k}]$.

\begin{remark}\label{imp1}
For implementation purposes, it is not necessary to explicitly
construct $\eta $.  Rather, one may work directly with \eqref{tfm}.
\end{remark}

\subsection{}\label{interior.ms}
%
Why do we expect $\xi ' $ to satisfy $\|\xi '\| < \|\xi \|$?  First of
all, in the right hand side of \eqref{tfm} there are no submodular
symbols of the form $[v_{i},v_{j}]$.  In fact, any submodular symbol
involving a point $v_{i}$ also includes a reducing point for 
$[v_{i},v_{j}]$.

On the other hand, consider the submodular symbols in \eqref{tfm} of
the form $[w_{i},w_{j}]$.  Since there is no relationship among the
$w_{i}$, one has no reason to believe that these modular symbols are
closer to unimodularity than those in $\uu $.  Indeed, for certain
choices of reducing points it can happen that $\|[w_{i},w_{j}]\|\geq
\|\uu \|$.

The upshot is that some care must be taken in choosing reducing
points.  In \cite[Conjectures 3.5 and 3.6]{experimental} we describe
two methods for finding reducing points for modular symbols, one using
\Vor reduction and one using LLL-reduction.  Our experience is that
if one selects reducing points using either of these conjectures,
then $\|[w_{i},w_{j}]\|< \|\uu \|$ for each of the new modular symbols
$[w_{i},w_{j}]$.  In fact, in practice these symbols are trivial or
satisfy $\|[w_{i},w_{j}]\|=1$.

\subsection{}\label{sl2.stop}
%
In the previous discussion we assumed that no submodular symbols of
any $\uu \in \support \xi $ were unimodular.  Now we say what to
do if some are.  There are three cases to consider.  

First, all submodular symbols of $\uu $ may be unimodular.  In this
case there are no reducing points, and \eqref{tfm} becomes
\begin{equation}\label{tfm2}
[v_{1},v_{2},v_{3}]\longmapsto [v_{1},v_{2},v_{3}].
\end{equation}

Second, one submodular symbol of $\uu $ may be nonunimodular, say the
symbol $[v_{1},v_{2}]$.  In this case, to build $\eta$, we use a
tetrahedron $P'$
and put $\eta (\uu ) = [v_{1},v_{2},v_{3},w_{3}]$
(Figure~\ref{tetra.fig}).  Since $[v_{1}, v_{2}, w_{3}]$ vanishes
in the boundary of $\eta $ mod $\Gamma $, \eqref{tfm} becomes
\begin{equation}\label{tfm3}
[v_{1},v_{2},v_{3}]\mapsto -[v_{1},v_{3},w_{3}] + [v_{2}, v_{3},
w_{3}].
\end{equation}

\begin{figure}[ht]
\begin{center}
\psfrag{v1}{$v_{1}$}
\psfrag{v2}{$v_{2}$}
\psfrag{v3}{$v_{3}$}
\psfrag{w3}{$w_{3}$}
\includegraphics[scale=.5]{diagrams/tetra}
\end{center}
\caption{\label{tetra.fig}}
\end{figure}

Finally, two submodular symbols of $\uu $ may be nonunimodular, say
$[v_{1},v_{2}]$ and $[v_{1},v_{3}]$.  In this case we use the cone on
a square $P''$ (Figure~\ref{cone-on-square.fig}).  To construct $\eta
(\uu )$, we must choose a decomposition of $P''$ into tetrahedra.
Since $P''$ has a nonsimplicial face, this choice affects $\xi '$ (in
contrast to the previous cases).  If we subdivide $P''$ by connecting
the vertex labelled $v_{2}$ with the vertex labelled $w_{2}$, we
obtain
\begin{equation}\label{tfm4}
[v_{1},v_{2},v_{3}]\longmapsto[v_{2},w_{2},w_{3}]+[v_{2},v_{3},w_{2}]+
[v_{1},v_{3},w_{2}].
\end{equation}

\begin{figure}[ht]
\begin{center}
\psfrag{v1}{$v_{1}$}
\psfrag{v2}{$v_{2}$}
\psfrag{v3}{$v_{3}$}
\psfrag{w2}{$w_{2}$}
\psfrag{w3}{$w_{3}$}
\includegraphics[scale=.5]{diagrams/cone-on-square}
\end{center}
\caption{\label{cone-on-square.fig}}
\end{figure}

\subsection{}
%
Now consider general $n$.  The basic technique is the same, but the
combinatorics become more complicated.  Suppose $\uu = [v_{1},\dots ,v_{n+1}]$
satisfies $q (\uu )\not =0$ in a $1$-sharbly cycle $\xi$, and for
$i=1,\dots ,n+1$ let $\vv_{i} $ be the submodular symbol $[v_{1},\dots
,\widehat{v_{i}},\dots ,v_{n+1}]$.  Assume that all $\vv_{i}$ are
nonunimodular, and for each $i$ let $w_{i}$ be a reducing point for
$\vv_{i}$.

For any subset $I\subset\{1,\dots ,n+1 \}$, let $\uu_{I}$ be the
$1$-sharbly $[u_{1},\dots ,u_{n+1}]$, where $u_{i}=w_{i}$ if $i\in I$,
and $u_{i}=v_{i}$ otherwise.  The polytope $P$ used to build $\eta
(\uu)$ is the \defn{cross polytope}, which is the higher-dimensional
analogue of the octahedron \cite[\S4.4]{experimental}.  We suppress the
details and give the final answer:
\eqref{tfm} becomes
\begin{equation}\label{tfmrelation}
\uu \longmapsto -\sum _{I} (-1)^{\#I}\uu_{I},
\end{equation}
where the sum is taken over all subsets $I\subset \{1,\dotsc ,n+1 \}$ 
of cardinality \emph{at least $2$}.  

More generally, if some $\vv_{i}$ happen to be unimodular, then the
polytope used to build $\eta$ is an iterated cone on a
lower-dimensional cross polytope.  This is already visible for $n=2$:
\begin{itemize}
\item The 2-dimensional cross polytope is a
square, and the polytope $P''$ is
a cone on a square.
\item The 1-dimensional cross polytope is an interval, and the
polytope $P'$ is a double cone on an interval.
\end{itemize}
Altogether there are $n+1$ relations generalizing
\eqref{tfm2}--\eqref{tfm4}.

\subsection{}\label{ss:sl4computations}
%
Now we describe how these computations are carried out in practice,
focusing on $\Gamma = \Gamma_{0} (N)\subset \SL_{4} (\Z)$ and $H^{5}
(\Gamma; \C)$.  Besides discussing technical details,
we also have to slightly modify some aspects of the construction in
Section~\ref{sl2.start}, since $\Gamma$ is not torsion-free.

Let $W$ be the well-rounded retract.  We can represent
a cohomology class $\beta \in H^{5} (\Gamma ; \C)$ as $ \beta = \sum q
(\sigma )\sigma $, where $\sigma $ denotes a codimension $1$ cell in
$W$.  In this case there are three types of codimension $1$ cells in
$W$.  Under the bijection $W \leftrightarrow \sC$, these cells
correspond to the \Vor cells indexed by the columns of the matrices 
\begin{equation}\label{eq:4onesharbs}
\left(\begin{array}{ccccc}
1&0&0&0&1\\
0&1&0&0&1\\
0&0&1&0&1\\
0&0&0&1&1
\end{array} \right),\quad 
\left(\begin{array}{ccccc}
1&0&0&0&1\\
0&1&0&0&1\\
0&0&1&0&1\\
0&0&0&1&0
\end{array} \right),\quad 
\left(\begin{array}{ccccc}
1&0&0&0&1\\
0&1&0&0&1\\
0&0&1&0&0\\
0&0&0&1&0
\end{array} \right).
\end{equation}
Thus each $\sigma$ in $W$ modulo $\Gamma$ corresponds to an $\SL_{4}
(\Z)$-translate of one of the matrices in \eqref{eq:4onesharbs}.  
These translates
determine basis $1$-sharblies $\uu$ (by taking the points $u_{i}$ to
be the columns), and hence we can represent $\beta$ by a 1-sharbly
chain $ \xi = \sum q
(\uu) \uu \in S_{1}$ that is a cycle in the complex of coinvariants
$(S_{*,\Gamma}, \partial_{\Gamma})$.

To make later computations more efficient, we precompute more data
attached to $\xi$.  Given a $1$-sharbly $\uu = [u_{1},\dotsc
,u_{n+1}]$, a \emph{lift} $M (\uu)$ of $\uu$ is defined to be an
integral matrix with primitive columns $M_{i}$ such that $\uu =
[M_{1},\dotsc ,M_{n+1}]$.  Then we encode $\xi $, once and for all, by
a finite collection $\Phi$ of $4$-tuples 
\[
(\uu , n (\uu ), \{\vv \}, \{M (\vv
) \}),
\]
where
\begin{enumerate}
\item $\uu$ ranges over the support of $\xi$,
\item $n (\uu )\in \C $ is the coefficient of $\uu$ in $\xi$,
\item $\{\vv \}$ is the set of submodular symbols appearing in the
boundary of $\uu$, and 
\item $\{M (\vv ) \}$ is a set of lifts for $\{\vv \}$.\label{ifour}  
\end{enumerate}
Moreover, the lifts in (\ref{ifour}) are chosen to satisfy the
following $\Gamma $-equivariance condition.  Suppose that for $\uu
,\uu '\in \support \xi $ we have $\vv \in \support (\partial \uu ) $
and $\vv '\in \support (\partial \uu ')$ satisfying $\vv = \gamma
\cdot \vv '$ for some $\gamma \in \Gamma $.  Then we require $M (\vv
)= \gamma M (\vv ')$.  This is possible since $\xi$ is a cycle modulo
$\Gamma$, although there is one complication since $\Gamma$ has torsion:
it can happen that some submodular symbol
$\vv$ of a $1$-sharbly $\uu$ is identified to \emph{itself} by an
element of $\Gamma$.  
This means that in constructing $\{M (\vv)
\}$ for $\uu$, we must somehow choose more than one lift for $\vv$.
To deal with this, let $M
(\vv)$ be any lift of $\vv$, and let $\Gamma (\vv)\subset \Gamma $ be
the stabilizer of $\vv$.  Then in $\xi$, we replace $q (\uu ) \uu$ by
\[
\frac{1}{\#\Gamma (\vv )}\sum_{\gamma \in \Gamma (\vv)} q (\uu) \uu_{\gamma},
\]
where $\uu_{\gamma }$ has the same data as $\uu$, 
except\footnote{In fact, we can be slightly
more clever than this and only introduce denominators 
that are powers of $2$.} 
that we give
$\vv$ the lift $\gamma M (\vv)$.


Next we compute and store the 1-sharbly transformation laws generalizing
\eqref{tfm2}--\eqref{tfm4}.  As a part of this we fix triangulations
of certain cross polytopes as in \eqref{tfm4}.  

We are now ready to begin the actual reduction algorithm.  We take a
Hecke operator $T (l,k)$ and build the coset representatives $\Omega $
as in \eqref{how.hecke.acts}.  For each $h\in \Omega $ and each
$1$-sharbly $\uu$ in the support of $\xi$, we obtain a non-reduced
$1$-sharbly $\uu_{h} := h\cdot\uu$.  Here $h$ acts on all the data
attached to $\uu$ in the list $\Phi$.  In particular, we replace each
lift $M (\vv)$ with $h\cdot M (\vv)$, where the dot means matrix
multiplication.

Now we check the submodular symbols of $\uu_{h}$ and choose reducing
points for the nonunimodular symbols.  This is where the lifts come in
handy.  Recall that reduction points must be chosen
$\Gamma$-equivariantly over the entire cycle.  Instead of explicitly
keeping track of the identifications between modular symbols, we do
the following trick:
\begin{enumerate}
\item Construct the \defn{Hermite normal form} $M_{\her}(\vv)$ of the lift $M
(\vv)$ (see \cite[\S2.4]{cohen:course_ant} and \exref{ch:linalg}{ex:hnf}).  Record the transformation matrix $U \in \GL_{4} (\Z)$ such
that $UM(\vv ) = M_{\her} (\vv)$.
\item Choose a reducing point $u$ for $M_{\her} (\vv)$.
\item Then the reducing point for $M (\vv)$ is $U^{-1}u$.
\end{enumerate}
This guarantees $\Gamma$-equivariance: if $\vv$, $\vv '$ are submodular
symbols of $\xi$ with $\gamma \cdot \vv = \vv '$ and with reducing
points $u, u'$, we have $\gamma u = u'$.  The reason is that the
Hermite normal form $M_{\her} (\vv)$ is a \emph{uniquely determined}
representative of the $\GL_{4}(\Z)$-orbit of
$M (\vv)$ \cite{cohen:course_ant}.  Hence if $\gamma M
(\vv) = M (\vv ')$, then $M_{\her} (\vv) = M_{\her} (\vv
')$.

After computing all reducing points, we apply the appropriate
transformation law.  The result will be a chain of $1$-sharblies, each
of which has (conjecturally) smaller norm than the original
$1$-sharbly $\uu$.  We output these $1$-sharblies if they are reduced;
otherwise they are fed into the reduction algorithm again.  Eventually
we obtain a reduced $1$-sharbly cycle $\xi '$ homologous to the
original cycle $\xi$.

The final step of the algorithm is to rewrite $\xi '$ as a cocycle on
$W$.  This is easy to do since the relevant cells of $W$ are in
bijection with the reduced $1$-sharblies.  There are some nuisances in
keeping orientations straight, but the computation is not difficult.
We refer to \cite{computation} for details.

\subsection{}
%
We now give some examples, taken from \cite{computation}, of Hecke
eigenclasses in $H^{5} (\Gamma_{0} (N); \C)$ for various levels
$N$.
Instead of giving a table of eigenvalues, we give the \defn{Hecke
polynomials}.  If $\beta$ is an eigenclass with $T (l,k) (\beta) = a
(l,k)\beta$, then we define
\[
H (\beta , l) = \sum_k (-1)^k l^{k(k-1)/2} a(l, k) X^k \in \C [X].
\]
For almost all $l$, after putting $X = l^{-s}$ where $s$ is a complex
variable, the function $H
(\beta , s)$ is the inverse of the local factor at $l$ of the
automorphic representation attached to $\beta$.
\begin{example}
Suppose $N=11$.  Then the cohomology $H^{5} (\Gamma_{0} (11); \C)$ is
2-dimensional.  There are two Hecke eigenclasses $u_{1}, u_{2}$, each
with rational Hecke eigenvalues.

\medskip
\begin{center}
\begin{tabular}{|p{40pt}|p{40pt}|p{200pt}|}
\hline
$u_{1}$&$T_2$&$ (1-4 X)(1-8 X)(1 +2 X + 2 X^2)$\\
&$T_3$&$ (1-9 X)(1-27 X)(1 + X + 3 X^2) $\\
&$T_5$&$ (1-25 X)(1-125 X)(1 - X + 5 X^2) $\\
&$T_7$&$ (1-49 X)(1-343 X)(1 +2 X + 7 X^2) $\\
\hline
$u_{2}$&$T_2$&$ (1- X)(1-2 X)(1 +8 X + 32 X^2)$\\
&$T_3$&$ (1- X)(1-3 X)(1 +9 X + 243 X^2) $\\
&$T_5$&$ (1- X)(1-5 X)(1 -25 X + 3125 X^2) $\\
&$T_7$&$ (1- X)(1-7 X)(1 +98 X + 16807 X^2) $\\
\hline
\end{tabular}
\end{center}

  
\end{example}
 
\begin{example}
Suppose $N=19$.  Then the cohomology $H^{5}
(\Gamma_{0} (19); \C)$ is 3-dimensional.  There are three Hecke
eigenclasses $u_{1}, u_{2}, u_{3}$, each with rational Hecke eigenvalues.

\medskip
\begin{center}
\begin{tabular}{|p{40pt}|p{40pt}|p{200pt}|}
\hline
$u_{1}$&$T_2$&$ (1-4 X)(1-8 X)(1  + 2 X^2)$\\
&$T_3$&$ (1-9 X)(1-27 X)(1 +2 X + 3 X^2) $\\
&$T_5$&$ (1-25 X)(1-125 X)(1 -3 X + 5 X^2) $\\
\hline
$u_{2}$&$T_2$&$ (1- X)(1-2 X)(1  + 32 X^2)$\\
&$T_3$&$ (1- X)(1-3 X)(1 +18 X + 243 X^2) $\\
&$T_5$&$ (1- X)(1-5 X)(1 -75 X + 3125 X^2) $\\
\hline
$u_{3}$&$T_2$&$ (1- 2 X)(1-4 X)(1 +3 X + 8 X^2)$\\
&$T_3$&$ (1- 3 X)(1-9 X)(1 +5 X + 27 X^2) $\\
&$T_5$&$ (1- 5 X)(1-25 X)(1 +12 X + 125 X^2) $\\
\hline
\end{tabular}
\end{center}

\end{example}

In these examples, the cohomology is completely accounted for by the
Eisenstein summand of \eqref{franke.decomp}.  In fact, let
$\Gamma'_{0} (N)\subset \SL_{2} (\Z)$ be the usual Hecke congruence
subgroup of matrices upper-triangular modulo $N$.  Then the cohomology
classes above actually come from classes in $H^{1} (\Gamma_{0}' (N))$,
that is from holomorphic modular forms of level $N$.

For $N=11$, the space of weight two cusp forms $S_{2} (11)$ is
1-dimensional.  This cusp form $f$ lifts in two different ways to
$H^{5} (\Gamma_{0} (11); \C)$, which can be seen from the quadratic
part of the Hecke polynomials for the $u_i$.  Indeed, for $u_{i}$ the
quadratic part is exactly the inverse of the local factor for the
$L$-function attached to $f$, after the substitution $X = l^{-s}$.
For $u_{2}$, we see that the lift is also twisted by the square of the
cyclotomic character.  (In fact the linear terms of the Hecke
polynomials come from powers of the cyclotomic character.)

For $N=19$, the space of weight two cusp forms $S_{2} (19)$ is again
1-dimensional.  The classes $u_{1}$ and $u_{2}$ are lifts of this
form, exactly as for $N=11$.  The class $u_{3}$, on the other hand,
comes from $S_{4} (19)$, the space of weight $4$ cusp forms on
$\Gamma_{0}' (19)$.  In fact, $\dim S_{4} (19) = 4$, with one Hecke
eigenform defined over $\Q$ and another defined over a totally real
cubic extension of $\Q$.  Only the rational weight four eigenform
contributes to $H^{5} (\Gamma_{0} (19); \C)$.  One can show that
whether or not a weight four cuspidal eigenform $f$ contributes to the
cohomology of $\Gamma_{0} (N)$ depends only on the sign of the
functional equation of $L (f,s)$ \cite{weselman}.  This phenomenon
is typical of what one encounters when studying Eisenstein cohomology.

In addition to the lifts of weight 2 and weight 4 cusp forms, for
other levels one finds lifts of Eisenstein series of weights 2 and 4
and lifts of cuspidal cohomology classes from subgroups of $\SL_{3}
(\Z)$.  For some levels one finds cuspidal classes that appear to be
lifts from the group of symplectic similitudes $\GSp (4)$.  More
details can be found in \cite{computation, computationII}.

\subsection{}
%
Here are some notes on the reduction algorithm and its implementation:
\begin{itemize}
\item Some additional care must be taken when selecting reducing
points for the submodular symbols of $\uu$.  In particular, in
practice one should choose $w$ for $\vv$ such that $\sum \| \vv_{i}
(w)\|$ is minimized.  Similar remarks apply when choosing a
subdivision of the crosspolytopes in Section~\ref{sl2.stop}.
\item In practice, the reduction algorithm has \emph{always} terminated with
a reduced $1$-sharbly cycle $\xi '$ homologous to $\xi$.  However, at
the moment we cannot prove that this will always happen.
\item Experimentally, the efficiency of the reduction step appears to
be comparable to that of Theorem \ref{thm:modsymbalg}.  In other words
the depth of the ``reduction tree'' associated to a given $1$-sharbly
$\uu$ seems to be bounded by a polynomial in $\log \log \| \uu \| $.  Hence
computing the Hecke action using this algorithm is extremely
efficient.

On the other hand, computing Hecke operators on $\SL_{4}$ is still a
much bigger computation---relative to the level---than on $\SL_{2}$
and $\SL_{3}$.  For example, the size of the full retract $W$ modulo
$\Gamma_{0} (p)$ is roughly $O (p^{6})$, which grows rapidly with $p$.
The portion of the retract corresponding to $H^{5}$ is much smaller,
around $p^{3}/10$, but this still grows quite quickly.  This makes
computing with $p>100$ out of reach at the moment.

The number of Hecke cosets grows rapidly as well, e.g., the number of
coset representatives of $T (l,2)$ is $l^{4}+l^{3}+2l^{2}+l+1$.  Hence
it is only feasible to compute Hecke operators for small $l$; for
large levels only $l=2$ is possible.

Here are some numbers to give an idea of the size of these
computations.  For level $73$, the rank of $H^{5}$ is 20.  There are
39504 cells of codimension~$1$ and 4128 top-dimensional cells in $W$
modulo $\Gamma_{0} (73)$.  The computational techniques in
\cite{computation} used at this level (a Lanczos scheme over a large
finite field) tend to produce sharbly cycles supported on almost all
the cells.  Computing $T (2,1)$ requires a reduction tree of 
depth $1$
and produces as many as 26 reduced $1$-sharblies for each of the 15
nonreduced Hecke images.  Thus one cycle produces a cycle supported on
as many as 15406560 $1$-sharblies, all of which must be converted to
an appropriate cell of $W$ modulo $\Gamma$.  Also this is just what
needs to be done for \emph{one} cycle; do not forget that the rank of
$H^{5}$ is 20.

In practice the numbers are slightly better, since the reduction step
produces fewer $1$-sharblies on average and since the support of the
initial cycle has size less than $39504$.  Nevertheless the orders of
magnitude are correct.  

\item Using lifts is a convenient way to encode the global
$\Gamma$-identifications in the cycle $\xi$, since it means we do not
have to maintain a big data structure keeping track of the identifications on
$\partial \xi$.  However, there is a certain expense in computing the
Hermite normal form.  This is balanced by the benefit that working
with the data $\Phi$ associated to $\xi$ allows us to reduce the
supporting $1$-sharblies $\uu$ \emph{independently}.  This means we
can cheaply parallelize our computation: each $1$-sharbly, encoded as
a $4$-tuple $(\uu , n (\uu ), \{\vv \}, \{M (\vv ) \})$, can be
handled by a separate computer.  The results of all these individual
computations can then be collated at the end, when producing a
$W$-cocycle.
\end{itemize}

%%%%%%%%%%%%%%%%%%%
%                 %
%  open problems  %
%                 %
%%%%%%%%%%%%%%%%%%%

\section{Complements and Open Problems}\label{apps}

\subsection{}
%
We conclude this appendix by giving some complements and describing
some possible directions for future work, both theoretical and
computational.  Since a full explanation of the material in this
section would involve many more pages, we will be brief and will
provide many references.

\subsection{Perfect Quadratic Forms over Number Fields and Retracts}
%
Since \Vor 's pioneering work \cite{vor1}, it has been the goal of
many to extend his results from $\Q$ to a general algebraic number
field $F$.  Recently Coulangeon \cite{coul}, building on work of Icaza
and Baeza \cite{icaza, baeza.icaza}, has found a good notion of
\defn{perfection} for quadratic forms over number
fields\footnote{Such forms are called \defn{Humbert forms} in the
literature.}.  One of the key ideas in \cite{coul} is that the correct
notion of equivalence between Humbert forms involves not only the
action of $\GL_{n} (\sO_{F})$, where $\sO_{F}$ is the ring of integers
of $F$, but also the action of a certain continuous group $U$ related
to the units $\sO_{F}^{\times }$.  One of Coulangeon's basic results
is that there are finitely many equivalence classes of perfect Humbert
forms modulo these actions.

On the other hand, Ash's original construction of retracts
\cite{ash77} introduces a geometric notion of perfection.  Namely he
generalizes the \Vor polyhedron $\Pi$ and defines a quadratic form to
be perfect if it naturally indexes a facet of $\Pi$.  What is the
connection between these two notions?  Can one use Coulangeon's
results to construct cell complexes to be used in cohomology
computations?  One tempting possibility is to try to use the group $U$ 
to collapse the \Vor cells of \cite{ash77} into a
cell decomposition of the symmetric space associated to $\SL_{n} (F)$.

\subsection{The Modular Complex}
%
In his study of multiple $\zeta$-values, Goncharov has recently defined
the \defn{modular complex} $M^{*}$ \cite{gonch1, gonch2}.  This is an
$n$-step complex of $\GL_{n} (\Z)$-modules closely related both to the
properties of multiple polylogarithms evaluated at $\mu_{N}$, the
$N$th roots of unity, and to the action of $G_{\Q}$ on $\pi_{1,N} =
\pi_{1}^{l} (\Proj^{1}\smallsetminus \{0,\infty ,\mu_{N} \})$, the
pro-$l$ completion of the algebraic fundamental group of
$\Proj^{1}\smallsetminus \{0,\infty ,\mu_{N} \}$.  

Remarkably, the modular complex is very closely related to the \Vor
decomposition $\sV$.  In fact, one can succinctly describe the modular
complex by saying that it is the chain complex of the cells coming
from the top-dimensional \Vor cone of type $A_{n}$.  This is all of
the \Vor decomposition for $n=2,3$, and Goncharov showed that the
modular complex is quasi-isomorphic to the full \Vor complex for
$n=4$.  Hence there is a precise relationship among multiple
polylogarithms, the Galois action on $\pi_{1,N}$, and the cohomology
of level $N$ congruence subgroups of $\SL_{n} (\Z)$.

The question then arises, how much of the cohomology of congruence
subgroups is captured by the modular complex for all $n$?  Table
\ref{vortable} indicates that asymptotically very little of the \Vor
decomposition comes from the $A_n$ cone, but this says nothing about
the cohomology.  The first interesting case to consider is $n=5$.


\subsection{Retracts for Other Groups}
%
The most general construction of retracts $W$ known \cite{ash.wrr}
applies only to \emph{linear} symmetric spaces.\index{linear symmetric spaces}  
The most familiar
example of such a space is $\SL_{n} (\R)/\SO (n)$; other examples are
the symmetric spaces associated to $\SL_{n}$ over number fields and
division algebras.

Now let $\Gamma\subset \bG (\Q )$ be an arithmetic group, and let $X =
G/K$ be the associated symmetric space.  What can one say about cell
complexes that can be used to compute $H^{*} (\Gamma; \MMM)$?  The
theorem of Borel--Serre mentioned in Section~\ref{ss:vcd} implies the
vanishing of $H^{k}(\Gamma; \MMM )$ for $k>\nu := \dim X - q$, where
$q$ is the \defn{\protect{$\Q$-rank}} of $\Gamma$.  For example, for the split
form of $\SL_{n}$, the $\Q$-rank is $n-1$.  For the split symplectic
group $\SP_{2n}$, the $\Q$-rank is $n$.  Moreover, this bound is
sharp: there will be coefficient modules $\MMM$ for which $H^{\nu}
(\Gamma ; \MMM) \not = 0$.  Hence any minimal cell complex used to compute the
cohomology of $\Gamma$ should have dimension $\nu$. 

Ideally one would like to see such a complex realized as a subspace of~$X$ 
and would like to be able to treat all finite index subgroups of
$\Gamma $ simultaneously.  This leads to the following question: is
there a $\Gamma$-equivariant deformation retraction of $X$ onto a
regular cell complex $W$ of dimension $\nu$?

For $\bG = \SP_{4}$, McConnell and MacPherson showed that the answer is
yes.  Their construction begins by realizing the symplectic symmetric
space $X_{\SP}$ as a subspace of the special linear symmetric space
$X_{\SL}$.  They then construct subsets of $X_{\SP}$ by intersecting
the \Vor cells in $X_{\SL}$ with $X_{\SP}$.  Through explicit
computations in coordinates they prove that these intersections are
cells and give a cell decomposition of $X_{\SP}$.  By taking an
appropriate dual complex (as suggested by Figures \ref{tess} and
\ref{retract2.fig} and as done in \cite{ash77}), they construct the
desired cell complex $W$.

Other progress has been recently made by Bullock \cite{bullock},
Bullock and Connell \cite{bullock.connell}, and Yasaki \cite{yasaki1,
yasaki2} in the case of groups of $\Q$-rank 1.  In particular, Yasaki
uses the \defn{tilings} of Saper \cite{saper} to construct an explicit
retract for the unitary group $\SU (2,1)$ over the Gaussian integers.
His method also works for Hilbert modular groups, although further
refinement may be needed to produce a regular cell complex.  Can one
generalize these techniques to construct retracts for groups of
arbitrary $\Q$-rank?  Is there an analogue of the \Vor decomposition
for these retracts (i.e., a dual cell decomposition of the symmetric
space)?  If so, can one generalize ideas in
Sections~\ref{mod.symb.section}--\ref{other.coho} and use 
that generalization to compute the
action of the Hecke operators on the cohomology?



\subsection{Deeper Cohomology Groups}
%
The algorithm in Section~\ref{other.coho} can be used
to compute the Hecke action on
$H^{\nu -1} (\Gamma)$.  For $n>4$, this group no longer contains
cuspidal cohomology classes.  Can one generalize this algorithm to
compute the Hecke action on deeper cohomology groups?  The first
practical case is $n=5$.  Here $\nu = 10$, and the highest degree in
which cuspidal cohomology can live is $8$.  This case is also
interesting since the cohomology of full level has been studied
\cite{gangl}.

Here are some indications of what one can expect.  The general
strategy is the same: for a $k$-sharbly $\xi$ representing a
class in $H^{\nu -k} (\Gamma)$, begin by $\Gamma$-equivariantly
choosing reducing points for the nonunimodular submodular symbols of
$\xi$.  This data can be packaged into a new $k$-sharbly cycle as in
Section~\ref{ss:buildingxiprime}ff, but the crosspolytopes must be replaced
with \defn{hypersimplices}.  By definition, the hypersimplex $\Delta
(n,k)$ is the convex hull in $\R^{n}$ of the points $\{\sum_{i\in I}
e_{i} \}$, where $I$ ranges over all order $k$ subsets of $\{1,\dotsc
,n \}$ and $e_{1},\dotsc ,e_{n}$ denotes the standard 
basis of $\R ^{n}$.

The simplest example is $n=2$, $k=2$.  From the point of view of
cohomology, this is even less interesting than $n=2$, $k=1$, since now
we are computing the Hecke action on $H^{-1} (\Gamma)$!  Nevertheless,
the geometry here illustrates what one can expect in general.

Each $2$-sharbly in the support of $\xi$ can be written as $
[v_{1},v_{2},v_{3},v_{4}]$ and determines six submodular symbols, of
the form $[v_{i}, v_{j}]$, $i\not =j$.  Assume for simplicity that all these
submodular symbols are nonunimodular.  Let $w_{ij}$ be the reducing
point for $[v_{i},v_{j}]$.  Then use the ten points $v_{i}, w_{ij}$ to
label the vertices of the hypersimplex $\Delta (5,2)$ as in Figure
\ref{hy1simp.fig} (note that $\Delta (5,2)$ is $4$-dimensional).

\begin{figure}[htb]
\psfrag{04}{$v_{1}$}
\psfrag{14}{$v_{2}$}
\psfrag{24}{$v_{3}$}
\psfrag{34}{$v_{4}$}
\psfrag{01}{$w_{12}$}
\psfrag{02}{$w_{13}$}
\psfrag{03}{$w_{14}$}
\psfrag{12}{$w_{23}$}
\psfrag{13}{$w_{24}$}
\psfrag{23}{$w_{34}$}
\begin{center}
\includegraphics[scale=0.4]{diagrams/hy1simp_2}
\end{center}
\caption{\label{hy1simp.fig}}
\end{figure}


The boundary of this hypersimplex gives the analogue of
\eqref{tfm}.  Which $2$-sharblies will appear in $\xi '$?  The boundary
$\partial \Delta (5,2)$ is a union of five tetrahedra and five octahedra.
The outer tetrahedron will not appear in $\xi '$, since that is the
analogue of the left side of \eqref{tfm}.  The four octahedra sharing a
triangular face with the outer tetrahedron also will not appear, since
they disappear when considering $\xi '$ modulo $\Gamma$.  The
remaining four tetrahedra and the central octahedron survive to $\xi '$
and constitute the right side of the analogue of \eqref{tfm}.  Note that we
must choose a simplicial subdivision of the central octahedron to write
the result as a $2$-sharbly cycle and that this must be done with
care since it introduces a new submodular symbol.

If some submodular symbols are unimodular, then again one must
consider iterated cones on hypersimplices, just as in
Section~\ref{sl2.stop}.  The analogues of these steps become more
complicated, since there are now many simplicial subdivisions of a
hypersimplex\footnote{Indeed, computing all simplicial subdivisions of
$\Delta (n,k)$ is a difficult problem in convex geometry.}. There is
one final complication: in general we cannot use reduced $k$-sharblies
alone to represent cohomology classes.  Thus one must terminate the
algorithm when $\|\xi \|$ is less than some predetermined bound.

\subsection{Other Linear Groups}
%
Let $F$ be a number field, and let $\bG = \bR_{F/\Q} (\SL_{n})$
(Example \ref{ex:resofscalars}).  Let $\Gamma \subset \bG (\Q)$ be an
arithmetic subgroup.  Can one compute the action of the
Hecke operators on $H^{*} (\Gamma)$?  

There are two completely different approaches to this problem.  The
first involves the generalization of the modular symbols method.  One
can define the analogue of the sharbly complex, and can try to extend
the techniques of Sections~\ref{mod.symb.section}--\ref{other.coho}.

This technique has been extensively used when 
$F$ is \emph{imaginary quadratic} and $n=2$.  We have $X = \SL_{2}
(\C)/\SU (2)$, which is isomorphic to $3$-dimensional hyperbolic
space $\fh_{3}$.  The arithmetic groups $\Gamma \subset \SL_{2}
(\sO_{F})$ are known as \defn{Bianchi groups}.  The retracts and
cohomology of these groups have been well studied; as a representative
sample of works we mention \cite{mendoza, egm, vogtmann,
grune.sch}.

Such groups have $\Q$-rank 1 and thus have cohomological dimension
$2$.  One can show that the cuspidal classes live in degrees $1$ and
$2$.  This means that we can use modular symbols to investigate the
Hecke action on cuspidal cohomology.  This was done by Cremona
\cite{crem} for \emph{euclidean} fields $F$.  In that case Theorem
\ref{thm:modsymbalg} works with no trouble (the euclidean algorithm is
needed to construct reducing points).  For noneuclidean fields
further work has been done by Whitley \cite{whit}, Cremona and Whitely
\cite{crem.whit} (both for principal ideal domains), Bygott
\cite{bygott} (for $F=\Q (\sqrt{-5})$ and any field with class group
an elementary abelian $2$-group), and Lingham \cite{lingham} (any
field with odd class number).  Putting all these ideas together allows
one to generalize the modular symbols method to \emph{any} imaginary
quadratic field \cite{crem.letter}.

For $F$ imaginary quadratic and $n>2$, very little has been studied.
The only related work to the best of our knowledge is that of
Staffeldt \cite{staffeldt}.  He determined the structure of the \Vor
polyhedron in detail for $\bR_{F/\Q} (\SL_{3})$, where $F = \Q
(\sqrt{-1})$.  We have $\dim X = 8$ and $\nu =6$.  The cuspidal
cohomology appears in degrees $3,4,5$, so one could try to use the
techniques of Section~\ref{other.coho} to investigate it.

Similar remarks apply to $F$ \emph{real quadratic} and $n=2$.  The
symmetric space $X \simeq \fh \times \fh$ has dimension $4$ and the
$\Q$-rank is 1, which means $\nu = 3$.  Unfortunately the cuspidal
cohomology appears only in degree $2$, which means modular symbols
cannot see it.  On the other hand, 1-sharblies can see it, and so one
can try to use ideas in Section~\ref{other.coho} here to compute the Hecke
operators.  The data needed to build the retract $W$ already
(essentially) appears in the literature for certain fields; see for
example \cite{ong}.

The second approach shifts the emphasis from modular symbols and the
sharbly complex to the \Vor fan and its cones.  For this approach we
must assume that the group $\Gamma$ is associated to a
\defn{self-adjoint homogeneous cone} over $\Q$.  (cf. \cite{ash77}).
This class of groups includes arithmetic subgroups of $\bR_{F/\Q}
(\SL_{n})$, where $F$ is a totally real or CM field.  Such groups have
all the nice structures in Section~\ref{vcd.section}.  For example, we have
a cone $C$ with a $G$-action.  We also have an analogue of the \Vor
polyhedron $\Pi$.  There is a natural compactification $\tilde{C}$ of
$C$ obtained by adjoining certain self-adjoint homogeneous cones of
lower rank.  The quotient $\Gamma \backslash \tilde{C}$ is singular in
general, but it can still be used to compute $H^{*} (\Gamma ; \C)$.  The
polyhedron $\Pi$ can be used to construct a fan $\sV$ that gives a
$\Gamma$-equivariant decomposition of all of $\tilde{C}$.  But the
most important structure we have is the \Vor reduction algorithm:
given any point $x\in \tilde{C}$, we can determine the unique \Vor
cone containing $x$.

Here is how this setup can be used to compute the Hecke action.  Full
details are in \cite{msa, ccalg}.  We define two chain complexes
$\bC^{V}_{*}$ and $\bC^{R}_{*}$.  The latter is essentially the chain
complex generated by all simplicial rational polyhedral cones in
$\tilde{C}$; the former is the subcomplex generated by the \Vor cones.
These are the analogues of the sharbly complex and the chain complex
associated to the retract $W$, and one can show that either can be used
to compute $H^{*} (\Gamma ; \C)$.  Take a cycle $\xi \in \bC^{V}_{*}$
representing a cohomology class in $H^{*} (\Gamma ; \C)$ and act on it
by a Hecke operator $T$.  We have $T (\xi) \in \bC^{R}_{*}$, and we
must push $T (\xi)$ back to $\bC^{V}_{*}$.

To do this, we use the linear structure on $\tilde{C}$ to subdivide $T
(\xi)$ very finely into a chain $\xi '$.  For each $1$-cone $\tau$ in
$\support \xi '$, we choose a $1$-cone $\rho_{\tau}\in
\tilde{C}\smallsetminus C$ and assemble them using the combinatorics
of $\xi '$ into a polyhedral chain $\xi ''$ homologous to $\xi'$.
Under certain conditions involved in the construction of $\xi '$, this
chain $\xi ''$ will lie in $\bC^{V}_{*}$.  

We illustrate this process for the split group $\SL_{2}$; more details
can be found in \cite{msa}.  We work modulo homotheties, so that the
three-dimensional cone $\tilde{C}$ becomes the extended upper
half plane $\fh^{*}:= \fh \cup \Q \cup \{\infty \} $, with $\partial
\tilde{C}$ passing to the cusps $\fh ^{*}\smallsetminus \fh $.  As
usual top-dimensional \Vor cones become the triangles of the Farey
tessellation, and the cones $\rho_{\tau}$ become cusps.  Given any
$x\in \fh$, let $R (x)$ be the set of cusps of the unique triangle or
edge containing $x$ (this can be computed using the \Vor reduction
algorithm).  Extend $R$ to a function on $\fh^{*}$ by putting $R (u) =
\{u \}$ for any cusp $u$.

In $\fh$, the support of $T (\xi)$ becomes a geodesic $\mu $ between
two cusps $u$, $u'$, in other words the support of a modular symbol
$[u,u' ]$ (Figure
\ref{partition.fig}).  Subdivide $\mu$ by choosing points
$x_{0},\dotsc ,x_{n}$ such that $x_{0}=u$, $x_{n} = u'$, and $R
(x_{i})\cap R (x_{i+1})\not =\emptyset$.  (This is easily done, for
example by repeatedly barycentrically subdividing $\mu$.)
For each $i<n$ choose a
cusp $q_{i}\in R (x_{i})\cap R (x_{i+1})$, and put $q_{n}=u'$.  Then
we have a relation in $H^{1}$:
\begin{equation}\label{eq:sahc.rel}
[u,u']=[q_{0},q_{1}]+\dotsb +[q_{n-1},q_{n}].
\end{equation}
Moreover, each $[q_{i},q_{i+1}]$ is unimodular, since $q_{i}$ and
$q_{i+1}$ are both vertices of a triangle containing $x_{i+1}$.
Upon lifting \eqref{eq:sahc.rel} back to $\bC^{R}_{*}$, the cusps
$q_{i}$ become the $1$-cones $\rho_{\tau}$ and give us a relation
$T (\xi) = \xi ''\in \bC^{V}_{*}$.

\begin{figure}[htb]
\psfrag{mu}{$\mu $}
\psfrag{u1}{$u$}
\psfrag{u2}{$u'$}
\centerline{\includegraphics[scale = .5]{diagrams/partition}}
\caption{A subdivision of $\mu $; the solid dots are the $x_{i}$.
Since the $x_{i}$ lie in the same or adjacent \Vor \ cells, we can
assign cusps to them to construct a homology to a cycle in $\bC ^{V}_{*}$.\label{partition.fig}}
\end{figure}

\subsection{The Sharbly Complex for General Groups}
%
In \cite{sympms} we generalized Theorem \ref{thm:modsymbalg} (without
the complexity statement) to the symplectic group $\SP_{2n}$.  Using
this algorithm and the symplectic retract \cite{mmc1, mmc2}, one can
compute the action of the Hecke operators on the top-degree cohomology
of subgroups of $\SP_{4} (\Z)$. 

More recently, Toth has investigated modular symbols for other groups.
He showed that the unimodular symbols generate the top-degree
cohomology groups for $\Gamma$ an arithmetic subgroup of a split
classical group or a split group of type $E_{6}$ or $E_{7}$
\cite{toth}.  His technique of proof is completely different from that
of \cite{sympms}.  In particular he does not give an analogue of the
Manin trick.  Can one extract an algorithm from Toth's proof that can
be used to explicitly compute the action of the Hecke operators on
cohomology?

The proof of the main result of \cite{sympms} uses a description of
the relations among the modular symbols.  These relations were
motivated by the structure of the cell complex in \cite{mmc1, mmc2}.
The modular symbols and these relations are analogues of the groups
$S_{0}$ and $S_{1}$ in the sharbly complex.  Can one 
extend these combinatorial constructions to form a \defn{symplectic
sharbly complex}?  What about for general groups $\bG$?  

Already for $\SP_{4}$, resolution of this question would have immediate
arithmetic applications.  Indeed, Harder has a beautiful conjecture
about certain congruences between holomorphic modular forms and Siegel
modular forms of full level \cite{harder-arbeit}.  Examples of these
congruences were checked numerically in \cite{harder-arbeit} using techniques of
\cite{faber-van.der.geer} to compute the Hecke action.

However, to investigate higher levels, one needs a different technique.
The relevant cohomology classes live in $H^{\nu -1} (\Gamma ; \MMM)$,
so one only needs to understand the first three terms of the complex
$S_{0}\leftarrow S_{1}\leftarrow S_{2}$.  We understand $S_{0}$,
$S_{1}$ from \cite{sympms}; the key is understanding $S_{2}$, which
should encode relations among elements of $S_{1}$.  If one could do
this and then could generalize the techniques of \cite{experimental},
one would have a way to investigate Harder's conjecture.

\subsection{Generalized Modular Symbols}
%
We conclude this appendix by discussing a geometric approach to
modular symbols.  This complements the algebraic approaches presented
in this book and leads to many new interesting phenomena and
problems.

Suppose $\bH$ and $\bG$ are connected semisimple algebraic groups over
$\Q$ with an injective map $f\colon \bH \rightarrow \bG$.  Let $K_{H}$
be a maximal compact subgroup of $H = \bH (\R)$, and suppose $K\subset
G$ is a maximal compact subgroup containing $f (K_{H})$.  Let $X =
G/K$ and $Y = H/K_{H}$.

Now let $\Gamma \subset \bG (\Q)$ be a torsion-free arithmetic
subgroup.  Let $\Gamma_{H} = f^{-1} (\Gamma)$.  We get a map
$\Gamma_{H}\backslash Y \rightarrow \Gamma \backslash X$, and we
denote the image by $S (H,\Gamma)$.  Any
compactly supported cohomology class $\xi \in H_{c}^{\dim Y} (\Gamma
\backslash X; \C)$ can be pulled back via $f$ to $\Gamma_{H}\backslash
Y$ and integrated to obtain a complex number.  Hence $S (H,\Gamma)$
defines a linear form on $H_{c}^{\dim Y} (\Gamma
\backslash X; \C)$.  By Poincar\'e duality, this linear form
determines a class $[S (H,\Gamma )] \in H^{\dim X - \dim Y} (\Gamma
\backslash X; \C)$, called a \defn{generalized
modular symbol}.  Such classes have been considered by many authors,
for example \cite{ash.borel, venky.speh, harder.modsym,
ash.ginzburg.rallis}.

As an example, we can take $\bG$ to be the split form of $\SL_{2}$,
and we can take $f \colon \bH \rightarrow \bG $ to be the inclusion of
connected component of the diagonal subgroup.  Hence $H \simeq
\R_{>0}$.  In this case $K_{H}$ is trivial.  The image of $Y$ in $X$
is the ideal geodesic from $0$ to $\infty$.  One way to vary $f$ is by
taking an $\SL_{2} (\Q)$-translate of this geodesic, which gives a
geodesic between two cusps.  Hence we can obtain the support of any
modular symbol this way.  This example generalizes to $\SL_{n}$ to
yield the modular symbols in Section~\ref{mod.symb.section}.  Here $H \simeq
(\R >0)^{n-1}$.  Note that $\dim Y = n-1$, so the cohomology classes
we have constructed live in the top degree $H^{\nu} (\Gamma \backslash
X; \C)$.

Another family of examples is provided by taking $\bH$ to be a Levi
factor of a parabolic subgroup; these are the modular symbols studied
in \cite{ash.borel}.

There are many natural questions to study for such objects.  Here are
two:

\begin{itemize}
\item Under what conditions on $\bG , \bH , \Gamma$ is $[S
(H,\Gamma)]$ nonzero?  This question is connected to relations between periods
of automorphic forms and functoriality lifting.  There are a variety
of partial results known; see for example \cite{venky.speh, ash.ginzburg.rallis}.
\item We know the usual modular symbols span the top-degree cohomology
for any arithmetic group $\Gamma$.  Fix a class of generalized modular
symbols by fixing the pair $\bG , \bH $ and fixing some class of maps
$f$.  How much of the cohomology can one span for a general arithmetic
group $\Gamma \subset \bG (\Q)$?

A simple example is given by the Ash--Borel construction for $\bG =
\SL_{3}$ and $\bH$ a Levi factor of a rational parabolic subgroup
$\bP$ of type $(2,1)$.  In this case $H \simeq \SL_{2} (\R) \times
\R_{>0}$ and sits inside $G$ via
\[
g\left( 
\begin{array}{cc}
\alpha^{-1}M&0\\
0&\alpha
\end{array}\right)g^{-1}, \quad M\in \SL_{2} (\R), \quad \alpha \in
\R_{>0}, \quad g \in \SL_{3} (\Q).
\]
For $\Gamma \subset \SL_{3} (\Z)$ these symbols define a subspace 
\[
S_{(2,1)}\subset H^{2} (\Gamma
\backslash X; \C).
\]
Are there $\Gamma $ for which $S_{(2,1)}$ equals
the full cohomology space?  For general $\Gamma$ how much is
captured?  Is there
a nice combinatorial way to write down the relations among these
classes?  Can one cook up a generalization of Theorem
\ref{thm:modsymbalg} for these classes and use it to compute Hecke
eigenvalues?

\end{itemize}




\backmatter

\bibliography{biblio}

\printindex

\end{document}

%-----------------------------------------------------------------------------
% End of main.tex
%-----------------------------------------------------------------------------
